% Propietario conservado: /Users/ruben/Documents/New project/output/EXTREMA_MEDIA_RAZON_RECIPROCIDAD_ES_EN_20260918/ARTICULO/ES/source/nuclear/antecedentes_fisicos.tex
\chapter{Realización canónica de la covarianza}
\label{x_reciprocidad:nuclear:covarianza-antecedente}

% BEGIN OWNER /Users/ruben/Documents/New project/output/EXTREMA_MEDIA_RAZON_RECIPROCIDAD_ES_EN_20260918/ARTICULO/ES/source/antecedentes_ii/sections/07c_covarianza_accion.tex
% Propuesta de inserción después de 07b_geometria_elipse.tex.
% RESULTADO_RECUPERADO: c52_06_01_incertidumbre.tex, líneas 10--91.
% Dependencias: c52_04_01_ley_areal.tex, líneas 137--198;
% c52_06_02_intercambio_constitutivo.tex, líneas 5--29.
% Explicitación local: dominio operatorio, prueba de covarianza,
% transporte recíproco y distinción dimensional posición--momento.
% Este fragmento no introduce una constante ni un generador adicional.

\section{Covarianza canónica y reciprocidad de la elipse de acción}
\label{x_reciprocidad:sec:ii-area-covarianza}

La construcción precedente determina dos invariantes complementarios:
el producto de los semiejes fija el área de acción y su cociente conserva
la anisotropía del par angular. Precisamos ahora la realización canónica
en la que esos semiejes son dos desviaciones estándar. La estructura de
partida es la elipse ya obtenida de las publicaciones angulares y de
acción de APP--TRIT--TPK; la realización operatoria se establece después
de esa construcción.

En este apartado escribimos \(\eta=\eta_{\rm el}
=\operatorname{artanh}(C^*/A)\), con \(|C^*/A|<1\) y
\(\hbar>0\). El cociente angular utiliza una misma unidad para sus dos
componentes. En las coordenadas equilibradas \((Q,P)\), ambas de
dimensión raíz de acción, definimos
\begin{equation}
 V_\eta:=\frac14
 \begin{pmatrix}a_S^2&0\\0&b_S^2\end{pmatrix}
 =\frac{\hbar}{2}
 \begin{pmatrix}e^\eta&0\\0&e^{-\eta}\end{pmatrix}.
 \label{x_reciprocidad:eq:ii-covarianza-geometrica}
\end{equation}
Esta matriz positiva tiene entradas de dimensión acción. Su interpretación
como covarianza cuántica queda especificada por la realización siguiente.

\begin{proposition}[Realización gaussiana de la covarianza de acción]
\label{x_reciprocidad:prop:ii-covarianza-gaussiana}
Considérese la representación canónica en \(L^2(\mathbb R,dQ)\), con
\(\widehat Qf(Q)=Qf(Q)\) y
\(\widehat Pf(Q)=-i\hbar f'(Q)\), inicialmente sobre el espacio de
Schwartz \(\mathcal S(\mathbb R)\). El estado normalizado
\begin{equation}
 \psi_\eta(Q)=\bigl(\pi\hbar e^\eta\bigr)^{-1/4}
 \exp\!\left(-\frac{Q^2}{2\hbar e^\eta}\right)
 \label{x_reciprocidad:eq:ii-gaussiana-accion}
\end{equation}
tiene medias nulas y matriz de covarianza \(V_\eta\). En particular,
\begin{equation}
 (\Delta Q)^2=\frac{\hbar e^\eta}{2},\qquad
 (\Delta P)^2=\frac{\hbar e^{-\eta}}{2},\qquad
 \operatorname{Cov}(Q,P)=0,\qquad
 \det V_\eta=\frac{\hbar^2}{4}.
 \label{x_reciprocidad:eq:ii-varianzas-accion}
\end{equation}
Esta realización satura la desigualdad de Robertson--Schrödinger.
\end{proposition}

\begin{proof}
Los operadores anteriores son esencialmente autoadjuntos sobre
\(\mathcal S(\mathbb R)\), que constituye un núcleo común invariante;
sobre este espacio se cumple
\([\widehat Q,\widehat P]=i\hbar I\). Para
\(\widehat Y=(\widehat Q,\widehat P)\), la covarianza simétrica es
\[
 V_{jk}=\frac12\left\langle
 (\widehat Y_j-\langle\widehat Y_j\rangle)
 (\widehat Y_k-\langle\widehat Y_k\rangle)
 +(\widehat Y_k-\langle\widehat Y_k\rangle)
 (\widehat Y_j-\langle\widehat Y_j\rangle)
 \right\rangle.
\]
Con \(s=\hbar e^\eta>0\), sea
\(I_s=\int_{\mathbb R}e^{-Q^2/s}\,dQ\). El producto de dos
integrales iguales, evaluado en coordenadas polares, da
\(I_s^2=2\pi\int_0^\infty r e^{-r^2/s}\,dr=\pi s\).
Además, la integral de
\(\frac{d}{dQ}(Qe^{-Q^2/s})\) es cero, por el decaimiento en ambos
extremos; de ahí
\(\int_{\mathbb R}Q^2e^{-Q^2/s}\,dQ=sI_s/2\).
Estas identidades y la paridad prueban la normalización, la media nula y
\(\langle\widehat Q^2\rangle=s/2\). Además,
\(\widehat P\psi_\eta=i e^{-\eta}Q\psi_\eta\).
Por tanto, \(\langle\widehat P\rangle=0\),
\(\|\widehat P\psi_\eta\|^2=\hbar e^{-\eta}/2\) y
\(\operatorname{Re}\langle\widehat Q\psi_\eta,
\widehat P\psi_\eta\rangle=0\).

Para cualquier estado normalizado de ese dominio, Cauchy--Schwarz aplicada
a \((\widehat Q-\langle\widehat Q\rangle)\psi\) y
\((\widehat P-\langle\widehat P\rangle)\psi\) da
\[
 (\Delta Q)^2(\Delta P)^2
 \geq \operatorname{Cov}(Q,P)^2
       +\frac14\left|\langle[\widehat Q,\widehat P]\rangle\right|^2
 =\operatorname{Cov}(Q,P)^2+\frac{\hbar^2}{4}.
\]
Las varianzas calculadas alcanzan la igualdad. Equivalentemente, el
operador
\[
 a_\eta=\frac{e^{-\eta/2}\widehat Q
                  +i e^{\eta/2}\widehat P}{\sqrt{2\hbar}}
\]
satisface \([a_\eta,a_\eta^\dagger]=I\) y
\(a_\eta\psi_\eta=0\), por sustitución directa.
\end{proof}

\begin{proposition}[Área, dispersión y transporte recíproco]
\label{x_reciprocidad:prop:ii-area-covarianza-reciprocidad}
La elipse de acción es exactamente el subnivel de covarianza
\begin{equation}
 \mathcal E_\eta
 =\left\{y=(Q,P)^{\mathsf T}:y^{\mathsf T}V_\eta^{-1}y\leq4\right\}
 =\left\{\frac{Q^2}{a_S^2}+\frac{P^2}{b_S^2}\leq1\right\}.
 \label{x_reciprocidad:eq:ii-elipse-covarianza}
\end{equation}
Su área es \(h=2\pi\hbar\), sus semiejes son
\(a_S=2\Delta Q\), \(b_S=2\Delta P\) y
\begin{equation}
 \frac{\Delta P}{\Delta Q}=e^{-\eta}
 =\frac{b_S}{a_S}=R_{\rm el}^2.
 \label{x_reciprocidad:eq:ii-covarianza-forma}
\end{equation}
La rotación canónica \(J_2\) de la sección~\ref{x_reciprocidad:iv:ellipse:section} transporta
\(V_\eta\) a \(V_{-\eta}\) y conserva la forma simpléctica.
\end{proposition}

\begin{proof}
La inversión de la matriz diagonal
\eqref{x_reciprocidad:eq:ii-covarianza-geometrica} da
\eqref{x_reciprocidad:eq:ii-elipse-covarianza}. Sus semiejes son las raíces positivas de
cuatro veces sus varianzas, y
\[
 \operatorname{Area}(\mathcal E_\eta)
 =4\pi\sqrt{\det V_\eta}=2\pi\hbar=h,
\]
en concordancia con la ley areal~\eqref{x_reciprocidad:iv:ellipse:area-fase}.
La razón de las raíces positivas prueba
\eqref{x_reciprocidad:eq:ii-covarianza-forma}.
Si \(M_\eta=\operatorname{diag}(e^{\eta/2},e^{-\eta/2})\), entonces
\[
 V_\eta=M_\eta V_0M_\eta^{\mathsf T},\qquad
 M_\eta^{\mathsf T}J_2M_\eta=J_2,\qquad
 J_2V_\eta J_2^{\mathsf T}=V_{-\eta}.
\]
Estas identidades se obtienen por multiplicación. En términos operatorios,
\((\widehat Q',\widehat P')=(-\widehat P,\widehat Q)\) conserva
\([\widehat Q',\widehat P']=i\hbar I\). La permutación
\(S_2:(Q,P)\mapsto(P,Q)\) también intercambia las varianzas, pero invierte
el signo del conmutador y de la forma simpléctica. Así, el transporte de
covarianza respeta la distinción entre reciprocidad canónica e inversión
de orientación ya establecida en la sección~\ref{x_reciprocidad:iv:ellipse:section}.
\end{proof}

\subsubsection{Posición, momento y transducción dimensional.}
Sea \(\lambda_{xp}>0\) un factor de la carta dimensional, de dimensión
momento por unidad de longitud. La transformación
\[
 Q=\sqrt{\lambda_{xp}}\,x,\qquad
 P=\frac{p_x}{\sqrt{\lambda_{xp}}}
\]
realiza el par equilibrado en una posición \(x\) y su momento conjugado
\(p_x\); conserva \(dQ\wedge dP=dx\wedge dp_x\) y
\([\widehat x,\widehat p_x]=i\hbar I\). Por transformación de la
covarianza,
\begin{equation}
 V_\eta^{(x,p_x)}=\frac{\hbar}{2}
 \begin{pmatrix}e^\eta/\lambda_{xp}&0\\
                 0&\lambda_{xp}e^{-\eta}\end{pmatrix},\qquad
 \Delta x\,\Delta p_x=\frac{\hbar}{2},\qquad
 \frac{\Delta p_x}{\Delta x}=\lambda_{xp}e^{-\eta}.
 \label{x_reciprocidad:eq:ii-covarianza-dimensiones}
\end{equation}
El cociente de dispersión adquiere aquí la dimensión momento por longitud.
En cambio, \eqref{x_reciprocidad:eq:ii-covarianza-forma} pertenece a la carta equilibrada.
Para la realización inductivo--capacitiva ya definida,
\eqref{x_reciprocidad:hmtII:eq:ii-elipse-impedancia} aporta la identidad adimensional
\(Z_\eta/Z_*=e^{-\eta}=\Delta P/\Delta Q\).
Cada lectura utiliza su transducción declarada y conserva el mismo
determinante de acción.

\subsubsection{Energía del contorno y energía media del estado.}
La realización inductivo--capacitiva tiene el Hamiltoniano
\[
 H_\eta(Q,P)=\frac{\omega}{2}
 \left(e^{-\eta}Q^2+e^\eta P^2\right),\qquad \omega>0.
\]
Sobre \(\partial\mathcal E_\eta\), su valor es
\(H_\eta=\hbar\omega\), por las ecuaciones de los semiejes. En la
realización cuántica de la proposición~\ref{x_reciprocidad:prop:ii-covarianza-gaussiana},
\begin{equation}
 \langle\widehat H_\eta\rangle_{\psi_\eta}
 =\frac{\omega}{2}
 \left(e^{-\eta}\frac{\hbar e^\eta}{2}
       +e^\eta\frac{\hbar e^{-\eta}}{2}\right)
 =\frac{\hbar\omega}{2}.
 \label{x_reciprocidad:eq:ii-energia-contorno-covarianza}
\end{equation}
El contorno de dos desviaciones estándar tiene, por tanto, el doble de la
energía media de este estado gaussiano. La ley areal fija la normalización
de acción; el par angular fija su distribución entre las dos coordenadas;
la representación canónica y el estado explícito realizan esa distribución
como covarianza. La saturación de incertidumbre queda así acreditada en
la realización especificada por sus operadores y su estado.

% END OWNER


% BEGIN OWNER /Users/ruben/Documents/New project/output/EXTREMA_MEDIA_RAZON_RECIPROCIDAD_ES_EN_20260918/ARTICULO/ES/source/antecedentes_iii/sections/03_hojas_y_orientacion.tex
\chapter[Retorno de hoja, incidencia areal--volumétrica y orientación constitutiva]{Retorno de hoja, incidencia areal--volumétrica\\y orientación constitutiva}
\label{x_reciprocidad:iii:sec:hojas}
\label{x_reciprocidad:nuclear:constitutiva-antecedente}

La estructura de hojas interviene antes de la evaluación de las magnitudes
electromagnéticas. Su función consiste en distinguir el cierre de una proyección,
la restitución de una orientación y la evolución del registro transportado.
Estas operaciones actúan sobre objetos relacionados, pero de distinto tipo: una
trayectoria proyectada puede haberse cerrado mientras su levantamiento permanece
en la hoja opuesta; la orientación puede haberse restituido mientras la memoria
conserva las dos vueltas efectuadas. La construcción constitutiva debe recibir
esta información con la estructura que la determina.

\section{Retorno proyectado y restitución de la cubierta orientada}

Sea $\ell$ un lazo de retorno de una extensión superviviente y sea
$\varepsilon:\langle\ell\rangle\to\{+1,-1\}$ el carácter de su
levantamiento, con $\varepsilon(\ell)=-1$. El fibrado de intervalos asociado
tiene la presentación
\begin{equation}
 \mathcal M=([0,1]\times[-1,1])/{(0,u)\sim(1,-u)}.
 \label{x_reciprocidad:iii:eq:mobius}
\end{equation}
El dato decisivo es el signo del transporte sobre el lazo especificado. La
identificación de los extremos realiza geométricamente ese carácter.

\begin{proposition}[Dos órdenes de clausura]
Una vuelta del lazo cierra la base e invierte la orientación de la fibra. Dos
vueltas restituyen la orientación. En la parametrización angular del lector de
borde, estos retornos corresponden, respectivamente, al cierre areal $2\pi$ y
al cierre volumétrico $4\pi$ de la cubierta orientada.
\end{proposition}
\begin{proof}
La relación de equivalencia de \eqref{x_reciprocidad:iii:eq:mobius} identifica la fibra terminal
con la inicial mediante $u\mapsto-u$. La composición de dos transportes es
$u\mapsto-(-u)=u$. Parametrizar una vuelta de la base por $2\pi$ asigna $4\pi$
a su cuadrado. La afirmación se refiere a este levantamiento y a esta lectura.
\end{proof}

La memoria del TPK conserva, además, la trayectoria recorrida y las actualizaciones
del registro. El cuadrado del carácter de signo es la identidad, pero no impone
que se anulen esas actualizaciones. La representación espinorial, cuando se
utilice, deberá transportar este dato mediante su aplicación propia.

\section{Descenso de la incidencia desde el calendario trítico}

En el bloque primitivo $(+1,-1,0)$, el TPK actualiza el modo aritmético mediante
$+1:m\mapsto\Sigma$, $-1:m\mapsto\Pi$ y $0:m\mapsto m$. Su órbita cíclica
es $(\Sigma,\Pi,\Pi)$. El lector de hojas asigna $\Sigma$ a la hoja areal y
$\Pi$ a la volumétrica. En el orden de bases $(e_{\rm ar},e_{\rm vol})$, la
incidencia de aristas queda representada por
\begin{equation}
 F_{\rm av}=\begin{pmatrix}0&1\\1&1\end{pmatrix}.
 \label{x_reciprocidad:iii:eq:fav}
\end{equation}
Los tres pares consecutivos son $\Sigma\to\Pi$, $\Pi\to\Pi$ y
$\Pi\to\Sigma$, cada uno con multiplicidad uno. Los calendarios de longitudes
$9$, $27$ y $108$ contienen $3$, $9$ y $36$ copias del bloque primitivo. Sus
matrices de incidencia son, por tanto, esos múltiplos de $F_{\rm av}$.

Esta construcción conserva la diferencia entre frecuencia cronológica y medida
del espacio de continuaciones. La primera asigna pesos $(1/3,2/3)$ al ciclo
de modos. La segunda se define sobre la envolvente simbólica mínima de memoria
uno que admite exactamente los pares anteriores. Su matriz de adyacencia es
$F_{\rm av}$; los recuentos de trayectorias pertenecen a esa envolvente.

\Needspace{10\baselineskip}
\begin{proposition}[Pesos de retorno de las dos hojas]
Sean $F_0=0$, $F_1=1$ y $F_{n+1}=F_n+F_{n-1}$. Para $n\geq1$,
\begin{equation}
 F_{\rm av}^{n}=
 \begin{pmatrix}F_{n-1}&F_n\\F_n&F_{n+1}\end{pmatrix}.
 \label{x_reciprocidad:iii:eq:fav-powers}
\end{equation}
La medida de Parry de esta envolvente asigna a las hojas pesos proporcionales
a $(1,\varphi^2)$, y la razón de recuentos diagonales converge a $\varphi^{-2}$.
\end{proposition}
\begin{proof}
La fórmula de potencias se prueba por inducción utilizando la recurrencia.
El cuadrado de $F_{\rm av}$ es estrictamente positivo. Su vector positivo de
Perron es $(1,\varphi)$, con autovalor $\varphi$, pues
$\varphi^2=\varphi+1$. Al ser simétrica la matriz, los pesos de Parry son
proporcionales a los cuadrados de las componentes de ese vector. La otra raíz
del polinomio característico es $-\varphi^{-1}$; su módulo es estrictamente
menor que $\varphi$, lo que da el límite de los cocientes diagonales.
\end{proof}

La identificación espectral de esta incidencia es posterior a su construcción
desde el calendario. La coordenada regional $\varphi$ conserva su genealogía
coinductiva anterior: el reconocimiento del autovalor no selecciona el calendario.

\section{Marcado de la frontera y razón areal--volumétrica}

Sea $\Pi_{\rm HMT}(X_\infty)$ el valor del lector regional de clausura ya
construido. Con los proyectores diagonales de hoja, se define
\begin{equation}
 D_\pi=\Pi_{\rm HMT}(X_\infty)P_{\rm ar}+P_{\rm vol},\qquad
 \Theta_n=
 \frac{\langle e_{\rm ar},D_\pi F_{\rm av}^{n}e_{\rm ar}\rangle}
 {\langle e_{\rm vol},F_{\rm av}^{n}e_{\rm vol}\rangle}.
 \label{x_reciprocidad:iii:eq:marked-return}
\end{equation}
La normalización volumétrica es unitaria; el cierre regional marca una sola vez
el retorno areal. Aplicando \eqref{x_reciprocidad:iii:eq:fav-powers},
\begin{equation}
 \Theta_n=\Pi_{\rm HMT}(X_\infty)\frac{F_{n-1}}{F_{n+1}}
 \longrightarrow\frac{\pi}{\varphi^2}.
 \label{x_reciprocidad:iii:eq:av-limit}
\end{equation}
La razón de retorno queda determinada por la incidencia, la medida y el marcado.
El cociente $2\pi/4\pi$ describe, en cambio, los dos órdenes de clausura de la
cubierta. Las dos relaciones participan de la lectura de hojas sin sustituirse.

\section{Transporte de la orientación a los canales constitutivos}

En una representación de dos canales, sea $\mathcal R$ una involución
autoadjunta y $\mathcal S$ un operador unitario de intercambio que satisface
$\mathcal S\mathcal R\mathcal S^{-1}=-\mathcal R$. Las coordenadas angulares
ya construidas se escriben $x=A_{\rm rad}$ e $y=C^*_{\rm rad}$, con $x>|y|$.
Entonces
\begin{equation}
 B=xI+y\mathcal R,\qquad T=e^{-B},\qquad
 \mathcal ST(y)\mathcal S^{-1}=T(-y).
 \label{x_reciprocidad:iii:eq:channel-conjugation}
\end{equation}
En efecto, la conjugación transforma $B(y)$ en $B(-y)$ y conmuta con su serie
exponencial convergente. Los autovalores $q_\pm=e^{-x\mp y}$ se intercambian.
Para fijar el orden de las respuestas se utiliza la cámara $x>y>0$, donde
$0<q_+<q_-<1$. La aplicación $(x,y)\mapsto(x,-y)$ la transporta a la
cámara opuesta $x>-y>0$, donde $0<q_-<q_+<1$. Ambas pertenecen al dominio
contractivo $x>|y|$. En la frontera $y=0$ los dos canales coinciden.
La conjugación intercambia las cámaras, no conserva el orden estricto de
sus etiquetas; sobre los proyectores actúa mediante
$\mathcal SP_\pm\mathcal S^{-1}=P_\mp$, con
$P_\pm=(I\pm\mathcal R)/2$.
Esta aplicación concreta transmite la orientación al operador constitutivo que
se desarrolla a continuación.

% END OWNER


% BEGIN OWNER /Users/ruben/Documents/New project/output/EXTREMA_MEDIA_RAZON_RECIPROCIDAD_ES_EN_20260918/ARTICULO/ES/source/antecedentes_iii/sections/03b_vacancias_prefactor.tex
% Ampliación propia de III; no modifica el núcleo común copiado de I.
% Residencia propuesta: antes de Retorno de hoja, incidencia areal--volumétrica.
% Procedencia y controles: VACANCIAS_PREFACTOR_PROCEDENCIA.md, en esta carpeta.
\chapter{Polarización discreta y modo trítico de las proyecciones aritméticas}
\label{x_reciprocidad:iii:sec:vacancias-prefactor}

El refinamiento y la comparación entre las hojas de APP proporcionan dos
observaciones complementarias. La primera sigue la discrepancia acumulada
entre el número de transiciones y la capacidad adquirida; la segunda conserva
el modo del selector trítico en el emparejamiento de las fibras aditivas y
multiplicativas. La polarización discreta y el prefactor electrónico se
obtienen de estas operaciones respectivas. Su articulación con el vacío
requiere mantener tanto la información temporal como los espacios sobre los
que actúan los operadores.

\section{Discrepancia del refinamiento y convención de extremos}

Los cardinales del bloque ternario y de su representación decimal fijan
\[
 \rho=\log_{1000}729=\log_{10}9,\qquad
 \delta=1-\rho=\log_{10}(10/9).
\]
Se tiene $0<\delta<1$. Además, $\delta$ es irracional: una igualdad
$\delta=p/q$, con $q>0$, implicaría $10^{q-p}=9^q$, incompatible con
la descomposición en factores primos. Este desajuste ya interviene en la
codificación del refinamiento. Aquí se explicita su lectura centrada,
conservando el origen de fase y la convención de extremos.

\begin{definition}[Codificaciones inferior y superior]
Para $\theta\in\mathbb R$ y $n\in\mathbb Z$, se definen
\begin{align}
 z_n^{\downarrow}(\theta)
 &=\lfloor(n+1)\delta+\theta\rfloor-\lfloor n\delta+\theta\rfloor,
 \label{x_reciprocidad:iii:vp:lower}\\
 z_n^{\uparrow}(\theta)
 &=\lceil(n+1)\delta+\theta\rceil-\lceil n\delta+\theta\rceil.
 \label{x_reciprocidad:iii:vp:upper}
\end{align}
Ambas toman valores en $\{0,1\}$. La flecha distingue exclusivamente la
convención de extremos de los intervalos; la orientación trítica conserva
su definición anterior.
\end{definition}

La codificación inferior coincide con la utilizada en el desarrollo
del refinamiento. La superior permite recuperar la formulación de la
polarización con función techo. Las dos expresiones se relacionan por
un término de frontera exactamente determinado.

\begin{proposition}[Cambio de convención y término de frontera]
Sea $b_n(\theta)=\mathbf1_{\mathbb Z}(n\delta+\theta)$. Entonces
\begin{equation}
 z_n^{\uparrow}(\theta)-z_n^{\downarrow}(\theta)
 =b_n(\theta)-b_{n+1}(\theta).
 \label{x_reciprocidad:iii:vp:endpoint}
\end{equation}
En particular, para $N\geq0$,
\[
 \sum_{n=0}^{N-1}(z_n^{\uparrow}-z_n^{\downarrow})
 =b_0-b_N.
\]
\end{proposition}
\begin{proof}
Para todo real $x$, $\lceil x\rceil=\lfloor x\rfloor+1-
\mathbf1_{\mathbb Z}(x)$. Aplicar esta identidad a los dos extremos
de cada incremento da \eqref{x_reciprocidad:iii:vp:endpoint}; su suma es telescópica.
La irracionalidad de $\delta$ implica que, fijada $\theta$, como máximo
un índice entero satisface $n\delta+\theta\in\mathbb Z$.
\end{proof}

Con $\theta=0$, el primer símbolo superior es $1$ y el inferior es $0$;
para $n\geq1$ coinciden. El término inicial pertenece al registro y se
conserva cuando se cambia la convención. Las densidades y las cotas de
discrepancia se transportan con ese término, sin modificar la actualización
del TPK ni reiniciar la historia en cada retorno de fase.

\section{Polarización acotada y balance temporal exacto}

\begin{definition}[Polarización discreta centrada]
La polarización de la codificación superior, respecto de la densidad
$\delta$, es
\begin{equation}
 P_N^{\uparrow}(\theta)
 =\sum_{n=0}^{N-1}\bigl(z_n^{\uparrow}(\theta)-\delta\bigr),
 \qquad P_0^{\uparrow}=0.
 \label{x_reciprocidad:iii:vp:polarization}
\end{equation}
Su incremento temporal es
\begin{equation}
 J_N^{\uparrow}(\theta)
 =P_{N+1}^{\uparrow}(\theta)-P_N^{\uparrow}(\theta)
 =z_N^{\uparrow}(\theta)-\delta.
 \label{x_reciprocidad:iii:vp:current}
\end{equation}
\end{definition}

El centrado separa el crecimiento uniforme de densidad $\delta$ de la
información de fase. Así, $P_N^{\uparrow}$ mide un desequilibrio acumulado,
no el número total de vacancias. El incremento $J_N^{\uparrow}$ conserva
el pulso local y el fondo respecto del cual se evalúa.

\begin{theorem}[Cota uniforme y continuidad discreta temporal]
Para $N\geq0$ y toda fase $\theta$,
\begin{equation}
 P_N^{\uparrow}(\theta)
 =\lceil N\delta+\theta\rceil-\lceil\theta\rceil-N\delta,
 \qquad |P_N^{\uparrow}(\theta)|<1.
 \label{x_reciprocidad:iii:vp:bound}
\end{equation}
Para cualesquiera enteros $0\leq a<b$,
\begin{equation}
 \sum_{n=a}^{b-1}J_n^{\uparrow}(\theta)
 =P_b^{\uparrow}(\theta)-P_a^{\uparrow}(\theta),
 \qquad
 \left|\sum_{n=a}^{b-1}J_n^{\uparrow}(\theta)\right|<1.
 \label{x_reciprocidad:iii:vp:continuity}
\end{equation}
La densidad de la sucesión de vacancias es $\delta$.
\end{theorem}
\begin{proof}
La suma de los incrementos de techo de \eqref{x_reciprocidad:iii:vp:upper} produce la
primera identidad. Escribiendo $r(x)=\lceil x\rceil-x\in[0,1)$,
se obtiene $P_N^{\uparrow}=r(N\delta+\theta)-r(\theta)$; su valor
absoluto es menor que uno. Sobre el intervalo $[a,b)$ se obtiene
análogamente $r(b\delta+\theta)-r(a\delta+\theta)$, lo que prueba
la segunda cota, más precisa que la suma de dos cotas individuales.
Finalmente,
$N^{-1}\sum_{n=0}^{N-1}z_n^{\uparrow}
=\delta+N^{-1}P_N^{\uparrow}\to\delta$.
\end{proof}

La codificación inferior admite las mismas demostraciones sustituyendo
$r(x)$ por $\lfloor x\rfloor-x\in(-1,0]$. En particular,
$P_N^{\uparrow}-P_N^{\downarrow}=b_0-b_N$. La discrepancia acotada
utilizada anteriormente en las lecturas de Dirichlet y la polarización
presente conservan, por tanto, una relación explícita de origen y extremos.
La aperiodicidad tampoco introduce una deriva secular de este desequilibrio:
la historia puede prolongarse indefinidamente mientras la discrepancia
centrada permanece uniformemente acotada.

\section{Realización en las rectas de carga y tiempo}

Para realizar dimensionalmente el balance se especifican una recta de carga
$\mathcal L_Q$ con sección no nula $u_Q$ y una sucesión creciente de
instantes $t_N$ en una recta temporal orientada. Con
$\Delta t_N=t_{N+1}-t_N>0$, la carga centrada y la corriente media de
cada intervalo quedan definidas por
\begin{equation}
 Q_N=P_N^{\uparrow}u_Q,\qquad
 I_N=\frac{Q_{N+1}-Q_N}{\Delta t_N}
     =\frac{u_Q}{\Delta t_N}\bigl(z_N^{\uparrow}-\delta\bigr).
 \label{x_reciprocidad:iii:vp:dimensional}
\end{equation}
Así, $Q_N\in\mathcal L_Q$ e
$I_N\in\mathcal L_Q\otimes\mathcal L_T^{-1}$. La sección $u_Q$ y
los intervalos temporales son datos de esta realización; la construcción
aritmética precedente ya ha determinado $P_N^{\uparrow}$ y
$J_N^{\uparrow}$.

\begin{proposition}[Conservación del balance bajo realización dimensional]
Para $a<b$,
\begin{equation}
 Q_b-Q_a=\sum_{n=a}^{b-1}\Delta t_n I_n,
 \qquad
 \left|\frac{Q_b-Q_a}{u_Q}\right|<1.
 \label{x_reciprocidad:iii:vp:dimensional-balance}
\end{equation}
Si todos los intervalos tienen longitud $t_0$, las dos corrientes posibles
son $-\delta u_Q/t_0$ y $(1-\delta)u_Q/t_0$.
\end{proposition}
\begin{proof}
Multiplicar cada identidad de \eqref{x_reciprocidad:iii:vp:dimensional} por
$\Delta t_n$ y sumar telescopa las cargas. La cota es
\eqref{x_reciprocidad:iii:vp:continuity}; los dos valores locales se obtienen de
$z_n^{\uparrow}\in\{0,1\}$.
\end{proof}

La continuidad demostrada es temporal y está referida a la carga centrada.
Una realización espacial incorpora además la localización de las cargas,
la incidencia entre celdas y la identificación de los flujos de frontera.
En una descripción constitutiva por canales, el acoplamiento se especifica
mediante una aplicación de excitación
$\iota:\mathcal L_Q\to\mathcal H_{\rm ch}\otimes\mathcal L_Q$
y una aplicación de observación
$\ell:\mathcal H_{\rm ch}\otimes\mathcal L_Q\to\mathcal L_Q$,
donde $\mathcal H_{\rm ch}$ es una realización real del espacio de canales.
Si ambas son
lineales e independientes del índice temporal, la respuesta de un operador
$\mathcal V$ fijo se transporta por
\[
 Q_N^{\rm resp}=\ell(\mathcal V\otimes I_{\mathcal L_Q})\iota(Q_N),
 \qquad
 I_N^{\rm resp}=\frac{Q_{N+1}^{\rm resp}-Q_N^{\rm resp}}{\Delta t_N}.
\]
La linealidad conserva el balance telescópico. Esta composición hace
explícitos los datos de una identificación constitutiva: la secuencia
escalar determina el pulso, mientras $\iota$, $\mathcal V$ y $\ell$
determinan su excitación, transporte y observación por canales. Las
identidades \eqref{x_reciprocidad:iii:vp:polarization}--\eqref{x_reciprocidad:iii:vp:dimensional-balance}
establecen la primera construcción con independencia de esa elección de
realización. El operador del vacío se construye en la
sección~\ref{x_reciprocidad:iii:sec:constitutiva} sobre sus canales y proyectores propios.

\section{Emparejamiento de las fibras aditivas y multiplicativas}

La segunda observación procede de la realización tridimensional de APP,
$\mathcal A_3=D_9^3$, con medida uniforme y
$D_9=\{1,\ldots,9\}$. Se conservan los observables
\[
 \Sigma(x)=\rho_9(x_1+x_2+x_3),\qquad
 \Pi(x)=\rho_9(x_1x_2x_3),
\]
donde $\rho_9(n)=1+[(n-1)\bmod9]$ para $n>0$. En
$\ell^2(\mathcal A_3)$, sean $P_{\Sigma,3}$ y $P_{\Pi,3}$ las
esperanzas condicionales sobre sus fibras. Por ejemplo,
\[
 (P_{\Sigma,3}f)(x)=
 \frac1{|\Sigma^{-1}(\Sigma(x))|}
 \sum_{y:\,\Sigma(y)=\Sigma(x)}f(y).
\]
Son proyecciones ortogonales definidas sobre el mismo soporte. La matriz
de incidencias $N_{ab}=|\Sigma^{-1}(a)\cap\Pi^{-1}(b)|$ resulta de
enumerar los $729$ triples, sumando y multiplicando sus coordenadas antes
de aplicar $\rho_9$:
\begin{equation}
 N=9\begin{pmatrix}
 0&1&1&0&1&1&0&1&4\\
 1&0&1&1&0&1&1&0&4\\
 1&0&2&0&1&2&0&0&3\\
 0&1&1&0&1&1&0&1&4\\
 1&0&1&1&0&1&1&0&4\\
 0&0&2&1&0&2&0&1&3\\
 0&1&1&0&1&1&0&1&4\\
 1&0&1&1&0&1&1&0&4\\
 0&1&2&0&0&2&1&0&3
 \end{pmatrix}.
 \label{x_reciprocidad:iii:vp:joint-matrix}
\end{equation}
Las filas suman $81$ y las sumas de columna son
\[
 (m_1,\ldots,m_9)=(36,36,108,36,36,108,36,36,297).
\]
Por ello, el emparejamiento de las indicatrices normalizadas de las fibras
es $C_{ab}=N_{ab}/\sqrt{81m_b}$. La normalización recoge las
multiplicidades producidas por APP, no una elección de valores singulares.

\begin{proposition}[Modo singular del selector trítico]
En la carta canónica, el vector del selector es
$m_{\rm ph}=(1,-1,0)^3$. Se cumple
\begin{equation}
 Cm_{\rm ph}=C^{\mathsf T}m_{\rm ph}=-\tfrac12m_{\rm ph}.
 \label{x_reciprocidad:iii:vp:trit-mode}
\end{equation}
Así, el valor singular correspondiente es $s=1/2$.
\end{proposition}
\begin{proof}
En las seis columnas donde $m_{\rm ph}$ es no nulo, $m_b=36$ y
$C_{ab}=N_{ab}/54$. Multiplicar las filas de
\eqref{x_reciprocidad:iii:vp:joint-matrix} por $(1,-1,0)^3$ produce
$-m_{\rm ph}/2$. Para el producto transpuesto, las sumas firmadas
de las columnas $3,6,9$ son cero; las otras seis dan los mismos
valores $-m_{{\rm ph},b}/2$. Esto prueba ambas identidades y, por
composición, $CC^{\mathsf T}m_{\rm ph}=m_{\rm ph}/4$.
\end{proof}

\section{Norma del conmutador y prefactor electrónico}

\begin{theorem}[Norma exacta de incompatibilidad aritmética]
Las dos proyecciones satisfacen
\begin{equation}
 \bigl\|[P_{\Sigma,3},P_{\Pi,3}]\bigr\|=\frac{\sqrt3}{4}.
 \label{x_reciprocidad:iii:vp:prefactor}
\end{equation}
El máximo se alcanza en el plano principal correspondiente al modo
trítico de \eqref{x_reciprocidad:iii:vp:trit-mode}.
\end{theorem}
\begin{proof}
La matriz racional $M=CC^{\mathsf T}$ tiene entradas
$M_{ac}=\sum_bN_{ab}N_{cb}/(81m_b)$. Su espectro completo se
comprueba mediante los siguientes subespacios independientes: los
vectores $(1,\ldots,1)$, $(1,-1,0)^3$ y $(1,1,-2)^3$ tienen
autovalores $1$, $1/4$ y $7/132$, respectivamente; el subespacio de
vectores soportados en $\{3,6,9\}$ cuya suma es cero tiene dimensión
$2$ y autovalor $1/18$; los vectores de suma cero soportados en
$\{1,4,7\}$ o en $\{2,5,8\}$ forman un núcleo de dimensión $4$.
La sustitución en la matriz verifica estas acciones y las dimensiones
suman $9$. En consecuencia,
\[
 \operatorname{spec}M=
 \{1,1/4,7/132,1/18,1/18,0,0,0,0\}.
\]
Para dos proyecciones ortogonales, un valor singular $s\in(0,1)$ del
emparejamiento determina un plano principal en el que sus matrices son
\[
 P=\begin{pmatrix}1&0\\0&0\end{pmatrix},\qquad
 Q=\begin{pmatrix}s^2&s\sqrt{1-s^2}\\
 s\sqrt{1-s^2}&1-s^2\end{pmatrix}.
\]
Esta forma se obtiene tomando un vector unitario $u$ de la primera
imagen con $PQu=s^2u$ y completándolo con
$(Qu-s^2u)/(s\sqrt{1-s^2})$. Su conmutador es
$s\sqrt{1-s^2}\bigl(\begin{smallmatrix}0&1\\-1&0\end{smallmatrix}\bigr)$.
Los sectores de valor singular $0$ o $1$ contribuyen con cero.
Por tanto, la norma total es el máximo de
$\sqrt{\lambda(1-\lambda)}$ sobre el espectro de $M$. Dicho máximo
es $\sqrt{3/16}$, alcanzado en $\lambda=1/4$.
\end{proof}

La norma es el prefactor que interviene en la composición electrónica.
Su obtención conserva el soporte tridimensional, las fibras y el modo
del selector; la evaluación escalar no sustituye esa estructura. En
particular, conocer $\sqrt3/4$ permite conocer la magnitud máxima de
la incompatibilidad, mientras los proyectores conservan las direcciones
en que actúa.

\clearpage
\section{Defecto cuadrático y respuesta de los sectores 90/120}

El modo trítico determina simultáneamente
\begin{equation}
 1-s^2=\frac34=\frac{90}{120},\qquad s=\frac12.
 \label{x_reciprocidad:iii:vp:defect-ratio}
\end{equation}
La primera expresión es el defecto cuadrático del plano principal de APP.
La última conserva las longitudes de los dos sectores del calendario.
La respuesta constitutiva utiliza esos sectores sobre otro objeto: un
canal contractivo $q\in(0,1)$. Su evaluación es
\[
 r(q)=\frac{1-q^{90}}{1-q^{120}}.
\]
La relación entre el cociente de sectores y esta respuesta puede
precisarse sin identificar sus operadores.

\begin{proposition}[Límite del cociente de respuesta]
Para $0<q<1$,
\begin{equation}
 \frac34<r(q)<1,\qquad
 \lim_{q\uparrow1}r(q)=\frac34=1-s^2.
 \label{x_reciprocidad:iii:vp:response-limit}
\end{equation}
\end{proposition}
\begin{proof}
Con $u=q^{30}\in(0,1)$, la factorización de $1-u^3$ y
$1-u^4$ da
$r(q)=(1+u+u^2)/(1+u+u^2+u^3)$. El límite en $u=1$ es
$3/4$. Además,
$4(1+u+u^2)-3(1+u+u^2+u^3)
=(1-u)(3u^2+2u+1)>0$, y el denominador supera al numerador
en $u^3>0$. Se obtienen las dos desigualdades.
\end{proof}

La coincidencia exacta concierne al defecto del modo singular y al límite
de la respuesta de los sectores. En el interior contractivo, la respuesta
depende de $q$ y es estrictamente mayor que $3/4$. El desarrollo
constitutivo conserva sus dos canales $q_\pm$ y sus proyectores
$P_\pm$; los proyectores $P_{\Sigma,3}$ y $P_{\Pi,3}$ conservan,
por su parte, las fibras aritméticas que determinaron el prefactor.
Una equivalencia operatoria entre ambas realizaciones exige una aplicación
entre sus espacios que entrelace las acciones correspondientes.
Las ecuaciones \eqref{x_reciprocidad:iii:vp:trit-mode}, \eqref{x_reciprocidad:iii:vp:prefactor} y
\eqref{x_reciprocidad:iii:vp:response-limit} precisan la información compartida y
mantienen visibles los datos de cada realización.

% END OWNER


% BEGIN OWNER /Users/ruben/Documents/New project/output/EXTREMA_MEDIA_RAZON_RECIPROCIDAD_ES_EN_20260918/ARTICULO/ES/source/antecedentes_iii/sections/04_respuesta_constitutiva.tex
\chapter{Operador constitutivo y reconstrucción de los canales del vacío}
\label{x_reciprocidad:iii:sec:constitutiva}

Las magnitudes constitutivas se obtienen mediante una operación sobre el
transporte angular. La información de partida comprende los canales, sus
proyectores y la orientación que los distingue. Esta organización permite
seguir por separado la construcción del operador, la evaluación de sus
funciones espectrales y su realización dimensional posterior.

\section{Respuesta positiva sobre dos hojas}

En la cámara orientada $x>y>0$ de la sección~\ref{x_reciprocidad:iii:sec:hojas}, sean
$P_\pm=(I\pm\mathcal R)/2$ y
$T=q_+P_++q_-P_-$, con $0<q_+<q_-<1$. Las incidencias construidas por el
calendario fijan los sectores $90$ y $120$. La respuesta constitutiva es
\begin{equation}
 \mathcal V=(I-T^{90})(I-T^{120})^{-1}.
 \label{x_reciprocidad:iii:eq:vacuum-op}
\end{equation}
Esta definición recibe el transporte procedente del TPK y conserva las
asignaciones de hoja. La regla de respuesta y su dominio son parte explícita
de la construcción.

El cociente procede de comparar las dos incidencias del calendario sobre
un mismo transporte, no de ajustar por separado cuatro magnitudes. Si
$S=T^{30}$ y $\Sigma_j(S)=I+S+\cdots+S^{j-1}$, la factorización de
$I-S^j$ da
\begin{equation}
 \mathcal V\,\Sigma_4(S)=\Sigma_3(S),\qquad
 \mathcal V=\Sigma_3(S)\Sigma_4(S)^{-1}.
 \label{x_reciprocidad:iii:eq:completion-response}
\end{equation}
Los factores $\Sigma_3$ y $\Sigma_4$ expresan respectivamente los
sectores $90=3\cdot30$ y $120=4\cdot30$. Como el espectro de $S$ está
contenido en $(0,1)$, $\Sigma_4(S)$ es invertible: la ecuación de
completación determina una única respuesta para el transporte fijado.
Esta unicidad concierne a la regla constitutiva declarada y a sus
incidencias. La identificación electromagnética añade la asignación
orientada de las hojas y las rectas dimensionales que se precisan más
abajo.

\begin{proposition}[Descomposición y positividad]
La respuesta está bien definida y
\begin{equation}
 \mathcal V=r_+P_++r_-P_-,\qquad
 r_\pm=\frac{1-q_\pm^{90}}{1-q_\pm^{120}}.
 \label{x_reciprocidad:iii:eq:rchannels}
\end{equation}
Sus autovalores satisfacen $3/4<r_-<r_+<1$ en esta cámara.
\end{proposition}
\begin{proof}
Cada coeficiente $1-q_\pm^{120}$ es positivo. La inversa de $I-T^{120}$
es la suma de los proyectores divididos por esos coeficientes, lo que prueba
\eqref{x_reciprocidad:iii:eq:rchannels}. Con $s=q^{30}$ se obtiene
\begin{equation}
 f(s)=\frac{1+s+s^2}{1+s+s^2+s^3},\qquad
 f'(s)=-\frac{s^2(3+2s+s^2)}{(1+s+s^2+s^3)^2}<0.
 \label{x_reciprocidad:iii:eq:monotone}
\end{equation}
Los límites en $0$ y $1$ son $1$ y $3/4$, respectivamente. La monotonía
y el orden de los canales concluyen la prueba.
\end{proof}

\section{Cuatro coordenadas de una respuesta común}

Se definen las evaluaciones normalizadas
\begin{equation}
 \widehat\varepsilon=r_+^2,\qquad
 \widehat\mu=r_-^2,\qquad
 \widehat Z=\frac{r_-}{r_+},\qquad
 \widehat c=\frac1{r_+r_-}.
 \label{x_reciprocidad:iii:eq:four}
\end{equation}
La primera pareja conserva las respuestas cuadráticas asignadas a cada hoja.
El producto recoge el modo común; el cociente distingue su orientación relativa.
Se verifican exactamente
\begin{equation}
 \widehat\mu\widehat\varepsilon=\widehat c^{-2},\qquad
 \frac{\widehat\mu}{\widehat\varepsilon}=\widehat Z^2,
 \qquad
 r_+=\frac1{\sqrt{\widehat Z\widehat c}},\quad
 r_-=\sqrt{\frac{\widehat Z}{\widehat c}}.
 \label{x_reciprocidad:iii:eq:four-identities}
\end{equation}
Las raíces positivas fijan la inversión. Bajo la conjugación de hojas,
leída respecto de las marcas $P_\pm$ fijas, se
intercambian $\widehat\varepsilon$ y $\widehat\mu$, se invierte
$\widehat Z$ y se conserva $\widehat c$. La respuesta se transporta entonces
a la cámara $x>-y>0$: las fórmulas siguen siendo válidas, con los órdenes
$q_-<q_+$ y $r_+<r_-$ invertidos. Por consiguiente, el modo común
y la orientación son datos recuperables mediante observables complementarios.

\begin{proposition}[Determinación de las evaluaciones cuadráticas]
Sea $W=r_+P_++r_-P_-$ una respuesta positiva sobre las dos hojas.
Fijadas las lecturas de producto y cociente
\[
 \widehat c^{-1}=\det W=r_+r_-,
 \qquad \widehat Z=r_-/r_+,
\]
existe una única pareja positiva $(\widehat\varepsilon,\widehat\mu)$
que satisface
$\widehat\mu\widehat\varepsilon=\widehat c^{-2}$ y
$\widehat\mu/\widehat\varepsilon=\widehat Z^2$.
Es precisamente $(r_+^2,r_-^2)$.
\end{proposition}
\begin{proof}
La positividad permite tomar las raíces sin ambigüedad:
\[
 \widehat\mu=\frac{\widehat Z}{\widehat c}
 =\frac{r_-}{r_+}r_+r_-=r_-^2,\qquad
 \widehat\varepsilon=\frac1{\widehat Z\widehat c}
 =\frac{r_+}{r_-}r_+r_-=r_+^2.
\]
Las expresiones verifican las dos relaciones y son las únicas raíces
positivas de ellas.
\end{proof}

Así, los cuadrados no constituyen dos elecciones independientes añadidas
a los canales: quedan determinados por las lecturas multiplicativa y
orientada, una vez fijadas las relaciones constitutivas. Esta proposición
establece la unicidad dentro de ese sistema de reglas. La interpretación
de las reglas como respuesta electromagnética conserva su contenido
adicional de representación física.

La lectura determinantal del transporte y la respuesta constitutiva
también deben conservar su orden. Para $T=e^{-xI-y\mathcal R}$ se tiene
$\det T=e^{-2x}$, mientras que
\[
 \det\mathcal V=f(q_+^{30})f(q_-^{30})=\widehat c^{-1}.
\]
Ambas proceden del mismo operador de dos hojas, pero evalúan funciones
diferentes de su espectro. En general, aplicar primero el determinante
y después una función racional no equivale a tomar el determinante de
esa función del operador. Mantener los dos autovalores hasta completar
la respuesta conserva la información de orientación que un determinante
aislado ya no distingue.

\section{Inversión cúbica y recuperación angular}

\begin{theorem}[Inversión constitutiva en el dominio contractivo]
Para cada $r\in(3/4,1)$ existe una única $s\in(0,1)$ tal que $f(s)=r$.
Esta $s$ es la raíz única en dicho intervalo de
\begin{equation}
 r s^3+(r-1)(1+s+s^2)=0.
 \label{x_reciprocidad:iii:eq:cubic}
\end{equation}
La respuesta etiquetada reconstruye $s_\pm$, $q_\pm=s_\pm^{1/30}$ y
\begin{equation}
 x=-\frac12\log(q_+q_-),\qquad
 y=\frac12\log\frac{q_-}{q_+}.
 \label{x_reciprocidad:iii:eq:recover-angular}
\end{equation}
\end{theorem}
\begin{proof}
La derivada y los límites de \eqref{x_reciprocidad:iii:eq:monotone} prueban que $f$ es una
biyección decreciente entre los intervalos indicados. Multiplicar $f(s)=r$
por su denominador produce \eqref{x_reciprocidad:iii:eq:cubic}. La raíz positiva de orden
$30$ recupera el canal. Finalmente, suma y diferencia de
$\log q_\pm=-x\mp y$ dan \eqref{x_reciprocidad:iii:eq:recover-angular}.
\end{proof}

La reconstrucción conserva las etiquetas de hoja. Con el operador completo
$\mathcal V$ se conservan también los proyectores y se aplica la función
inversa por cálculo espectral. Recuperar dos coordenadas angulares es una
operación definida sobre esta representación; la traza histórica completa
del TPK conserva datos adicionales.

El intervalo $r\in(3/4,1)$ y la función inversa de este teorema pertenecen
a la normalización interna de \eqref{x_reciprocidad:iii:eq:four}. Si se dispone de
coordenadas expresadas en otra carta de unidades, se recupera primero
esta normalización mediante la transformación de la
ecuación~\eqref{x_reciprocidad:iii:eq:dimensions}; después se aplica la inversión cúbica.

\section{Coordenadas logarítmicas y refinamiento}

Sean $\mathcal E=\log(r_+r_-)$ y $\mathcal B=\log(r_-/r_+)$. Entonces
\begin{equation}
 \log\mathcal V=\tfrac12\mathcal E I-\tfrac12\mathcal B\mathcal R.
 \label{x_reciprocidad:iii:eq:logvac}
\end{equation}
Esta descomposición exhibe la componente invariante y la componente orientada.
En la amplificación $\mathcal V_m=\mathcal V\otimes I_{9^m}$, las
expectativas parciales normalizadas satisfacen
$E_m(\mathcal V_{m+1})=\mathcal V_m$. Para todo polinomio $p$,
$p(\mathcal V_m)=p(\mathcal V)\otimes I_{9^m}$; de ahí
$E_m(p(\mathcal V_{m+1}))=p(\mathcal V_m)$. La continuidad uniforme
extiende la identidad a las funciones continuas del espectro. Se conservan
así los proyectores, las potencias y el logaritmo en la torre indicada.

\section{Tipos dimensionales de las cuatro evaluaciones}

Con las rectas de acción, carga y velocidad, provistas de secciones
$u_S,u_Q,u_C$, la realización dimensional se escribe
\begin{align}
 Z_{\rm vac}&=\widehat Z u_Su_Q^{-2},&
 c_{\rm vac}&=\widehat c u_C,\\
 \mu_{\rm vac}&=\widehat\mu u_Su_C^{-1}u_Q^{-2},&
 \varepsilon_{\rm vac}&=\widehat\varepsilon u_S^{-1}u_C^{-1}u_Q^2.
 \label{x_reciprocidad:iii:eq:dimensions}
\end{align}
Multiplicando y dividiendo estas expresiones se preservan
$\mu_{\rm vac}\varepsilon_{\rm vac}=c_{\rm vac}^{-2}$ y
$\mu_{\rm vac}/\varepsilon_{\rm vac}=Z_{\rm vac}^2$.
La carta SI evalúa después las secciones. La comparación física requiere
exponer su construcción y normalización junto a la evaluación; las cuatro
coordenadas normalizadas de \eqref{x_reciprocidad:iii:eq:four} conservan su tipo adimensional.

% END OWNER


% BEGIN OWNER /Users/ruben/Documents/New project/output/EXTREMA_MEDIA_RAZON_RECIPROCIDAD_ES_EN_20260918/ARTICULO/ES/source/antecedentes_iii/sections/04b_reconstruccion_forma_lc.tex
% INSERCIÓN FOCAL AUTORIZADA: no introduce un capítulo ni el concepto de radión.
% Incluir dentro de la sección constitutiva, después de la inversión cúbica
% y de la declaración de las cartas dimensionales. Etiquetas fuera de grupos
% de prefijado automático. Conserva íntegramente el bloque geométrico anterior.
\section{Reconstrucción de la forma de acción desde la respuesta constitutiva}
\label{x_reciprocidad:iii:subsec:puente-forma-LC}

La respuesta constitutiva y la realización inductiva--capacitiva reciben
el mismo par angular construido desde APP--TRIT--TPK. La primera evalúa
funciones racionales de los dos canales; la segunda utiliza la razón de
sus exponentes para determinar la forma de la elipse de acción. La
inversión cúbica permite componer ambas lecturas sin identificar sus
cocientes numéricos.

Sean $x=A_{\rm rad}$ e $y=C^*_{\rm rad}$, con $x>|y|$, las coordenadas
levantadas ya construidas, y consérvense sus etiquetas de hoja. Escribimos
\begin{equation}
 q_\pm=e^{-(x\pm y)},\qquad s_\pm=q_\pm^{30},\qquad
 r_\pm=f(s_\pm),\qquad
 f(s)=\frac{1+s+s^2}{1+s+s^2+s^3}.
 \label{x_reciprocidad:iii:eq:forma-canales}
\end{equation}
El dominio contractivo da $s_\pm\in(0,1)$ y
$r_\pm\in(3/4,1)$. En la cámara $x>y>0$ se cumplen
$s_+<s_-$ y $r_+>r_-$; el intercambio de hojas transporta esta cámara
a la orientación opuesta.

La imagen de las dos coordenadas constitutivas normalizadas es el dominio
\begin{equation}
 \mathcal U=\left\{(z,v)\in(0,\infty)^2:
   (zv)^{-1/2}\in(3/4,1),\quad
   (z/v)^{1/2}\in(3/4,1)\right\}.
 \label{x_reciprocidad:iii:eq:forma-dominio}
\end{equation}
Aquí $(z,v)=(\widehat Z,\widehat c)$. Las restricciones expresan la
pertenencia de los canales reconstruidos al dominio del teorema de
inversión; no constituyen una elección ulterior de sus valores.

\Needspace{23\baselineskip}
\begin{proposition}[Composición constitutiva de la forma y de la impedancia LC]
\label{x_reciprocidad:iii:prop:forma-LC}
Para $(\widehat Z,\widehat c)\in\mathcal U$, sean $s_\pm\in(0,1)$ las
raíces únicas de
\begin{equation}
 r_\pm s_\pm^3+(r_\pm-1)(1+s_\pm+s_\pm^2)=0,
 \qquad
 r_+=\frac1{\sqrt{\widehat Z\widehat c}},\quad
 r_-=\sqrt{\frac{\widehat Z}{\widehat c}}.
 \label{x_reciprocidad:iii:eq:forma-inversion}
\end{equation}
La razón angular, la anisotropía de la elipse de acción y la impedancia
LC relativa se reconstruyen exactamente mediante
\begin{align}
 \frac{y}{x}
 &=\frac{\log s_+-\log s_-}{\log s_++\log s_-},
 \label{x_reciprocidad:iii:eq:forma-razon}\\
 \eta_{\rm el}
 &=\frac12\log\frac{\log s_+}{\log s_-},
 \label{x_reciprocidad:iii:eq:forma-eta}\\
 \frac{Z_{\rm LC}}{Z_*}
 &=e^{-\eta_{\rm el}}
  =\sqrt{\frac{\log s_-}{\log s_+}}.
 \label{x_reciprocidad:iii:eq:forma-LC}
\end{align}
En la última identidad $Z_{\rm LC}=Z_{\eta_{\rm el}}$ es la impedancia
de la carta LC de \eqref{x_reciprocidad:hmtII:eq:ii-elipse-lc}.
\end{proposition}
\begin{proof}
La inversión cúbica de \eqref{x_reciprocidad:iii:eq:cubic} determina $s_\pm$ de manera
unívoca, porque $f$ es una biyección decreciente de $(0,1)$ sobre
$(3/4,1)$. Sus raíces positivas de orden $30$ recuperan $q_\pm$.
Las identidades
\[
 \log s_+=-30(x+y),\qquad \log s_-=-30(x-y)
\]
prueban \eqref{x_reciprocidad:iii:eq:forma-razon} por diferencia y suma. Ambos logaritmos
son negativos; su cociente es positivo y su suma no es nula. Por ello,
\[
 \frac12\log\frac{\log s_+}{\log s_-}
   =\frac12\log\frac{x+y}{x-y}
   =\operatorname{artanh}(y/x)=\eta_{\rm el}.
\]
La relación $Z_{\eta_{\rm el}}/Z_*=e^{-\eta_{\rm el}}$, demostrada en
\eqref{x_reciprocidad:hmtII:eq:ii-elipse-impedancia}, da \eqref{x_reciprocidad:iii:eq:forma-LC}.
\end{proof}

La composición conserva la información necesaria para cada reconstrucción.
Con $\hbar$ se recuperan
$a_S=\sqrt{2\hbar}\,e^{\eta_{\rm el}/2}$ y
$b_S=\sqrt{2\hbar}\,e^{-\eta_{\rm el}/2}$; con la carta
$(\omega,Z_*)$ se recuperan la inductancia y la capacitancia ya definidas.
La acción, el reloj y la referencia dimensional acompañan la forma como
datos tipados. La inversión de dos coordenadas adimensionales no
reconstruye sus unidades ni toda la memoria del estado.

La normalización interviene antes de resolver las ecuaciones cúbicas.
Si se reciben $Z_{\rm vac}$ y $c_{\rm vac}$ en una carta dimensional,
las relaciones \eqref{x_reciprocidad:iii:eq:dimensions} determinan primero
\[
 \widehat Z=\frac{Z_{\rm vac}}{u_Su_Q^{-2}},\qquad
 \widehat c=\frac{c_{\rm vac}}{u_C}.
\]
Esas son las coordenadas a las que se aplica
\eqref{x_reciprocidad:iii:eq:forma-inversion}. La referencia $Z_*$ pertenece a la
realización LC y se conserva expresamente en \eqref{x_reciprocidad:iii:eq:forma-LC}.

Los dos cocientes están relacionados mediante los canales, pero evalúan
funciones diferentes:
\begin{equation}
 \widehat Z=\frac{f(s_-)}{f(s_+)},\qquad
 \frac{Z_{\rm LC}}{Z_*}
       =\sqrt{\frac{\log s_-}{\log s_+}}.
 \label{x_reciprocidad:iii:eq:forma-dos-cocientes}
\end{equation}
La reconstrucción anterior proporciona el mapa entre ambas lecturas.
En particular, intercambiar las hojas invierte ambos cocientes, cambia
$\eta_{\rm el}$ por $-\eta_{\rm el}$ y conserva $\widehat c$.
En el límite simétrico $y=0$, los canales coinciden y ambos cocientes
valen $1$. Fuera de ese caso, la igualdad de origen no impone igualdad
de las dos funciones. La forma de acción es recuperable desde la
respuesta constitutiva conservando orientación y normalización.

% END OWNER


% BEGIN OWNER /Users/ruben/Documents/New project/output/EXTREMA_MEDIA_RAZON_RECIPROCIDAD_ES_EN_20260918/ARTICULO/ES/source/antecedentes_iii/sections/04c_velocidad_y_reloj.tex
% RESULTADO_RECUPERADO: funtor dimensional integral, eq:cZ-internos,
% eq:longitud-ruta; composición reunida en Artículo V REV02, sección 69.
% Incorporación autónoma a III: mismo transporte, mismas hojas y misma carta.
\section{Velocidad constitutiva y realización cinemática del reloj}
\label{x_reciprocidad:iii:sec:enlace-velocidad-constitutiva}

La respuesta constitutiva y el lector cinemático realizan una misma
sección de velocidad. Su identificación conserva el orden de construcción:
las hojas de APP, su orientación TRIT y el transporte angular publicado
por el TPK anteceden a las lecturas de velocidad y a la elección de unidades.
En la representación de dos hojas de la
sección~\ref{x_reciprocidad:iii:sec:hojas}, con $x>|y|$ y proyectores de rango uno
$P_\pm$, el transporte es
\[
 T=e^{-xI-y\mathcal R}=q_+P_++q_-P_-,
 \qquad q_\pm=e^{-x\mp y},\qquad 0<q_\pm<1.
\]
Sus bloques temporales de longitudes $90$ y $120$ producen la respuesta
$\mathcal V$ de \eqref{x_reciprocidad:iii:eq:vacuum-op}. Las coordenadas $x,y$ y los
canales son salidas del desarrollo angular previo; esta composición
aplica sus lectores a ese transporte ya construido.

El lector logarítmico simétrico del transporte temporal es
\begin{equation}
 \mathcal E=
 \log\!\bigl[(1-q_+^{90})(1-q_-^{90})\bigr]
 -\log\!\bigl[(1-q_+^{120})(1-q_-^{120})\bigr].
 \label{x_reciprocidad:iii:eq:velocidad-lector-simetrico}
\end{equation}
En la recta dimensional de velocidad, con sección positiva $u_C$, el
lector cinemático publica $c_{\mathrm{int}}=\exp(-\mathcal E)u_C$.
La respuesta constitutiva publica
$c_{\mathrm{vac}}=\widehat c u_C$ mediante \eqref{x_reciprocidad:iii:eq:four}.

\begin{proposition}[Identidad constitutiva y cinemática]
\label{x_reciprocidad:iii:prop:identidad-velocidad}
En la misma carta dimensional se verifica
\begin{equation}
 c_{\mathrm{int}}=c_{\mathrm{vac}}=\widehat c\,u_C,
 \qquad \widehat c=(r_+r_-)^{-1}=\exp(-\mathcal E).
 \label{x_reciprocidad:iii:eq:identidad-velocidad}
\end{equation}
La identificación conserva el intercambio de hojas.
\end{proposition}
\begin{proof}
Todos los factores de \eqref{x_reciprocidad:iii:eq:velocidad-lector-simetrico} son
positivos. La ley del logaritmo y \eqref{x_reciprocidad:iii:eq:rchannels} dan
\[
 \mathcal E=\log(r_+r_-)=\log\det\mathcal V,
 \qquad e^{-\mathcal E}=(\det\mathcal V)^{-1}.
\]
El determinante corresponde a la realización de dos hojas de rango uno.
Las dos definiciones publican, por tanto, la misma sección. El intercambio
de hojas permuta $r_+$ y $r_-$ y conserva su producto. El determinante
$\det T=e^{-2x}$ corresponde al transporte elemental, mientras que
$\det\mathcal V=r_+r_-$ corresponde a su respuesta temporal;
\eqref{x_reciprocidad:iii:eq:identidad-velocidad} utiliza el segundo.
\end{proof}

\subsection{Longitud de trayectoria y duración elemental}

El marco dimensional positivo $(u_S,u_C,u_T,u_Q)$ induce la base de
longitud $u_L=u_Cu_T$. Para una trayectoria enriquecida $\gamma$ de
$n=|\gamma|$ transiciones y canal constante, los lectores aditivos son
\begin{equation}
 T_{\mathrm{int}}(\gamma)=nu_T,
 \qquad L_{\mathrm{int}}(\gamma)=n\widehat c\,u_L.
 \label{x_reciprocidad:iii:eq:velocidad-longitud-reloj}
\end{equation}
La concatenación suma el número de transiciones y sus realizaciones.
Las rutas enriquecidas conservan, además, su orientación y su memoria.

\begin{corollary}[Normalización de la longitud elemental]
Para $n>0$, $L_{\mathrm{int}}(\gamma)/T_{\mathrm{int}}(\gamma)
=c_{\mathrm{vac}}$. En la normalización $t_0=u_T$,
\begin{equation}
 \ell_0=c_{\mathrm{int}}t_0=\widehat c\,u_L.
 \label{x_reciprocidad:iii:eq:longitud-elemental-constitutiva}
\end{equation}
En una carta temporal $t_0=\tau_0u_T$, con $\tau_0>0$, se obtiene
$\ell_0=\tau_0\widehat c\,u_L$.
\end{corollary}
\begin{proof}
Dividir las dos expresiones de \eqref{x_reciprocidad:iii:eq:velocidad-longitud-reloj}
cancela $n$ y utiliza $u_L/u_T=u_C$. Multiplicar la sección de velocidad
por $t_0$ prueba ambas fórmulas elementales.
\end{proof}

Cuando el canal depende de la arista, el lector conserva la forma aditiva
$L_{\mathrm{int}}(\gamma)=\sum_{e\in\gamma}\widehat c(e)u_L$.
Su cociente por $nu_T$ es la media de las velocidades de esas aristas.
La expresión de canal constante es la especialización homogénea de
esta lectura.

\subsection{Cambio de carta y normalización adaptada}

\begin{proposition}[Covariancia de la identificación]
Sean $u_C'=d u_C$ y $u_T'=\eta u_T$, con $d,\eta>0$.
Las coordenadas de las mismas secciones satisfacen
\begin{equation}
 \widehat c'=d^{-1}\widehat c,
 \qquad \tau_0'=\eta^{-1}\tau_0,
 \qquad u_L'=d\eta u_L.
 \label{x_reciprocidad:iii:eq:velocidad-cambio-carta}
\end{equation}
Si otra publicación escribe $c'=\widehat c' u_C'$ con $u_C'=d u_C$,
la igualdad de secciones equivale exactamente a $d\widehat c'=\widehat c$.
\end{proposition}
\begin{proof}
La igualdad de cada sección en las dos bases da
$\widehat c'u_C'=\widehat c u_C$ y $\tau_0'u_T'=\tau_0u_T$.
En particular,
\[
 \tau_0'\widehat c'u_L'
 =\frac{\tau_0}{\eta}\frac{\widehat c}{d}\,d\eta u_L
 =\tau_0\widehat c u_L=\ell_0.
\]
La última equivalencia procede de la misma sustitución para la velocidad.
\end{proof}

La elección adaptada $d=\widehat c$ da $u_C'=c_{\mathrm{vac}}$ y
$\widehat c'=1$: expresa la sección construida en una base adaptada.
El coeficiente de la carta interna sigue siendo $(r_+r_-)^{-1}$, y esa
es la normalización utilizada por la inversión cúbica. Manteniendo las
bases de acción y carga, \eqref{x_reciprocidad:iii:eq:dimensions} produce
\[
 \widehat\mu'=d\widehat\mu,\qquad
 \widehat\varepsilon'=d\widehat\varepsilon,\qquad
 \widehat\mu'\widehat\varepsilon'=(\widehat c')^{-2}.
\]
En la carta adaptada se tiene $\widehat\mu'=r_-/r_+$ y
$\widehat\varepsilon'=r_+/r_-$, por sustitución de
$\widehat\mu=r_-^2$ y $\widehat\varepsilon=r_+^2$.

\subsection{Conservación del lector bajo amplificación}

Para el refinamiento $\mathcal V_m=\mathcal V\otimes I_{9^m}$, la
traza normalizada $\tau_m=\operatorname{Tr}/(2\cdot9^m)$ satisface
\begin{equation}
 \exp[-2\tau_m(\log\mathcal V_m)]
 =\exp[-\log r_+-\log r_-]=\widehat c.
 \label{x_reciprocidad:iii:eq:velocidad-traza-normalizada}
\end{equation}
En efecto, cada autovalor aparece $9^m$ veces; esa multiplicidad se
cancela con el denominador de la traza normalizada. El determinante
ordinario ampliado es $(r_+r_-)^{9^m}$. La normalización de dos hojas,
equivalentemente expresada por \eqref{x_reciprocidad:iii:eq:velocidad-traza-normalizada},
conserva el lector constitutivo durante la amplificación. Esta
realización conserva el transporte enriquecido que la origina.

% END OWNER


% BEGIN OWNER /Users/ruben/Documents/New project/output/EXTREMA_MEDIA_RAZON_RECIPROCIDAD_ES_EN_20260918/ARTICULO/ES/source/antecedentes_iii/sections/07_carga_y_cuantos.tex
\chapter{Sección de carga y relaciones de cuantización eléctrica}
\label{x_reciprocidad:iii:sec:carga}

La respuesta constitutiva establece una relación entre las dos hojas y una
escala relativa de propagación. La sección de carga se obtiene al componer
esta respuesta con la constante de estructura fina y la acción, ambas
construidas en los capítulos antecedentes. El orden tiene una consecuencia
precisa: la carga no determina retrospectivamente el registro dodecafásico.
Es el registro, junto con la acción y la respuesta del vacío, el que
determina la sección de carga en la recta correspondiente.

\section{Determinación positiva de la sección de carga}

Consideremos las secciones positivas $\hbar_a$, con
$a\in\{\mathrm{pre},\mathrm{ret}\}$, y $h_a=2\pi\hbar_a$.
La impedancia $Z_{\rm vac}$ ya está construida mediante
\eqref{x_reciprocidad:iii:eq:dimensions}. El cociente $h_a/Z_{\rm vac}$ pertenece a
la recta cuadrática de carga. Esta compatibilidad dimensional permite
formular la siguiente definición sin elegir todavía una carta SI.

\begin{definition}[Sección positiva de carga]
La sección de carga correspondiente a la orientación de acción $a$ es
\begin{equation}
 e_a=\left(\frac{2\alpha h_a}{Z_{\rm vac}}\right)^{1/2}.
 \label{x_reciprocidad:iii:eq:charge}
\end{equation}
La raíz se toma en la semirrecta positiva de carga. Su sección conjugada
es $-e_a$; el signo de esta última no se identifica con el índice de
retorno de la acción.
\end{definition}

\begin{proposition}[Existencia, unicidad y compatibilidad del acoplamiento]
La ecuación $Z_{\rm vac}e_a^2=2\alpha h_a$ determina una única sección
positiva. En la realización constitutiva se cumple también
\begin{equation}
 \alpha=\frac{Z_{\rm vac}e_a^2}{2h_a}
 =\frac{e_a^2}{4\pi\varepsilon_{\rm vac}\hbar_a c_{\rm vac}}.
 \label{x_reciprocidad:iii:eq:coupling}
\end{equation}
\end{proposition}
\begin{proof}
La positividad de $\alpha$, $h_a$ y $Z_{\rm vac}$ garantiza una raíz
positiva única. Las identidades constitutivas dan
$\varepsilon_{\rm vac}c_{\rm vac}=Z_{\rm vac}^{-1}$.
Sustituir esta igualdad y $h_a=2\pi\hbar_a$ en la primera expresión
de \eqref{x_reciprocidad:iii:eq:coupling} produce la segunda.
\end{proof}

La última expresión es la forma habitual del acoplamiento electromagnético
en una carta física. Aquí aparece como identidad posterior de las secciones
construidas. La prueba no vuelve a calcular $\alpha$ a partir de una carga
medida: establece la coherencia algebraica del resultado anterior con la
sección que acaba de definirse. La identificación física exige, además,
mantener la carta de unidades y el régimen del acoplamiento empleados en la
comparación. Una identidad entre magnitudes y una comparación experimental
tienen funciones diferentes en ese reconocimiento.

\section{Resistencia, conductancia, flujo y frecuencia}

Las siguientes cuatro magnitudes utilizan la misma sección $e_a$. Su
dependencia puede resolverse por completo antes de evaluar ningún decimal.

\begin{theorem}[Composición de los cuantos eléctricos]
Las secciones
\begin{align}
 R_K&=\frac{h_a}{e_a^2}=\frac{Z_{\rm vac}}{2\alpha},
 &G_0&=\frac{2e_a^2}{h_a}=\frac{4\alpha}{Z_{\rm vac}},
 \label{x_reciprocidad:iii:eq:rk-g0}\\
 \Phi_{0,a}&=\frac{h_a}{2e_a}
 =\sqrt{\frac{h_aZ_{\rm vac}}{8\alpha}},
 &K_{J,a}&=\frac{2e_a}{h_a}=\Phi_{0,a}^{-1}
 \label{x_reciprocidad:iii:eq:flux-josephson}
\end{align}
satisfacen exactamente
\begin{equation}
 R_KG_0=2,\qquad G_0Z_{\rm vac}=4\alpha,\qquad
 K_{J,a}\Phi_{0,a}=1,\qquad
 K_{J,a}^{2}R_K=\frac4{h_a}.
 \label{x_reciprocidad:iii:eq:quanta-identities}
\end{equation}
\end{theorem}
\begin{proof}
La sustitución de $e_a^2=2\alpha h_a/Z_{\rm vac}$ en las dos primeras
definiciones cancela $h_a$ y da \eqref{x_reciprocidad:iii:eq:rk-g0}. Al elevar al
cuadrado la sección positiva $h_a/(2e_a)$ se obtiene
$h_aZ_{\rm vac}/(8\alpha)$; la positividad fija la raíz en
\eqref{x_reciprocidad:iii:eq:flux-josephson}. Su inversa es $2e_a/h_a$.
Los tres primeros productos de \eqref{x_reciprocidad:iii:eq:quanta-identities} se
reducen, respectivamente, a $2$, $4\alpha$ y $1$. Finalmente,
$(4e_a^2/h_a^2)(h_a/e_a^2)=4/h_a$.
\end{proof}

La realización física reconoce en estas secciones la resistencia de von
Klitzing, el cuanto de conductancia, el cuanto de flujo y la constante de
Josephson. El factor $2$ de la conductancia pertenece a la normalización
de la magnitud aquí definida; su identificación con una multiplicidad de
canales de un sistema material exige especificar esa representación.
Las cuatro ecuaciones no constituyen cuatro generadores independientes:
comparten $\alpha$, la acción y la respuesta del vacío. Precisamente esa
dependencia produce las identidades cruzadas y permite comprobar una
evaluación sin alterar las operaciones que la originaron.

\section{Retorno de la acción y transporte de las secciones}

El cociente $R_{\rm act}=h_{\rm pre}/h_{\rm ret}$, obtenido antes del
par angular, conserva la diferencia entre las dos secciones de acción.
No cambia las asignaciones de hoja de la respuesta constitutiva. Son dos
operaciones distintas: retorno de la acción e intercambio de los canales.

\begin{proposition}[Covariancia bajo el retorno]
Con $\alpha$ y $Z_{\rm vac}$ comunes, se tiene
\begin{align}
 \frac{e_{\rm pre}}{e_{\rm ret}}&=R_{\rm act}^{1/2},
 &R_K^{\rm pre}&=R_K^{\rm ret},
 &G_0^{\rm pre}&=G_0^{\rm ret},\label{x_reciprocidad:iii:eq:charge-return}\\
 \frac{\Phi_{0,\rm pre}}{\Phi_{0,\rm ret}}&=R_{\rm act}^{1/2},
 &\frac{K_{J,\rm pre}}{K_{J,\rm ret}}&=R_{\rm act}^{-1/2}.
 \label{x_reciprocidad:iii:eq:flux-return}
\end{align}
\end{proposition}
\begin{proof}
La primera identidad sigue de dividir las dos versiones de
\eqref{x_reciprocidad:iii:eq:charge} y tomar la raíz positiva. En
\eqref{x_reciprocidad:iii:eq:rk-g0} la acción ya se ha cancelado. Las otras dos
identidades siguen de las potencias $h_a^{1/2}$ y $h_a^{-1/2}$ en
\eqref{x_reciprocidad:iii:eq:flux-josephson}.
\end{proof}

Resistencia y conductancia conservan el cociente constitutivo, mientras
que flujo y frecuencia retienen también la orientación de acción. Esta
separación da una lectura estructural de sus dependencias: unas magnitudes
son invariantes del cociente de hojas y otras miden, además, la sección de
acción utilizada. El contraste metrológico debe declarar qué sección evalúa;
la proximidad numérica de las dos acciones no autoriza identificarlas.

% END OWNER


% BEGIN OWNER /Users/ruben/Documents/New project/output/EXTREMA_MEDIA_RAZON_RECIPROCIDAD_ES_EN_20260918/ARTICULO/ES/source/antecedentes_vii/sections/v_funcional_barbero.tex
% RESULTADO_RECUPERADO desde II REV09. Véase MAPA_PROCEDENCIA.json.
% Selección de líneas y cambios mecánicos explícitos; ninguna prueba nueva.
\chapter{Incidencia hexádico--octádica y funcional de Barbero--Immirzi}
\label{x_reciprocidad:v:ant:sec:barbero}
Sea $H\hookrightarrow O$ la inclusión marcada de una hexada en una octada,
$|H|=6$, $|O|=8$, construida mediante el residuo binario de la
sección~\ref{x_reciprocidad:exc:residuo-binario}. Se conservan la hexada, la dodecada
y la carta regional que determinan esa inclusión.
Los objetos de incidencia
\[
\mathcal F_6=H\times\binom H2,\qquad
\mathcal F_8=O\times\binom H2
\]
tienen cardinales $90$ y $120$. Su diferencia normalizada da
\[
\frac{|\mathcal F_6|-|\mathcal F_8|}{24\binom62}=-\frac1{12}.
\]
\begin{figure}[htbp]
\centering
\begin{tikzpicture}[>=Latex,every node/.style={font=\small}]
\draw[rounded corners,draw=ink] (-1.8,-1.35) rectangle (1.8,1.35);
\foreach \a in {0,60,...,300}{
\fill[ink] (\a:0.83) circle(2pt);}
\node at (0,-1.7) {$H:\ 6$ puntos};
\draw[->,thick,ink] (2.1,0)--(3.6,0);
\node[above] at (2.85,1.5) {inclusión marcada};
\draw[rounded corners,draw=ink] (4,-1.35) rectangle (7.6,1.35);
\foreach \a in {0,60,...,300}{
\fill[ink] ($(5.8,0)+(\a:0.83)$) circle(2pt);}
\fill[accent] (4.5,0) circle(2pt);
\fill[accent] (7.1,0) circle(2pt);
\node at (5.8,-1.7) {$O:\ 8$ puntos};
\node at (0,-2.35) {$6\binom62=90$};
\node at (5.8,-2.35) {$8\binom62=120$};
\end{tikzpicture}
\caption{Incidencia marcada de seis a ocho puntos. El dibujo representa
los conjuntos y la inclusión; no pretende convertir las caras o vértices
de un cubo en bloques de Steiner sin su mapa de incidencia.}
\end{figure}

\section{Cadena fundamental dodecafásica y coeficiente de polo}
En el módulo libre de sucesos orientados se fija la cadena
fundamental de la vuelta
\[
c_{12}^{+}=
\sum_{j\in\Z/12\Z}[j,\mathcal F_6]
-[\partial_{\rm ret},\mathcal F_8].
\]
La multiplicidad doce corresponde a los sectores y la multiplicidad uno
a la costura orientada. La involución cambia la orientación de la cadena
completa. Sea
\[
S_m(q)=\frac{q^m}{1-q^{3m}},\qquad 0<q<1.
\]
El lector lineal asigna $S_{90}$ a cada sector y $S_{120}$ a la costura;
en consecuencia,
\[
\mathscr R(c_{12}^{+})=12S_{90}-S_{120}.
\]

\begin{definition}[Evaluación orientada de la incidencia hexádico--octádica]
\label{x_reciprocidad:paper:barbero-definicion}
El parámetro \(\gamma_{\HMT}\) es la evaluación conjugada de la cadena
\(c_{12}^{+}\), completada por su componente angular. En el dominio
\(0<q_\pm<1\), se define mediante
\[
\gamma_{\HMT}=\frac{\pi AC^*}{180}
+[\mathscr R(c_{12}^{+})](q_-)-[\mathscr R(c_{12}^{+})](q_+).
\]
La incidencia marcada determina los espacios de sucesos; la cadena
orientada fija las multiplicidades, y los canales de la
sección~\ref{x_reciprocidad:iv:angular:inversion} determinan la evaluación.
\end{definition}

\begin{theorem}[Coeficiente de polo y evaluación conjugada]
La imagen de la cadena fundamental tiene coeficiente $1/24$
respecto de $(1-q)^{-1}$ y residuo $-1/24$ respecto de $(q-1)^{-1}$.
Su evaluación en los canales previamente generados determina
\begin{equation}
\gamma_{\HMT}=
\frac{\pi AC^*}{180}
+12\bigl[S_{90}(q_-)-S_{90}(q_+)\bigr]
-\bigl[S_{120}(q_-)-S_{120}(q_+)\bigr].
\label{x_reciprocidad:v:ant:eq:gamma}
\end{equation}
\end{theorem}
\begin{proof}
La factorización $1-q^{3m}=(1-q)\sum_{j=0}^{3m-1}q^j$ implica
$\lim_{q\uparrow1}(1-q)S_m(q)=1/(3m)$. Por linealidad,
\[
\frac{12}{270}-\frac1{360}=\frac1{24}.
\]
La coordenada $q-1$ invierte el signo del residuo.
La evaluación conjugada y la componente angular de la definición~\ref{x_reciprocidad:paper:barbero-definicion}
producen \eqref{x_reciprocidad:v:ant:eq:gamma}.
\end{proof}

El certificado diofántico $4a-3b=45$ admite los pares
$(12+3t,1+4t)$, $t\ge0$; certifica la minimalidad de $(12,1)$,
pero no reemplaza el origen de las multiplicidades en la cadena.
La evaluación de \eqref{x_reciprocidad:v:ant:eq:gamma} es
\[
\gamma_{\HMT}
=0.27400767269445927474725526077398041940\ldots.
\]
Esta cifra es una salida del funcional orientado; no es un valor externo
usado para seleccionar $A$, $C^*$ o sus coeficientes.

\section{Paridad y sensibilidad estructural}
Al invertir $C^*$ se intercambian $q_+$ y $q_-$ y cambia de signo
la componente angular. Por ello
\[
\gamma_{\HMT}(A,-C^*)=-\gamma_{\HMT}(A,C^*).
\]
El funcional conserva orientación, mientras que otros lectores,
como $q_+q_-$, son pares. Esta diferencia precede a su interpretación física.

% END OWNER


% BEGIN OWNER /Users/ruben/Documents/New project/output/EXTREMA_MEDIA_RAZON_RECIPROCIDAD_ES_EN_20260918/ARTICULO/ES/source/antecedentes_vii/sections/v_unidad_areal.tex
% RESULTADO_RECUPERADO desde II REV09. Véase MAPA_PROCEDENCIA.json.
% Selección de líneas y cambios mecánicos explícitos; ninguna prueba nueva.
\chapter{Acción, gravedad e información térmica}
\label{x_reciprocidad:v:ant:sec:gravity}
\section{Período dodecafásico y realización circular de la acción}
El calendario ya construido determina $T_{108}=108t_0$.
La velocidad constitutiva, la duración y la longitud de trayectoria se
obtienen mediante las identidades de la
proposición~\ref{x_reciprocidad:iii:prop:identidad-velocidad} y de las
ecuaciones~\eqref{x_reciprocidad:iii:eq:velocidad-longitud-reloj} y
\eqref{x_reciprocidad:iii:eq:longitud-elemental-constitutiva}.
En su realización circular,
\[
\omega_{108}=\frac{2\pi}{108t_0},\qquad
R_{108}=\frac{c_{\rm int}}{\omega_{108}}
=\frac{54}{\pi}\ell_0,\qquad \ell_0=c_{\rm int}t_0.
\]
La unidad $\ell_0$ y las cartas de tiempo y velocidad se mantienen
explícitas; la carta SI no selecciona el ciclo.
Para cada sección orientada de acción $a\in\{\mathrm{pre},\mathrm{ret}\}$,
\[
E_{108,a}=\hbar_a\omega_{108},\qquad
m_{108,a}=\frac{E_{108,a}}{c_{\rm int}^2}.
\]
La longitud reducida y la completa son respectivamente
\[
\frac{\hbar_a}{m_{108,a}c_{\rm int}}=R_{108},
\qquad
\frac{h_a}{m_{108,a}c_{\rm int}}=2\pi R_{108}.
\]
La misma acción y la misma duración intervienen en ambas lecturas.

\begin{proposition}[Selector de coincidencia de radios]
En la familia dimensional
\[
G_a(\vartheta)=\vartheta\frac{c_{\rm int}^5t_0^2}{\hbar_a},
\qquad \vartheta>0,
\]
la condición de la realización circular
$G_a(\vartheta)m_{108,a}/c_{\rm int}^2=R_{108}$
tiene la única solución
\[
\vartheta^\circ_{108}=(54/\pi)^2.
\]
\end{proposition}
\begin{proof}
El primer radio es
$r_g(\vartheta)=\vartheta(\pi/54)\ell_0$.
Compararlo con $(54/\pi)\ell_0$, con $\ell_0>0$, reduce la
igualdad a una ecuación lineal estrictamente creciente.
\end{proof}
La condición de coincidencia es parte explícita de esta realización;
el teorema resuelve su coeficiente sin usar un valor metrológico de $G$.
El radio aquí empleado es $Gm/c^2$, por lo que no se traslada
inadvertidamente el factor dos de la convención $2Gm/c^2$.

\begin{proposition}[Acción y gravedad recíprocas, área conservada]
En las dos cartas anteriores,
\[
\frac{G_{\rm pre}}{G_{\rm ret}}=R_{\rm act}^{-1},\qquad
\hbar_aG_a=\vartheta^\circ_{108}c_{\rm int}^5t_0^2,
\qquad
\ell_P^2=\frac{\hbar_aG_a}{c_{\rm int}^3}
=\vartheta^\circ_{108}\ell_0^2.
\]
\end{proposition}
\begin{proof}
La definición de $G_a$ contiene $\hbar_a^{-1}$ y el resto de factores
es común a las dos cartas. Tomar el cociente y el producto da
las tres igualdades.
\end{proof}
Esta reciprocidad explica por qué la información de orientación puede
variar sin que cambie el área. Se dispone así de una relación precisa
entre las secciones recíprocas de acción, la escala de Planck y el operador areal
de la sección~\ref{x_reciprocidad:v:ant:sec:area}.

% END OWNER


% BEGIN OWNER /Users/ruben/Documents/New project/output/EXTREMA_MEDIA_RAZON_RECIPROCIDAD_ES_EN_20260918/ARTICULO/ES/source/antecedentes_vii/sections/v_area_informacion.tex
% RESULTADO_RECUPERADO desde II REV09. Véase MAPA_PROCEDENCIA.json.
% Selección de líneas y cambios mecánicos explícitos; ninguna prueba nueva.
\chapter{Operador de área y reconstrucción de la información sectorial}
\label{x_reciprocidad:v:ant:sec:area}
El reagrupamiento de seis trits construye una biyección entre
$(\Z/9\Z)^3$ y $(\Z/27\Z)^2$, ambos de cardinal $729$.
La biyección no se presenta como isomorfismo aditivo. Sobre un bloque
de borde $9\times9$, la fuente descompone 36 enlaces en dos
conjuntos de 18, etiquetados por los canales $90$ y $120$.

\section{Reagrupamiento trítico e incidencia del borde}
\label{x_reciprocidad:v:ant:sec:ii-area-incidencia-explicita}

La correspondencia conserva una palabra ordenada, además de su cardinal.
Escribamos cada coordenada \(x_j\in\mathbb Z/9\mathbb Z\), para
\(j=1,2,3\), mediante sus representantes posicionales únicos
\(x_j=r_{2j-1}+3r_{2j}\), con \(r_k\in\{0,1,2\}\). El mapa es
\begin{equation}
 \beta_{729}(x_1,x_2,x_3)
 =\bigl(r_1+3r_2+9r_3,\ r_4+3r_5+9r_6\bigr)
 \in(\mathbb Z/27\mathbb Z)^2.
 \label{x_reciprocidad:v:ant:eq:ii-area-beta729}
\end{equation}

\begin{proposition}[Correspondencia posicional entre interior y borde]
\label{x_reciprocidad:v:ant:prop:ii-area-beta729}
La aplicación \(\beta_{729}\) es biyectiva y conserva el orden de los
seis trits. Su inversa extrae las dos ternas de dígitos y las reagrupa
en tres parejas.
\end{proposition}
\begin{proof}
La expansión posicional de los representantes \(0,\ldots,8\) y
\(0,\ldots,26\) es única. Extraer los seis dígitos, reagruparlos y
extraerlos de nuevo devuelve la misma palabra. La operación inversa
restituye así cada una de las tres coordenadas iniciales.
\end{proof}

El acarreo conserva su posición en la palabra enriquecida. Las operaciones
aditivas de los dos empaquetados tienen fronteras diferentes: por ejemplo,
\(\beta_{729}(0,2,0)=(18,0)\) y
\(\beta_{729}(0,1,0)=(9,0)\), cuya suma en el toro de orden
veintisiete es \((0,0)\), mientras
\(\beta_{729}(0,3,0)=(0,1)\). Esta distinción permite reutilizar
la correspondencia posicional manteniendo tipadas las operaciones.

En el toro \((\mathbb Z/27\mathbb Z)^2\), el bloque
\(B_\partial=\{0,\ldots,8\}^2\) tiene los conjuntos de enlaces
orientados
\begin{align}
 \partial B_{90}
 &=\{((-1,j),(0,j)),\ ((8,j),(9,j)):0\le j\le8\},\\
 \partial B_{120}
 &=\{((i,-1),(i,0)),\ ((i,8),(i,9)):0\le i\le8\}.
 \label{x_reciprocidad:v:ant:eq:ii-area-enlaces}
\end{align}
Todas las coordenadas se leen módulo veintisiete. Cada familia contiene
dieciocho enlaces y ambas son disjuntas. Sus etiquetas conservan los
canales de incidencia; el cardinal de enlaces es \(18+18\).

\begin{theorem}[Capacidad mínima del corte cúbico]
\label{x_reciprocidad:v:ant:thm:ii-area-corte-minimo}
En el grafo cúbico de vértices \(\{0,\ldots,N-1\}^3\), con
\(N\ge2\), enlaces entre vecinos y capacidad unitaria por enlace,
el corte mínimo entre dos caras opuestas tiene capacidad \(N^2\).
Para \(N=9\), un estado producto máximamente mezclado sobre los
trits de una capa mínima tiene información
\begin{equation}
 \operatorname{cap}_{\min}=81,\qquad s_{\rm corte}=81\log3.
 \label{x_reciprocidad:v:ant:eq:ii-area-corte-informacion}
\end{equation}
\end{theorem}
\begin{proof}
Las \(N^2\) rectas paralelas a la normal de las caras forman caminos
disjuntos por enlaces. Todo corte separador intercepta cada camino, de
modo que contiene al menos \(N^2\) enlaces. Los enlaces que atraviesan
una capa transversal constituyen un corte de ese cardinal. Cada trit
uniforme aporta \(\log3\) nats; la entropía es aditiva para el estado
producto de los ochenta y un trits de la capa.
\end{proof}

La capacidad del corte cúbico y el estado de los treinta y seis enlaces
de \eqref{x_reciprocidad:v:ant:eq:ii-area-enlaces} pertenecen a observables definidos sobre
cortes diferentes. La genealogía conserva ambos dominios y sus medidas.

\section{Transporte de los modos colectivos de incidencia}
\label{x_reciprocidad:v:ant:sec:ii-area-transporte-colectivo}

Los espacios \(\mathcal F_6\) y \(\mathcal F_8\) de la
sección~\ref{x_reciprocidad:v:ant:sec:barbero} proporcionan los vectores uniformes
\[
 u_{90}=\frac1{\sqrt{90}}\sum_{f\in\mathcal F_6}(\delta_f,0),
 \qquad
 u_{120}=\frac1{\sqrt{120}}\sum_{f\in\mathcal F_8}(0,\delta_f)
\]
en \(\ell^2(\mathcal F_6)\oplus\ell^2(\mathcal F_8)\).
En \(\ell^2(\partial B_{90})\oplus\ell^2(\partial B_{120})\)
se construyen
\[
 v_{90}=\frac1{\sqrt{18}}\sum_{e\in\partial B_{90}}(\delta_e,0),
 \qquad
 v_{120}=\frac1{\sqrt{18}}\sum_{e\in\partial B_{120}}(0,\delta_e).
\]
Las dos parejas son ortonormales. Denotemos por
\(\mathscr H_{\rm inc}^{\rm coll}\) y
\(\mathscr H_\partial^{\rm coll}\) sus respectivos planos generados.

\begin{theorem}[Entrelazamiento colectivo de las etiquetas sectoriales]
\label{x_reciprocidad:v:ant:thm:ii-area-entrelazamiento-colectivo}
Existe una única isometría lineal sobreyectiva entre esos planos tal que
\begin{equation}
 \mathcal J_{6\to8}u_m=v_m\quad(m=90,120).
 \label{x_reciprocidad:v:ant:eq:ii-area-j-colectivo}
\end{equation}
Para los operadores
\[
 \begin{aligned}
 D_{\rm inc}&=90|u_{90}\rangle\langle u_{90}|
             +120|u_{120}\rangle\langle u_{120}|,\\
 D_\partial&=90|v_{90}\rangle\langle v_{90}|
             +120|v_{120}\rangle\langle v_{120}|,
 \end{aligned}
\]
se cumple
\begin{equation}
 \mathcal J_{6\to8}D_{\rm inc}
 =D_\partial\mathcal J_{6\to8}.
 \label{x_reciprocidad:v:ant:eq:ii-area-j-entrelaza}
\end{equation}
Se transportan también los proyectores sectoriales y su combinación
orientada de coeficientes \(12:-1\).
\end{theorem}
\begin{proof}
La imagen de una base fija la aplicación lineal única. Las normas y los
productos internos se conservan porque ambas bases son ortonormales.
La identidad \eqref{x_reciprocidad:v:ant:eq:ii-area-j-entrelaza} vale sobre cada vector
\(u_m\), pues sus dos miembros son \(m v_m\). Los proyectores se
expresan como \((120I-D)/30\) y \((D-90I)/30\); su cálculo funcional
lineal y cualquier combinación de ambos conservan el entrelazamiento.
\end{proof}

Este transporte actúa sobre los dos modos uniformes y sus etiquetas
espectrales. Las incidencias individuales y las fluctuaciones dentro de
cada sector permanecen en sus espacios completos. El enlace conserva
exactamente el dato colectivo que utilizará la ponderación areal.

\section{Carácter angular y operador de área}
\label{x_reciprocidad:v:ant:sec:ii-area-caracter}

Sea $N_e=\operatorname{diag}(0,1,2)$ en cada enlace y
$N_{90},N_{120}$ las sumas por sector. Fijada la escala
\[
\theta_{\rm clk}=\frac{C^*}{6}\frac{\pi}{180},\qquad
a_0=\frac{1+2\cos(10^\circ)}3\theta_{\rm clk},
\]
el factor angular es el carácter real normalizado del triplete
\begin{equation}
 \chi_{10}=\frac13\Tr\operatorname{diag}
       (1,e^{i\pi/18},e^{-i\pi/18})
 =\frac{1+2\cos(10^\circ)}3,
 \qquad a_0=\chi_{10}\theta_{\rm clk}.
 \label{x_reciprocidad:v:ant:eq:ii-area-caracter-triplete}
\end{equation}
La suma de las fases conjugadas prueba la igualdad. En la rama positiva
considerada, \(\theta_{\rm clk}>0\), \(\chi_{10}>0\) y
\(\gamma_{\HMT}>0\); por tanto, \(a_0\) y la escala areal son
positivos. El triplete angular y su normalización preceden a la
distribución entrópica sobre los enlaces. Con estas cantidades,
el operador areal normalizado es
\[
\widehat{\mathcal A}/\ell_P^2
=\gamma_{\HMT}a_0(N_{90}+N_{120}).
\]
Su espectro tiene valores $n\gamma_{\HMT}a_0$, $0\le n\le72$,
con multiplicidades dadas por $(1+z+z^2)^{36}$.
La prueba consiste en sumar los autovalores $0,1,2$ de los
36 factores tensoriales.

\begin{theorem}[Recuperación de los sectores de borde]
Sean
\[
N=N_{90}+N_{120},\qquad D=90N_{90}+120N_{120}.
\]
Entonces
\begin{equation}
N_{90}=4N-\frac D{30},\qquad
N_{120}=\frac D{30}-3N.
\label{x_reciprocidad:v:ant:eq:recover}
\end{equation}
Por tanto, el total y el lector ponderado recuperan conjuntamente
las dos ocupaciones.
\end{theorem}
\begin{proof}
La matriz $\left(\begin{smallmatrix}1&1\\90&120\end{smallmatrix}\right)$
tiene determinante $30$. Su inversión da \eqref{x_reciprocidad:v:ant:eq:recover}.
En el dominio de ocupaciones enteras, la imagen satisface
$D\in30\Z$ y $90N\le D\le120N$.
\end{proof}

El Hamiltoniano modular de la fuente es
$K_A=-\log\rho_A=cI+\Delta_{\rm mod}D$.
Para $\Delta_{\rm mod}\ne0$ y normalización conocida,
$D=(K_A-cI)/\Delta_{\rm mod}$ publica el segundo observable.
Este operador describe la estructura espectral del estado reducido;
no se identifica por notación con el generador de evolución temporal.
El área conserva la ocupación total y la pareja añade la procedencia.
Por ejemplo, $(N_{90},N_{120})=(1,0)$ y $(0,1)$ comparten $N=1$,
pero tienen $D=90$ y $D=120$.


% BEGIN OWNER /Users/ruben/Documents/New project/output/EXTREMA_MEDIA_RAZON_RECIPROCIDAD_ES_EN_20260918/ARTICULO/ES/source/antecedentes_vii/sections/v_estado_modular.tex
% RESULTADO_RECUPERADO desde II REV09. Véase MAPA_PROCEDENCIA.json.
% Selección de líneas y cambios mecánicos explícitos; ninguna prueba nueva.
\section{Normalización del estado y significado del Hamiltoniano modular}
El segundo observable se explicita mediante el estado que lo define.
Se mantiene la descomposición tensorial anterior: dieciocho qutrits
por sector, cada uno con operador número $\operatorname{diag}(0,1,2)$;
$N_m$ es la suma de esos operadores y los sectores actúan en factores
distintos. Sea
\[
 z_m=1+e^{-m\Delta_{\rm mod}}+e^{-2m\Delta_{\rm mod}},
 \qquad m\in\{90,120\}.
\]
La familia de estados de la fuente tiene la realización producto
\[
 \rho_A=Z^{-1}e^{-\Delta_{\rm mod}(90N_{90}+120N_{120})},
 \qquad Z=z_{90}^{18}z_{120}^{18}.
\]
La conmutatividad de los operadores número y la factorización tensorial
prueban la expresión de $Z$. El estado es de rango completo para
$\Delta_{\rm mod}$ real y finito, y su logaritmo espectral da
\[
 K_A=-\log\rho_A=(\log Z)I+\Delta_{\rm mod}D.
\]
Para $\Delta_{\rm mod}\ne0$, $K_A$ pesa de forma distinta las excitaciones
de los dos sectores. La lectura conjunta con $N$ recupera las dos ocupaciones
por la ecuación~\eqref{x_reciprocidad:v:ant:eq:recover}.

La necesidad de la lectura conjunta se ve en dos ejemplos complementarios.
Los estados de ocupación $(1,0)$ y $(0,1)$ tienen igual total $N=1$
y diferente peso modular. A la inversa, $(4,0)$ y $(0,3)$ tienen
el mismo $D=360$ y diferentes totales. La pareja $(N,D)$ separa ambos
casos. Esta recuperación se refiere a ocupaciones sectoriales, no a
cada vector microscópico dentro de una degeneración de esas ocupaciones.

Cuando $\Delta_{\rm mod}=0$, el estado es proporcional a la identidad
y $K_A$ es escalar: esa lectura deja de distinguir los sectores.
Para $\Delta_{\rm mod}\ne0$, la normalización conocida permite recuperar
$D$ y componerlo con el área. Así queda definido el papel del Hamiltoniano
modular: caracteriza el estado reducido y su distribución de información.
Su inclusión en el artículo está motivada por esa función matemática,
mientras el Hamiltoniano de evolución pertenece al desarrollo dinámico
de la acción.

% END OWNER


\section{Estado reducido, descomposición de Schmidt y entropía de entrelazamiento}
Si $\rho_A=\sum_\nu p_\nu|\nu\rangle\langle\nu|$,
la purificación
\[
|\Psi\rangle=\sum_\nu\sqrt{p_\nu}\,
|\nu\rangle_A|\nu\rangle_B
\]
satisface $\Tr_B|\Psi\rangle\langle\Psi|=\rho_A$.
Así, $-\Tr(\rho_A\log\rho_A)$ es la entropía de entrelazamiento
de esta bipartición pura concreta. El estado ampliado declara
qué sistema porta esa interpretación.

Para una involución de dos hojas con un autovector de cada signo,
\[
\rho_y=\frac{e^{-yR}}{2\cosh y},\qquad R^2=I,
\]
la diagonalización proporciona
\[
S(\rho_y)=\log(2\cosh y)-y\tanh y,\qquad
D(\rho_y\Vert\rho_{-y})=2y\tanh y.
\]
La entropía es par en $y$; la distancia compara explícitamente
las dos orientaciones. No se identifica esta entropía de dos niveles
con la capacidad $81\log3$ del corte triádico.

% END OWNER


% BEGIN OWNER /Users/ruben/Documents/New project/output/EXTREMA_MEDIA_RAZON_RECIPROCIDAD_ES_EN_20260918/ARTICULO/ES/source/antecedentes_vii/sections/v_transduccion_termica.tex
% RESULTADO_RECUPERADO desde II REV09. Véase MAPA_PROCEDENCIA.json.
% Selección de líneas y cambios mecánicos explícitos; ninguna prueba nueva.
% Fragmento aditivo. No sustituye 10b_boltzmann.tex.
% Procedencia y distinción entre recuperación y formalización:
% 10c_boltzmann_procedencia.md. Etiquetas locales prefijadas ii-kb-aut.

\chapter{Memoria genealógica, reducción informacional y transducción térmica}
\label{x_reciprocidad:v:ant:sec:ii-kb-aut-desarrollo}

La construcción térmica utiliza dos publicaciones del mismo desarrollo
APP--TRIT--TPK. La primera conserva los estados, las rutas admisibles y
la información del registro; la segunda proporciona la acción orientada
y la duración del retorno. Su composición determina una energía por
decisión trítica. Las coordenadas de temperatura y energía se introducen
después mediante aplicaciones dimensionales explícitas. Esta secuencia
permite identificar qué información se conserva, qué información publica
un lector y qué magnitud convierte esa información en entropía dimensional.

\section{Lector, memoria y recuperación del estado}

Fijemos un nivel finito del desarrollo y denotemos por \(\Omega\) el
conjunto de estados enriquecidos admisibles. El registro conserva hoja,
orientación, residuo, cociente, acarreo, ruta, fase y memoria en los
dominios donde interviene cada coordenada. Sean
\[
 q:\Omega\longrightarrow Y,\qquad
 M:\Omega\longrightarrow\mathcal M
\]
la publicación visible y la memoria retenida, respectivamente. Para una
distribución \(p\) sobre \(\Omega\), escribimos
\(\mathsf H(p)=-\sum_xp_x\log p_x\), con \(0\log0=0\).
El logaritmo se aplica a probabilidades adimensionales.

\begin{definition}[Recuperación genealógica de una publicación]
El par \(\widehat q=(q,M)\) es un lector recuperable sobre el soporte
\(\Omega_p\) cuando existe
\[
 d:\widehat q(\Omega_p)\longrightarrow\Omega_p,
 \qquad d\circ\widehat q=\operatorname{id}_{\Omega_p}.
\]
La información de fibra de \(q\) es
\(\mathsf H_{\mathrm{fib}}(q;p)=\mathsf H(X\mid q(X))\), donde
\(X\) tiene ley \(p\).
\end{definition}

\begin{proposition}[Identidad de recuperación y coste de la reducción]
Para un lector recuperable,
\begin{align}
 \mathsf H(X\mid q(X),M(X))&=0,\\
 \mathsf H(X)&=\mathsf H(q(X))+\mathsf H(X\mid q(X)),\\
 \mathsf H(X\mid q(X))&=\mathsf H(M(X)\mid q(X)).
 \label{x_reciprocidad:v:ant:eq:ii-kb-aut-fibra}
\end{align}
\end{proposition}
\begin{proof}
La aplicación \(d\) recupera \(X\) del par \((q(X),M(X))\),
por lo que la primera entropía condicional es nula. Como \(q\) es
función de \(X\), la regla de la cadena da la segunda igualdad.
Para la tercera, se desarrolla de dos maneras
\(\mathsf H(X,M(X)\mid q(X))\). Una proporciona
\(\mathsf H(X\mid q(X))\), pues \(M\) es función de \(X\);
la otra proporciona
\(\mathsf H(M(X)\mid q(X))+\mathsf H(X\mid q(X),M(X))\).
\end{proof}

La división euclídea de las hojas de APP muestra el mecanismo
aritmético elemental. Con \(i,j\in\{1,\ldots,9\}\), se conserva
\[
 i+j=9q^+_{ij}+R^+_{ij},\qquad
 ij=9q^\times_{ij}+R^\times_{ij},\qquad
 1\leq R^\sigma_{ij}\leq9\quad(\sigma\in\{+,\times\}).
\]
El par residuo--cociente recupera el valor entero de la operación.
Al sumar las diferencias sobre las ochenta y una posiciones se obtiene
\[
 2025-810=54+9\cdot129.
\]
La memoria del cociente restaura la contribución que no publica la
reducción residual. La recuperabilidad del estado completo requiere,
además, conservar las coordenadas de posición y de ruta pertinentes;
el par residuo--cociente identifica aquí el resultado de la operación.

\section{Actualización reversible, decisión trítica y borrado}

Sea \(U:\Omega\to\Omega\) la actualización TPK en un dominio donde
el transporte enriquecido posee inversa. La ley transportada es
\(p'_y=p_{U^{-1}y}\). La biyección conserva el multiconjunto de
probabilidades y, por tanto,
\begin{equation}
 \mathsf H(U(X))=\mathsf H(X),\qquad
 \mathsf H(X\mid U(X))=0.
 \label{x_reciprocidad:v:ant:eq:ii-kb-aut-reversible}
\end{equation}
Si la salida observada es \(q(U(X))\), la información no publicada
se cuantifica por \(\mathsf H(X\mid q(U(X)))\). El estado completo,
su publicación y el reinicio de un registro son así tres mapas distintos.

Para los tres estados orientados del TRIT, sea
\(p=(p_{-1},p_0,p_{+1})\) una distribución normalizada. Su información es
\begin{equation}
 I_{\mathrm{TRIT}}(p)=-\sum_{a=-1}^{1}p_a\log p_a,
 \qquad 0\leq I_{\mathrm{TRIT}}(p)\leq\log3.
 \label{x_reciprocidad:v:ant:eq:ii-kb-aut-trit}
\end{equation}
El extremo superior corresponde a \(p_a=1/3\). En efecto,
\[
 D(p\Vert(1/3,1/3,1/3))
 =\sum_ap_a\log(3p_a)=\log3-I_{\mathrm{TRIT}}(p)\geq0;
\]
la desigualdad se obtiene de la convexidad de \(t\log t\), y la
igualdad caracteriza la distribución uniforme. El extremo inferior
corresponde a una distribución concentrada en un estado.

La realización ideal de Landauer utiliza una memoria con estados lógicos
energéticamente degenerados, un entorno isotermo de temperatura \(\Theta\)
y un balance que conserva la salida del mapa lógico y contabiliza los
registros auxiliares. En ese modelo, la disminución de entropía lógica es
\(k_B[\mathsf H(X)-\mathsf H(F(X))]
=k_B\mathsf H(X\mid F(X))\) para un mapa determinista \(F\).
El balance entrópico exige transferir al entorno al menos esa cantidad;
su expresión energética es \(Q\geq k_B\Theta\mathsf H(X\mid F(X))\).
La distinción entre evolución lógica reversible y borrado físico
corresponde a la realización de Landauer--Bennett \cite{x_reciprocidad:bennett2003}.
La actualización invertible
\eqref{x_reciprocidad:v:ant:eq:ii-kb-aut-reversible} tiene contribución lógica nula; el
reinicio de una terna uniforme requiere transferir \(\log3\) nats
al entorno y satisface \(Q\geq k_B\Theta\log3\).
La energía de control y la disipación de una realización física
concreta pertenecen a su balance termodinámico adicional.

\section{Acción, frecuencia y energía de decisión}

Sean \(\mathcal L_{\rm act}\) y \(\mathcal L_T\) las rectas
dimensionales de acción y duración; la recta energética es
\(\mathcal L_E=\mathcal L_{\rm act}\otimes\mathcal L_T^*\).
El desarrollo anterior aporta las secciones positivas
\(\hbar_a\in\mathcal L_{\rm act}\), con
\(a\in\{\mathrm{pre},\mathrm{ret}\}\), y la duración elemental
\(t_0\in\mathcal L_T\). El calendario completo proporciona
\[
 T_{108}=108t_0,\qquad
 \omega_{108}=\frac{2\pi}{T_{108}},\qquad
 h_a=2\pi\hbar_a.
\]
La contracción dimensional entre acción y frecuencia define
\begin{equation}
 \epsilon_a:=\hbar_a\omega_{108}
 =\frac{h_a}{108t_0}
 =\frac{\hbar_a}{t_P},\qquad
 t_P=\frac{54}{\pi}t_0.
 \label{x_reciprocidad:v:ant:eq:ii-kb-aut-energia}
\end{equation}
El último miembro emplea el mismo período circular reducido que la
sección gravitatoria. Las tres expresiones corresponden a una única
sección de energía, con una orientación \(a\) mantenida en todos
los factores. La información normalizada de una decisión uniforme
produce la sección de energía por nat
\begin{equation}
 \varepsilon_a:=\frac{\epsilon_a}{\log3}\in\mathcal L_E^+.
 \label{x_reciprocidad:v:ant:eq:ii-kb-aut-energia-nat}
\end{equation}
La acción previamente construida determina la escala de energía al
fijar la duración; el cardinal de alternativas del TRIT determina su
normalización informacional. El valor de \(\log3\) entra por el
conteo y la medida declarados en \eqref{x_reciprocidad:v:ant:eq:ii-kb-aut-trit}.

\section{Definición dimensional de Boltzmann y selección de coordenadas}

Sea \(\mathcal L_\Theta\) la recta orientada de temperatura.
El espacio dimensional de entropía es
\(\mathcal L_{\rm ent}=\mathcal L_E\otimes\mathcal L_\Theta^*\).
Una aplicación lineal positiva
\(k_B:\mathcal L_\Theta\to\mathcal L_E\) es, equivalentemente,
un elemento positivo de \(\mathcal L_{\rm ent}\). Por eso el mismo
objeto puede actuar sobre una temperatura o multiplicar una información
adimensional.

\begin{definition}[Normalización térmica de la decisión trítica]
Un par térmico compatible con la sección orientada \(a\) consta de
un transductor positivo \(k_B\) y una sección positiva
\(\Theta_{{\rm clk},a}\in\mathcal L_\Theta\) tales que
\begin{equation}
 k_B(\Theta_{{\rm clk},a})=\varepsilon_a,
 \qquad
 k_B\Theta_{{\rm clk},a}\log3=\epsilon_a.
 \label{x_reciprocidad:v:ant:eq:ii-kb-aut-par-termico}
\end{equation}
La constante de Boltzmann desempeña la función de transductor
energía--temperatura de la información trítica en esta realización.
\end{definition}

\begin{theorem}[Unicidad relativa a la sección térmica y covariancia dimensional]
Fijada una sección positiva \(\Theta_{{\rm clk},a}\), existe una
única aplicación lineal positiva que satisface
\eqref{x_reciprocidad:v:ant:eq:ii-kb-aut-par-termico}. Para bases positivas \(u_E,u_\Theta\),
si
\[
 \epsilon_a=\epsilon_a^{\rm num}u_E,\quad
 \Theta_{{\rm clk},a}=\theta_a u_\Theta,\quad
 k_B(u_\Theta)=b\,u_E,
\]
su coordenada es
\begin{equation}
 b=\frac{\epsilon_a^{\rm num}}{\theta_a\log3}.
 \label{x_reciprocidad:v:ant:eq:ii-kb-aut-coordenada}
\end{equation}
Bajo \(u'_E=s_Eu_E\), \(u'_\Theta=s_\Theta u_\Theta\), con
\(s_E,s_\Theta>0\), se cumple
\[
 b'=\frac{s_\Theta}{s_E}b,\qquad
 \theta'_a=\frac{\theta_a}{s_\Theta},\qquad
 (\epsilon_a^{\rm num})'=\frac{\epsilon_a^{\rm num}}{s_E}.
\]
\end{theorem}
\begin{proof}
Una sección no nula genera una recta. Prescribir su imagen determina
la aplicación lineal única; la positividad de ambas secciones prueba
su positividad. Sustituir las coordenadas en
\eqref{x_reciprocidad:v:ant:eq:ii-kb-aut-par-termico} da
\(b\theta_a\log3=\epsilon_a^{\rm num}\). Las tres leyes de
transformación se obtienen expresando las mismas secciones y la misma
aplicación en las bases nuevas.
\end{proof}

Una misma constante de Boltzmann debe actuar en las dos orientaciones.
Puesto que \(t_P\) es común, la compatibilidad simultánea equivale a
\begin{equation}
 \frac{\Theta_{{\rm clk},{\rm pre}}}{\Theta_{{\rm clk},{\rm ret}}}
 =\frac{\epsilon_{\rm pre}}{\epsilon_{\rm ret}}
 =\frac{\hbar_{\rm pre}}{\hbar_{\rm ret}}.
 \label{x_reciprocidad:v:ant:eq:ii-kb-aut-orientaciones}
\end{equation}
En efecto, un isomorfismo de rectas conserva cocientes de secciones.
Recíprocamente, si la igualdad anterior se cumple, el isomorfismo
determinado en una orientación satisface también la normalización de
la otra. Así, el retorno modifica la sección térmica del ciclo en la
misma proporción que la energía, mientras el transductor permanece
común a ambas lecturas.

La ecuación anterior explicita la elección requerida para publicar un
valor individual. Si la sección térmica todavía no ha sido seleccionada
por un lector independiente, para cada \(\lambda>0\) el par
\[
 (k_B,\Theta_{{\rm clk},a})
 \longmapsto(\lambda k_B,\lambda^{-1}\Theta_{{\rm clk},a})
\]
conserva \(\epsilon_a\). Por consiguiente, la presentación de una
coordenada individual internamente seleccionada exige dar el lector de
\(\Theta_{{\rm clk},a}\), las bases y su transporte dimensional.
Las ecuaciones de acción e información permanecen determinadas antes
de esa especificación. La identificación posterior de las bases con
joule y kelvin pertenece a la publicación metrológica del resultado.

\section{Entropía dimensional e información del estado reducido}

Para una distribución de rutas admitidas \(p\), o para un estado
reducido de matriz de densidad \(\rho\), sean
\[
 s(p)=-\sum_xp_x\log p_x,\qquad
 s(\rho)=-\operatorname{Tr}(\rho\log\rho).
\]
Sus publicaciones dimensionales son
\begin{equation}
 S(p)=k_Bs(p),\qquad S(\rho)=k_Bs(\rho)
 \quad\text{en }\mathcal L_{\rm ent}.
 \label{x_reciprocidad:v:ant:eq:ii-kb-aut-entropia}
\end{equation}
La evaluación en \(\Theta\) produce una energía
\(\Theta S=k_B\Theta\,s\). Para la terna uniforme a la temperatura
del ciclo, \eqref{x_reciprocidad:v:ant:eq:ii-kb-aut-par-termico} da exactamente
\[
 S_{\rm TRIT}=k_B\log3,\qquad
 \Theta_{{\rm clk},a}S_{\rm TRIT}=\epsilon_a.
\]
Si \(\rho_A\) es la reducción de la purificación bipartita
construida en la sección areal, \(s(\rho_A)\) mide el entrelazamiento
de esa bipartición. El transporte de unidades conserva sus
autovalores, su orientación y las ocupaciones sectoriales. La escala
areal \(\gamma a_0\ell_P^2\) convierte ocupaciones en área; el
transductor \(k_B\) convierte información en entropía dimensional.
Ambas lecturas utilizan el mismo registro mediante operadores distintos.

\section{Normalización espectral de las historias térmicas}

Sea \(\mathcal G\) un grafo finito de extensiones admisibles ya
construido, con matriz de adyacencia \(A\). Cada arista \(e\) conserva
sus coordenadas de ruta y recibe un incremento adimensional de acción
\(a_*(e)>0\) del lector correspondiente. Elegida la sección de energía
\(\epsilon_*\), pongamos \(t=\beta\epsilon_*\), donde
\(\beta=(k_B\Theta)^{-1}\). El operador térmico es
\begin{equation}
 (L_tf)(y)=\sum_{e:x\to y}\exp[-t a_*(e)]f(x).
 \label{x_reciprocidad:v:ant:eq:ii-kb-aut-operador}
\end{equation}
Todos sus exponentes son adimensionales. La condición
\(\rho(L_t)=1\) normaliza el crecimiento ponderado de las historias.

% REMATE_LOCAL_X_TERMICA
\Needspace{35\baselineskip}
\begin{proposition}[Existencia y unicidad de la presión normalizada]
Supongamos que \(A\) es irreducible, \(\rho(A)>1\), y que
\(0<m\leq a_*(e)\leq M\) para todas las aristas. Existe un único
\(t_*>0\) tal que \(\rho(L_{t_*})=1\), y
\begin{equation}
 \frac{\log\rho(A)}{M}\leq t_*\leq
 \frac{\log\rho(A)}{m}.
 \label{x_reciprocidad:v:ant:eq:ii-kb-aut-presion}
\end{equation}
\end{proposition}
\begin{proof}
La comparación entrada a entrada proporciona
\(e^{-tM}A\leq L_t\leq e^{-tm}A\) para \(t\geq0\).
La monotonía del radio espectral de matrices no negativas implica
las cotas correspondientes para \(\rho(L_t)\). El radio espectral
es continuo, vale \(\rho(A)>1\) en cero y tiende a cero en infinito.
Para \(s>0\), además,
\(L_{t+s}\leq e^{-sm}L_t\), de modo que
\(\rho(L_{t+s})\leq e^{-sm}\rho(L_t)<\rho(L_t)\).
La continuidad y la disminución estricta prueban existencia y
unicidad. Sustituir el valor uno en las cotas iniciales demuestra
\eqref{x_reciprocidad:v:ant:eq:ii-kb-aut-presion}.
\end{proof}

La realización local de tres alternativas equiponderadas, representada
por un vértice con tres lazos y \(a_*(e)=1\), satisface
\(\rho(L_t)=3e^{-t}\). Su normalización da \(t_*=\log3\).
Con \(\epsilon_*=\epsilon_a\), resulta
\(\beta_{\rm clk}=\log3/\epsilon_a\), exactamente la relación
\eqref{x_reciprocidad:v:ant:eq:ii-kb-aut-par-termico}. Este cálculo identifica la
confluencia local trítica. Para un grafo de historias con restricciones
adicionales, la confluencia con el mismo reloj exige comprobar
\[
 \rho\!\left(L_{(\log3/\epsilon_a)\epsilon_*}\right)=1
\]
con sus incrementos efectivos. La normalización local y la presión
del sistema completo conservan así sus respectivos dominios.

La relación obtenida reúne estructura, valor y significado. El registro
determina la información recuperable; las probabilidades sobre las
continuaciones determinan su lectura entrópica; la acción y el período
determinan \(\epsilon_a\); el transductor energía--temperatura
publica la entropía dimensional. El área y el entrelazamiento se
relacionan a través del estado reducido y de sus operadores de borde,
manteniendo visible qué dato produce cada magnitud.

% END OWNER


% BEGIN OWNER /Users/ruben/Documents/New project/output/EXTREMA_MEDIA_RAZON_RECIPROCIDAD_ES_EN_20260918/ARTICULO/ES/source/antecedentes_vii/28_portador_tensorial.tex
% FORMALIZACION_REUNIDA: representación pentacomponente de II,
% con antecedente Paley y demostración directa de los momentos del marco.
% El retorno transporta la carta; no se identifica un ciclo de orden doce con A5.
\chapter{Incidencia orientada y representación tensorial gravitatoria}
\label{x_reciprocidad:vii:sec:pentacomponente}
La publicación excepcional del estado conserva la incidencia que acompaña
a su lectura aritmética. La matriz ternaria \(A_W\) de
\eqref{x_reciprocidad:exc:aw} proporciona aquí un antecedente material preciso. Su
levantamiento balanceado \(C=\operatorname{bal}(A_W)\), con representantes
\(0,1,-1\), satisface
\[
 C=C^{\mathsf T},\qquad C^2=5I_6,\qquad \operatorname{Tr}C=0.
\]
Las seis posiciones tienen el orden \((\infty,0,1,2,3,4)\).
Esta incidencia produce primero un espacio tridimensional; su lectura
cuadrática produce después el espacio de cinco componentes sobre el que
actúa el acoplamiento gravitatorio.

\section{Recubrimiento firmado y transporte del calendario}
Definimos
\[
 E_C=\frac12(I_6+C/\sqrt5),\quad
 \mathcal H_C=\operatorname{im}E_C,\quad
 v_{u,\sigma}=\sigma\sqrt2E_Ce_u,\quad \sigma\in\{1,-1\}.
\]
\begin{proposition}[Doce posiciones orientadas]
\label{x_reciprocidad:vii:prop:doce-orientadas}
Los vectores \(v_{u,\sigma}\) son doce vectores unitarios distintos
en un espacio de dimensión tres. Su incidencia está determinada por
\(\langle v_{u,\sigma},v_{v,\tau}\rangle=1/\sqrt5\).
La inversión de la orientación actúa por antipodía.
\end{proposition}
\begin{proof}
La identidad cuadrática de \(C\) implica \(E_C^2=E_C=E_C^*\);
su traza es tres. Puesto que la diagonal de \(C\) es nula,
\[
 \langle v_{u,\sigma},v_{v,\tau}\rangle
 =\begin{cases}\sigma\tau,&u=v,\\
 \sigma\tau C_{uv}/\sqrt5,&u\ne v.
 \end{cases}
\]
Esto prueba normalización, distinción de los doce vectores e incidencia.
Cambiar \(\sigma\) por \(-\sigma\) cambia el signo del vector.
\end{proof}

Una expresión euclidiana explícita se obtiene poniendo
\(\nu=\sqrt{1+\varphi^2}\) y
\[
 B=\frac1\nu
 \begin{pmatrix}
 0&-1&-\varphi&0&\varphi&1\\
 -1&-\varphi&0&1&0&-\varphi\\
 -\varphi&0&-1&-\varphi&-1&0
 \end{pmatrix}.
\]
La relación \(\varphi^2=\varphi+1\), para la coordenada ya generada,
da \(B^*B=2E_C\).
En consecuencia \(B/\sqrt2|_{\mathcal H_C}\) es una isometría y lleva
\(v_{u,\sigma}\) a \(\sigma Be_u\). La imagen es
\[
 \mathcal V=\frac1\nu\bigl(
 \{0\}\times\{\pm1\}\times\{\pm\varphi\}\ \cup\
 \{\pm1\}\times\{\pm\varphi\}\times\{0\}\ \cup\
 \{\pm\varphi\}\times\{0\}\times\{\pm1\}\bigr).
\]
Así quedan relacionados la coordenada interna, la incidencia y sus
doce representantes geométricos.

La coordenada del calendario se escribe \((b,r)\in
\mathbb Z/12\mathbb Z\times\{0,\ldots,8\}\). Una ruta de longitud \(n\)
induce
\[
 k(r,n)=\left\lfloor\frac{r+n}{9}\right\rfloor,\qquad
 (b,r)\longmapsto
 \bigl(b+k(r,n)\bmod12,\ (r+n)\bmod9\bigr).
\]
El alineamiento explícito de sus posiciones \(p=b+1\) con la carta firmada es
\[
\begin{array}{c|rrrrrrrrrrrr}
p&1&2&3&4&5&6&7&8&9&10&11&12\\ \hline
u&\infty&0&0&\infty&1&3&3&1&2&4&4&2\\
\sigma&+&-&+&-&-&+&-&+&-&+&-&+
\end{array}
\]
Esta tabla fija una trivialización del portador. Cada cambio de posición
transporta conjuntamente los vectores y su matriz de incidencia.

\begin{proposition}[Covariancia del marco bajo retorno]
Si \(S e_p=e_{p+1\bmod12}\), la representación matricial de un
operador \(P\) del portador se transporta como \(P_k=S^kPS^{-k}\).
Este transporte conserva productos escalares, incidencias y composición
de rutas. Tras doce retornos vuelve la carta, conservándose la memoria
acumulada en el estado fuente.
\end{proposition}
\begin{proof}
\(S\) es ortogonal y \(S^{12}=I\); de aquí se siguen las propiedades
métricas y la periodicidad de la carta. Para la composición, con
\(r'=(r+n)\bmod9\), la división euclidiana da
\(k(r,n+m)=k(r,n)+k(r',m)\). Los exponentes de \(S\) se suman
exactamente con el acarreo. La operación actúa sobre la coordenada
del calendario; las actualizaciones de memoria de la ruta siguen
presentes en su dominio enriquecido.
\end{proof}

\section{Lectura cuadrática y espacio de cinco componentes}
Sea \(\mathbb R^{\mathcal V}\) el espacio de las doce coordenadas,
\(\mathsf A e_v=e_{-v}\) su involución antipodal y
\(\mathsf J=\mathbf1\mathbf1^{\mathsf T}\). Se definen
\[
 P_5=\frac12(I+\mathsf A)-\frac1{12}\mathsf J,\qquad
 Q_v=vv^{\mathsf T}-\frac13I_3,\qquad
 \Gamma_0x=\sum_{v\in\mathcal V}x_vQ_v.
\]
\begin{theorem}[Proyector pentacomponente e isometría tensorial]
\label{x_reciprocidad:vii:thm:portador-pentacomponente}
\(P_5\) es el proyector ortogonal de rango cinco sobre el espacio de
coordenadas pares bajo antipodía y de suma nula. Además,
\[
 \Gamma_0^*\Gamma_0=\frac85P_5,\qquad
 \Gamma_0\Gamma_0^*=\frac85I_{\operatorname{Sym}^2_0(\mathbb R^3)}.
\]
Por tanto
\[
 \Gamma_5=\sqrt{\frac58}\,\Gamma_0|_{\operatorname{im}P_5}
\]
es una isometría sobre el espacio de tensores simétricos sin traza.
\end{theorem}
\begin{proof}
El espacio par tiene una coordenada por cada par antipodal, por lo que
su dimensión es seis. El vector constante pertenece a él; restar su
proyector de rango uno da \(P_5\), de rango cinco.

Como
\(\langle Q_v,Q_w\rangle_F=(v^{\mathsf T}w)^2-1/3\), las entradas
de \(\Gamma_0^*\Gamma_0\) valen \(2/3\) para \(w=\pm v\) y
\(-2/15\) en los otros diez lugares de cada fila. Éstas son
exactamente las entradas de \((8/5)P_5\).

También puede comprobarse la identidad en el espacio tensorial sin
recurrir a una clasificación de representaciones. La lista de vértices da
\[
 \sum_vv_i^4=\frac{12}{5},\qquad
 \sum_vv_i^2v_j^2=\frac45\quad(i\ne j),
\]
y los momentos con alguna potencia impar se anulan por los cambios de
signo. En efecto,
\((1+\varphi^2)^2=5\varphi^2\) y \(1+\varphi^4=3\varphi^2\).
Para un tensor simétrico \(Q\), estas sumas dan
\[
 \sum_v(v^{\mathsf T}Qv)\,vv^{\mathsf T}
 =\frac85Q+\frac45(\operatorname{tr}Q)I.
\]
Con \(\operatorname{tr}Q=0\), y usando
\(\sum_vv^{\mathsf T}Qv=4\operatorname{tr}Q=0\), resulta
\(\Gamma_0\Gamma_0^*Q=(8/5)Q\). Las dos identidades prueban
la isometría normalizada y su sobreyectividad.
\end{proof}

La sección gravitatoria de \(\ref{x_reciprocidad:v:ant:sec:gravity}\) determina el
operador \(G_aP_5\). Su espectro tiene \(G_a\) con multiplicidad
cinco y cero con multiplicidad siete. La orientación se conserva en
el portador y en su transporte; \(G_a\) fija el acoplamiento común.

\section{Proyección transversal y helicidades}
Para \(n\in S^2\), pongamos \(\Pi=I-nn^{\mathsf T}\) y
\[
 \Lambda_nQ=\Pi Q\Pi-\frac12\operatorname{tr}(\Pi Q\Pi)\Pi.
\]
\begin{proposition}[Reducción transversal de rango dos]
\(\Lambda_n\) es el proyector ortogonal sobre
\(\{Q\in\operatorname{Sym}^2_0(\mathbb R^3):Qn=0\}\), de dimensión dos.
\end{proposition}
\begin{proof}
\(\Pi n=0\) da transversalidad y \(\operatorname{tr}\Pi=2\) da
traza nula. Un tensor transversal sin traza queda fijo.
La fórmula es autoadjunta para el producto de Frobenius.
En un marco \((e_1,e_2,n)\), la imagen tiene base ortonormal
\[
 E_+=\frac{e_1e_1^{\mathsf T}-e_2e_2^{\mathsf T}}{\sqrt2},\qquad
 E_\times=\frac{e_1e_2^{\mathsf T}+e_2e_1^{\mathsf T}}{\sqrt2}.
\]
\end{proof}
Los otros tres vectores de una base ortonormal del portador son
\[
 E_1=\frac{e_1n^{\mathsf T}+ne_1^{\mathsf T}}{\sqrt2},\quad
 E_2=\frac{e_2n^{\mathsf T}+ne_2^{\mathsf T}}{\sqrt2},\quad
 E_0=\frac{2nn^{\mathsf T}-e_1e_1^{\mathsf T}-e_2e_2^{\mathsf T}}{\sqrt6}.
\]
Al rotar el marco transversal un ángulo \(\theta\), las parejas
\((E_+,E_\times)\) y \((E_1,E_2)\) giran \(2\theta\) y
\(\theta\), respectivamente, mientras \(E_0\) queda fijo. La
complejificación tiene así los pesos \(2,-2,1,-1,0\).
\(\Lambda_n\) conserva los dos primeros y anula los restantes.
Queda establecida la cadena
\[
 \mathbb R^{12}\xrightarrow{P_5}\operatorname{im}P_5
 \xrightarrow{\Gamma_5}\operatorname{Sym}^2_0(\mathbb R^3)
 \xrightarrow{\Lambda_n}\mathcal H_{\rm TT}(n),
 \qquad 12\longrightarrow5\longrightarrow5\longrightarrow2.
\]
La representación pentacomponente y su reducción son resultados
algebraicos anteriores a la elección de ecuaciones de campo.
En la realización gravitatoria sin masa, las restricciones de calibre
seleccionan la imagen transversal. Esta distinción conserva tanto las
cinco coordenadas estructurales como las dos helicidades radiativas.


% END OWNER


% BEGIN OWNER /Users/ruben/Documents/New project/output/EXTREMA_MEDIA_RAZON_RECIPROCIDAD_ES_EN_20260918/ARTICULO/ES/source/antecedentes_vii/29_retorno_y_memoria_traslacional.tex
% Artículo VII, recuperación aditiva, revisión 04, 10-09-2026.
% RESULTADO_RECUPERADO: S15, 11c_gravedad_holonomia_rev6.tex:16--126.
% El cociclo canónico c27=(1,18) se conserva como antecedente del cómputo
% de rutas; este fragmento prueba su composición y representación.
% La memoria traslacional finita y la amplitud cosmológica s0 tienen
% tipos distintos. El cociclo no evalúa por sí solo una densidad local.

\chapter{Retorno de fase y representación de la memoria acumulada}
\label{x_reciprocidad:vii:sec:retorno-memoria-gravedad}

La cronología de las rutas enriquecidas distingue el retorno de posición,
el retorno de orientación y el retorno de fase. La publicación observable
identifica ciertos estados de llegada, mientras el registro acumulativo
conserva la secuencia recorrida. La representación geométrica debe
transportar ambos datos conjuntamente.

\section{Cociclo de inversión y cambio de hoja}
\label{x_reciprocidad:vii:subsec:cociclo-memoria}

Sea \(F_{\rm obs}\) la publicación cronológica del transporte TPK.
Sus bloques canónicos tienen la jerarquía
\[
\begin{array}{c|l}
27&\text{retorno posicional con inversión de orientación},\\
54&\text{retorno posicional y orientacional},\\
108&\text{retorno completo de la fase observable}.
\end{array}
\]
El registro cuenta las inversiones de orientación \(g(a)\) y los
cambios efectivos de hoja \(s(a)\). Para una ruta,
\begin{equation}
 c(\gamma)=\sum_{a\subset\gamma}(g(a),s(a)),\qquad
 c(\gamma_2\gamma_1)=c(\gamma_2)+c(\gamma_1).
 \label{x_reciprocidad:vii:eq:cociclo-ruta}
\end{equation}
El cómputo canónico del bloque de longitud \(27\) publica
\(c_{27}=(1,18)\). La composición conserva esos incrementos en
los cuatro bloques que completan el período observable.
En rutas orientadas hacia adelante el registro es acumulativo;
su extensión al grupoide de rutas utiliza \(\mathbb Z^2\) y
\(c(\bar\gamma)=-c(\gamma)\).

\begin{proposition}[Memoria de la vuelta observable completa]
\label{x_reciprocidad:vii:prop:memoria-vuelta-completa}
Escribiendo \(\mathcal K_{27}=F_{\rm obs}^{27}\), la elevación
al registro, para \(x\in\mathcal O_{108}\) y \(z\in\mathbb Z^2\), satisface
\[
 \widetilde{\mathcal K}_{27}(x,z)
   =(\mathcal K_{27}x,z+(1,18)),\qquad
 \widetilde{\mathcal K}_{27}^{\,4}(x,z)
   =(x,z+(4,72)).
\]
Tras \(N\) vueltas completas, el incremento es \(N(4,72)\).
\end{proposition}

\begin{proof}
La proyección observable de \(\mathcal K_{27}\) tiene orden cuatro.
El cómputo canónico anterior fija \(c_{27}=(1,18)\) de manera constante
sobre esa órbita observable. La aditividad de
\eqref{x_reciprocidad:vii:eq:cociclo-ruta} suma los cuatro incrementos y da
\(c_{108}=4(1,18)=(4,72)\).
La repetición aplica la misma ley de concatenación y produce
\(Nc_{108}\). El retorno de la primera coordenada conserva el
incremento de la segunda.
\end{proof}

\section{Representación afín del cociclo}
\label{x_reciprocidad:vii:subsec:representacion-afin-cociclo}

La realización discreta utiliza un elemento \(R_4\) de orden cuatro en
\(\operatorname{Spin}^+(1,3)\), un vector temporal unitario \(u\)
fijado por su acción vectorial y una unidad de longitud \(\ell_0>0\).
La escala \(\ell_0\) conserva su procedencia en la sección dimensional.
Definimos
\begin{equation}
 \rho_{\rm C}^{\rm disc}(\gamma)=
 \bigl(R_4^{c_g(\gamma)},\,\ell_0c_s(\gamma)u\bigr)
 \in\operatorname{Spin}^+(1,3)\ltimes\mathbb R^{1,3}.
 \label{x_reciprocidad:vii:eq:representacion-discreta-memoria}
\end{equation}
La potencia actúa sobre el vector mediante la representación vectorial
asociada del grupo de espín.

\begin{theorem}[Composición y amplitud traslacional del retorno]
\label{x_reciprocidad:vii:thm:traslacion-memoria}
La aplicación \(\rho_{\rm C}^{\rm disc}\) conserva la identidad,
la concatenación y la inversión. En los retornos completos,
\begin{equation}
 \rho_{\rm C}^{\rm disc}(\gamma_{108})=(I,72\ell_0u),\qquad
 \rho_{\rm C}^{\rm disc}(\gamma_{108}^{\,N})=(I,72N\ell_0u).
 \label{x_reciprocidad:vii:eq:traslacion-retorno-108}
\end{equation}
La inversión cambia el signo de la traslación.
\end{theorem}

\begin{proof}
El producto semidirecto es
\((R,a)(S,b)=(RS,a+Rb)\). Como \(R_4u=u\),
\[
\begin{split}
 \rho_{\rm C}^{\rm disc}(\gamma_2)
 \rho_{\rm C}^{\rm disc}(\gamma_1)
 &=
 \bigl(R_4^{c_g(\gamma_2)+c_g(\gamma_1)},
 \ell_0(c_s(\gamma_2)+c_s(\gamma_1))u\bigr)\\
 &=\rho_{\rm C}^{\rm disc}(\gamma_2\gamma_1).
\end{split}
\]
El cociclo de la identidad es cero y el de la inversa es el opuesto.
Para la vuelta completa se sustituyen \(c_{108}=(4,72)\) y
\(R_4^4=I\). La concatenación de \(N\) vueltas mantiene la parte
rotacional identidad y suma sus traslaciones.
\end{proof}

La amplitud \(72\ell_0\) es una longitud asociada a una ruta finita.
Una realización mediante cotetrada y conexión la transporta a su
holonomía afín. La contribución local a las ecuaciones gravitatorias
se obtiene después mediante la curvatura de esa conexión, la
variación de la acción material y la normalización de su densidad.
La representación conserva así el contenido acumulativo exacto que
debe intervenir en esa realización.

% END OWNER


% BEGIN OWNER /Users/ruben/Documents/New project/output/EXTREMA_MEDIA_RAZON_RECIPROCIDAD_ES_EN_20260918/ARTICULO/ES/source/antecedentes_vii/30_pantalla_y_respuesta.tex
% Artículo VII. Incorporación de resultados preexistentes, 10-09-2026.
% Propietarios íntegros preservados en fuentes_conservadas:
% 17_pantalla_cubica_soldadura.tex (P), 18_clausura_energetica.tex (E).
% P: incidencia, cociente, soldadura celular y refinamiento.
% E: respuesta, unicidad, extensión mínima y expectativa tracial.
% Dependencias anteriores de VII: núcleo APP--TRIT--TPK; realización
% cúbica de la ruta, cociclo de transporte y compresión q_vis=5/9.
% Este fragmento no constituye por sí solo el artículo autónomo.
% La identidad de pantalla se distingue de la contracción material de espín.

\chapter{Soldadura de frontera y respuesta energética}
\label{x_reciprocidad:vii:sec:pantalla-respuesta}

La realización cúbica del transporte conserva seis incidencias orientadas.
Su cociente geométrico, la forma lorentziana y la respuesta energética
se obtienen de la misma matriz de incidencias. El transporte procede de
la ruta enriquecida APP--TRIT--TPK: la cara, el signo, el acarreo y la
memoria acompañan a cada arista. En esta sección se desarrolla la
operación lineal que actúa sobre esa publicación de frontera.

\section{Incidencia cúbica y cociente de rango cuatro}
\label{x_reciprocidad:vii:sec:cociente-pantalla}

Sea \(H_F=\mathbb R^F\), con
\(F=\{(i,\epsilon):1\leq i\leq3,\ \epsilon\in\{-1,1\}\}\),
su producto escalar euclidiano y su base \(e_{i,\epsilon}\).
La oposición de caras es la involución ortogonal
\(Re_{i,\epsilon}=e_{i,-\epsilon}\). Definimos
\[
 u=\sum_{i,\epsilon}e_{i,\epsilon},\qquad
 P_u=\frac{uu^{\mathsf T}}6,\qquad P_-=\frac{I-R}{2},\qquad
 P_0=\frac{I+R}{2}-P_u,\qquad P_\square=P_u+P_-.
\]
Los vértices son \(V=\{-1,1\}^3\). La matriz de incidencia
\(M:\mathbb R^V\longrightarrow H_F\) queda determinada por
\[
 M_{(i,\epsilon),\nu}=\mathbf1_{\{\nu_i=\epsilon\}}.
\]

\begin{proposition}[Descomposición de la incidencia]
\label{x_reciprocidad:vii:prop:incidencia}
Se tiene
\[
 MM^{\mathsf T}=12P_u+4P_-,
 \qquad
 \ker M^{\mathsf T}=P_0H_F.
\]
Por consiguiente, el cociente
\(Q=H_F/P_0H_F\), identificado ortogonalmente con
\(P_\square H_F\), tiene dimensión cuatro.
\end{proposition}

\begin{proof}
Una cara contiene cuatro vértices. Dos caras opuestas carecen de vértices
comunes y dos caras de ejes diferentes tienen dos vértices comunes.
Escribiendo \(\mathbf J=uu^{\mathsf T}\), resulta
\[
 MM^{\mathsf T}=4I+2(\mathbf J-I-R)
 =2(I-R)+2\mathbf J=4P_-+12P_u.
\]
Los tres proyectores \(P_u,P_-,P_0\) son ortogonales y suman \(I\);
sus rangos son, respectivamente, uno, tres y dos. La identidad
\(\|M^{\mathsf T}x\|^2=\langle x,MM^{\mathsf T}x\rangle\)
identifica el núcleo anunciado. El rango complementario es cuatro.
\end{proof}

Sobre \(Q\), la forma
\[
 B=P_--P_u
\]
tiene signatura \((3,1)\). La orientación temporal se fija por \(u\).
Con
\[
 t=\frac{u}{\sqrt6},\qquad
 d_i=\frac{e_{i,+}-e_{i,-}}{\sqrt2},\qquad
 W=\sqrt6P_u+\sqrt2P_-,
\]
se cumple \(We_{i,\epsilon}=t+\epsilon d_i\). Esos seis vectores
son nulos para \(B\), mientras que
\[
 3t+\nu_1d_1+\nu_2d_2+\nu_3d_3
\]
tiene norma cuadrática \(-6\). La misma incidencia satisface
\(MM^{\mathsf T}=2W^*W\) sobre \(Q\). Las rotaciones propias del cubo
conmutan con los proyectores, preservan ambas formas y fijan \(t\).
Estas identidades determinan la métrica de la publicación cúbica.

\section{Soldadura celular y transporte de memoria}
\label{x_reciprocidad:vii:sec:soldadura-celular}

Para cada arista orientada \(a\), la publicación de la ruta proporciona
la cara \(\lambda(\underline a)\), el signo
\(\sigma(a)\), las fibras de origen y destino y el transporte
\(U(a):Q_{s(a)}\longrightarrow Q_{t(a)}\).
La inversión verifica \(U(\bar a)=U(a)^{-1}\).
Las semiaristas de frontera conservan su inversión formal y su incidencia
en el borde. Con \(\ell_m=\ell_0\,9^{-m}>0\), la soldadura es
\begin{equation}
 \beta_m(a)=\ell_m\sigma_m(a)
              n_{\lambda(\underline a)}^{s(a)},\qquad
 n_{i,\epsilon}=t+\epsilon d_i.
 \label{x_reciprocidad:vii:eq:soldadura-arista}
\end{equation}
La cara y el signo se transportan junto con la memoria de la ruta.
La convención de inversión exige
\begin{equation}
 \beta_m(\bar a)=-U_m(a)\beta_m(a).
 \label{x_reciprocidad:vii:eq:inversion-soldadura}
\end{equation}
Aplicando dos veces esta relación se recupera \(\beta_m(a)\);
el registro de acarreo conserva, por separado, la composición ordenada.

\begin{proposition}[Naturalidad de la soldadura por refinamiento]
\label{x_reciprocidad:vii:prop:refinamiento-soldadura}
Sea \(a=a_1\cdots a_9\) un refinamiento compatible de una arista.
Supóngase que las nueve incidencias hijas, transportadas a la fibra del
padre, conservan la cara y el signo, y que
\(\ell_{m+1}=\ell_m/9\). Sea
\(\mathcal U_j:Q_{m,s(a)}\longrightarrow Q_{m+1,s(a_j)}\)
el transporte acumulado que incluye la identificación por refinamiento
y el trayecto hasta el origen del hijo \(j\). Entonces
\[
 \beta_m(a)=
 \sum_{j=1}^{9}\mathcal U_j^{-1}\beta_{m+1}(a_j).
\]
La igualdad es compatible con la inversión y con la concatenación
del registro de memoria.
\end{proposition}

\begin{proof}
Cada sumando, expresado en la fibra del origen del padre, es
\((\ell_m/9)\sigma_m(a)n_{\lambda(\underline a)}\).
La suma de nueve términos reproduce la soldadura del padre.
La inversión invierte el orden de los transportes y sustituye cada uno
por su inverso; la relación \eqref{x_reciprocidad:vii:eq:inversion-soldadura} proporciona
el signo global. La memoria se concatena en ese mismo orden, por lo
que el refinamiento conserva tanto la publicación vectorial como su
registro enriquecido.
\end{proof}

\section{Respuesta positiva, unicidad y extensión mínima}
\label{x_reciprocidad:vii:sec:respuesta-positiva}

La soldadura lineal de pantalla es
\(\bar\beta_m=\ell_mW|_Q\). Es invertible, con
\[
 \bar\beta_m^{-1}
 =\ell_m^{-1}\left(\frac1{\sqrt6}P_u+\frac1{\sqrt2}P_-\right).
\]
El autovalor \(5/9\) del sector antisimétrico del bloque normalizado
\(B_c/9\), construido en la sección~\ref{x_reciprocidad:vii:sec:bloque-compresion},
determina el coeficiente de
compresión \(q_{\mathrm{vis}}=5/9\). Reservamos esta notación para
distinguirlo del cociente euclídeo \(q_9(n)\) del núcleo aritmético.
Su realización sobre la pantalla utiliza una isometría
\(V_9:Q\longrightarrow\mathcal K\). La contribución visible
\(C_9=q_{\mathrm{vis}}V_9\) y su complemento
\(D_9=\sqrt{1-q_{\mathrm{vis}}^2}\,V_9\) satisfacen
\[
 C_9^*C_9+D_9^*D_9=I_Q,\qquad D_9^*D_9=\frac{56}{81}I_Q.
\]
El coeficiente \(56\) es \(9^2-5^2\); pertenece al defecto cuadrático
de ese lector de compresión.

\begin{theorem}[Respuesta energética determinada por la soldadura]
\label{x_reciprocidad:vii:thm:respuesta-unica}
Existe un único operador \(\Lambda_m:Q\longrightarrow Q\) tal que
\(\bar\beta_m^*\Lambda_m\bar\beta_m=D_9^*D_9\). Es positivo definido y
\begin{equation}
 \Lambda_m=\frac{28}{243\ell_m^2}(P_u+3P_-)
 =\frac{112}{81\ell_m^2}
       (MM^{\mathsf T}|_Q)^{-1}.
 \label{x_reciprocidad:vii:eq:lambda-pantalla}
\end{equation}
Sobre \(H_F\), la identidad correspondiente tiene como segundo miembro
\((56/81)P_\square\).
\end{theorem}

\begin{proof}
La congruencia por la soldadura invertible determina necesariamente
\[
 \Lambda_m=\bar\beta_m^{-*}\frac{56}{81}I_Q\bar\beta_m^{-1}.
\]
Sustituyendo la inversa, el coeficiente de \(P_u\) es
\(56/(486\ell_m^2)=28/(243\ell_m^2)\); el de \(P_-\) es
\(28/(81\ell_m^2)\). Ambos son positivos. Por la
proposición \ref{x_reciprocidad:vii:prop:incidencia},
\((MM^{\mathsf T}|_Q)^{-1}=P_u/12+P_-/4\),
lo que da la segunda expresión. La extensión a seis caras anula
\(P_0H_F\); por ello la identidad se extiende mediante \(P_\square\).
\end{proof}

\begin{proposition}[Caracterización por mínima extensión]
\label{x_reciprocidad:vii:prop:minima-extension}
Para todo \(b\in Q\),
\[
 \langle b,\Lambda_m b\rangle
 =\frac{112}{81\ell_m^2}
   \min_{Mx=b}\|x\|^2.
\]
La extensión minimizante es única:
\(x_{\min}=M^{\mathsf T}(MM^{\mathsf T}|_Q)^{-1}b\).
\end{proposition}

\begin{proof}
El vector anunciado satisface \(Mx_{\min}=b\) y pertenece a
\(\operatorname{im}M^{\mathsf T}=(\ker M)^\perp\).
Toda otra solución es \(x_{\min}+z\), con \(z\in\ker M\); por
ortogonalidad, su norma al cuadrado es \(\|x_{\min}\|^2+\|z\|^2\).
Además,
\[
 \|x_{\min}\|^2
 =\langle b,(MM^{\mathsf T}|_Q)^{-1}b\rangle.
\]
La fórmula de \(\Lambda_m\) concluye la prueba.
\end{proof}

El operador \(MM^{\mathsf T}|_Q\) realiza la respuesta
Dirichlet--a--Neumann de esta formulación de incidencia; su inversa
es la respuesta Neumann--a--Dirichlet. La energía anterior corresponde
a esta última, con la normalización de la soldadura y del defecto de
compresión. El complemento de Schur de
\(\left(\begin{smallmatrix}I&M^{\mathsf T}\\M&0\end{smallmatrix}\right)\)
es \(-MM^{\mathsf T}\), lo que ofrece la misma eliminación variacional.

\section{Conservación de la densidad por refinamiento}
\label{x_reciprocidad:vii:sec:densidad-refinamiento}

Consideremos el sector de transportes cúbicos, ortogonales para el producto
escalar energético y compatibles con la forma \(B\). Fijamos un
espacio de etiquetas común \(Q^{\mathrm{lab}}\), identificado
isométricamente con la pantalla del padre. Las soldaduras tienen
los tipos
\(\bar\beta_m:Q^{\mathrm{lab}}\longrightarrow Q_m\) y
\(\bar\beta_{m+1,j}:Q^{\mathrm{lab}}\longrightarrow Q_{m+1,j}\);
los operadores \(\Lambda\) actúan en sus respectivas fibras de llegada.
Si
\(U_j:Q_j\longrightarrow Q\) transporta la fibra hija a la del padre,
entonces
\[
 \bar\beta_{m+1,j}=\frac19U_j^{-1}\bar\beta_m,\qquad
 \Lambda_{m+1,j}=81U_j^{-1}\Lambda_mU_j.
\]
El factor \(81\) procede de \(\ell_{m+1}^{-2}=81\ell_m^{-2}\).

\begin{theorem}[Invariancia de la densidad normalizada]
\label{x_reciprocidad:vii:thm:densidad}
Cada hijo conserva la respuesta cuadrática del padre:
\[
 \bar\beta_{m+1,j}^*\Lambda_{m+1,j}\bar\beta_{m+1,j}
 =\frac{56}{81}I_{Q^{\mathrm{lab}}}.
\]
Para \(r\) refinamientos, la respuesta conjunta, expresada en las fibras
del padre, es
\[
 E_{m+r}=\frac{56}{81}I_{Q^{\mathrm{lab}}}\otimes I_{9^r}.
\]
La expectativa con traza normalizada
\(\tau_{9^r}(A)=9^{-r}\operatorname{Tr}(A)\) satisface
\[
 (\operatorname{id}\otimes\tau_{9^r})(E_{m+r})
 =\frac{56}{81}I_{Q^{\mathrm{lab}}}
\]
a toda profundidad finita.
\end{theorem}

\begin{proof}
La congruencia de un hijo contiene los factores \(1/9,81,1/9\);
su producto es uno. La ortogonalidad cancela los transportes
intermedios. La suma directa de nueve respuestas iguales identifica
el primer refinamiento con \((56/81)I_Q\otimes I_9\).
La iteración produce el tensor con \(I_{9^r}\).
Finalmente \(\tau_{9^r}(I_{9^r})=1\). Las expectativas se componen
por la identidad
\(\tau_{9^{r+s}}=\tau_{9^r}\otimes\tau_{9^s}\).
\end{proof}

Las inclusiones unitales \(A\mapsto A\otimes I_9\) son compatibles con
esa familia de trazas. En el límite inductivo de las álgebras matriciales,
la respuesta es el elemento central
\((56/81)I_Q\otimes\mathbf1\). La suma extensiva sin normalizar produce,
en cambio, \(9^r(56/81)I_Q\). Esta distinción especifica qué magnitud
se conserva: la densidad cuadrática de pantalla.

\section{Realización variacional de la respuesta}
\label{x_reciprocidad:vii:sec:realizacion-respuesta}

La realización material usa una isometría
\(J_{\mathrm{scr}}:Q\longrightarrow\mathcal H_\partial\)
compatible con la acción de espín y con el refinamiento.
La cotetrada de frontera es \(e=J_{\mathrm{scr}}\bar\beta_m\).
Sea \(\Lambda_{\mathrm{CE}}\) la respuesta de frontera obtenida
por segunda variación de la acción material y geométrica, con sus
condiciones de contorno y la eliminación de las variables interiores.
El defecto de compatibilidad es
\begin{equation}
 \mathcal R_{\mathrm{grav}}
 =\bar\beta_m^*
   (J_{\mathrm{scr}}^*\Lambda_{\mathrm{CE}}J_{\mathrm{scr}}-\Lambda_m)
   \bar\beta_m
 =e^*\Lambda_{\mathrm{CE}}e-D_9^*D_9.
 \label{x_reciprocidad:vii:eq:defecto-respuesta}
\end{equation}
Por la invertibilidad de \(\bar\beta_m\),
\[
 \mathcal R_{\mathrm{grav}}=0
 \quad\Longleftrightarrow\quad
 J_{\mathrm{scr}}^*\Lambda_{\mathrm{CE}}J_{\mathrm{scr}}=\Lambda_m.
\]
La isometría \(J_{\mathrm{scr}}\) transporta los estados de frontera;
la corriente variacional de espín es una forma de grado tres con valores
en bivectores. Cada una interviene con su dominio propio en la
realización de la respuesta.

La construcción determina una forma energética positiva, su extensión
interior minimizante y la densidad conservada por refinamiento.
La amplitud de una contribución material concreta se obtiene al evaluar
la corriente de esa realización en su forma cuadrática variacional.
Así quedan diferenciados el operador que fija la respuesta, la corriente
sobre la que actúa y el observable gravitatorio que publica la evaluación.

% END OWNER


% BEGIN OWNER /Users/ruben/Documents/New project/output/EXTREMA_MEDIA_RAZON_RECIPROCIDAD_ES_EN_20260918/ARTICULO/ES/source/antecedentes_vii/30d_realizacion_geometrica.tex
% RESULTADO_RECUPERADO: S15, 11c_gravedad_holonomia_rev6.tex:128--212;
% S01, 02_dinamica_tres_hojas.tex:155--219; pantalla y respuesta de30.
% Pruebas de curvatura, covariancia y Bianchi reunidas para autonomía.
% La realización suave conserva sus condiciones de holonomía explícitas.

\chapter{Realización geométrica del transporte y de la soldadura}
\label{x_reciprocidad:vii:sec:realizacion-cartan}

La representación discreta del retorno y la soldadura de pantalla
determinan dos aspectos del transporte: la acumulación de memoria en
las rutas y la comparación de sus fibras. Su realización geométrica
conserva ambas operaciones. En esta sección se fijan el espacio de
llegada, la normalización dimensional de la conexión y las identidades
que utilizará su variación.

\section{Pantalla, cotétrada y conexión}

El cociente de pantalla \(Q_\square\), su forma lorentziana y la
soldadura \(\beta_m\) proceden de
\ref{x_reciprocidad:vii:sec:pantalla-respuesta}. En cada celda, una realización de
pantalla \(J_{{\rm scr},m}\) transporta la forma y la soldadura
a la fibra geométrica correspondiente. Su imagen celular es
\[
 e_m=J_{{\rm scr},m}\beta_m.
\]
La escala \(\ell_m\) permanece dentro de la soldadura. Los
cambios de marco actúan sobre \(J_{{\rm scr},m}\), la conexión
y el transporte de las celdas conjuntamente.

Una realización suave de esos datos sobre una variedad orientada
\(\mathcal U\) de dimensión cuatro consiste en una cotétrada
no degenerada \(e=(e^I)\), una conexión de Lorentz
\(\omega=(\omega^I{}_J)\) y una realización de rutas
\(\mathfrak R_{\rm C}\) que satisfacen
\begin{equation}
 \operatorname{Hol}_{\mathcal A_{\ell_0}}
       (\mathfrak R_{\rm C}(\gamma))
 =\rho_{\rm C}(\operatorname{Hol}_{\rm TPK}(\gamma)).
 \label{x_reciprocidad:vii:eq:compatibilidad-holonomia-geometrica}
\end{equation}
En esta igualdad matricial, la representación se normaliza mediante
\[
 N_{\ell_0}(R,a)=(\Lambda(R),a/\ell_0),\qquad
 \rho_{\rm C}=N_{\ell_0}\circ\rho_{\rm C}^{\rm disc},
\]
donde \(\Lambda:\operatorname{Spin}^+(3,1)\to SO^+(3,1)\)
es la representación vectorial. La aplicación \(N_{\ell_0}\)
preserva el producto semidirecto: dividir la traslación
\(a+\Lambda(R)b\) por \(\ell_0\) equivale a componer las
dos traslaciones normalizadas. Por ello la vuelta completa publica
\(72u\) en la conexión normalizada y \(72\ell_0u\) en la
carta dimensional. El levantamiento de espín conserva, por separado,
\(R_4^{c_g(\gamma)}\); la representación vectorial tiene el
núcleo central \(\{1,-1\}\), de modo que se retiene ese
levantamiento al utilizar campos espinoriales.
Identidad, concatenación, inversión y subdivisión se mantienen en la
realización correspondiente. El emparejamiento de
pantalla y la condición de holonomía especifican los datos de la
realización utilizada en las ecuaciones de campo.

Con \(\eta=\operatorname{diag}(-1,1,1,1)\), se define
\(g=\eta_{IJ}e^I\otimes e^J\). Para una escala elemental
positiva \(\ell_0\), constante en la carta, la conexión afín es
\begin{equation}
 \mathcal A_{\ell_0}
 =\begin{pmatrix}\omega&\ell_0^{-1}e\\0&0\end{pmatrix}.
 \label{x_reciprocidad:vii:eq:conexion-afin-realizada}
\end{equation}
La cotétrada tiene dimensión de longitud; el factor
\(\ell_0^{-1}\) homogeneiza su componente traslacional.

\begin{proposition}[Curvatura y torsión de la conexión realizada]
\label{x_reciprocidad:vii:prop:curvatura-afin}
La curvatura de \eqref{x_reciprocidad:vii:eq:conexion-afin-realizada} es
\begin{equation}
 \mathcal F_{\ell_0}
 =\mathrm d\mathcal A_{\ell_0}
   +\mathcal A_{\ell_0}\wedge\mathcal A_{\ell_0}
 =\begin{pmatrix}F_\omega&\ell_0^{-1}T\\0&0\end{pmatrix},
 \quad F_\omega=\mathrm d\omega+\omega\wedge\omega,
 \quad T=\mathrm de+\omega\wedge e.
 \label{x_reciprocidad:vii:eq:curvatura-afin-realizada}
\end{equation}
\end{proposition}
\begin{proof}
El producto de matrices de formas produce en el bloque diagonal
\(\omega\wedge\omega\) y en el traslacional
\(\ell_0^{-1}\omega\wedge e\). La derivada exterior aporta
\(\mathrm d\omega\) y \(\ell_0^{-1}\mathrm de\), pues
\(\ell_0\) es constante. La suma demuestra la igualdad.
\end{proof}

La curvatura describe el defecto de transporte de marcos y la torsión
el de la cotétrada. Ambas pertenecen a la misma conexión realizada;
sus componentes conservan dominios tensoriales distintos.

\begin{proposition}[Covariancia e identidades de Bianchi]
\label{x_reciprocidad:vii:prop:bianchi-cartan}
Bajo un cambio de marco \(g_L:\mathcal U\to SO^+(3,1)\),
\[
 e'=g_Le,\qquad
 \omega'=g_L\omega g_L^{-1}-\mathrm dg_L\,g_L^{-1},
\]
se cumple \(T'=g_LT\) y \(F_{\omega'}=g_LF_\omega g_L^{-1}\).
Además,
\begin{equation}
 D_\omega T=F_\omega\wedge e,\qquad D_\omega F_\omega=0.
 \label{x_reciprocidad:vii:eq:bianchi-realizacion}
\end{equation}
\end{proposition}
\begin{proof}
En \(\mathrm d(g_Le)+\omega'\wedge g_Le\), los términos
que contienen \(\mathrm dg_L\wedge e\) se cancelan.
La misma sustitución en la curvatura da su transformación por
conjugación. Para la primera identidad de Bianchi, se expande
\[
 D_\omega(\mathrm de+\omega\wedge e)
 =(\mathrm d\omega+\omega\wedge\omega)\wedge e.
\]
En \(\mathrm dF_\omega+\omega\wedge F_\omega-F_\omega\wedge\omega\),
los términos que contienen \(\mathrm d\omega\) y los productos
cúbicos se cancelan por pares. Esto prueba la segunda identidad.
\end{proof}

\section{Las tres realizaciones cuadráticas de curvatura}

La hoja \(\tau\in\{+1,0,-1\}\) del TRIT se representa en
el álgebra de Cartan mediante generadores \(J_{ab}=-J_{ba}\)
y \(P_a^{(\tau)}\), con relaciones
\begin{align*}
 [J_{ab},J_{cd}]&=
 \eta_{bc}J_{ad}-\eta_{ac}J_{bd}
 -\eta_{bd}J_{ac}+\eta_{ad}J_{bc},\\
 [J_{ab},P_c^{(\tau)}]&=
 \eta_{bc}P_a^{(\tau)}-\eta_{ac}P_b^{(\tau)},\\
 [P_a^{(\tau)},P_b^{(\tau)}]&=-\tau J_{ab}.
\end{align*}
Para \(\ell>0\) constante, se escribe
\[
 \mathcal A_\tau^{\rm Cart}
 =\tfrac12\omega^{ab}J_{ab}+\ell^{-1}e^aP_a^{(\tau)}.
\]

\begin{proposition}[Curvatura de la realización tritificada]
\label{x_reciprocidad:vii:prop:curvatura-tritificada}
Con las relaciones anteriores,
\begin{equation}
 \mathcal F_\tau^{\rm Cart}
 =\tfrac12(F_\omega^{ab}-\tau\ell^{-2}e^a\wedge e^b)J_{ab}
     +\ell^{-1}T^aP_a^{(\tau)}.
 \label{x_reciprocidad:vii:eq:curvatura-tritificada}
\end{equation}
\end{proposition}
\begin{proof}
Se expande \(\mathrm d\mathcal A+\tfrac12[\mathcal A,\mathcal A]\).
El corchete del sector de Lorentz produce \(F_\omega\), el
mixto produce \(D_\omega e\), y el traslacional aporta
\(-\tfrac12\tau\ell^{-2}e^a\wedge e^bJ_{ab}\).
\end{proof}

En la hoja central se recupera la conexión afín; las hojas extremas
conservan los dos signos del término de curvatura constante. La
representación mantiene así el régimen trítico de la ruta antes de
aplicar las ecuaciones variacionales. La igualdad de holonomía
\eqref{x_reciprocidad:vii:eq:compatibilidad-holonomia-geometrica} corresponde a la
hoja central. En cada hoja extrema se utiliza la representación
\(\rho_{{\rm C},\tau}\) en su grupo de Cartan y su condición
de holonomía respectiva. La corriente y el tensor métrico utilizan
los campos y la representación de la misma realización.

% END OWNER


% BEGIN OWNER /Users/ruben/Documents/New project/output/EXTREMA_MEDIA_RAZON_RECIPROCIDAD_ES_EN_20260918/ARTICULO/ES/source/antecedentes_vii/30b_variacion_energia_memoria.tex
% Artículo VII. Desarrollo aditivo, revisión 04, 10-09-2026.
% ARQUITECTURA_AUTORAL_PREEXISTENTE: transporte enriquecido, energía
% D_9^* W D_9 y compatibilidad por refinamiento, propietario S01,
% 02_dinamica_tres_hojas.tex, líneas 221--250.
% FORMALIZACION_NUEVA / CERTIFICADO_NUEVO: diferencial explícito de ese
% funcional, representante antihermítico, composición e inversión ponderada.
% Las variaciones pertenecen al espacio de realizaciones unitarias del
% transporte. No se presume que toda variación continua sea una ruta TPK.
% Su realización como corriente material de espín se trata por separado.
% \mathsf W_m es el peso energético; W de la sección anterior es soldadura.

\chapter{Respuesta variacional de la energía de memoria}
\label{x_reciprocidad:vii:sec:variacion-memoria}

El transporte de una ruta identifica las fibras de sus extremos y conserva
la orientación y la memoria que intervienen en esa identificación. La energía
asociada compara el campo final con el campo inicial transportado. Su
diferencial cuantifica cómo responde esa comparación a una variación del
transporte. Esta sección obtiene la respuesta directamente del funcional
discreto y desarrolla su composición, covariancia e inversión.

\section{Funcional discreto y espacio de variaciones}
\label{x_reciprocidad:vii:subsec:funcional-memoria}

En un nivel finito \(m\), sean \(\mathcal V_m\) los vértices de la
publicación de rutas y \(\mathcal E_m^+\) una orientación de cada arista.
Cada vértice lleva una fibra hermítica finito-dimensional \(H_v\).
Definimos
\[
 \mathcal H_m=\bigoplus_{v\in\mathcal V_m}H_v,\qquad
 \mathcal K_m=\bigoplus_{a\in\mathcal E_m^+}H_{t(a)}.
\]
El producto hermítico es antilineal en la primera variable.
Para \(U_a:H_{s(a)}\to H_{t(a)}\) unitario, el operador de diferencias es
\begin{equation}
 (D_mf)_a=f_{t(a)}-U_af_{s(a)},\qquad
 E_m(f,U)=\langle D_mf,\mathsf W_mD_mf\rangle,\qquad
 \mathsf W_m=\mathsf W_m^*>0.
 \label{x_reciprocidad:vii:eq:energia-memoria-discreta}
\end{equation}
Es la energía del término \(D_m^*\mathsf W_mD_m\) del generador discreto.
El peso \(\mathsf W_m\) puede acoplar distintas aristas. El operador \(W\)
de la soldadura de la sección~\ref{x_reciprocidad:vii:sec:pantalla-respuesta} y este peso
energético tienen dominios y funciones diferentes.

El transporte \(U\) es una publicación del estado enriquecido. El campo
\(f\) y el peso \(\mathsf W_m\) permanecen explícitos como argumentos
del funcional; cada aplicación material conserva las reglas que los
determinan. Se consideran familias diferenciables de las realizaciones
unitarias del transporte,
\[
 \dot U_a=A_aU_a,\qquad A_a^*=-A_a.
\]
Las variaciones físicamente utilizadas forman el subespacio que conserve
las condiciones de la realización correspondiente. El diferencial se
calcula primero en el espacio tangente unitario y después se restringe
a ese subespacio.

\begin{proposition}[Covector de respuesta al transporte]
\label{x_reciprocidad:vii:prop:covector-memoria}
Escribamos
\[
 v_a=U_af_{s(a)},\qquad r=D_mf,\qquad q=\mathsf W_mr.
\]
A campo y peso fijos, el diferencial de la energía es
\begin{equation}
 \mathfrak j_m(A)=-2\operatorname{Re}
       \sum_{a\in\mathcal E_m^+}\langle q_a,A_av_a\rangle .
 \label{x_reciprocidad:vii:eq:covector-memoria}
\end{equation}
Su representante para el producto real de Hilbert--Schmidt es
\begin{equation}
 \mathcal J_a=v_aq_a^*-q_av_a^*,\qquad
 \mathcal J_a^*=-\mathcal J_a,\qquad
 \mathfrak j_m(A)=
 \sum_a\operatorname{Re}\operatorname{Tr}(\mathcal J_a^*A_a).
 \label{x_reciprocidad:vii:eq:representante-corriente-discreta}
\end{equation}
Para variaciones simultáneas de campo, transporte y peso se tiene
\begin{equation}
 \dot E_m=
 2\operatorname{Re}\sum_a
 \left\langle q_a,\dot f_{t(a)}-U_a\dot f_{s(a)}-A_av_a\right\rangle
 +\langle r,\dot{\mathsf W}_mr\rangle .
 \label{x_reciprocidad:vii:eq:variacion-completa-memoria}
\end{equation}
\end{proposition}

\begin{proof}
La actualización de primer orden, escrita en el álgebra de números
duales con \(\epsilon^2=0\), da
\[
 r+\epsilon\dot r,\qquad
 \dot r_a=\dot f_{t(a)}-U_a\dot f_{s(a)}-A_av_a.
\]
El coeficiente de \(\epsilon\) en la energía es
\[
 \langle\dot r,\mathsf W_mr\rangle+
 \langle r,\mathsf W_m\dot r\rangle+
 \langle r,\dot{\mathsf W}_mr\rangle.
\]
La autoadjunción del peso produce
\eqref{x_reciprocidad:vii:eq:variacion-completa-memoria}. Al fijar campo y peso queda
\eqref{x_reciprocidad:vii:eq:covector-memoria}. Además,
\[
 \operatorname{Tr}(\mathcal J_a^*A_a)
 =v_a^*A_aq_a-q_a^*A_av_a.
\]
La antihermiticidad de \(A_a\) implica
\(v_a^*A_aq_a=-\overline{q_a^*A_av_a}\), de donde se obtiene
el representante anunciado.
\end{proof}

Así queda evaluada la respuesta: en cada arista se transporta el campo,
se calcula su diferencia con el campo final y se aplica el peso energético.
El producto antisimetrizado de esos dos vectores determina
\(\mathcal J_a\in\mathfrak u(H_{t(a)})\). Este operador representa,
mediante el producto real de Hilbert--Schmidt, el covector de respuesta
sobre las variaciones unitarias de la fibra de destino.

\section{Composición y transporte de la respuesta}
\label{x_reciprocidad:vii:subsec:composicion-respuesta}

\begin{proposition}[Regla de composición sobre una ruta]
\label{x_reciprocidad:vii:prop:composicion-respuesta}
Para una ruta formada por \(N\) transportes, sea
\(U_\gamma=U_N\cdots U_1\) y
\(B_k=U_N\cdots U_{k+1}\), con \(B_N=I\).
Si \(\dot U_k=A_kU_k\), entonces
\begin{equation}
 \dot U_\gamma U_\gamma^{-1}
 =\sum_{k=1}^N B_kA_kB_k^{-1}.
 \label{x_reciprocidad:vii:eq:variacion-transporte-compuesto}
\end{equation}
Por tanto, un representante \(\mathcal J_\gamma\) en la fibra final
induce en el factor \(k\) la contribución
\(B_k^*\mathcal J_\gamma B_k\). Cuando un factor pertenece a varias
rutas del funcional, sus contribuciones se suman con las incidencias
que éste utiliza.
\end{proposition}

\begin{proof}
La regla del producto da
\[
 \dot U_\gamma=\sum_k
 U_N\cdots U_{k+1}A_kU_k\cdots U_1.
\]
Multiplicar por \(U_\gamma^{-1}\) elimina los factores a la derecha
de \(A_k\) y produce \eqref{x_reciprocidad:vii:eq:variacion-transporte-compuesto}.
Por unitariedad y ciclicidad de la traza,
\[
 \operatorname{Re}\operatorname{Tr}
   (\mathcal J_\gamma^*B_kA_kB_k^{-1})
 =\operatorname{Re}\operatorname{Tr}
   ((B_k^*\mathcal J_\gamma B_k)^*A_k).
\]
La linealidad del diferencial demuestra la suma de contribuciones.
\end{proof}

Esta fórmula conserva la posición de cada operación dentro de la ruta.
El transporte posterior a una incidencia determina cómo se representa
su contribución en la fibra final; la respuesta local se recupera
transportándola a la fibra de esa incidencia.

\begin{proposition}[Covariancia bajo cambios unitarios de marco]
\label{x_reciprocidad:vii:prop:covariancia-corriente-memoria}
Sean \(g_v\) cambios unitarios de marco, independientes del parámetro
de variación, y \(G_{\mathcal E}=\bigoplus_a g_{t(a)}\).
Bajo las transformaciones
\[
 f'_v=g_vf_v,\quad U'_a=g_{t(a)}U_ag_{s(a)}^{-1},\quad
 \mathsf W'_m=G_{\mathcal E}\mathsf W_mG_{\mathcal E}^{-1},\quad
 A'_a=g_{t(a)}A_ag_{t(a)}^{-1},
\]
se tiene
\begin{equation}
 \mathcal J'_a=g_{t(a)}\mathcal J_ag_{t(a)}^{-1},
 \qquad \mathfrak j'_m(A')=\mathfrak j_m(A).
 \label{x_reciprocidad:vii:eq:covariancia-respuesta}
\end{equation}
\end{proposition}

\begin{proof}
Se obtiene \(r'=G_{\mathcal E}r\), \(q'=G_{\mathcal E}q\) y
\(v'_a=g_{t(a)}v_a\). La sustitución en
\eqref{x_reciprocidad:vii:eq:representante-corriente-discreta} demuestra la primera
identidad. La invariancia del producto hermítico, o la ciclicidad de
la traza, demuestra la segunda.
\end{proof}

Cuando los cambios de marco dependan del parámetro, sus derivadas
actúan también sobre el campo y el peso. La expresión covariante
correspondiente es entonces el diferencial completo
\eqref{x_reciprocidad:vii:eq:variacion-completa-memoria}.

\section{Inversión orientada y transporte del peso}
\label{x_reciprocidad:vii:subsec:inversion-respuesta}

La inversión puede verse explícitamente en el caso de una energía
local por arista, \(\mathsf W_m=\bigoplus_a W_a\).
Para \(a:s\to t\), su inversa lleva
\[
 U_{\bar a}=U_a^{-1},\qquad
 r_{\bar a}=-U_a^{-1}r_a,\qquad
 W_{\bar a}=U_a^{-1}W_aU_a.
\]
Estas reglas conservan exactamente la energía de la arista.

\begin{proposition}[Compatibilidad diferencial con la inversión]
\label{x_reciprocidad:vii:prop:inversion-respuesta}
Si \(W_a\) y los campos de ambos extremos permanecen fijos, entonces
\[
 A_{\bar a}=-U_a^{-1}A_aU_a,\qquad
 \dot W_{\bar a}=U_a^{-1}[W_a,A_a]U_a.
\]
La variación completa de la energía de la arista invertida coincide con
la original. En particular, si \([W_a,A_a]=0\), el covector parcial
de transporte satisface
\begin{equation}
 \mathfrak j_{\bar a}(U_a^{-1}A_aU_a)
 =-\mathfrak j_a(A_a).
 \label{x_reciprocidad:vii:eq:imparidad-transporte}
\end{equation}
\end{proposition}

\begin{proof}
La derivada de \(U_a^{-1}\) es \(-U_a^{-1}A_a\), de donde resultan
las dos fórmulas. Escribiendo \(v_a=U_af_s\) y \(f_t=r_a+v_a\),
la parte debida exclusivamente al transporte invertido vale
\[
 \mathfrak j_{\bar a}(A_{\bar a})
 =-2\operatorname{Re}\langle r_a,W_aA_af_t\rangle
 =\mathfrak j_a(A_a)
   -2\operatorname{Re}\langle r_a,W_aA_ar_a\rangle.
\]
La variación del peso aporta
\[
 \langle r_{\bar a},\dot W_{\bar a}r_{\bar a}\rangle
 =\langle r_a,[W_a,A_a]r_a\rangle
 =2\operatorname{Re}\langle r_a,W_aA_ar_a\rangle.
\]
Ambos términos se compensan. Si el conmutador es nulo, esa contribución
desaparece; la linealidad en \(A_a\) da
\eqref{x_reciprocidad:vii:eq:imparidad-transporte}.
\end{proof}

La energía se suma sobre una orientación por arista. Una suma sobre ambas
orientaciones lleva el factor \(1/2\). Esta convención conserva el mismo
funcional y, por tanto, la misma respuesta.

\section{Compatibilidad con el refinamiento}
\label{x_reciprocidad:vii:subsec:refinamiento-respuesta}

\begin{theorem}[Transporte del diferencial por refinamiento]
\label{x_reciprocidad:vii:thm:refinamiento-respuesta}
Sean \(J_m:\mathcal H_m\to\mathcal H_{m+1}\) y
\(K_m:\mathcal K_m\to\mathcal K_{m+1}\) aplicaciones fijas.
Para una familia \(C^1\), supongamos las identidades de refinamiento
\begin{equation}
 D_{m+1}(t)J_m=K_mD_m(t),\qquad
 K_m^*\mathsf W_{m+1}(t)K_m=\mathsf W_m(t).
 \label{x_reciprocidad:vii:eq:familia-refinamiento-memoria}
\end{equation}
Entonces \(E_{m+1}(J_mf)=E_m(f)\).
Los covectores de transporte coinciden sobre las variaciones enlazadas
por la primera identidad. Cuando varían los pesos, sus contribuciones
al diferencial también coinciden.
\end{theorem}

\begin{proof}
Por sustitución,
\[
 \langle D_{m+1}J_mf,\mathsf W_{m+1}D_{m+1}J_mf\rangle
 =\langle D_mf,K_m^*\mathsf W_{m+1}K_mD_mf\rangle=E_m(f).
\]
Diferenciando la primera identidad de
\eqref{x_reciprocidad:vii:eq:familia-refinamiento-memoria}, se obtiene
\(\dot D_{m+1}J_m=K_m\dot D_m\). Por consiguiente,
\[
 2\operatorname{Re}
 \langle D_{m+1}J_mf,\mathsf W_{m+1}\dot D_{m+1}J_mf\rangle
 =2\operatorname{Re}\langle D_mf,\mathsf W_m\dot D_mf\rangle.
\]
La derivada de la segunda identidad es
\(K_m^*\dot{\mathsf W}_{m+1}K_m=\dot{\mathsf W}_m\), que demuestra
la coincidencia de los términos de variación del peso.
\end{proof}

El enunciado utiliza la compatibilidad de la familia de transportes.
Si \(J_m\) depende del parámetro, la regla de la cadena conserva
también el término
\[
 2\operatorname{Re}
 \langle D_{m+1}J_mf,\mathsf W_{m+1}D_{m+1}\dot J_mf\rangle.
\]
La naturalidad obtenida transporta conjuntamente el funcional, sus
incidencias y su diferencial.

\section{Realización material de la respuesta}
\label{x_reciprocidad:vii:subsec:realizacion-respuesta-discreta}

El resultado de esta sección es el covector explícito
\(\mathfrak j_m\), con su representante \(\mathcal J_a\).
La realización material compone ese covector con la variación del
transporte inducida por la conexión y con la normalización de la acción.
Si \(\mathcal L_m\) denota ese mapa diferencial, el covector sobre
variaciones de conexión es \(\mathcal L_m^*\mathfrak j_m\).
La regla de composición
\eqref{x_reciprocidad:vii:eq:variacion-transporte-compuesto} evalúa su parte
correspondiente a productos finitos de transportes.

En la sección~\ref{x_reciprocidad:vii:sec:corriente-respuesta-cartan}, la corriente
\(\sigma_{IJ}\) es una tres-forma definida por la variación de la acción
material. La identificación entre ambas realizaciones conserva el campo,
la representación, el peso, la normalización dimensional y el paso al
límite utilizados por esa acción. Estos datos permiten evaluar la
contracción que interviene en la densidad torsional. El desarrollo
presente aporta el diferencial discreto que se compone en ese cálculo.

% END OWNER


% BEGIN OWNER /Users/ruben/Documents/New project/output/EXTREMA_MEDIA_RAZON_RECIPROCIDAD_ES_EN_20260918/ARTICULO/ES/source/antecedentes_vii/30c_composicion_corriente_conexion.tex
% Artículo VII, incorporación aditiva, revisión 05.
% Arquitectura autoral: energía de memoria APP--TRIT--TPK, S01.
% Antecedente recuperado: diferencial y refinamiento de 30b.
% Formalización reunida: nota COMPOSICION_VARIACIONAL_TRIT_CONEXION.md.
% Se conserva íntegra 30b; su representante antihermítico es la
% restricción unitaria del diferencial que aquí se desarrolla.
% Esta sección evalúa la corriente de una acción y una realización
% declaradas. No selecciona retrospectivamente su campo ni su peso.

\chapter{Realización de la corriente en coordenadas de conexión}
\label{x_reciprocidad:vii:sec:composicion-corriente-conexion}

La energía de memoria asigna a cada ruta una respuesta que conserva sus
incidencias y el orden de los transportes. Para insertarla en una acción
geométrica se necesitan sus componentes respecto de las variaciones de
conexión. El paso se obtiene diferenciando el transporte que ya realiza
la ruta: la corriente de conexión es la composición de ese diferencial
con el covector energético de la sección anterior.

Los regímenes elíptico, parabólico e hiperbólico del TRIT requieren
conservar el diferencial completo antes de restringirlo a una
representación particular. Las rotaciones unitarias constituyen un
dominio de esa construcción. La ampliación siguiente conserva también
las direcciones de dilatación y transporte que admite la realización
geométrica utilizada.

\section{Diferencial sobre transportes invertibles}

Se mantienen las fibras, los campos y el peso de
\eqref{x_reciprocidad:vii:eq:energia-memoria-discreta}. Consideremos ahora transportes
invertibles y familias diferenciables con
\(\dot U_a=A_aU_a\), donde \(A_a\) pertenece al espacio tangente
de la representación elegida. Escribimos
\[
 v_a=U_af_{s(a)},\qquad r_a=f_{t(a)}-v_a,\qquad
 q=\mathsf W_mr,\qquad E_m=\langle r,\mathsf W_mr\rangle.
\]

\begin{proposition}[Representante completo del diferencial]
\label{x_reciprocidad:vii:prop:diferencial-invertible}
A campo y peso fijos, se cumple
\begin{equation}
 \mathrm d_UE_m(A)=-2\operatorname{Re}\sum_a
       \langle q_a,A_av_a\rangle
 =\sum_a\operatorname{Re}\operatorname{Tr}(\mathcal G_a^*A_a),
 \qquad \mathcal G_a=-2q_av_a^*.
 \label{x_reciprocidad:vii:eq:gradiente-transporte-completo}
\end{equation}
Las proyecciones antihermítica y hermítica de ese representante son
\begin{equation}
 \frac{\mathcal G_a-\mathcal G_a^*}{2}
   =v_aq_a^*-q_av_a^*=\mathcal J_a,
 \qquad
 \frac{\mathcal G_a+\mathcal G_a^*}{2}
   =-(q_av_a^*+v_aq_a^*).
 \label{x_reciprocidad:vii:eq:proyecciones-gradiente-memoria}
\end{equation}
La restricción a generadores antihermíticos recupera exactamente
\eqref{x_reciprocidad:vii:eq:representante-corriente-discreta}.
\end{proposition}

\begin{proof}
La identidad \(\dot r_a=-A_av_a\) y la autoadjunción del peso dan
\(\dot E_m=2\operatorname{Re}\langle q,\dot r\rangle\).
Además,
\[
 \operatorname{Tr}(\mathcal G_a^*A_a)
 =-2\operatorname{Tr}(v_aq_a^*A_a)=-2q_a^*A_av_a.
\]
Las dos proyecciones se obtienen tomando el adjunto. Para el producto
real de Hilbert--Schmidt, los subespacios hermítico y antihermítico son
ortogonales; por ello la parte hermítica tiene contracción nula con
un generador antihermítico.
\end{proof}

La fórmula permite pesos que acoplan aristas: su acción se calcula en
\(q=\mathsf W_mr\) antes de separar las contribuciones de cada
incidencia. Si varían campo o peso, se conserva la expresión completa
\eqref{x_reciprocidad:vii:eq:variacion-completa-memoria}, válida también para los
transportes invertibles aquí considerados.

\section{Representaciones cuadráticas del TRIT}

En una realización matricial bidimensional de
\(\mathbb A_\tau\), pueden tomarse
\[
 J_{+}=\begin{pmatrix}0&1\\-1&0\end{pmatrix},\qquad
 J_{-}=\begin{pmatrix}0&1\\1&0\end{pmatrix},\qquad
 J_{0}=\begin{pmatrix}0&1\\0&0\end{pmatrix}.
\]
Por multiplicación directa,
\[
 J_+^2=-I,\qquad J_-^2=I,\qquad J_0^2=0.
\]
Sus exponenciales realizan, respectivamente, los regímenes elíptico,
hiperbólico y parabólico. El primero es ortogonal para el producto
positivo fijado. El segundo preserva la forma
\(B=\operatorname{diag}(-1,1)\), pues
\(J_-^{\mathsf T}B+BJ_-=0\). El tercero representa el álgebra
de números duales. Estas representaciones fijan ejemplos de los tipos
cuadráticos; su uso en una pantalla de otra dimensión conserva el mapa
de representación correspondiente.

Un cálculo distingue las componentes del diferencial. En una arista,
tómense \(U=I\), \(f_s=(1,1)^{\mathsf T}\),
\(f_t=2f_s\) y \(\mathsf W=I\). Entonces \(q=v=f_s\),
\(\mathcal J=0\), mientras
\begin{equation}
 \mathrm dE(J_-)=-4,\qquad \mathrm dE(J_0)=-2.
 \label{x_reciprocidad:vii:eq:control-regimenes-corriente}
\end{equation}
El representante completo conserva ambas respuestas. La distinción
identifica el espacio de variaciones antes de efectuar la contracción
material.

\section{Composición de rutas e inversión orientada}

\begin{proposition}[Transporte contragrediente del diferencial]
\label{x_reciprocidad:vii:prop:pullback-invertible}
Para \(U_\gamma=U_N\cdots U_1\), sean
\(B_k=U_N\cdots U_{k+1}\) y \(B_N=I\). Se cumple
\[
 \dot U_\gamma U_\gamma^{-1}
   =\sum_kB_kA_kB_k^{-1}.
\]
El representante \(\mathcal G_\gamma\) en la fibra final induce
en la incidencia \(k\)
\begin{equation}
 \mathcal G_k=B_k^*\mathcal G_\gamma(B_k^{-1})^*.
 \label{x_reciprocidad:vii:eq:transporte-contragrediente}
\end{equation}
\end{proposition}

\begin{proof}
La regla del producto, seguida de multiplicación por
\(U_\gamma^{-1}\), demuestra la primera igualdad. La ciclicidad
de la traza da
\[
 \operatorname{Re}\operatorname{Tr}
  (\mathcal G_\gamma^*B_kA_kB_k^{-1})
 =\operatorname{Re}\operatorname{Tr}
  ((B_k^*\mathcal G_\gamma(B_k^{-1})^*)^*A_k).
\]
Para \(B_k\) unitario se recupera la conjugación de
\ref{x_reciprocidad:vii:prop:composicion-respuesta}.
\end{proof}

\begin{proposition}[Inversión por congruencia energética]
\label{x_reciprocidad:vii:prop:inversion-congruencia}
Para una energía local por arista, la inversión se representa por
\begin{equation}
 U_{\bar a}=U_a^{-1},\qquad
 r_{\bar a}=-U_a^{-1}r_a,\qquad
 W_{\bar a}=U_a^*W_aU_a.
 \label{x_reciprocidad:vii:eq:inversion-peso-general}
\end{equation}
Conserva la energía. Si \(W_a\) permanece fijo,
\[
 A_{\bar a}=-U_a^{-1}A_aU_a,\qquad
 \dot W_{\bar a}=U_a^*(A_a^*W_a+W_aA_a)U_a.
\]
La variación completa de la energía invertida coincide con la original.
\end{proposition}

\begin{proof}
La sustitución en \(r_{\bar a}^*W_{\bar a}r_{\bar a}\)
produce \(r_a^*W_ar_a\). La derivada de la inversa es
\(\dot U_a^{-1}=-U_a^{-1}A_a\); la del peso se obtiene por
la regla del producto. La igualdad de energías vale para toda la
familia, de modo que sus derivadas coinciden, incluyendo el término
\(\langle r_{\bar a},\dot W_{\bar a}r_{\bar a}\rangle\).
\end{proof}

La congruencia conserva positividad para todo transporte invertible.
En la representación unitaria se reduce a la regla de inversión de
\ref{x_reciprocidad:vii:prop:inversion-respuesta}. La suma utiliza una orientación
por arista, o la mitad de la suma sobre ambas orientaciones.

\section{Evaluación de las componentes de corriente}

Sea \(\theta=(\theta_b)\) una variación real de las coordenadas
de conexión de la realización utilizada. El diferencial efectivo del
transporte define los operadores
\begin{equation}
 A_a(\theta)=\dot U_aU_a^{-1}
     =\sum_bB_{ab}\theta_b.
 \label{x_reciprocidad:vii:eq:diferencial-conexion-transporte}
\end{equation}
Para una composición finita, sus contribuciones se obtienen por
\eqref{x_reciprocidad:vii:eq:transporte-contragrediente}. Los operadores
\(B_{ab}\) son las derivadas de la representación de conexión que
realiza la ruta.

\begin{theorem}[Corriente de conexión de la acción de memoria]
\label{x_reciprocidad:vii:thm:corriente-conexion-evaluada}
Para un término material \(S_{\rm mem}=\mathfrak a_m E_m\),
con \(\mathfrak a_m\) su normalización en la recta de acción,
a campo, peso y normalización fijos se tiene
\begin{equation}
 \delta S_{\rm mem}=-\frac12\sum_b\theta_b s_b,\qquad
 s_b=4\mathfrak a_m\operatorname{Re}
        \sum_a\langle q_a,B_{ab}v_a\rangle.
 \label{x_reciprocidad:vii:eq:componentes-corriente-conexion}
\end{equation}
Si el emparejamiento celular entre las bases de conexión y corriente
es una matriz invertible \(M\), definida por
\(\delta S_{\rm mem}=-\tfrac12\theta^{\mathsf T}M\sigma\),
las coordenadas de la corriente son
\begin{equation}
 \sigma=M^{-1}s.
 \label{x_reciprocidad:vii:eq:coordenadas-corriente-material}
\end{equation}
\end{theorem}

\begin{proof}
Se sustituye \eqref{x_reciprocidad:vii:eq:diferencial-conexion-transporte} en
\eqref{x_reciprocidad:vii:eq:gradiente-transporte-completo}, se multiplica por
\(\mathfrak a_m\) y se agrupan los coeficientes de
\(\theta_b\). La convención variacional \(-1/2\) fija el
factor \(4\). Comparar ambas expresiones para toda variación
\(\theta\) da \(M\sigma=s\); la invertibilidad del
emparejamiento demuestra existencia y unicidad de esas coordenadas.
\end{proof}

La fórmula realiza explícitamente la composición mencionada en
\ref{x_reciprocidad:vii:subsec:realizacion-respuesta-discreta}. Para evaluar cada
componente se transporta el campo, se calcula el residuo, se aplica el
peso y se contrae con la variación de conexión de la misma realización.
El emparejamiento \(M\) conserva orientación y volumen de las celdas.
Un emparejamiento diagonal da la correspondiente división por sus
volúmenes; otras bases conservan la matriz completa.

Si la conexión interviene también en \(f\), \(\mathsf W_m\)
o \(\mathfrak a_m\), se aplica
\eqref{x_reciprocidad:vii:eq:variacion-completa-memoria} y se añade
\(E_m\delta\mathfrak a_m\). Las contribuciones de todos los
términos materiales se suman antes de componer la respuesta de Cartan.

\begin{corollary}[Compatibilidad de la corriente con el refinamiento]
\label{x_reciprocidad:vii:cor:refinamiento-corriente-conexion}
Para las familias compatibles de
\ref{x_reciprocidad:vii:thm:refinamiento-respuesta}, sea \(P_m\) el mapa fijo
que lleva variaciones de conexión del nivel \(m\) al nivel
\(m+1\). Si las acciones de memoria coinciden sobre esos campos
y transportes refinados, sus covectores satisfacen
\[
 s_m=P_m^{\mathsf T}s_{m+1},\qquad
 M_m\sigma_m=P_m^{\mathsf T}M_{m+1}\sigma_{m+1}.
\]
\end{corollary}

\begin{proof}
Diferenciar la igualdad de acciones en la dirección
\(P_m\theta\) da
\(-\tfrac12\theta^{\mathsf T}s_m
 =-\tfrac12\theta^{\mathsf T}P_m^{\mathsf T}s_{m+1}\).
La arbitrariedad de \(\theta\) y las representaciones celulares
de ambos covectores prueban las dos identidades. La composición de
estos mapas itera la igualdad a todo número finito de refinamientos.
\end{proof}

\section{Inserción en la ecuación de Cartan}

La sección~\ref{x_reciprocidad:vii:sec:corriente-respuesta-cartan} normaliza la
tres-forma material mediante
\(\delta_\omega S_m=-\tfrac12\int\delta\omega^{IJ}\wedge
\sigma_{IJ}\). El teorema anterior proporciona sus componentes
celulares para la acción de memoria y el emparejamiento declarados.
La respuesta de Cartan utiliza esa corriente junto a la cotétrada,
el parámetro de Barbero--Immirzi y el acoplamiento gravitatorio de la
misma realización.

La dependencia material respecto de la contorsión determina cómo se
resuelve esa ecuación. Para materia afín, la corriente es independiente
de la contorsión y se aplica la inversión lineal y la eliminación
cuadrática demostradas en la sección siguiente. Para una acción de
holonomía, el diferencial se evalúa en la conexión actual y se conserva
esa dependencia al resolver la ecuación. Por ejemplo, una arista con
\(f_s=f_t=1\), peso uno y \(U(t)=e^{it}\) tiene
\[
 E(t)=2-2\cos t,\qquad E'(t)=2\sin t.
\]
La corriente depende aquí de \(t\). La fórmula completa mantiene
esa dependencia en vez de sustituirla por su valor en una conexión
particular.

Finalmente, la densidad gravitatoria se obtiene de la variación
métrica de la acción efectiva, con su campo y su normalización.
La representación de la corriente, la contorsión que produce y la
densidad resultante son las sucesivas operaciones de esa composición.
La ley de dilución posterior transporta la amplitud así evaluada en
la sección material correspondiente.

% END OWNER


% BEGIN OWNER /Users/ruben/Documents/New project/output/EXTREMA_MEDIA_RAZON_RECIPROCIDAD_ES_EN_20260918/ARTICULO/ES/source/antecedentes_vii/30e_corriente_extension_minima.tex
% Artículo VII. Incorporación aditiva, 10-09-2026.
% RESULTADO_RECUPERADO: fuentes_conservadas/18_clausura_energetica.tex,
% líneas 357--485: Gramiano, peso, extensión mínima y eliminación de Schur;
% líneas 149--180 y 235--295: normalización y unicidad de la respuesta.
% Residencia anterior en VII: 30_pantalla_y_respuesta.tex,
% vii:prop:incidencia, vii:eq:lambda-pantalla, vii:prop:minima-extension.
% FORMALIZACION_NUEVA / CERTIFICADO_NUEVO: diferencial del mínimo con
% incidencia, frontera y normalización variables; incidencia transportada
% por bloques y composición con coordenadas de conexión.
% La familia de rutas que realiza cada incidencia se mantiene explícita.
% No se atribuyen al propietario histórico ni esa familia diferenciable
% completa ni la evaluación de una amplitud material de espín.
% Dependencias: 30, 30c y la convención variacional de 31. No modifica
% esos fragmentos ni constituye por sí solo una realización material.

\chapter{Diferencial de la extensión mínima de frontera}
\label{x_reciprocidad:vii:sec:corriente-extension-minima}

La incidencia cúbica determina simultáneamente el campo de extensión
mínima y el peso de su respuesta. Para un dato de frontera, ambos se
calculan mediante el inverso del Gramiano de incidencia. Se desarrolla
ahora el diferencial de esa misma operación, conservando las variaciones
de la frontera, del transporte y de la escala de normalización. De este
modo, el campo que interviene en la respuesta queda seleccionado por
el problema variacional de pantalla.

\section{Campo minimizante y peso de respuesta}
\label{x_reciprocidad:vii:subsec:extension-minima-parametrizada}

Sean \(H\) y \(Q\) espacios hermíticos de dimensión finita, con
productos fijos y antilineales en la primera variable. El caso real se
obtiene sustituyendo los adjuntos por transpuestas. En un abierto de
parámetros reales \(\Omega\), consideremos familias de clase \(C^1\)
\[
 \mathsf C(\theta):H\longrightarrow Q,\qquad
 b(\theta)\in Q,\qquad \chi(\theta)>0,
 \qquad \operatorname{Ran}\mathsf C(\theta)=Q.
\]
La sobreyectividad equivale a rango fila completo en bases ortonormales.
Definimos el Gramiano, el potencial conjugado y la extensión
\begin{equation}
 \mathsf L=\mathsf C\mathsf C^*,\qquad
 y=\mathsf L^{-1}b,\qquad f_{\min}=\mathsf C^*y,
 \qquad \Lambda=\chi\mathsf L^{-1}.
 \label{x_reciprocidad:vii:eq:gramiano-extension-minima}
\end{equation}
La incidencia \(\mathsf C\) y el emparejamiento celular \(M\) de
\eqref{x_reciprocidad:vii:eq:coordenadas-corriente-material} son operadores diferentes.

\begin{proposition}[Selección variacional de la extensión]
\label{x_reciprocidad:vii:prop:minimo-parametrizado}
La extensión \(f_{\min}\) es la única minimizante de la norma entre
los campos \(f\in H\) que satisfacen \(\mathsf C f=b\). Su energía es
\begin{equation}
 E(\theta)=\chi\min_{\mathsf C f=b}\|f\|^2
 =\chi\langle b,\mathsf L^{-1}b\rangle
 =\langle b,\Lambda b\rangle.
 \label{x_reciprocidad:vii:eq:energia-extension-minima}
\end{equation}
El campo \(f_{\min}\), el potencial \(y\) y la energía \(E\)
dependen de manera \(C^1\) de los parámetros.
\end{proposition}

\begin{proof}
La sobreyectividad de \(\mathsf C\) hace inyectivo a \(\mathsf C^*\),
por lo que \(\langle z,\mathsf Lz\rangle=\|\mathsf C^*z\|^2>0\)
para \(z\ne0\). Así, \(\mathsf L\) es invertible y
\(\mathsf C f_{\min}=\mathsf L y=b\). Toda solución es
\(f_{\min}+z\), con \(z\in\ker\mathsf C\). Como
\(f_{\min}\in\operatorname{Ran}\mathsf C^*=(\ker\mathsf C)^\perp\),
\[
 \|f_{\min}+z\|^2=\|f_{\min}\|^2+\|z\|^2,
 \qquad
 \|f_{\min}\|^2=\langle y,\mathsf L y\rangle
 =\langle b,\mathsf L^{-1}b\rangle.
\]
Estas identidades prueban el mínimo, su unicidad y su valor. La inversión
es una operación diferenciable en el abierto de matrices invertibles;
las expresiones de \eqref{x_reciprocidad:vii:eq:gramiano-extension-minima} proporcionan
la regularidad afirmada.
\end{proof}

La realización cúbica de \ref{x_reciprocidad:vii:prop:incidencia} corresponde a
\(H=\mathbb R^{\{-1,1\}^3}\), al cociente de pantalla
\(Q=P_\square\mathbb R^6\) y a la incidencia \(\mathsf C=M:H\to Q\).
Su Gramiano y su normalización son
\begin{equation}
 \mathsf L_\square=12P_u+4P_-,\qquad
 \chi_m=\frac{112}{81\ell_m^2},\qquad
 \Lambda_m=\frac{28}{243\ell_m^2}(P_u+3P_-).
 \label{x_reciprocidad:vii:eq:normalizacion-minimo-cubico}
\end{equation}
El factor \(112/81\) procede de la soldadura y del defecto cuadrático
de compresión de \eqref{x_reciprocidad:vii:eq:lambda-pantalla}. Si \(b\) tiene dimensión
de longitud, también la tienen \(y\) y \(f_{\min}\), la incidencia es
adimensional y \(E\) es adimensional. La conversión a acción utiliza
la normalización en la recta de acción que se especifica más adelante.

\Needspace{15\baselineskip}
\section{Diferencial completo del mínimo}
\label{x_reciprocidad:vii:subsec:diferencial-minimo}

\begin{theorem}[Respuesta a la variación de incidencia y frontera]
\label{x_reciprocidad:vii:thm:diferencial-minimo}
Para toda dirección tangente real \(\delta\theta\), se tiene
\begin{equation}
 \delta E
 =\delta\chi\,\langle b,y\rangle
   +2\chi\operatorname{Re}
       \langle y,\delta b-(\delta\mathsf C)f_{\min}\rangle.
 \label{x_reciprocidad:vii:eq:diferencial-minimo-completo}
\end{equation}
Si \(\chi=112/(81\ell_m^2)\), con \(\ell_m>0\) variable,
\begin{equation}
 \delta E=-2E\frac{\delta\ell_m}{\ell_m}
 +2\chi_m\operatorname{Re}
       \langle y,\delta b-(\delta\mathsf C)f_{\min}\rangle.
 \label{x_reciprocidad:vii:eq:diferencial-minimo-escala}
\end{equation}
\end{theorem}

\begin{proof}
Diferenciar \(\mathsf L\mathsf L^{-1}=I_Q\) da
\[
 \delta\mathsf L^{-1}
 =-\mathsf L^{-1}(\delta\mathsf L)\mathsf L^{-1},\qquad
 \delta\mathsf L=(\delta\mathsf C)\mathsf C^*
                      +\mathsf C(\delta\mathsf C)^*.
\]
Aplicando la regla del producto a \eqref{x_reciprocidad:vii:eq:energia-extension-minima}
y usando la autoadjunción de \(\mathsf L^{-1}\), resulta
\[
 \delta E=\delta\chi\,\langle b,y\rangle
 +2\chi\operatorname{Re}\langle y,\delta b\rangle
 -\chi\langle y,(\delta\mathsf L)y\rangle.
\]
Los dos términos del último producto son conjugados entre sí y
\(\mathsf C^*y=f_{\min}\); por tanto,
\[
 \langle y,(\delta\mathsf L)y\rangle
 =2\operatorname{Re}\langle y,(\delta\mathsf C)f_{\min}\rangle.
\]
Se obtiene \eqref{x_reciprocidad:vii:eq:diferencial-minimo-completo}.
Finalmente, \(\delta\chi_m=-2\chi_m\delta\ell_m/\ell_m\)
produce \eqref{x_reciprocidad:vii:eq:diferencial-minimo-escala}.
\end{proof}

La fórmula incorpora la variación del campo minimizante mediante la
ecuación que lo determina. También se comprueba directamente:
\(\mathsf C\delta f_{\min}=\delta b-(\delta\mathsf C)f_{\min}\)
y \(f_{\min}=\mathsf C^*y\) implican
\[
 \delta\|f_{\min}\|^2
 =2\operatorname{Re}\langle f_{\min},\delta f_{\min}\rangle
 =2\operatorname{Re}
     \langle y,\delta b-(\delta\mathsf C)f_{\min}\rangle.
\]
Por ejemplo, si \(\mathsf C\) es fijo y \(b=\ell_m b_0\), con
\(b_0\) fijo, los dos términos de
\eqref{x_reciprocidad:vii:eq:diferencial-minimo-escala} se cancelan y la energía
permanece constante. Este control conserva conjuntamente la dimensión
de frontera y su normalización.

\section{Cociente fijo y cambios de marco}
\label{x_reciprocidad:vii:subsec:marcos-minimo}

Las derivadas anteriores están tomadas en espacios \(H,Q\) y productos
hermíticos fijos. En la presentación de seis caras, la inversa siempre
actúa sobre \(Q=P_\square H_F\). Para una familia de cocientes
\(Q(\theta)\), se fija primero una trivialización isométrica
\(V_Q(\theta):Q\to Q(\theta)\), y análogamente para el dominio.
Las derivadas de esos mapas se incluyen al diferenciar las expresiones
transportadas a los espacios fijos. Las fórmulas se aplican en cada
abierto donde se conserve el rango. Los cambios del producto energético
requieren diferenciar además ese producto; una métrica lorentziana
y el producto positivo usado para minimizar conservan aquí sus tipos
distintos.

\begin{proposition}[Covariancia del mínimo]
\label{x_reciprocidad:vii:prop:covariancia-minimo}
Sean \(g_H,g_Q\) unitarios. Bajo
\[
 \mathsf C'=g_Q\mathsf C g_H^*,\qquad b'=g_Qb,\qquad \chi'=\chi,
\]
se cumple \(y'=g_Qy\), \(f'_{\min}=g_Hf_{\min}\) y \(E'=E\).
Cuando los cambios de marco dependen del parámetro, el diferencial
completo conserva esa igualdad. En particular, una variación pura de
marco tiene \(\delta E=0\).
\end{proposition}

\begin{proof}
El Gramiano transformado es \(\mathsf L'=g_Q\mathsf Lg_Q^*\),
de donde se obtienen las transformaciones de \(y,f_{\min}\) y
la igualdad de energías. Para una variación pura de marco, en un punto
con marcos identidad, sean \(A_Q^*=-A_Q\) y \(A_H^*=-A_H\)
sus generadores. Entonces
\[
 \delta\mathsf C=A_Q\mathsf C-\mathsf C A_H,\qquad
 \delta b=A_Qb,\qquad
 \delta b-(\delta\mathsf C)f_{\min}=\mathsf C A_Hf_{\min}.
\]
Su contribución a \eqref{x_reciprocidad:vii:eq:diferencial-minimo-completo} es
\(2\chi\operatorname{Re}\langle f_{\min},A_Hf_{\min}\rangle=0\).
La prueba incluye, así, los términos debidos a los marcos móviles.
\end{proof}

\section{Incidencia transportada y componentes de conexión}
\label{x_reciprocidad:vii:subsec:incidencia-conexion-minima}

La dependencia de conexión se especifica antes de efectuar la
contracción. Una realización explícita de la incidencia por bloques
se obtiene como sigue. Sean \(\mathcal K\) una fibra hermítica común,
\(F=\{(i,\epsilon):1\leq i\leq3,\ \epsilon=\pm1\}\) y
\(V=\{-1,1\}^3\). Para cada incidencia \(\nu_i=\epsilon\),
la realización utilizada declara una ruta \(\gamma_{i,\epsilon,\nu}\)
desde la fibra del vértice hasta la de la cara, con transporte
\(U_{i,\epsilon,\nu}\). Las trivializaciones identifican esas fibras
con \(\mathcal K\). Definimos
\begin{equation}
 \begin{split}
 (\mathsf B_Uf)_{i,\epsilon}
   &=\sum_{\nu_i=\epsilon}U_{i,\epsilon,\nu}f_\nu,\\
 \mathsf C_U&=(q_\square\otimes I_{\mathcal K})\mathsf B_U:
   \bigoplus_{\nu\in V}\mathcal K\longrightarrow Q\otimes\mathcal K,
 \qquad q_\square=P_\square:H_F\to Q.
 \end{split}
 \label{x_reciprocidad:vii:eq:incidencia-transportada-minima}
\end{equation}
Esta es una continuación variacional explícita de la incidencia cúbica
en la representación de rutas declarada. Para transportes identidad,
\(\mathsf C_U=M\otimes I_{\mathcal K}\) y se recupera
\eqref{x_reciprocidad:vii:eq:normalizacion-minimo-cubico}, tensorizada con la fibra.
El Gramiano sigue siendo positivo en un entorno de esa realización,
pues la positividad definida es una condición abierta. En cada dominio
de rango completo, el peso y el campo se fijan por
\eqref{x_reciprocidad:vii:eq:gramiano-extension-minima}. Las rutas que realizan las
incidencias, sus orientaciones y sus memorias son datos de esa
realización del transporte APP--TRIT--TPK y acompañan a la construcción.

Si \(\omega=(\omega_j)\) son coordenadas reales de conexión,
la derivada de \(\mathsf C_U\) se calcula sustituyendo en cada bloque
\begin{equation}
 \partial_j U_\gamma
 =\sum_{k=1}^{N}
    U_N\cdots U_{k+1}(\partial_jU_k)U_{k-1}\cdots U_1,
 \qquad U_\gamma=U_N\cdots U_1.
 \label{x_reciprocidad:vii:eq:derivada-incidencia-rutas}
\end{equation}
Los productos vacíos son identidades. La fórmula procede de la regla
del producto y conserva el lugar de cada operación en la ruta.
Así, \(\mathsf C_j:=\partial_j\mathsf C_U\) se obtiene de una
suma finita de matrices conocidas en la realización elegida, seguida
de la proyección \(q_\square\otimes I\). Si se utiliza un cociente
móvil, su derivada se incorpora conforme a
\ref{x_reciprocidad:vii:subsec:marcos-minimo}. La matriz cúbica entera fija corresponde
al caso \(\mathsf C_j=0\); la respuesta a conexión procede de la
dependencia de los transportes efectivamente declarada en
\eqref{x_reciprocidad:vii:eq:incidencia-transportada-minima}.

% REMATE_LOCAL_X_EXTENSION_MINIMA
\Needspace{34\baselineskip}
\begin{theorem}[Componentes variacionales de la extensión mínima]
\label{x_reciprocidad:vii:thm:componentes-minimo-conexion}
Para la acción de pantalla
\(S_{\min}=\mathfrak a_m E\), donde \(\mathfrak a_m\) es su
normalización en la recta de acción, escribamos
\[
 b_j=\partial_jb,\qquad \chi_j=\partial_j\chi,
 \qquad a_j=\partial_j\mathfrak a_m.
\]
Con la convención \(\delta S_{\min}=-\tfrac12\sum_j s_j\delta\omega_j\),
las componentes son
\begin{equation}
 \begin{split}
 s_j={}&4\mathfrak a_m\chi\operatorname{Re}
             \langle y,\mathsf C_jf_{\min}-b_j\rangle\\
      &-2\mathfrak a_m\chi_j\langle b,y\rangle-2a_jE.
 \end{split}
 \label{x_reciprocidad:vii:eq:componentes-minimo-conexion}
\end{equation}
Para la normalización cúbica, el segundo término equivale a
\(4\mathfrak a_m E\,\partial_j\ell_m/\ell_m\).
Si el emparejamiento celular de esta misma acción tiene matriz
invertible \(\mathsf M_{\rm cel}\), sus coordenadas de corriente
son \(\sigma=\mathsf M_{\rm cel}^{-1}s\).
\end{theorem}

\begin{proof}
La regla del producto da
\(\partial_jS_{\min}=a_jE+\mathfrak a_m\partial_jE\).
Sustituir \eqref{x_reciprocidad:vii:eq:diferencial-minimo-completo} y multiplicar
por \(-2\) prueba \eqref{x_reciprocidad:vii:eq:componentes-minimo-conexion}.
La derivada de \(\chi_m\) produce la expresión de escala.
Finalmente, comparar
\(\delta S_{\min}=-\tfrac12\delta\omega^{\mathsf T}s\)
con \(-\tfrac12\delta\omega^{\mathsf T}\mathsf M_{\rm cel}\sigma\)
para toda variación da \(\mathsf M_{\rm cel}\sigma=s\),
y su inversa determina las coordenadas.
\end{proof}

La construcción recupera un campo y un peso seleccionados conjuntamente
por la incidencia, y evalúa su respuesta a las variaciones de la misma
realización. El dato de frontera \(b\), la familia de rutas, la escala
\(\ell_m\), la normalización \(\mathfrak a_m\) y el emparejamiento
celular conservan sus reglas de publicación. La interpretación como
corriente material utiliza esos objetos y la acción completa. La
respuesta de Cartan y su densidad métrica posterior se componen con
ella en el dominio de materia correspondiente, conservando cualquier
dependencia de la corriente respecto de la conexión.

% END OWNER

\begin{HMTCompactSection}

% BEGIN OWNER /Users/ruben/Documents/New project/output/EXTREMA_MEDIA_RAZON_RECIPROCIDAD_ES_EN_20260918/ARTICULO/ES/source/antecedentes_vii/31_corriente_y_respuesta_cartan.tex
% Artículo VII. Fragmento de cuerpo: corriente y respuesta Cartan--Holst.
% RESULTADO_RECUPERADO: acción, corriente variacional, inversa de Holst,
% torsión algebraica, respuesta cuadrática y normalización métrica.
% Fuentes leídas íntegramente:
% [II10] output/ARTICULO_II_REV10_EDICION_INTEGRADA_20260910/manuscrito/
%        sections/10_gravedad.tex, líneas 109--155.
% [G11c] integral activo 2249, fuente/manuscrito/sections/md/
%        11c_gravedad_holonomia_rev6.tex, líneas 214--334.
% [B028] integral activo, fuente/manuscrito/integracion_83/parte_iv/body/
%        B028_ch68_realizacion_cartan_holst_barbero_7df2afd9f675_11_gravedad_cartan_holst.tex,
%        líneas 76--218; copia idéntica conservada en II10/fuentes/propietarios_II/G03/.
%        SHA-256 e1e1f3325399b3ef09c331ed8268e3883f7bcc561cd7a560cf80d1be7eb1e91c.
% CERTIFICADO_NUEVO de resultados recuperados: inversión explícita del mapa
% torsión -> D(e wedge e), prueba de unicidad de contorsión y eliminación
% homogénea del término cuadrático. Las dos inversas se comprobaron sobre
% las 24 bases correspondientes mediante aritmética racional exacta.
% Precisión de dominio: la respuesta lineal en corriente y su eliminación
% cuadrática se enuncian para materia afín en la contorsión. Se conserva
% la dependencia del coeficiente respecto de la acción material concreta.
%
% DEPENDENCIAS DE INTEGRACIÓN (no constituyen remisiones bibliográficas):
% 1. vii:sec:realizacion-cartan: cotetrada, conexión y realización del
%    transporte enriquecido; incluir su prueba en el cuerpo precedente.
% 2. vii:sec:acoplamientos: genealogía de gamma_HMT, G_int, c_int,
%    homogeneidad dimensional y elección de la carta, con pruebas.
% 3. Para evaluar una especie material: acción S_m concreta, representación
%    de sus campos y cálculo de su corriente sigma. Este fragmento recibe
%    esa acción declarada y demuestra su respuesta variacional.
% 4. La identificación de una corriente de pantalla con sigma, su
%    contracción física y su normalización material tienen residencia propia.
%    Aquí no se introduce ni se selecciona una amplitud s_0.
% 5. El balance métrico utiliza la covariancia de la acción material y
%    sus ecuaciones de movimiento; no supone conservación separada del
%    tensor de espín y del tensor material de Levi--Civita.

\chapter{Corriente de espín y respuesta algebraica de Cartan--Holst}
\label{x_reciprocidad:vii:sec:corriente-respuesta-cartan}

La realización del transporte enriquecido de la
sección~\ref{x_reciprocidad:vii:sec:realizacion-cartan} proporciona la cotetrada y la
conexión; las secciones dimensionales y el parámetro de incidencia de la
secciones~\ref{x_reciprocidad:v:ant:sec:barbero} y~\ref{x_reciprocidad:v:ant:sec:gravity} fijan sus acoplamientos. Su composición
determina un problema variacional de primer orden. La corriente material
ocupa en él una posición precisa: es la derivada de la acción respecto de
la conexión y determina la respuesta torsional con la misma normalización.

\section{Dominio geométrico y normalización de la corriente}
\label{x_reciprocidad:vii:subsec:dominio-corriente-cartan}

Trabajamos en una variedad lisa orientada \(M\) de dimensión cuatro,
con fibrado lorentziano interno \(E\), métrica
\(\eta=\operatorname{diag}(-1,1,1,1)\) y cotetrada no degenerada
\(e\in\Omega^1(M;E)\). En un marco local escribimos
\(e^I\), con \(I=0,1,2,3\). La conexión métrica
\(\omega^{IJ}=-\omega^{JI}\) actúa mediante \(D_\omega\).
Cuando los campos \(\Psi\) sean espinoriales, se conserva una estructura
de espín y la representación material correspondiente. Definimos
\begin{equation}
 \Sigma^{IJ}=e^I\wedge e^J,\qquad
 F^{IJ}=\mathrm d\omega^{IJ}+\omega^I{}_{K}\wedge\omega^{KJ},
 \qquad T^I=D_\omega e^I.
 \label{x_reciprocidad:vii:eq:campos-cartan}
\end{equation}
El operador \(\star\) actúa sobre los índices bivectoriales internos,
\begin{equation}
 (\star X)^{IJ}=\tfrac12\epsilon^{IJ}{}_{KL}X^{KL},
 \qquad \star^2=-I,\qquad
 \mathsf P_\gamma=\star+\gamma^{-1}I.
 \label{x_reciprocidad:vii:eq:operador-holst}
\end{equation}
Se distingue así esta dualidad de la dualidad exterior de formas sobre
\(M\). Los índices internos se elevan y contraen con \(\eta\).

Fijamos una carta con \(G_{\rm int}>0\), \(c_{\rm int}>0\),
\(\gamma=\gamma_{\rm HMT}\in\mathbb R\setminus\{0\}\) y
\(\Lambda_{\rm int}\) constantes sobre ella. Abreviamos
\(\kappa=8\pi G_{\rm int}/c_{\rm int}^4\) y \(c=c_{\rm int}\).
Estas secciones permanecen fijas durante la variación. La acción es
\begin{equation}
 \begin{split}
 S[e,\omega,\Psi]={}&
 \frac1{2\kappa c}\int_M\Sigma_{IJ}\wedge
                         (\mathsf P_\gamma F)^{IJ}
 -\frac{\Lambda_{\rm int}}{\kappa c}
                         \int_M\operatorname{vol}_e
 +S_{\rm m}[e,\omega,\Psi].
 \end{split}
 \label{x_reciprocidad:vii:eq:accion-cartan-holst}
\end{equation}
La convención de volumen y orientación es la que da
\(\Sigma_{IJ}\wedge(\star F)^{IJ}=R\operatorname{vol}_e\)
en el sector de Levi--Civita. El tipo dimensional previo hace de cada
sumando una sección de acción.

Para la variación de la conexión se utilizan variaciones de soporte
compacto en el interior, o se fija
\(\delta\omega|_{\partial M}=0\). Puede emplearse también una
completación de borde cuyo problema variacional anule el mismo término
fronterizo. A cotetrada y campos materiales fijos, la corriente de espín
es la forma de grado tres, con valores en bivectores internos duales,
determinada por
\begin{equation}
 \delta_\omega S_{\rm m}
 =-\tfrac12\int_M\delta\omega^{IJ}\wedge\sigma_{IJ}.
 \label{x_reciprocidad:vii:eq:corriente-variacional}
\end{equation}
El emparejamiento con una variación arbitraria de soporte compacto fija
\(\sigma\) de manera única. Esta definición conserva la unidad de
acción de la corriente y el factor de normalización que interviene en
su respuesta geométrica.

\section{Variación de la conexión e inversión de Holst}
\label{x_reciprocidad:vii:subsec:variacion-holst}

\begin{theorem}[Ecuación variacional de Cartan--Holst]
\label{x_reciprocidad:vii:thm:ecuacion-cartan-holst}
En el dominio anterior, la estacionariedad respecto de la conexión
equivale a
\begin{equation}
 D_\omega(\mathsf P_\gamma\Sigma)^{IJ}
       =\kappa c\,\sigma^{IJ}.
 \label{x_reciprocidad:vii:eq:ecuacion-cartan-holst}
\end{equation}
\end{theorem}
\begin{proof}
La variación de la curvatura es
\(\delta F=D_\omega\delta\omega\). La dualidad interna es paralela
y autoadjunta para la contracción bivectorial; como \(\gamma\) es
constante, \(D_\omega\mathsf P_\gamma=\mathsf P_\gamma D_\omega\).
Poniendo \(B_{IJ}=(\mathsf P_\gamma\Sigma)_{IJ}\), se obtiene
\[
 B_{IJ}\wedge D_\omega\delta\omega^{IJ}
 =\mathrm d(B_{IJ}\wedge\delta\omega^{IJ})
   -D_\omega B_{IJ}\wedge\delta\omega^{IJ}.
\]
La última forma cambia de signo al permutar sus factores de grados tres
y uno. Las condiciones de borde eliminan la integral del término
exacto. El término cosmológico es independiente de \(\omega\).
Por tanto,
\[
 \delta_\omega S=
 \frac1{2\kappa c}\int_M\delta\omega^{IJ}\wedge
       \bigl[D_\omega B_{IJ}-\kappa c\,\sigma_{IJ}\bigr].
\]
La arbitrariedad de la variación prueba la ecuación, incluido su signo
y su coeficiente.
\end{proof}

\begin{lemma}[Inversa real del operador de Holst]
\label{x_reciprocidad:vii:lem:inversa-holst}
Para \(\gamma\in\mathbb R\setminus\{0\}\),
\begin{equation}
 \mathsf P_\gamma^{-1}
 =\frac{\gamma^2}{1+\gamma^2}
                      (-\star+\gamma^{-1}I).
 \label{x_reciprocidad:vii:eq:inversa-holst}
\end{equation}
\end{lemma}
\begin{proof}
Los dos factores son polinomios en \(\star\) y conmutan. Su producto
antes de normalizar es
\[
 (\star+\gamma^{-1}I)(-\star+\gamma^{-1}I)
 =-\star^2+\gamma^{-2}I=(1+\gamma^{-2})I.
\]
El denominador \(1+\gamma^2\) es positivo en el dominio real fijado.
La multiplicación por \(\gamma^2/(1+\gamma^2)\) da la identidad.
\end{proof}

La corriente determina así la forma bivectorial de grado tres
\begin{equation}
 Y^{IJ}:=\kappa c\,
   \frac{\gamma^2}{1+\gamma^2}
   (-\star+\gamma^{-1}I)\sigma^{IJ},
 \qquad D_\omega\Sigma^{IJ}=Y^{IJ}.
 \label{x_reciprocidad:vii:eq:corriente-dualizada}
\end{equation}

\section{Reconstrucción puntual de la torsión y de la contorsión}
\label{x_reciprocidad:vii:subsec:reconstruccion-torsion}

La regla de Leibniz da
\begin{equation}
 D_\omega\Sigma^{IJ}
 =T^I\wedge e^J-e^I\wedge T^J
 =:\mathsf L_e(T)^{IJ}.
 \label{x_reciprocidad:vii:eq:mapa-torsion}
\end{equation}
En cada punto, \(\mathsf L_e\) actúa de
\(\Lambda^2T^*M\otimes E\) en
\(\Lambda^3T^*M\otimes\Lambda^2E\). Ambos espacios tienen dimensión
veinticuatro. Sea \(E_I\) el marco dual de la cotetrada, de modo que
\(\iota_{E_I}e^J=\delta_I^J\).

\begin{theorem}[Inversión algebraica de la ecuación torsional]
\label{x_reciprocidad:vii:thm:inversion-torsion}
Para una cotetrada no degenerada, \(\mathsf L_e\) es un isomorfismo.
Su inversa sobre la forma \(Y\) se escribe
\begin{equation}
 B^I=\sum_J\iota_{E_J}Y^{IJ},\qquad
 b=\tfrac14\sum_I\iota_{E_I}B^I,\qquad
 T^I=B^I-e^I\wedge b.
 \label{x_reciprocidad:vii:eq:torsion-explicita}
\end{equation}
La ecuación~\eqref{x_reciprocidad:vii:eq:ecuacion-cartan-holst} determina por tanto
la torsión puntualmente a partir de \(e\) y de la corriente que figura
en ella.
\end{theorem}
\begin{proof}
Sea primero \(Y=\mathsf L_e(T)\) y pongamos
\(t=\sum_J\iota_{E_J}T^J\). La contracción de
\eqref{x_reciprocidad:vii:eq:mapa-torsion}, utilizando que \(T^I\) tiene grado dos
y que hay cuatro elementos en la cotetrada, da
\[
 \sum_J\iota_{E_J}(T^I\wedge e^J)=2T^I,
 \qquad
 \sum_J\iota_{E_J}(e^I\wedge T^J)=T^I-e^I\wedge t.
\]
De aquí \(B^I=T^I+e^I\wedge t\). Una segunda contracción produce
\[
 \sum_I\iota_{E_I}B^I=t+3t=4t.
\]
Se recupera \(t=b\) y después \(T^I=B^I-e^I\wedge b\).
En particular, el núcleo de \(\mathsf L_e\) es nulo. La igualdad
de dimensiones prueba su sobreyectividad y hace válida la fórmula
para toda forma \(Y\) del codominio.
\end{proof}

La conexión se descompone de manera única como
\begin{equation}
 \omega^{IJ}=\mathring\omega^{IJ}(e)+K^{IJ},\qquad
 K^{IJ}=-K^{JI},\qquad T^I=K^I{}_{J}\wedge e^J,
 \label{x_reciprocidad:vii:eq:contorsion}
\end{equation}
donde \(\mathring\omega(e)\) es la conexión de Levi--Civita y
\(K\) la contorsión. Para hacer explícita la inversa, escribimos
\(K_{IJ}=K_{IJK}e^K\) y
\(T_I=\tfrac12T_{IJK}e^J\wedge e^K\). Entonces
\begin{equation}
 T_{IJK}=K_{IKJ}-K_{IJK},\qquad
 K_{IJK}=\tfrac12(T_{KIJ}-T_{IJK}-T_{JKI}).
 \label{x_reciprocidad:vii:eq:contorsion-componentes}
\end{equation}
La segunda igualdad se obtiene combinando cíclicamente la primera y
usando \(K_{IJK}=-K_{JIK}\); su sustitución recupera las componentes
de \(T\). La torsión y la contorsión son, por consiguiente, dos
presentaciones equivalentes de la respuesta algebraica con cotetrada fija.

\begin{corollary}[Sector sin corriente y respuesta lineal mínima]
\label{x_reciprocidad:vii:cor:desacoplamiento-torsional}
Si \(\sigma=0\), se obtiene \(T=0\) y
\(\omega=\mathring\omega(e)\). Para un acoplamiento material cuya
corriente \(\sigma[e,\Psi]\) sea independiente de \(K\),
las ecuaciones~\eqref{x_reciprocidad:vii:eq:corriente-dualizada},
\eqref{x_reciprocidad:vii:eq:torsion-explicita} y
\eqref{x_reciprocidad:vii:eq:contorsion-componentes} proporcionan una respuesta única
\(K_*(e,\Psi)\), lineal en \(\sigma\).
\end{corollary}
\begin{proof}
Las tres aplicaciones de reconstrucción son lineales con \(e\),
\(\gamma\) y \(\kappa c\) fijos. La corriente nula produce
\(Y=0\), después \(T=0\) y finalmente \(K=0\).
\end{proof}

En la acción mínima, la ecuación de conexión contiene esta respuesta
puntual. Su evolución se obtiene al evolucionar \(e\) y \(\Psi\):
\(K_*=K_*(e,\Psi)\) cambia con ellos. Si una acción material conserva
dependencia algebraica adicional en \(K\), la ecuación de Cartan
mantiene su forma variacional y debe resolverse con esa corriente;
la respuesta lineal anterior corresponde al dominio mínimo especificado.

\section{Eliminación de la contorsión y contribución cuadrática}
\label{x_reciprocidad:vii:subsec:respuesta-cuadratica}

Consideremos ahora el acoplamiento mínimo afín en \(K\), expresado por
\begin{equation}
 S_{\rm m}[e,\mathring\omega+K,\Psi]
 =S_{\rm m}[e,\mathring\omega,\Psi]
   -\tfrac12\int_M K^{IJ}\wedge\sigma_{IJ}[e,\Psi].
 \label{x_reciprocidad:vii:eq:materia-afin-contorsion}
\end{equation}
Esta condición identifica la acción material a la que se aplica la
eliminación cuadrática; la corriente sigue siendo la definida por
\eqref{x_reciprocidad:vii:eq:corriente-variacional}.

La expansión de la curvatura es
\begin{equation}
 F(\mathring\omega+K)
 =F(\mathring\omega)+D_{\mathring\omega}K+K\wedge K,
 \qquad (K\wedge K)^{IJ}=K^I{}_{L}\wedge K^{LJ}.
 \label{x_reciprocidad:vii:eq:curvatura-contorsion}
\end{equation}
Como \(D_{\mathring\omega}\Sigma=0\), el término lineal es la
derivada exterior de \(\Sigma_{IJ}\wedge(\mathsf P_\gamma K)^{IJ}\).
La contribución de borde correspondiente se mantiene explícita:
\begin{equation}
 S_\partial[e,K]=\frac1{2\kappa c}
     \int_{\partial M}\Sigma_{IJ}\wedge(\mathsf P_\gamma K)^{IJ}.
 \label{x_reciprocidad:vii:eq:borde-contorsion}
\end{equation}
En un dominio interior, o con la completación de borde compatible,
la densidad de acción que depende de la contorsión es
\begin{equation}
 \mathcal Q_{e,\gamma}(K)-\tfrac12K^{IJ}\wedge\sigma_{IJ},
 \qquad
 \mathcal Q_{e,\gamma}(K)
  =\frac1{2\kappa c}\Sigma_{IJ}\wedge
                         \mathsf P_\gamma(K\wedge K)^{IJ}.
 \label{x_reciprocidad:vii:eq:forma-cuadratica-contorsion}
\end{equation}
La densidad \(\mathcal Q_{e,\gamma}\) es una forma diferencial de
grado cuatro, homogénea de grado dos en \(K\).

\begin{theorem}[Respuesta efectiva cuadrática en la corriente]
\label{x_reciprocidad:vii:thm:respuesta-cuadratica}
Bajo \eqref{x_reciprocidad:vii:eq:materia-afin-contorsion}, la sustitución de la
solución algebraica \(K_*\) produce la densidad efectiva de espín
\begin{equation}
 \mathcal L_{\rm spin}^{\rm eff}
 =-\tfrac14 K_*^{IJ}\wedge\sigma_{IJ}
 =-\mathcal Q_{e,\gamma}(K_*).
 \label{x_reciprocidad:vii:eq:accion-efectiva-spin}
\end{equation}
Con cotetrada y acoplamientos fijos, es homogénea de grado dos en
la corriente. La acción efectiva conserva además
\(S_\partial[e,K_*]\), o su completación de borde declarada.
\end{theorem}
\begin{proof}
Por el corolario~\ref{x_reciprocidad:vii:cor:desacoplamiento-torsional},
\(K_*\) es lineal en \(\sigma\). La ecuación de conexión es
precisamente la estacionariedad de
\eqref{x_reciprocidad:vii:eq:forma-cuadratica-contorsion}. Como ésta es algebraica,
se aplica puntualmente la variación radial \(\delta K=K_*\);
también puede multiplicarse por una función de soporte compacto.
La identidad de Euler para una forma cuadrática da
\[
 2\mathcal Q_{e,\gamma}(K_*)
             -\tfrac12K_*^{IJ}\wedge\sigma_{IJ}=0.
\]
La sustitución en la densidad de acción prueba
\eqref{x_reciprocidad:vii:eq:accion-efectiva-spin}. Si
\(\sigma\mapsto\lambda\sigma\), entonces
\(K_*\mapsto\lambda K_*\) y la densidad efectiva se multiplica
por \(\lambda^2\). La expansión
\eqref{x_reciprocidad:vii:eq:curvatura-contorsion} identifica separadamente el
término de borde que acompaña a esa eliminación.
\end{proof}

La corriente concreta, su contracción lorentziana, la dualidad y las
convenciones de la acción material determinan el signo y el coeficiente
de esta contribución. La identidad cuadrática proporciona la regla de
composición que permite evaluarlos una vez especificados esos datos.

\section{Fuente métrica efectiva y balance covariante}
\label{x_reciprocidad:vii:subsec:fuente-metrica-efectiva}

La normalización de la fuente métrica se fija en la misma recta de
acción. Tras eliminar \(K\), definimos
\begin{align}
 \delta_g S_{\rm m}[e,\mathring\omega,\Psi]
 &=-\tfrac12\int_M\operatorname{vol}_e\,
                   \mathcal T^{S,\rm LC}_{\mu\nu}\,\delta g^{\mu\nu},
 \\
 \delta_g S_{\rm spin}^{\rm eff}
 &=-\tfrac12\int_M\operatorname{vol}_e\,
                   \mathcal U^{S,\rm spin}_{\mu\nu}\,\delta g^{\mu\nu}.
 \label{x_reciprocidad:vii:eq:fuentes-normalizadas-accion}
\end{align}
Aquí se consideran variaciones interiores, con los términos de borde
tratados conforme al problema variacional fijado. Para el acoplamiento
afín, la segunda fuente es la variación métrica de
\eqref{x_reciprocidad:vii:eq:accion-efectiva-spin}. La variación incluye la dependencia
de la corriente material respecto de \(e\) y de \(\Psi\).

\begin{proposition}[Ecuación métrica y conservación de la fuente total]
\label{x_reciprocidad:vii:prop:fuente-metrica-balance}
Con las convenciones anteriores, las ecuaciones métricas son
\begin{equation}
 G_{\mu\nu}[g]+\Lambda_{\rm int}g_{\mu\nu}
 =\kappa c\bigl(\mathcal T^{S,\rm LC}_{\mu\nu}
                    +\mathcal U^{S,\rm spin}_{\mu\nu}\bigr)
 =\kappa\bigl(T^{\rm phys,LC}_{\mu\nu}
                    +U^{\rm phys,spin}_{\mu\nu}\bigr),
 \label{x_reciprocidad:vii:eq:ecuacion-metrica-efectiva}
\end{equation}
donde \(T^{\rm phys,LC}=c\mathcal T^{S,\rm LC}\) y
\(U^{\rm phys,spin}=c\mathcal U^{S,\rm spin}\) cuando la
cotetrada toma valores en la recta de longitud. Toda solución satisface
\begin{equation}
 \mathring\nabla^\mu
   (T^{\rm phys,LC}_{\mu\nu}+U^{\rm phys,spin}_{\mu\nu})=0.
 \label{x_reciprocidad:vii:eq:balance-metrico-total}
\end{equation}
\end{proposition}
\begin{proof}
La eliminación de \(K\) conserva las ecuaciones de las variables
restantes: en la regla de la cadena, el término que multiplica
\(\delta K_*\) se anula por la ecuación de conexión. En el sector
de Levi--Civita, la primera identidad de Bianchi anula el término
\(\Sigma_{IJ}\wedge F^{IJ}\), y la parte gravitatoria se escribe
\((2\kappa c)^{-1}\int_M(R-2\Lambda_{\rm int})\operatorname{vol}_e\).
Su variación interior es
\[
 \frac1{2\kappa c}\int_M\operatorname{vol}_e\,
        (G_{\mu\nu}+\Lambda_{\rm int}g_{\mu\nu})\delta g^{\mu\nu}.
\]
La suma con \eqref{x_reciprocidad:vii:eq:fuentes-normalizadas-accion} da
\eqref{x_reciprocidad:vii:eq:ecuacion-metrica-efectiva}. La identidad de Bianchi
contraída para \(\mathring\nabla\), la compatibilidad métrica y
la constancia de \(\kappa\) y \(\Lambda_{\rm int}\) en la carta
producen \eqref{x_reciprocidad:vii:eq:balance-metrico-total}.
\end{proof}

Antes de eliminar la conexión, las identidades geométricas conservan
la forma
\begin{equation}
 D_\omega T^I=F^I{}_{J}\wedge e^J,\qquad D_\omega F^{IJ}=0.
 \label{x_reciprocidad:vii:eq:bianchi-cartan}
\end{equation}
La primera es \(D_\omega^2e=F\wedge e\); la segunda se obtiene
expandiendo \(D_\omega(\mathrm d\omega+\omega\wedge\omega)\)
y cancelando los términos. En una acción material covariante, el
balance métrico total coincide con la identidad de Noether sobre las
ecuaciones materiales. Geometría y materia conservan así una fuente
conjunta; el reparto entre su contribución de Levi--Civita y su
contribución de espín queda determinado por la acción especializada.

% END OWNER

\end{HMTCompactSection}

% BEGIN OWNER /Users/ruben/Documents/New project/output/EXTREMA_MEDIA_RAZON_RECIPROCIDAD_ES_EN_20260918/ARTICULO/ES/source/antecedentes_vii/32_eliminacion_torsional_no_lineal.tex
% FORMALIZACION_NUEVA / CERTIFICADO_NUEVO, 10-09-2026.
% Composición de las operaciones efectivas ya reunidas en 30e y 31.
% No modifica la selección de rutas, la frontera, los acoplamientos ni
% la normalización de la acción; no evalúa una amplitud cosmológica s_0.
% Propietarios y alcance: desarrollo/COMPOSICION_METRICA_REV06.md.
% Control independiente: pruebas/verificar_eliminacion_torsional_no_lineal.py.

\chapter{Eliminación torsional y respuesta métrica de la extensión mínima}
\label{x_reciprocidad:vii:sec:eliminacion-torsional-no-lineal}

La extensión mínima proporciona un campo, un peso de respuesta y una
corriente mediante una sola operación variacional. Su incidencia
transportada depende de la conexión a través de los productos de ruta de
\eqref{x_reciprocidad:vii:eq:derivada-incidencia-rutas}. Esa dependencia se conserva al
componer la corriente con la ecuación de Cartan--Holst. La eliminación
de la conexión se realiza, por tanto, sobre la acción completa de esa
extensión. El caso afín de
\ref{x_reciprocidad:vii:thm:respuesta-cuadratica} queda contenido como una
especialización exacta de la construcción siguiente.

\section{Composición variacional en un nivel celular}
\label{x_reciprocidad:vii:subsec:composicion-torsional-celular}

Fijemos un nivel finito, sus rutas, la representación de transporte y
el emparejamiento celular de
\ref{x_reciprocidad:vii:thm:componentes-minimo-conexion}. Sean \(p\) coordenadas
reales de la cotetrada y de los campos materiales en una trivialización
local declarada, y \(k\in\mathbb R^d\) las coordenadas de la
contorsión \(K\). La familia \(p\mapsto e(p)\) conserva la
no degeneración de la cotetrada. La conexión que actúa en cada ruta es
\begin{equation}
 \omega(p,k)=\mathring\omega(e(p))+K(k),\qquad
 F(p,k)=S_{\min}[e(p),\omega(p,k),\Psi(p)]
       =\mathfrak a(p,k)E(p,k).
 \label{x_reciprocidad:vii:eq:accion-minima-en-contorsion}
\end{equation}
Los espacios y las bases de esta expresión se transportan a espacios
fijos antes de derivar. Se trabaja en un abierto donde la incidencia
\(\mathsf C(p,k)\) es sobreyectiva y las familias son \(C^2\).
Esta regularidad adicional permite examinar el Hessiano de la acción.
Los acoplamientos \(G_{\rm int},c_{\rm int},\gamma_{\rm HMT}\)
conservan la carta y la convención de variación de
\ref{x_reciprocidad:vii:subsec:dominio-corriente-cartan}.

Sea \(Q(p,k)\) la integral, o la realización celular declarada, de
la forma cuadrática
\eqref{x_reciprocidad:vii:eq:forma-cuadratica-contorsion}. En bases reales,
\begin{equation}
 Q(p,k)=\tfrac12 k^{\mathsf T}\mathsf A(p)k,
 \qquad \mathsf A(p)^{\mathsf T}=\mathsf A(p).
 \label{x_reciprocidad:vii:eq:cuadratica-celular-contorsion}
\end{equation}
El operador \(\mathsf A\) procede de la acción gravitatoria y de
su emparejamiento, mientras que el Gramiano
\(\mathsf C\mathsf C^*\) procede del problema de extensión.
La positividad del segundo garantiza la existencia del campo
minimizante; la forma gravitatoria conserva su propia signatura.
Los términos de borde se mantienen según
\eqref{x_reciprocidad:vii:eq:borde-contorsion}; las fórmulas interiores siguientes
se aplican con variaciones admisibles que anulan su contribución, o
con la completación de borde correspondiente.

Definamos el covector de respuesta mediante
\begin{equation}
 d_kF(p,k)[h]= -\tfrac12 h^{\mathsf T}s(p,k),
 \qquad s(p,k)=\mathsf M_{\rm cel}(p,k)\sigma(p,k).
 \label{x_reciprocidad:vii:eq:covector-torsional-minimo}
\end{equation}
Aquí \(s\) y \(\sigma\) conservan los tipos distintos de
\eqref{x_reciprocidad:vii:eq:componentes-minimo-conexion}: el primero contiene los
coeficientes del diferencial y el segundo las coordenadas de la
corriente en el emparejamiento material. Si se cambia de base en
\(k\), también se transporta ese emparejamiento. Al diferenciar
\(s=\mathsf M_{\rm cel}\sigma\), se incluyen las derivadas de
ambos factores cuando el emparejamiento depende del parámetro.

\begin{proposition}[Ecuación compuesta de la extensión y la contorsión]
\label{x_reciprocidad:vii:prop:ecuacion-torsional-no-lineal}
La estacionariedad interior de la acción respecto de la contorsión es
\begin{equation}
 \mathcal R(p,k):=\mathsf A(p)k-\tfrac12s(p,k)=0.
 \label{x_reciprocidad:vii:eq:ecuacion-torsional-no-lineal}
\end{equation}
Cada componente de \(s\) se calcula a partir de la misma extensión:
\begin{equation}
 \begin{split}
 s_j={}&4\mathfrak a\chi\operatorname{Re}
       \langle y,(\partial_{k_j}\mathsf C)f_{\min}
                         -\partial_{k_j}b\rangle\\
 &-2\mathfrak a(\partial_{k_j}\chi)\langle b,y\rangle
       -2(\partial_{k_j}\mathfrak a)E.
 \end{split}
 \label{x_reciprocidad:vii:eq:seleccion-variacional-corriente-k}
\end{equation}
Cuando el emparejamiento celular realiza la ecuación de conexión de
\ref{x_reciprocidad:vii:thm:ecuacion-cartan-holst}, esta ecuación equivale a componer
\(\sigma(p,k)\) con la inversa de Holst, la inversa de
\(\mathsf L_e\) y la reconstrucción de contorsión:
\begin{equation}
 K=\mathsf C_e^{\rm tor}\,\mathsf L_e^{-1}
           \bigl(\kappa c\,\mathsf P_\gamma^{-1}
             \sigma[e,\mathring\omega(e)+K,\Psi]\bigr).
 \label{x_reciprocidad:vii:eq:composicion-implicita-cartan}
\end{equation}
El símbolo \(\mathsf C_e^{\rm tor}\) designa exclusivamente la
aplicación \(T\mapsto K\) de
\eqref{x_reciprocidad:vii:eq:contorsion-componentes}, con cotetrada fija.
\end{proposition}

\begin{proof}
El diferencial de \(Q+F\), en la dirección arbitraria \(h\), es
\[
 d_k(Q+F)[h]
   =h^{\mathsf T}\bigl(\mathsf A k-\tfrac12s\bigr).
\]
Su anulación para toda dirección prueba
\eqref{x_reciprocidad:vii:eq:ecuacion-torsional-no-lineal}.
El teorema~\ref{x_reciprocidad:vii:thm:componentes-minimo-conexion}, aplicado a
\(\partial_{k_j}\), da
\eqref{x_reciprocidad:vii:eq:seleccion-variacional-corriente-k}. Con \(p\) fijo,
\(\delta\omega=\delta K\). La ecuación de Cartan se transforma,
sucesivamente, mediante \eqref{x_reciprocidad:vii:eq:inversa-holst},
\eqref{x_reciprocidad:vii:eq:torsion-explicita} y
\eqref{x_reciprocidad:vii:eq:contorsion-componentes}; cada una de estas operaciones
es invertible en su dominio. Esto prueba la equivalencia con
\eqref{x_reciprocidad:vii:eq:composicion-implicita-cartan} cuando ambas expresiones
usan la misma realización material y celular.
\end{proof}

La composición conserva la conexión en los dos miembros. Las rutas
que recorren varias celdas pueden acoplar las componentes de la
corriente entre ellas. La reconstrucción puntual de \(T\) a partir
de una corriente dada y la resolución conjunta de
\eqref{x_reciprocidad:vii:eq:ecuacion-torsional-no-lineal} son, en ese caso,
operaciones sucesivas. Esta formulación finita mantiene explícita
la representación de rutas que interviene en el funcional.

\section{Existencia local de una rama estacionaria}
\label{x_reciprocidad:vii:subsec:rama-torsional-local}

\begin{theorem}[Continuación local de la solución torsional]
\label{x_reciprocidad:vii:thm:rama-torsional-local}
Supongamos que \((p_0,k_0)\) pertenece al abierto anterior,
\(\mathcal R(p_0,k_0)=0\), y que el Hessiano total
\begin{equation}
 \mathsf H_0
 =\mathsf A(p_0)+\partial_k^2F(p_0,k_0)
 =\mathsf A(p_0)-\tfrac12\partial_ks(p_0,k_0)
 \label{x_reciprocidad:vii:eq:hessiano-total-torsional}
\end{equation}
es invertible. Existen entornos de \(p_0\) y \(k_0\) en los que
una única rama \(C^1\), \(k=k_*(p)\), satisface
\eqref{x_reciprocidad:vii:eq:ecuacion-torsional-no-lineal} y pasa por \(k_0\).
En esa rama,
\begin{equation}
 \partial_{p_i}k_*
 =-\bigl[\mathsf A+\partial_k^2F\bigr]^{-1}
     \left[(\partial_{p_i}\mathsf A)k_*
                      +\partial_{p_i}\partial_kF\right].
 \label{x_reciprocidad:vii:eq:derivada-rama-torsional}
\end{equation}
\end{theorem}

\begin{proof}
La inversión del Gramiano es \(C^2\) en el abierto de rango
completo. Por ello \(F\) es \(C^2\) y \(\mathcal R\) es
\(C^1\). Su derivada respecto de \(k\), en \((p_0,k_0)\), es
\(\mathsf H_0\). El teorema de la función implícita proporciona
la rama, su regularidad y su unicidad en un entorno del punto.
También puede verse la existencia mediante la aplicación
\[
 \mathcal T_p(k)=k-\mathsf H_0^{-1}\mathcal R(p,k).
\]
Su derivada respecto de \(k\) se anula en \((p_0,k_0)\).
En una bola suficientemente pequeña tiene norma a lo sumo
\(q<1\). Reduciendo el entorno de \(p_0\), se consigue
\(\|\mathcal T_p(k_0)-k_0\|<(1-q)r\), donde \(r\) es el
radio de la bola. Así, \(\mathcal T_p\) lleva la bola cerrada
en sí misma y es contractiva. Su punto fijo es la solución local.
Diferenciar \(\mathcal R(p,k_*(p))=0\) y despejar
\(\partial_{p_i}k_*\) da
\eqref{x_reciprocidad:vii:eq:derivada-rama-torsional}.
\end{proof}

La hipótesis se verifica sobre la acción compuesta: el rango del
Gramiano asegura el mínimo de pantalla y el Hessiano total controla
la rama torsional. La unicidad afirmada es local alrededor de la
solución especificada; otras ramas pueden coexistir en el mismo
dominio de rango completo.

\section{Acción efectiva y contribución no afín}
\label{x_reciprocidad:vii:subsec:accion-efectiva-no-afin}

\begin{theorem}[Identidad exacta de eliminación]
\label{x_reciprocidad:vii:thm:eliminacion-torsional-no-lineal}
Sea \(k_*(p)\) una rama estacionaria. Supongamos, además, que el
segmento \(t\mapsto (p,tk_*)\), \(0\leq t\leq1\), permanece
en el abierto donde está definido \(F\). La corrección de acción
respecto de la realización de Levi--Civita es
\begin{equation}
 \begin{split}
 \Delta S(p)
  &:=Q(p,k_*)+F(p,k_*)-F(p,0)\\
  &=\tfrac14 k_*^{\mathsf T}s(p,k_*)
     -\tfrac12\int_0^1
           k_*^{\mathsf T}s(p,tk_*)\,dt.
 \end{split}
 \label{x_reciprocidad:vii:eq:accion-efectiva-no-lineal}
\end{equation}
Equivalentemente,
\begin{equation}
 \Delta S=-\tfrac14 k_*^{\mathsf T}s(p,k_*)
  +\tfrac12\int_0^1 k_*^{\mathsf T}
       \bigl[s(p,k_*)-s(p,tk_*)\bigr]\,dt.
 \label{x_reciprocidad:vii:eq:correccion-no-afin-integrada}
\end{equation}
Estas identidades se refieren a la acción integrada o celular.
En una realización con densidad local, se aplican a la densidad
con su emparejamiento de formas correspondiente. La acción efectiva
conserva separadamente el término de borde declarado.
\end{theorem}

\begin{proof}
Por la regla de la cadena y
\eqref{x_reciprocidad:vii:eq:covector-torsional-minimo},
\[
 \frac{d}{dt}F(p,tk_*)
      =-\tfrac12 k_*^{\mathsf T}s(p,tk_*).
\]
Integrar entre \(0\) y \(1\) da la diferencia de acciones
materiales. La ecuación de estacionariedad, evaluada en la
dirección \(k_*\), y la homogeneidad cuadrática de \(Q\) dan
\[
 2Q(p,k_*)=\tfrac12 k_*^{\mathsf T}s(p,k_*).
\]
La suma de ambas identidades demuestra
\eqref{x_reciprocidad:vii:eq:accion-efectiva-no-lineal}; sumar y restar
\(\tfrac12k_*^{\mathsf T}s(p,k_*)\) proporciona
\eqref{x_reciprocidad:vii:eq:correccion-no-afin-integrada}.
\end{proof}

\begin{corollary}[Especialización afín]
\label{x_reciprocidad:vii:cor:recuperacion-eliminacion-afin}
Si \(s(p,k)=s(p)\), entonces
\[
 \Delta S=-\tfrac14 k_*^{\mathsf T}s(p)=-Q(p,k_*).
\]
Es la realización integrada de
\eqref{x_reciprocidad:vii:eq:accion-efectiva-spin}.
\end{corollary}
\begin{proof}
La diferencia dentro de la integral de
\eqref{x_reciprocidad:vii:eq:correccion-no-afin-integrada} se anula. La segunda
igualdad procede de la identidad radial de estacionariedad.
\end{proof}

\section{Diferencial métrico de la acción eliminada}
\label{x_reciprocidad:vii:subsec:diferencial-metrico-no-lineal}

\begin{theorem}[Respuesta de las variables geométricas]
\label{x_reciprocidad:vii:thm:envolvente-metrica}
Sobre la rama del teorema~\ref{x_reciprocidad:vii:thm:rama-torsional-local},
\begin{equation}
 \partial_{p_i}\Delta S
 =\tfrac12 k_*^{\mathsf T}(\partial_{p_i}\mathsf A)k_*
  +(\partial_{p_i}F)(p,k_*)-(\partial_{p_i}F)(p,0),
 \label{x_reciprocidad:vii:eq:envolvente-metrica-no-lineal}
\end{equation}
donde las derivadas parciales mantienen fijas las coordenadas \(k\).
En productos hermíticos fijos, cada derivada material se calcula por
\begin{equation}
 \begin{split}
 \partial_{p_i}F={}&(\partial_{p_i}\mathfrak a)E
       +\mathfrak a(\partial_{p_i}\chi)\langle b,y\rangle\\
   &+2\mathfrak a\chi\operatorname{Re}
       \langle y,\partial_{p_i}b
                    -(\partial_{p_i}\mathsf C)f_{\min}\rangle.
 \end{split}
 \label{x_reciprocidad:vii:eq:diferencial-geometrico-minimo}
\end{equation}
Si \(p_i\) varía la cotetrada, \(\partial_{p_i}\mathsf C\)
incluye la variación de
\(\mathring\omega(e(p))+K(k)\), de las rutas realizadas y de
las trivializaciones utilizadas, según sus dependencias declaradas.
\end{theorem}

\begin{proof}
La derivada total de \(Q(p,k_*)+F(p,k_*)\) contiene, además
de sus derivadas parciales, el término
\[
 \bigl[\partial_kQ(p,k_*)+\partial_kF(p,k_*)\bigr]
                      \partial_{p_i}k_*.
\]
El corchete es cero por la ecuación de conexión. Restar la derivada
de \(F(p,0)\) prueba
\eqref{x_reciprocidad:vii:eq:envolvente-metrica-no-lineal}.
La regla del producto y
\eqref{x_reciprocidad:vii:eq:diferencial-minimo-completo} prueban
\eqref{x_reciprocidad:vii:eq:diferencial-geometrico-minimo}. La composición
\eqref{x_reciprocidad:vii:eq:accion-minima-en-contorsion} obliga a diferenciar
la conexión de Levi--Civita al variar \(e\) con \(k\) fijo.
\end{proof}

Cuando el producto positivo de extensión también varía, su
contribución se conserva de forma explícita. En coordenadas, sea
\(\mathsf h(p,k)>0\) su matriz hermítica y considérese el mínimo
de \(f^*\mathsf h f\) bajo \(\mathsf C f=b\). Entonces
\begin{equation}
 \mathsf L=\mathsf C\mathsf h^{-1}\mathsf C^*,\quad
 y=\mathsf L^{-1}b,\quad
 f_{\min}=\mathsf h^{-1}\mathsf C^*y,\quad
 E=\chi b^*y.
 \label{x_reciprocidad:vii:eq:minimo-producto-variable}
\end{equation}
Aquí \(y\) representa el covector de la restricción en bases
fijas. La matriz positiva \(\mathsf h\) conserva su papel
energético, distinto del de la métrica lorentziana.

\begin{lemma}[Variación del producto de extensión]
\label{x_reciprocidad:vii:lem:producto-variable-extension}
Con las convenciones anteriores,
\begin{equation}
 \delta E=(\delta\chi)b^*y
   +2\chi\operatorname{Re}\bigl[y^*(\delta b-\delta\mathsf C f_{\min})\bigr]
   +\chi f_{\min}^*(\delta\mathsf h)f_{\min}.
 \label{x_reciprocidad:vii:eq:diferencial-producto-variable}
\end{equation}
\end{lemma}
\begin{proof}
La fórmula \(\delta\mathsf h^{-1}
=-\mathsf h^{-1}(\delta\mathsf h)\mathsf h^{-1}\) da
\[
 \delta\mathsf L=(\delta\mathsf C)\mathsf h^{-1}\mathsf C^*
 +\mathsf C\mathsf h^{-1}(\delta\mathsf C)^*
 -\mathsf C\mathsf h^{-1}(\delta\mathsf h)
                              \mathsf h^{-1}\mathsf C^*.
\]
Al diferenciar \(E=\chi b^*\mathsf L^{-1}b\), el término
\(-\chi y^*(\delta\mathsf L)y\) produce los dos términos
con \(\delta\mathsf C\) y el término positivo con
\(\delta\mathsf h\). Esto prueba la identidad completa.
\end{proof}

En una realización material covariante cuya variación admita el
emparejamiento métrico de
\eqref{x_reciprocidad:vii:eq:fuentes-normalizadas-accion}, la corrección determina
\begin{equation}
 \delta_g\Delta S
 =-\tfrac12\int_M\operatorname{vol}_e\,
          \mathcal U^{S,\min}_{\mu\nu}\,\delta g^{\mu\nu},
 \qquad U^{\rm phys,\min}_{\mu\nu}
                      =c\mathcal U^{S,\min}_{\mu\nu}.
 \label{x_reciprocidad:vii:eq:fuente-metrica-minimo-no-lineal}
\end{equation}
Las ecuaciones \eqref{x_reciprocidad:vii:eq:envolvente-metrica-no-lineal} y
\eqref{x_reciprocidad:vii:eq:diferencial-geometrico-minimo}, completadas por
\eqref{x_reciprocidad:vii:eq:diferencial-producto-variable} cuando corresponde,
evalúan sus coeficientes en las variaciones métricas declaradas.
El mínimo de pantalla, la solución torsional y la derivada métrica
son tres operaciones consecutivas sobre la misma acción.
La acción reducida conserva los posibles acoplamientos entre
celdas presentes en las rutas originales.

La especialización homogénea de
\ref{x_reciprocidad:vii:sec:dilucion-rebote} utiliza además la realización del
tensor obtenido: su simetría espacial, el signo de la contracción
con el observador y la ley de escala. Cuando esa evaluación da
\(\rho_{\rm corr}(a)=-s_0a^{-6}\), la amplitud se lee en el
estado de referencia como
\(s_0=-a_{\rm ref}^{6}\rho_{\rm corr}(a_{\rm ref})\).
Esta lectura conserva el estado material y su normalización; el
teorema de eliminación proporciona la acción que debe contraerse.

\section{Modelo exacto de comprobación no afín}
\label{x_reciprocidad:vii:subsec:control-no-afin}

El siguiente modelo finito comprueba los coeficientes y la
dependencia no afín de las identidades. Sus parámetros pertenecen
al ejemplo algebraico y mantienen separado el problema de
realización material de las rutas.

\begin{proposition}[Eliminación exacta de una extensión de rango uno]
\label{x_reciprocidad:vii:prop:ejemplo-no-afin}
Sean \(H=\mathbb R^2\), \(Q=\mathbb R\),
\(\mathsf C(k)=(1\ \ k)\), \(b=1\), \(\chi=1\),
\(\mathfrak a=\lambda>0\), y \(Q_{\rm ex}(k)=k^2/2\).
Entonces
\begin{equation}
 f_{\min}=\frac{(1,k)^{\mathsf T}}{1+k^2},\quad
 F(\lambda,k)=\frac{\lambda}{1+k^2},\quad
 s(\lambda,k)=\frac{4\lambda k}{(1+k^2)^2},\quad
 \mathcal R(\lambda,k)
    =k\left(1-\frac{2\lambda}{(1+k^2)^2}\right).
 \label{x_reciprocidad:vii:eq:ejemplo-minimo-no-afin}
\end{equation}
Para \(\lambda=25/2\), las soluciones estacionarias son
\(0,2,-2\). En \(k_*=2\),
\begin{equation}
 \mathsf H_*=\frac{16}{5},\quad s_*=4,\quad
 Q_{\rm ex}(k_*)=2,\quad F(\lambda,k_*)=\frac52,
 \quad \Delta S=-8.
 \label{x_reciprocidad:vii:eq:valores-ejemplo-no-afin}
\end{equation}
La expresión afín \(-k_*s_*/4=-2\) difiere de la corrección
exacta. En la rama positiva,
\begin{equation}
 \partial_\lambda\Delta S=-\frac45,
 \qquad \partial_\lambda k_* =\frac1{20}
 \quad\text{en }\lambda=25/2.
 \label{x_reciprocidad:vii:eq:derivadas-ejemplo-no-afin}
\end{equation}
\end{proposition}

\begin{proof}
El Gramiano es \(1+k^2>0\), por lo que el mínimo tiene las
expresiones indicadas. Diferenciar \(F\) da
\(s=-2\partial_kF\) y después \(\mathcal R=k-s/2\).
Para \(\lambda=25/2\), la ecuación es
\(k[(1+k^2)^2-25]=0\); como \(1+k^2>0\), sus raíces
son precisamente \(0,\pm2\). El Hessiano total es
\[
 1-\frac{2\lambda(1-3k^2)}{(1+k^2)^3},
\]
que vale \(16/5\) en ambas raíces no nulas. La corrección es
\(2+5/2-25/2=-8\). En la identidad radial,
\[
 \int_0^1 k_*s(\lambda,tk_*)\,dt
    =\int_0^1\frac{200t}{(1+4t^2)^2}\,dt=20,
\]
pues una primitiva es \(-25/(1+4t^2)\).
Así, \(\Delta S=8/4-20/2=-8\), incluido el término no afín.
Para \(\lambda>1/2\), la rama positiva satisface
\(k_*^2=\sqrt{2\lambda}-1\) y
\(\Delta S=\sqrt{2\lambda}-1/2-\lambda\).
Sus derivadas dan \eqref{x_reciprocidad:vii:eq:derivadas-ejemplo-no-afin}.
La fórmula de envolvente proporciona el mismo resultado:
\((\partial_\lambda F)(\lambda,k_*)
 -(\partial_\lambda F)(\lambda,0)=1/5-1=-4/5\).
\end{proof}

% END OWNER


% BEGIN OWNER /Users/ruben/Documents/New project/output/EXTREMA_MEDIA_RAZON_RECIPROCIDAD_ES_EN_20260918/ARTICULO/ES/source/antecedentes_vii/33_corriente_espinorial_y_densidad.tex
% Artículo VII. Incorporación aditiva autorizada, 10-09-2026.
% RESULTADO_RECUPERADO: representación real de Clifford y CAR:
% TEORIA_HOLOGRAFICA_INTEGRAL_MASA_HMT_MD_REV2_2026-08-20/fuente/modulos/
% 18_cuantica_estadistica_y_termodinamica.md, líneas 131--181 y 218--255;
% acción de primer orden: 14_topologia_exoticos_y_gravedad.md, 459--517;
% distinción corriente media/segundo momento y dilución:
% HMT_OBRAS/OBRA_II/chapters/11_cartan_bianchi_cosmologia.tex, 201--230.
% PDF REV2: b56bb38d10b9a4f4d0de47b7da8010fb31a6e8e8684a3bfd4237e37d7a3dddfa.
% FORMALIZACION_NUEVA: acción espinorial afín explícita, contracción
% normalizada mediante 31, observable cuártico CAR y funcional de amplitud.
% CERTIFICADO_NUEVO: cálculo racional de la forma cuadrática en las cuatro
% componentes de una trivector y evaluación de la realización CAR finita.
% Convención explicitada: las matrices g tienen firma (-+++); A_D=-i g^0
% y la acción cinética simétrica lleva hbar/2. La combinación equivalente
% i hbar Gamma^I usa Gamma^I=-i g^I y firma opuesta. No se mezclan ambas.
% Dependencias: cotetrada/conexión de 30d; constantes internas heredadas;
% corriente e inversas de 31. Estado material, celda y normal orden son
% datos de la realización declarada, no constantes universales extraídas.
% La acción afín de este fragmento y la acción de extensión mínima de 30e
% conservan sus problemas variacionales propios. No se identifican sus
% corrientes. No se cambia 50: su dominio de amplitud positiva se concreta.

\chapter{Corriente espinorial y amplitud torsional del estado material}
\label{x_reciprocidad:vii:sec:corriente-espinorial-densidad}

El desarrollo precedente y la realización espinorial que sigue articulan
dos problemas variacionales con antecedentes geométricos comunes. La
acción de extensión mínima \(S_{\min}=\mathfrak a_m E\), construida en
la sección~\ref{x_reciprocidad:vii:sec:corriente-extension-minima}, incorpora la
dependencia de la incidencia transportada y de su mínimo respecto de la
conexión; su eliminación torsional conserva esa dependencia completa en
la sección~\ref{x_reciprocidad:vii:sec:eliminacion-torsional-no-lineal}. La acción
espinorial \(S_D\) de \eqref{x_reciprocidad:vii:eq:accion-dirac-material} acopla la
misma conexión a los campos de su representación y es afín en la
contorsión al fijar la cotetrada y los campos. Cada corriente se obtiene
variando su propia acción. Una preparación minimizante puede determinar
el estado material inicial de la realización espinorial; la sustitución
de campos dependientes de la conexión conserva, además, los términos
del diferencial compuesto. Esta distinción sitúa el cálculo de amplitud
siguiente en su acción y en su dominio variacional específicos.

La representación espinorial del transporte y la respuesta algebraica
de Cartan--Holst permiten calcular la intensidad torsional como un
funcional del estado material. Su evaluación conserva tres datos:
la representación de los campos, el segundo momento de su corriente y
el volumen respecto del cual se define la densidad. El coeficiente de
acoplamiento resulta de la acción y de su inversión; el estado inicial
determina la amplitud material.

\section{Representación espinorial y acción material afín}
\label{x_reciprocidad:vii:subsec:accion-espinorial-afin}

Conservamos la cotetrada, la orientación y los acoplamientos de la
sección~\ref{x_reciprocidad:vii:sec:corriente-respuesta-cartan}. La realización
espinorial utiliza el módulo complejo de dimensión cuatro obtenido al
complejificar las matrices reales
\begin{equation}
 \begin{gathered}
 J=\begin{pmatrix}0&-1\\1&0\end{pmatrix},\qquad
 R=\begin{pmatrix}0&1\\1&0\end{pmatrix},\qquad
 Z=\begin{pmatrix}1&0\\0&-1\end{pmatrix},\\
 g^0=J\otimes I_2,\qquad g^1=R\otimes I_2,\qquad
 g^2=Z\otimes R,\qquad g^3=Z\otimes Z.
 \end{gathered}
 \label{x_reciprocidad:vii:eq:clifford-material-explicito}
\end{equation}
La multiplicación da
\(\{g^I,g^J\}=2\eta^{IJ}I_4\), con
\(\eta=\operatorname{diag}(-1,1,1,1)\). Las matrices se distinguen
del parámetro escalar \(\gamma\) del operador de Holst. Definimos
\begin{equation}
 g_I=\eta_{IJ}g^J,\quad
 \mathbb B_{IJ}=\tfrac14[g_I,g_J],\quad
 A_D=-i g^0,\quad \bar\psi=\psi^\dagger A_D,\quad
 \Omega=\tfrac12\omega^{IJ}\mathbb B_{IJ}.
 \label{x_reciprocidad:vii:eq:adjunto-conexion-espinorial}
\end{equation}
El producto de Clifford verifica
\[
 [\mathbb B_{IJ},g^K]=\delta_J^K g_I-\delta_I^K g_J.
\]
Así, la misma conexión métrica actúa en los campos mediante
\(D\psi=\mathrm d\psi+\Omega\psi\) y
\(D\bar\psi=\mathrm d\bar\psi-\bar\psi\Omega\).
El transporte sobre una ruta se obtiene integrando esta conexión en
su representación espinorial; conserva la conexión y la ruta de la
realización geométrica precedente.

Sea \(\eta_I=\iota_{E_I}\operatorname{vol}_e\), de modo que
\(e^J\wedge\eta_I=\delta_I^J\operatorname{vol}_e\). Para un
campo de tipo dimensional \([\psi]=L^{-3/2}\), tomamos la acción
material mínima
\begin{equation}
 S_D=\int_M\left\{
 \frac{\hbar}{2}
 \bigl[\bar\psi g^I D\psi-(D\bar\psi)g^I\psi\bigr]\wedge\eta_I
 -mc\,\bar\psi\psi\operatorname{vol}_e\right\}.
 \label{x_reciprocidad:vii:eq:accion-dirac-material}
\end{equation}
Aquí \(\hbar=\hbar_{\rm int}\), \(c=c_{\rm int}\) y la masa
\(m\) es la lectura de la especie material en la carta elegida.
Estos valores son salidas anteriores a esta realización. La matriz
\(A_D\) es hermítica y \(A_Dg^I\) es antihermítica. Por ello el
término cinético simétrico tiene la convención real correspondiente a
\eqref{x_reciprocidad:vii:eq:clifford-material-explicito}. En esta firma, la ecuación
sin torsión tiene término principal \(\hbar g^I D_I\); la notación
equivalente \(i\hbar\Gamma^I D_I\), con \(\Gamma^I=-i g^I\),
emplea la relación de Clifford de firma opuesta. Esta precisión fija
conjuntamente adjunto, firma y factor imaginario de la acción.

Los campos materiales y la cotetrada permanecen fijos en la variación
de la conexión. Una preparación procedente de la extensión mínima de
la sección~\ref{x_reciprocidad:vii:sec:corriente-extension-minima} puede determinar
el dato inicial del campo o del estado. Cada acción conserva su
variación propia: si se sustituye en una acción un campo minimizante
dependiente de la conexión, se aplica además su diferencial compuesto,
o el teorema de la envolvente bajo sus hipótesis de estacionariedad.
La acción~\eqref{x_reciprocidad:vii:eq:accion-dirac-material} realiza aquí un
acoplamiento afín en la contorsión.

\begin{proposition}[Corriente espinorial con normalización variacional]
\label{x_reciprocidad:vii:prop:corriente-dirac-explicita}
Sea
\(B^{IJK}=\bar\psi g^{[I}g^Jg^{K]}\psi\), donde los corchetes
designan la antisimetrización normalizada. La corriente definida en
\eqref{x_reciprocidad:vii:eq:corriente-variacional} por la acción
\eqref{x_reciprocidad:vii:eq:accion-dirac-material} es
\begin{equation}
 \sigma_{JK}
 =-\frac{\hbar}{2}\,
   \bar\psi\{g^I,\mathbb B_{JK}\}\psi\,\eta_I
 =-\frac{\hbar}{2}\,B^I{}_{JK}\eta_I.
 \label{x_reciprocidad:vii:eq:corriente-dirac-trivector}
\end{equation}
Es independiente de la contorsión cuando \(e\) y \(\psi\) están
fijos. La inversión de la sección~\ref{x_reciprocidad:vii:subsec:reconstruccion-torsion}
se aplica por tanto a esta corriente.
\end{proposition}
\begin{proof}
La variación de \(\Omega\) es
\(\delta\Omega=\frac12\delta\omega^{JK}\mathbb B_{JK}\).
Los dos sumandos cinéticos dan, respectivamente, el producto con
\(g^I\) a la izquierda y a la derecha. En consecuencia,
\[
 \delta_\omega S_D
 =\frac{\hbar}{4}\int_M\delta\omega^{JK}\wedge\eta_I\,
                 \bar\psi\{g^I,\mathbb B_{JK}\}\psi.
\]
La comparación con \eqref{x_reciprocidad:vii:eq:corriente-variacional} fija el
factor y el signo de \eqref{x_reciprocidad:vii:eq:corriente-dirac-trivector}.
Al expandir el anticonmutador y usar la relación de Clifford se
cancelan las contracciones de grado uno, y queda
\(\{g^I,\mathbb B_{JK}\}=g^{[I}g_Jg_{K]}\).
La acción es lineal en \(\Omega\), lo que prueba la última
afirmación.
\end{proof}

\section{Contracción efectiva y coeficiente del segundo momento}
\label{x_reciprocidad:vii:subsec:contraccion-espinorial}

La forma cuadrática lorentziana del trivector se escribe
\begin{equation}
 q(B)=\frac1{3!}B^{IJK}B_{IJK}
 =-(B^{012})^2-(B^{013})^2-(B^{023})^2+(B^{123})^2.
 \label{x_reciprocidad:vii:eq:cuadratica-trivector}
\end{equation}
Esta contracción conserva la firma del módulo espinorial. Su segundo
momento depende del estado y de la realización de los productos de
campos.

\begin{theorem}[Interacción efectiva de la corriente espinorial]
\label{x_reciprocidad:vii:thm:coeficiente-torsional-espinorial}
Con la acción material~\eqref{x_reciprocidad:vii:eq:accion-dirac-material}, la
eliminación algebraica de la contorsión da
\begin{equation}
 \mathcal L_{\rm spin}^{\rm eff}
 =C_\gamma q(B)\operatorname{vol}_e,\qquad
 C_\gamma=\frac{3\kappa c\hbar^2}{16}
                      \frac{\gamma^2}{1+\gamma^2}.
 \label{x_reciprocidad:vii:eq:coeficiente-cuadratico-espinorial}
\end{equation}
El coeficiente se obtiene de la corriente y de la inversa de Holst
con las normalizaciones de la acción precedente.
\end{theorem}
\begin{proof}
Por \eqref{x_reciprocidad:vii:eq:accion-efectiva-spin}, basta evaluar
\(-\frac14K_*^{IJ}\wedge\sigma_{IJ}\), con
\[
 Y^{IJ}=\kappa c\frac{\gamma^2}{1+\gamma^2}
           (-\star+\gamma^{-1}I)\sigma^{IJ},\qquad
 K_*=\mathsf C_e(Y).
\]
Aquí \(\mathsf C_e\) es la composición explícita de
\eqref{x_reciprocidad:vii:eq:torsion-explicita} con
\eqref{x_reciprocidad:vii:eq:contorsion-componentes}. La linealidad de esas dos
inversas reduce el cálculo a las partes \(-\star\) e \(I\).
Para dar la contracción completa, ordenamos las componentes del
trivector como \((012,013,023,123)\), extraemos los factores
\(\kappa c\hbar^2\gamma^2/(1+\gamma^2)\), y sustituimos
\(\sigma_{JK}=-\frac12B^I{}_{JK}\eta_I\). Las dos matrices
de la forma cuadrática restante son
\begin{equation}
 \mathcal B_{-\star}
 =\frac3{16}\operatorname{diag}(-1,-1,-1,1),\qquad
 \mathcal B_I=0_{4\times4}.
 \label{x_reciprocidad:vii:eq:matrices-contraccion-espinorial}
\end{equation}
Estas matrices se obtienen usando únicamente
\(e^I\wedge\eta_J=\delta_J^I\operatorname{vol}_e\): para
cada componente \(B^{abc}\), se forma \(\sigma\), se calcula
\(Y\), después
\(T^I=\sum_J\iota_{E_J}Y^{IJ}
 -\frac14e^I\wedge\sum_{L,J}\iota_{E_L}\iota_{E_J}Y^{LJ}\),
y finalmente
\(K_{IJK}=(T_{KIJ}-T_{IJK}-T_{JKI})/2\).
La contracción da las cuatro entradas diagonales mostradas; las
seis polarizaciones entre componentes distintas son cero para ambos
operadores. Es el cálculo de todos los coeficientes de las dos
formas cuadráticas. La parte \(\gamma^{-1}I\) se anula y la parte
\(-\star\) produce exactamente \eqref{x_reciprocidad:vii:eq:cuadratica-trivector},
con el coeficiente de \eqref{x_reciprocidad:vii:eq:coeficiente-cuadratico-espinorial}.
\end{proof}

\section{Realización fermiónica finita y segundo momento del estado}
\label{x_reciprocidad:vii:subsec:segundo-momento-car}

Fijamos una celda comóvil periódica, un marco ortonormal adaptado a su
observador y el modo espacial homogéneo de las cuatro componentes
espinoriales. Esta realización finita tiene espacio de estados
\(\mathcal F=\bigoplus_{n=0}^4\Lambda^n\mathbb C^4\).
Sean \(a_r,a_r^\dagger\), \(r=1,\ldots,4\), sus operadores
canónicos, con
\(\{a_r,a_s^\dagger\}=\delta_{rs}\) y
\(\{a_r,a_s\}=0\). El vacío de ocupación fija el orden
normal. Para una matriz \(M\), escribimos
\[
 \mathrm d\Gamma(M)=\sum_{r,s}a_r^\dagger M_{rs}a_s,
 \qquad \widehat N=\mathrm d\Gamma(I_4).
\]
La letra \(\Gamma\) de esta notación designa la segunda
cuantización; la holonomía nonádica conserva su notación propia.

Las cuatro matrices de los bilineales se calculan directamente:
\begin{equation}
 \begin{array}{c|cccc}
 abc&012&013&023&123\\ \hline
 M_{abc}=A_Dg^ag^bg^c
       &iJ\otimes R&iJ\otimes Z&iI_2\otimes J&iZ\otimes J
 \end{array}
 \label{x_reciprocidad:vii:eq:matrices-bilineales-espin}
\end{equation}
Todas son hermíticas, tienen traza cero y satisfacen
\(M_{abc}^2=I_4\). El observable cuártico correspondiente a
\eqref{x_reciprocidad:vii:eq:cuadratica-trivector}, en este orden normal declarado,
es
\begin{equation}
 \begin{split}
 \widehat{\mathscr Q}_{\rm sp}
 &=\sum_{a<b<c}\eta_{aa}\eta_{bb}\eta_{cc}
       :\!\mathrm d\Gamma(M_{abc})^2\!:\\
 &=\sum_{a<b<c}\eta_{aa}\eta_{bb}\eta_{cc}
       \mathrm d\Gamma(M_{abc})^2+2\widehat N.
 \end{split}
 \label{x_reciprocidad:vii:eq:observable-cuadratico-espin}
\end{equation}
El orden normal es parte de la realización finita especificada;
su definición hace explícita la sustracción de las contracciones de
una sola ocupación.

\begin{proposition}[Evaluación por las relaciones canónicas]
\label{x_reciprocidad:vii:prop:operador-espin-car}
El operador \(\widehat{\mathscr Q}_{\rm sp}\) es autoadjunto,
conserva cada sector de ocupación y tiene restricciones
\begin{equation}
 \widehat{\mathscr Q}_{\rm sp}|_{\Lambda^0\oplus\Lambda^1}=0,
 \qquad
 \widehat{\mathscr Q}_{\rm sp}|_{\Lambda^3}=4I_4,
 \qquad
 \widehat{\mathscr Q}_{\rm sp}|_{\Lambda^4}=8I_1.
 \label{x_reciprocidad:vii:eq:sectores-espin-car}
\end{equation}
En la base de \(\Lambda^2\mathbb C^4\) cuyos subconjuntos de
ocupación son \((12,13,23,14,24,34)\), su matriz es
\begin{equation}
 \begin{pmatrix}
 0&0&0&0&0&-4\\
 0&2&0&0&6&0\\
 0&0&2&-6&0&0\\
 0&0&-6&2&0&0\\
 0&6&0&0&2&0\\
 -4&0&0&0&0&0
 \end{pmatrix}.
 \label{x_reciprocidad:vii:eq:sector-dos-espin-car}
\end{equation}
\end{proposition}
\begin{proof}
La relación \(a_s a_t^\dagger=\delta_{st}-a_t^\dagger a_s\)
aplicada a \(\mathrm d\Gamma(M)^2\) separa su contracción de
una ocupación y da
\[
 :\!\mathrm d\Gamma(M)^2\!:
 =\mathrm d\Gamma(M)^2-\mathrm d\Gamma(M^2).
\]
Como los tres primeros signos de
\eqref{x_reciprocidad:vii:eq:matrices-bilineales-espin} son negativos y el último
positivo, la suma de las contracciones es \(-2\widehat N\).
Esto prueba \eqref{x_reciprocidad:vii:eq:observable-cuadratico-espin}.
La autoadjunción y la conservación de la ocupación se siguen de
esas mismas expresiones.

En \(\Lambda^0\) todos los sumandos se anulan. En \(\Lambda^1\),
\(\mathrm d\Gamma(M)=M\), y los cuadrados suman \(-2I_4\),
que se cancela con \(2\widehat N\). En \(\Lambda^4\),
\(\mathrm d\Gamma(M)=\operatorname{tr}M=0\), quedando ocho.
La identificación por complemento entre \(\Lambda^3\mathbb C^4\)
y el dual de \(\mathbb C^4\) lleva
\(\mathrm d\Gamma(M)\) a
\((\operatorname{tr}M)I-M^{\mathsf T}=-M^{\mathsf T}\).
Sus cuadrados vuelven a sumar \(-2I_4\), mientras
\(2\widehat N=6I_4\); por tanto la restricción es \(4I_4\).
Finalmente, la aplicación de las cuatro matrices explícitas a los
seis productos exteriores de la base indicada produce
\eqref{x_reciprocidad:vii:eq:sector-dos-espin-car}. Así queda evaluado el operador
en las dieciséis dimensiones de \(\mathcal F\).
\end{proof}

Para un operador de estado \(\varrho\geq0\),
\(\operatorname{Tr}\varrho=1\), definimos
\(q_\varrho=\operatorname{Tr}(\varrho\widehat{\mathscr Q}_{\rm sp})\).
Su dependencia de los segundos momentos es explícita:
\begin{equation}
 q_\varrho=\sum_{a<b<c}\eta_{aa}\eta_{bb}\eta_{cc}
 \left\{
  \langle\widehat B_{abc}\rangle_\varrho^2
  +\operatorname{Var}_\varrho(\widehat B_{abc})
  -\langle\widehat N\rangle_\varrho
 \right\},\qquad
 \widehat B_{abc}=\mathrm d\Gamma(M_{abc}).
 \label{x_reciprocidad:vii:eq:covarianza-espin-material}
\end{equation}
La covarianza y la sustracción normal se conservan incluso cuando
las corrientes medias se anulan. En particular, si \(P_3\) es la
proyección sobre \(\Lambda^3\mathbb C^4\),
\begin{equation}
 \varrho_3=\tfrac14P_3,\qquad
 \langle\widehat B_{abc}\rangle_{\varrho_3}=0,\qquad
 q_{\varrho_3}=4>0.
 \label{x_reciprocidad:vii:eq:testigo-positivo-espin}
\end{equation}
Es un estado positivo de traza uno y constituye un testigo explícito
del dominio positivo del funcional. El sector de dos ocupaciones
contiene también valores negativos; por ejemplo, el vector
\((|12\rangle+|34\rangle)/\sqrt2\) tiene valor \(-4\).
El signo de la intensidad procede así de la contracción y del
estado, y queda separado de la selección de condiciones cosmológicas.

\section{Densidad, presión y conservación de la amplitud comóvil}
\label{x_reciprocidad:vii:subsec:amplitud-espinorial-comovil}

En la realización espacialmente homogénea, fijamos
\[
 e^0=N(t)c\,\mathrm dt,\qquad e^j=a(t)\,\mathrm dx^j,
 \qquad \mathcal V(t)=a(t)^3V_c,
\]
con lapse \(N(t)>0\), factor de escala adimensional \(a(t)>0\)
y volumen comóvil \(V_c>0\). El marco y el modo periódico de la
subsección~\ref{x_reciprocidad:vii:subsec:segundo-momento-car} permanecen declarados.
La normalización del campo homogéneo es
\(\widehat\psi=\mathcal V^{-1/2}(a_1,a_2,a_3,a_4)^{\mathsf T}\).
Cada bilineal lleva un factor \(\mathcal V^{-1}\), y su producto
normal ordenado lleva \(\mathcal V^{-2}\). La forma temporal
cinética simétrica conserva esta normalización canónica: los términos
procedentes de derivar \(\mathcal V^{-1/2}\) se cancelan entre
sus dos sumandos.

\begin{theorem}[Amplitud torsional como funcional del estado]
\label{x_reciprocidad:vii:thm:amplitud-material-explicita}
Manteniendo fijos el estado canónico y la realización del producto
durante las variaciones métricas, el término efectivo de espín tiene
acción reducida
\begin{equation}
 S_{\rm spin}^{\rm hom}
 =\int\!\mathrm dt\,N(t)
       \frac{cC_\gamma}{a(t)^3V_c}\,q_\varrho.
 \label{x_reciprocidad:vii:eq:accion-homogenea-espin}
\end{equation}
Su densidad de energía y su presión isotrópica son
\begin{equation}
 \rho_{\rm spin}=p_{\rm spin}
 =-\frac{cC_\gamma}{V_c^2}\,q_\varrho a^{-6}.
 \label{x_reciprocidad:vii:eq:densidad-presion-espin}
\end{equation}
En cualquier evolución que conserve \(q_\varrho\), la amplitud
constante correspondiente es
\begin{equation}
 \boxed{\displaystyle
 s_0[\varrho,V_c]
 =\frac{3\kappa c^2\hbar^2}{16}
    \frac{\gamma^2}{1+\gamma^2}
    \frac{\operatorname{Tr}
       (\varrho\widehat{\mathscr Q}_{\rm sp})}{V_c^2}.}
 \label{x_reciprocidad:vii:eq:amplitud-s0-espinorial}
\end{equation}
\end{theorem}
\begin{proof}
La integración espacial de
\eqref{x_reciprocidad:vii:eq:coeficiente-cuadratico-espinorial} multiplica el
factor \(\mathcal V^{-2}\) por
\(Nc\mathcal V\,\mathrm dt\), lo que da
\eqref{x_reciprocidad:vii:eq:accion-homogenea-espin}. En variables canónicas,
la variación del lapse define la energía mediante
\(\delta_N S_{\rm m}=-\int E\,\delta N\,\mathrm dt\).
Por tanto,
\[
 E_{\rm spin}=-\frac{cC_\gamma q_\varrho}{a^3V_c},
 \qquad \rho_{\rm spin}=\frac{E_{\rm spin}}{a^3V_c}.
\]
La variación espacial isotrópica satisface
\(\delta_a S_{\rm m}=\int3Na^2V_c p\,\delta a\,\mathrm dt\).
Aplicada a \eqref{x_reciprocidad:vii:eq:accion-homogenea-espin}, da
\[
 3Na^2V_c p_{\rm spin}
 =-\frac{3NcC_\gamma q_\varrho}{a^4V_c}.
\]
Esto prueba simultáneamente el signo, la presión y la densidad de
\eqref{x_reciprocidad:vii:eq:densidad-presion-espin}. Al conservarse
\(q_\varrho\), el coeficiente de \(-a^{-6}\) es constante y
se obtiene \eqref{x_reciprocidad:vii:eq:amplitud-s0-espinorial}. La dimensionalidad
es consistente: \([cC_\gamma]=\text{energía}\,L^3\),
de modo que \(s_0\) tiene dimensión de densidad de energía.
\end{proof}

\begin{corollary}[Dominio conservado de amplitud positiva]
\label{x_reciprocidad:vii:cor:dominio-positivo-conservado}
Sea una evolución unitaria de \(\mathcal F\) que preserve
\(\Lambda^3\mathbb C^4\). Para todo estado inicial de traza uno
soportado en ese sector, se cumple \(q_{\varrho(t)}=4\), y
\begin{equation}
 s_0=\frac{3\kappa c^2\hbar^2}{4V_c^2}
                   \frac{\gamma^2}{1+\gamma^2}>0.
 \label{x_reciprocidad:vii:eq:amplitud-positiva-sector-tres}
\end{equation}
\end{corollary}
\begin{proof}
La invariancia del sector conserva el soporte del estado; en él,
\eqref{x_reciprocidad:vii:eq:sectores-espin-car} da
\(\widehat{\mathscr Q}_{\rm sp}=4I\). La conservación de la
traza fija entonces su esperanza en cuatro, independientemente de
la evolución interna de sus coherencias. La positividad de los
acoplamientos, \(V_c>0\) y \(\gamma\ne0\) prueba la fórmula
y su signo.
\end{proof}

Para estados generales, la conservación de
\(\langle\widehat N\rangle\) determina la dilución de la densidad
de ocupación, \(n(t)=\langle\widehat N\rangle/(a^3V_c)\).
La conservación del momento cuártico se comprueba por separado.
Si \(i\hbar\dot\varrho=[H(t),\varrho]\), entonces
\begin{equation}
 \frac{\mathrm d q_\varrho}{\mathrm dt}
 =\frac{i}{\hbar}\operatorname{Tr}
       \bigl(\varrho[H(t),\widehat{\mathscr Q}_{\rm sp}]\bigr).
 \label{x_reciprocidad:vii:eq:conservacion-momento-cuartico}
\end{equation}
La conmutación, la invariancia de un sector donde el observable es
escalar, o una conservación demostrada para el estado concreto
proporcionan las condiciones pertinentes. Si el momento evoluciona,
\eqref{x_reciprocidad:vii:eq:densidad-presion-espin} conserva su evaluación
instantánea; su balance incluye
\[
 \dot\rho_{\rm spin}+6H\rho_{\rm spin}
 =-\frac{cC_\gamma}{V_c^2}\,a^{-6}\dot q_\varrho,
 \qquad H=\dot a/a
\]
en tiempo propio. Este término registra el intercambio con los
restantes grados de libertad materiales.

La composición queda determinada: representación espinorial del
transporte, acción material, corriente variacional, inversión de
Cartan--Holst, segundo momento y densidad normalizada. El dominio
\(s_0>0\) de la sección~\ref{x_reciprocidad:vii:sec:dilucion-rebote} dispone así
de un funcional evaluable y de estados testigo explícitos. Su uso en
la dinámica cosmológica conserva además las condiciones materiales,
geométricas y de conservación enunciadas en aquella sección.

% END OWNER


% BEGIN OWNER /Users/ruben/Documents/New project/output/EXTREMA_MEDIA_RAZON_RECIPROCIDAD_ES_EN_20260918/ARTICULO/ES/source/antecedentes_vii/50_dilucion_y_rebote.tex
% Artículo VII: pruebas preexistentes incorporadas al cuerpo.
% Propietario íntegro: fuentes_conservadas/17_rebote_einstein_cartan.tex.
% Incluye la ley de escala y las pruebas de rebote y persistencia.
% La amplitud s_0 es la salida de la realización material antecedente;
% este fragmento no certifica por sí solo la evaluación de esa realización.

\chapter{Dilución de memoria y persistencia del rebote}
\label{x_reciprocidad:vii:sec:dilucion-rebote}

La evolución de la memoria y la expansión espacial son dos lecturas de
una misma sucesión de estados. El índice de refinamiento puede eliminarse
entre ambas; esa eliminación determina el exponente de dilución que la
realización cosmológica transmite a su contribución torsional.
La amplitud inicial depende de la evaluación de la corriente material
concreta en su respuesta variacional. La ley de escala y el teorema
de rebote que siguen conservan esa dependencia.

El teorema~\ref{x_reciprocidad:vii:thm:amplitud-material-explicita} determina esa
amplitud como funcional del estado material y del volumen comóvil.
El corolario~\ref{x_reciprocidad:vii:cor:dominio-positivo-conservado} proporciona
un dominio explícito de estados con amplitud positiva y conservada.
En ese dominio, la evaluación de \(s_0\) precede a la integración
cosmológica que se desarrolla a continuación. La densidad ordinaria
\(\rho_0>0\) y su ley barotrópica se especifican adicionalmente en la
sección~\ref{x_reciprocidad:vii:sec:umbral-rebote}: el testigo de positividad de la
amplitud no evalúa por sí mismo esa densidad material.

\section{Ley de dilución \texorpdfstring{\(a^{-6}\)}{a elevado a -6}}
\label{x_reciprocidad:vii:sec:dilucion}

\begin{proposition}[Eliminación del índice de memoria]
\label{x_reciprocidad:vii:prop:dilucion}
Sean \(a_{\mathrm{ref}}>0\), \(\varphi>1\) y los lectores de la ruta
\[
 \frac{a_n}{a_{\mathrm{ref}}}=\varphi^{n/3},\qquad
 M_n=\varphi^{-2n}.
\]
Entonces
\[
 M_n=\left(\frac{a_n}{a_{\mathrm{ref}}}\right)^{-6}
\]
para todo índice de la sucesión. Si la realización de la densidad
cuadrática es \(s_n^2=s_*^2M_n\), con \(s_*^2>0\), se obtiene
\[
 s_n^2=s_*^2\left(\frac{a_n}{a_{\mathrm{ref}}}\right)^{-6}.
\]
\end{proposition}

\begin{proof}
Elevando la primera igualdad a la potencia \(-6\) se obtiene
\(\varphi^{-2n}\), que es exactamente el segundo lector. Multiplicar
por la amplitud inicial reproduce la segunda afirmación. El exponente
queda determinado por el cociente entre los exponentes de ambos
lectores; la amplitud es la evaluación del estado de referencia.
\end{proof}

En la realización homogénea continua se conserva esta ley constitutiva:
la densidad torsional se escribe \(s_0a^{-6}\), donde la normalización
del factor \(a\), el coeficiente material y la amplitud de referencia
se han absorbido en \(s_0\). La ley de escala y la evaluación de
\(s_0\) son operaciones sucesivas de la construcción.

\section{Umbral homogéneo y transversalidad}
\label{x_reciprocidad:vii:sec:umbral-rebote}

Consideremos la carta homogénea, espacialmente plana, en unidades \(c=1\),
con \(G>0\) constante, densidades positivas \(\rho_0,s_0\) y parámetro
barotrópico constante \(w<1\). La realización especificada satisface
\(\Lambda_{\mathrm{eff}}=0\) y las ecuaciones
\begin{align}
 \rho_{\mathrm m}(a)&=\rho_0a^{-3(1+w)},&
 \rho_{\mathrm s}(a)&=s_0a^{-6},\\
 H^2&=\frac{8\pi G}{3}
        \bigl(\rho_{\mathrm m}-\rho_{\mathrm s}\bigr)
        =:F_0(a),&
 \dot H&=-4\pi G
        \bigl((1+w)\rho_{\mathrm m}-2\rho_{\mathrm s}\bigr).
 \label{x_reciprocidad:vii:eq:dinamica-homogenea}
\end{align}
Los acoplamientos pertenecen a la sección dimensional ya construida.
La expresión de \(\rho_{\mathrm s}\) recoge el signo de su contribución
en esta realización gravitatoria.

\begin{theorem}[Existencia y unicidad del umbral transversal]
\label{x_reciprocidad:vii:thm:rebote}
Existe un único \(a_{\mathrm b}>0\) con \(F_0(a_{\mathrm b})=0\):
\[
 a_{\mathrm b}^{\,3(1-w)}=\frac{s_0}{\rho_0}.
\]
La región permitida por \(H^2\geq0\) es \(a\geq a_{\mathrm b}\).
En el umbral, con
\(\rho_{\mathrm b}=\rho_{\mathrm m}(a_{\mathrm b})
                  =\rho_{\mathrm s}(a_{\mathrm b})>0\),
\[
 \dot H_{\mathrm b}=4\pi G(1-w)\rho_{\mathrm b}>0.
\]
La solución que atraviesa \(H=0\) presenta un mínimo estricto de \(a\),
con expansión local
\[
 a(t)=a_{\mathrm b}\left[
 1+2\pi G(1-w)\rho_{\mathrm b}(t-t_{\mathrm b})^2
 +O((t-t_{\mathrm b})^3)\right].
\]
\end{theorem}

\begin{proof}
La factorización
\[
 F_0(a)=\frac{8\pi G}{3}a^{-6}
             \bigl(\rho_0a^{3(1-w)}-s_0\bigr)
\]
reduce sus ceros al de una función estrictamente creciente de
\((0,\infty)\) sobre \((-s_0,\infty)\). Esto demuestra existencia,
unicidad y el signo de \(F_0\) a ambos lados del umbral.
La segunda ecuación de \eqref{x_reciprocidad:vii:eq:dinamica-homogenea}, evaluada
en la igualdad de densidades, da la fórmula de \(\dot H_{\mathrm b}\).
Como \(\dot a=aH\), en el umbral
\(\ddot a_{\mathrm b}=a_{\mathrm b}\dot H_{\mathrm b}>0\).
Las ecuaciones son suaves para \(a>0\), por lo que su desarrollo de
Taylor da la expansión indicada. El cambio de signo de \(H\) es
transversal: pasa de contracción a expansión.
\end{proof}

Para \(w=1\), ambos términos tienen el mismo exponente y la igualdad de
amplitudes determina un caso degenerado. Para \(w>1\), el signo de
\(\dot H\) en una coincidencia de densidades es negativo.
El dominio \(w<1\) del teorema expresa exactamente el orden entre los
dos exponentes de dilución.

\section{Persistencia local bajo perturbaciones controladas}
\label{x_reciprocidad:vii:sec:persistencia}

La persistencia del umbral utiliza dos estimaciones independientes:
el control del signo en los extremos y el de la derivada en el
intervalo. Sea \(I=[a_-,a_+]\subset(0,\infty)\), con
\(a_-<a_{\mathrm b}<a_+\), elegido de forma que
\[
 F_0(a_-)<0<F_0(a_+),\qquad
 m_I:=\inf_{a\in I}F_0'(a)>0.
\]
Tal intervalo existe porque
\[
 F_0'(a_{\mathrm b})
   =\frac{8\pi G}{a_{\mathrm b}}(1-w)\rho_{\mathrm b}>0.
\]
Definimos
\(\delta_I=\min\{-F_0(a_-),F_0(a_+)\}>0\).

\begin{theorem}[Persistencia local del rebote transversal]
\label{x_reciprocidad:vii:thm:persistencia}
Sea \(R\in C^1(I)\) una perturbación que satisface
\[
 \|R\|_{\infty,I}<\delta_I,\qquad
 \|R'\|_{\infty,I}<m_I.
\]
Entonces \(F_R=F_0+R\) tiene un único cero \(a_R\in(a_-,a_+)\) y
\(F_R'(a_R)>0\). La realización \(H^2=F_R(a)\) admite el rebote local
transversal, con
\[
 \dot H_R=\frac{a_R}{2}F_R'(a_R)>0
\]
en el umbral. Si \(R(a,\eta)\) es conjuntamente \(C^1\), los ceros
persistentes dependen localmente de manera \(C^1\) del parámetro \(\eta\).
\end{theorem}

\begin{proof}
La primera cota conserva
\(F_R(a_-)<0<F_R(a_+)\); el teorema del valor intermedio proporciona
un cero interior. La segunda da
\(F_R'(a)\geq m_I-\|R'\|_{\infty,I}>0\) en \(I\), por lo que el cero
es único y simple.

En \(a>a_R\) cercano al cero,
\(F_R(a)=F_R'(a_R)(a-a_R)+o(a-a_R)\).
La integral
\[
 t-t_{\mathrm b}=\int_{a_R}^{a}
          \frac{d\xi}{\xi\sqrt{F_R(\xi)}}
\]
es finita y estrictamente creciente. Su inversa define la rama
expansiva; el reflejo temporal define la rama contractiva.
Ambas se unen con \(\dot a(t_{\mathrm b})=0\) y
\[
 a(t)-a_R
 =\frac{a_R^2F_R'(a_R)}4(t-t_{\mathrm b})^2
   +o((t-t_{\mathrm b})^2).
\]
Fuera del umbral, la derivación de \(H^2=F_R(a)\) da
\(\dot H=aF_R'(a)/2\); su límite continuo da la expresión positiva
en \(a_R\). La regularidad \(C^1\) de \(F_R\) garantiza la continuidad
de la aceleración en esta unión. Finalmente,
\(\partial_aF_R(a_R)>0\) permite aplicar el teorema de la función
implícita cuando la perturbación depende conjuntamente de \((a,\eta)\).
\end{proof}

La primera cota controla la existencia; la segunda asegura unicidad
local y transversalidad. La estabilidad demostrada corresponde a
ese umbral en el intervalo seleccionado. La amplitud torsional
interviene en su localización, mientras que la relación entre
exponentes determina la orientación temporal del cruce.

% END OWNER


% BEGIN OWNER /Users/ruben/Documents/New project/output/EXTREMA_MEDIA_RAZON_RECIPROCIDAD_ES_EN_20260918/ARTICULO/ES/source/antecedentes_vii/51_solucion_homogenea.tex
% FORMALIZACION_NUEVA de la realización homogénea preexistente.
% Prioridad de la fórmula integrada: procedencia histórica pendiente de cotejo.
% Conserva la ecuación y el dominio del capítulo de rebote; no selecciona s_0.
\chapter{Solución homogénea explícita y completitud causal}
\label{x_reciprocidad:vii:sec:solucion-homogenea}
La determinación de la amplitud permite desarrollar una realización
completa del rebote, además del argumento local. Se conserva la carta
espacialmente plana de \eqref{x_reciprocidad:vii:eq:dinamica-homogenea}, en unidades
\(c=1\), con \(\rho_0,s_0,G>0\), y se especializa el sector material
a polvo, \(w=0\), con \(\Lambda_{\mathrm{eff}}=0\).
El coeficiente \(s_0\) es el obtenido del estado
material en la realización correspondiente; la integración siguiente
emplea ese coeficiente y conserva su dominio.

\begin{theorem}[Integración exacta del rebote de polvo]
\label{x_reciprocidad:vii:thm:rebote-polvo}
Pongamos
\[
 a_{\rm b}^3=\frac{s_0}{\rho_0},\qquad
 \tau^2=\frac{s_0}{6\pi G\rho_0^2},\qquad u=t-t_{\rm b}.
\]
La solución del sistema homogéneo con
\(a(t_{\rm b})=a_{\rm b}\) y \(H(t_{\rm b})=0\) es, para todo
\(t\in\mathbb R\),
\begin{equation}
 a(t)=a_{\rm b}\left(1+\frac{u^2}{\tau^2}\right)^{1/3},
 \qquad H(t)=\frac{2u}{3(\tau^2+u^2)},\qquad
 \dot H(t)=\frac{2(\tau^2-u^2)}{3(\tau^2+u^2)^2}.
 \label{x_reciprocidad:vii:eq:solucion-polvo}
\end{equation}
Tiene un único mínimo, en \(t_{\rm b}\); es contractiva antes y
expansiva después. La evolución es suave y satisface
\(a(t)\geq a_{\rm b}>0\).
\end{theorem}
\begin{proof}
Sea \(y=a^3\). La ecuación de Friedmann da
\[
 \dot y^{\,2}=24\pi G(\rho_0y-s_0).
\]
La función \(y=s_0/\rho_0+6\pi G\rho_0u^2\) verifica esa
identidad y las condiciones iniciales. Sus derivadas dan las
tres expresiones de \eqref{x_reciprocidad:vii:eq:solucion-polvo}. Para comprobar
también la ecuación de evolución, escribamos
\[
 \rho_{\rm b}=\frac{\rho_0^2}{s_0},\qquad
 x=\frac{u}{\tau}.
\]
Entonces
\[
 \rho_{\rm m}=\frac{\rho_{\rm b}}{1+x^2},\qquad
 \rho_{\rm s}=\frac{\rho_{\rm b}}{(1+x^2)^2},\qquad
 4\pi G\rho_{\rm b}=\frac{2}{3\tau^2}.
\]
Sustituir en \(\dot H=-4\pi G(\rho_{\rm m}-2\rho_{\rm s})\)
da exactamente la tercera fórmula. Se satisface, por tanto,
el sistema completo \((\dot a,\dot H)\), cuyo campo es suave
para \(a>0\). Su unicidad local, prolongada sobre la solución
explícita, fija la evolución por esas condiciones iniciales.
El denominador es positivo para todo \(t\), y las afirmaciones
sobre el mínimo y el signo siguen de \(H\).
\end{proof}

\begin{figure}[htbp]
\centering
\begin{tikzpicture}[x=1.45cm,y=1.15cm,>=Latex]
\draw[->] (-3.25,0)--(3.3,0) node[right] {\(\displaystyle u/\tau\)};
\draw[->] (0,0)--(0,2.8) node[above] {\(\displaystyle a/a_{\rm b}\)};
\draw[densely dashed,gray] (-3.1,1)--(3.1,1);
\draw[very thick,ink,domain=-3:3,samples=121,smooth]
 plot (\x,{(1+\x*\x)^(1/3)});
\fill[accent] (0,1) circle (1.6pt);
\node[anchor=north east] at (-.08,.98) {\(1\)};
\node[ink] at (-2,2.65) {Contracción};
\node[ink] at (2,2.65) {Expansión};
\foreach \t in {-3,-2,-1,1,2,3}
{\draw (\t,.045)--(\t,-.045) node[below] {\(\t\)};}
\end{tikzpicture}
\caption{Perfil adimensional de la solución demostrada:
\(a/a_{\rm b}=(1+(u/\tau)^2)^{1/3}\).
Los parámetros \(a_{\rm b}\) y \(\tau\) se calculan desde
\(\rho_0,s_0,G\); el dibujo representa toda la familia reescalada.}
\label{x_reciprocidad:vii:fig:rebote-polvo}
\end{figure}

\section{Curvatura finita y extensión de las geodésicas}
La métrica realizada en esta carta es
\(g=-dt^2+a(t)^2\,d\mathbf x^2\) sobre
\(\mathbb R\times\mathbb R^3\). Con la convención en que
el escalar de curvatura plano es \(R=6(\dot H+2H^2)\), la solución da
\[
 R(t)=\frac{4\tau^2+\frac43u^2}{(\tau^2+u^2)^2}.
\]
En particular, \(R(t_{\rm b})=4/\tau^2\) es finito.
En un marco ortonormal, las componentes de Riemann dependen de
\(H^2\) y de \(\dot H+H^2\), que son funciones acotadas en
\(\mathbb R\). Lo mismo ocurre con las contracciones polinómicas
de esas componentes; por ejemplo,
\[
 R_{\mu\nu\rho\sigma}R^{\mu\nu\rho\sigma}
 =12\bigl[(\dot H+H^2)^2+H^4\bigr].
\]

\begin{proposition}[Completitud de las geodésicas causales]
\label{x_reciprocidad:vii:prop:completitud-causal-polvo}
La realización métrica precedente es geodésicamente completa
para las geodésicas temporales y nulas de su conexión de
Levi--Civita, en ambas direcciones.
\end{proposition}
\begin{proof}
Las traslaciones espaciales proporcionan las constantes
\(p_i=a^2\,dx^i/d\lambda\). La normalización causal da
\[
 \left|\frac{dt}{d\lambda}\right|
 =\sqrt{\epsilon+\frac{|\mathbf p|^2}{a(t)^2}},
 \qquad \epsilon=1\ \text{en el caso temporal},\quad
 \epsilon=0\ \text{en el caso nulo}.
\]
Una geodésica nula no constante tiene \(|\mathbf p|>0\).
Como \(a\geq a_{\rm b}\), en el caso temporal
\[
 \left|\frac{d\lambda}{dt}\right|
 \geq\left(1+\frac{|\mathbf p|^2}{a_{\rm b}^2}\right)^{-1/2}>0,
\]
y en el nulo
\(\left|d\lambda/dt\right|=a/|\mathbf p|
 \geq a_{\rm b}/|\mathbf p|>0\).
Así, alcanzar \(t=\pm\infty\) requiere longitud afín infinita.
En un intervalo temporal compacto, \(a\) y sus derivadas son
suaves y \(a\) está separado de cero. Las ecuaciones explícitas
anteriores y \(dx^i/d\lambda=p_i/a^2\) permiten continuar la
geodésica a través de sus extremos finitos. Quedan cubiertas
las dos posibilidades de prolongación.
\end{proof}

El resultado de completitud pertenece a esta realización homogénea de
polvo, con sus condiciones de espacio, materia y amplitud. Su
persistencia local bajo perturbaciones del umbral conserva las cotas
del teorema~\ref{x_reciprocidad:vii:thm:persistencia}; esa propiedad local y la
completitud global de la solución explícita son conclusiones distintas.
La escala mínima procede de la relación entre densidad material y
respuesta torsional. La memoria interviene así en una evolución
geométrica verificable, con una realización regular de la contracción
y la expansión.

% END OWNER


% BEGIN OWNER /Users/ruben/Documents/New project/output/EXTREMA_MEDIA_RAZON_RECIPROCIDAD_ES_EN_20260918/ARTICULO/ES/source/antecedentes_vii/60_tensor_residual_y_balance.tex
% Artículo VII. Recuperación por dependencia, 10-09-2026.
% RESULTADO_RECUPERADO: integral activo 2249, fuente/colaboracion/
% partes_iv_v/tex/residencias/86_cosmologia_cerrada.tex (S04),
% líneas 2--44 (cartas seleccionadas), 174--248 (tensor y aceleración).
% Fuente S04 leída íntegramente, SHA-256:
% feb2582c3316c8b30829cfb0476bfa2f4e95889dabfba16fe43f2d835a226673.
% Se recuperan definición, unicidad, conservación, descomposición,
% intercambio y aceleración, con sus condiciones explícitas.
% FORMALIZACION_REUNIDA: relación con la fuente efectiva de 31;
% se distingue conservación total de conservación material separada.
% CERTIFICADO_NUEVO: expansión de las pruebas de traza, balance y signos;
% cálculo local de Ricci para las tres cartas ya seleccionadas por S04.
% No se modifica el rebote de 50 ni se determina aquí una amplitud s_0.
%
% DEPENDENCIAS DE INTEGRACIÓN:
% - vii:sec:realizacion-cartan: la métrica procede de la cotetrada realizada.
% - vii:sec:acoplamientos: G_int y el reloj t_0 del mismo estado HMT;
%   c_int=1 es una elección de unidades posterior a su construcción.
% - vii:eq:ecuacion-metrica-efectiva y vii:eq:balance-metrico-total, de 31:
%   se usan únicamente para identificar el residual sobre una solución.
% - T_b es el tensor publicado por el transductor material declarado;
%   su conservación separada es una condición expresamente enunciada.
% - El selector TRIT de la carta y t_0(x)>0 son antecedentes de la
%   especialización de Sitter, no datos cosmológicos observacionales.
% El presente fragmento no certifica la clausura del conjunto de VII.

\chapter{Tensor residual, descomposición de traza y balance covariante}
\label{x_reciprocidad:vii:sec:tensor-residual-balance}

La realización geométrica del estado enriquecido proporciona una métrica,
mientras el transductor material publica la fuente a la que se refiere la
lectura. El tensor residual conserva la diferencia entre ambas publicaciones
en un mismo espacio tensorial. Su descomposición permite distinguir el
sector proporcional a la métrica y el sector de traza nula, y determina
la corriente de intercambio que corresponde a ese reparto.

\section{Datos de la carta y construcción del tensor residual}
\label{x_reciprocidad:vii:subsec:datos-tensor-residual}

Fijemos un estado HMT \(x\) y una carta conexa \(\mathcal U\) de su
realización geométrica. Sobre \(\mathcal U\), la cotetrada no degenerada
de la sección~\ref{x_reciprocidad:vii:sec:realizacion-cartan} determina una métrica lisa
\(g=g_x\) de firma \((-+++)\). Denotamos por
\(\mathring\nabla\) su conexión de Levi--Civita y escribimos
\[
 \mathsf G_{\mu\nu}[g]
 :=\operatorname{Ric}_{\mu\nu}[g]-\tfrac12R[g]g_{\mu\nu}.
\]
Así se distingue el tensor de Einstein \(\mathsf G[g]\) del acoplamiento
gravitatorio \(G_{\rm int}\). La derivada utilizada en los balances
de esta sección es \(\mathring\nabla\), coherentemente con la ecuación
métrica efectiva de la sección~\ref{x_reciprocidad:vii:subsec:fuente-metrica-efectiva}.

Trabajamos en la carta dimensional \(c_{\rm int}=1\). Suponemos
\(G_{\rm int}(x)>0\) y \(\Lambda_{\rm ref}(x)\) constantes sobre
\(\mathcal U\), y ponemos \(\kappa_x=8\pi G_{\rm int}(x)\).
El símbolo \(x\) identifica el estado que suministra esos lectores;
la diferenciación covariante actúa sobre los puntos de \(\mathcal U\).
La genealogía de los acoplamientos permanece en la
sección~\ref{x_reciprocidad:v:ant:sec:gravity}.

Sea \(T_b\) el tensor simétrico liso publicado por el transductor material
de la carta; en la lectura bariónica, es su fuente bariónica. Definimos
\begin{equation}
 \mathcal E^{\rm ref}_{\mu\nu}
 :=\kappa_x^{-1}\bigl(\mathsf G_{\mu\nu}[g]
                            +\Lambda_{\rm ref}g_{\mu\nu}\bigr),
 \qquad
 T^{\rm res}_{\mu\nu}:=\mathcal E^{\rm ref}_{\mu\nu}-T_{b,\mu\nu}.
 \label{x_reciprocidad:vii:eq:residual-definicion}
\end{equation}
El tensor \(T^{\rm res}\) es el tensor residual efectivo, también
denotado \(T_{\rm osc}\) al interpretar el sector geométrico que queda
por publicar respecto de \(T_b\). La construcción conserva como
argumentos efectivos \(g_x,G_{\rm int},\Lambda_{\rm ref}\) y \(T_b\).

\begin{theorem}[Unicidad del residual y balance con la fuente material]
\label{x_reciprocidad:vii:thm:residual-unicidad-balance}
Para los datos anteriores existe un único tensor simétrico
\(T^{\rm res}\) que satisface
\(\mathcal E^{\rm ref}=T_b+T^{\rm res}\). Su divergencia es
\begin{equation}
 \mathring\nabla^\mu T^{\rm res}_{\mu\nu}
 =-\mathring\nabla^\mu T_{b,\mu\nu}.
 \label{x_reciprocidad:vii:eq:balance-residual-material}
\end{equation}
En particular, si la fuente material se conserva por separado,
\(\mathring\nabla^\mu T_{b,\mu\nu}=0\), también se conserva
\(T^{\rm res}\).
\end{theorem}
\begin{proof}
La diferencia de dos tensores que satisfagan la ecuación anunciada es
cero; la expresión~\eqref{x_reciprocidad:vii:eq:residual-definicion} proporciona su
existencia. La identidad de Bianchi contraída y la compatibilidad métrica
de Levi--Civita dan
\[
 \mathring\nabla^\mu\mathsf G_{\mu\nu}=0,
 \qquad \mathring\nabla g=0.
\]
Como \(\kappa_x\) y \(\Lambda_{\rm ref}\) son constantes sobre la
carta, \(\mathring\nabla^\mu\mathcal E^{\rm ref}_{\mu\nu}=0\).
Derivar la diferencia que define \(T^{\rm res}\) prueba
\eqref{x_reciprocidad:vii:eq:balance-residual-material}, incluida la afirmación bajo
conservación material separada.
\end{proof}

La constancia espacial y temporal de los acoplamientos en la carta fija
el dominio de este balance. Una familia de estados puede publicar
acoplamientos diferentes en cartas diferentes; cada identidad anterior
se aplica a la carta y al estado que la especifican.

\section{Descomposición canónica y corriente de intercambio}
\label{x_reciprocidad:vii:subsec:traza-corriente-residual}

Definimos la función de traza y los dos tensores
\begin{equation}
 \theta:=g^{\mu\nu}T^{\rm res}_{\mu\nu},\qquad
 T^{\Lambda}_{\mu\nu}:=\tfrac14\theta g_{\mu\nu},\qquad
 T^{0}_{\mu\nu}:=T^{\rm res}_{\mu\nu}-T^{\Lambda}_{\mu\nu}.
 \label{x_reciprocidad:vii:eq:descomposicion-traza-residual}
\end{equation}
El superíndice \(\Lambda\) identifica la componente proporcional a la
métrica; \(\Lambda_{\rm ref}\) es el parámetro de referencia de
\eqref{x_reciprocidad:vii:eq:residual-definicion}. Sus funciones son distintas.

\begin{proposition}[Proyectores de traza y unicidad de la descomposición]
\label{x_reciprocidad:vii:prop:proyectores-traza-residual}
Sobre los tensores simétricos covariantes de orden dos, las aplicaciones
\[
 P_{\rm tr}(S)=\tfrac14\operatorname{tr}_g(S)g,
 \qquad P_0(S)=S-P_{\rm tr}(S)
\]
satisfacen
\begin{equation}
 P_{\rm tr}^2=P_{\rm tr},\quad P_0^2=P_0,\quad
 P_{\rm tr}P_0=P_0P_{\rm tr}=0,\quad P_{\rm tr}+P_0=I.
 \label{x_reciprocidad:vii:eq:proyectores-traza-residual}
\end{equation}
En particular, \(T^{\rm res}=T^\Lambda+T^0\),
\(\operatorname{tr}_gT^0=0\), y esta descomposición es única.
\end{proposition}
\begin{proof}
En dimensión cuatro, \(g^{\mu\nu}g_{\mu\nu}=4\). Por tanto,
\(\operatorname{tr}_g(P_{\rm tr}S)=\operatorname{tr}_gS\), lo que
prueba la idempotencia de \(P_{\rm tr}\), la anulación de la traza
de \(P_0S\) y, por expansión, las otras identidades.
Si \(S=fg+Q\) con \(\operatorname{tr}_gQ=0\), tomar la traza
da \(f=\operatorname{tr}_gS/4\). Queda determinado también
\(Q=S-fg\), y con ello la unicidad.
\end{proof}

\begin{theorem}[Intercambio covariante entre los dos sectores]
\label{x_reciprocidad:vii:thm:intercambio-traza-residual}
En el dominio de \eqref{x_reciprocidad:vii:eq:residual-definicion}, sea
\(b_\nu:=\mathring\nabla^\mu T_{b,\mu\nu}\). La corriente
de intercambio de traza es la uno-forma
\begin{equation}
 \mathcal J^{\rm tr}_\nu
 :=\tfrac14\partial_\nu\theta
 =\mathring\nabla^\mu T^\Lambda_{\mu\nu}.
 \label{x_reciprocidad:vii:eq:corriente-intercambio-traza}
\end{equation}
Los balances completos son
\begin{equation}
 \mathring\nabla^\mu T^0_{\mu\nu}
       =-b_\nu-\mathcal J^{\rm tr}_\nu,
 \qquad
 \mathring\nabla^\mu(T^\Lambda+T^0)_{\mu\nu}=-b_\nu.
 \label{x_reciprocidad:vii:eq:balances-sectores-residuales}
\end{equation}
Si \(b=0\), ambos sectores intercambian corrientes exactamente opuestas.
En ese caso se conservan por separado si y sólo si \(\theta\) es
constante en la carta conexa.
\end{theorem}
\begin{proof}
La compatibilidad métrica permite calcular
\[
 \mathring\nabla^\mu\bigl(\tfrac14\theta g_{\mu\nu}\bigr)
 =\tfrac14g^{\mu\alpha}(\partial_\alpha\theta)g_{\mu\nu}
 =\tfrac14\partial_\nu\theta.
\]
Restar esta igualdad a \eqref{x_reciprocidad:vii:eq:balance-residual-material} produce
\eqref{x_reciprocidad:vii:eq:balances-sectores-residuales}. Con \(b=0\), la
conservación separada equivale a \(\mathrm d\theta=0\); una función
lisa con diferencial nulo es constante en cada componente conexa.
\end{proof}

La uno-forma \(\mathcal J^{\rm tr}\) mide el intercambio del reparto
tensorial. Conserva un tipo distinto del covector de transporte de la
sección~\ref{x_reciprocidad:vii:sec:variacion-memoria} y de la corriente de espín
\(\sigma_{IJ}\in\Omega^3(\mathcal U;\Lambda^2E^*)\) de
\eqref{x_reciprocidad:vii:eq:corriente-variacional}.

\section{Enlace con la fuente efectiva de Cartan--Holst}
\label{x_reciprocidad:vii:subsec:residual-fuente-cartan}

Sobre una solución de \eqref{x_reciprocidad:vii:eq:ecuacion-metrica-efectiva}, en las
mismas unidades \(c=1\), el residual tiene la expresión
\begin{equation}
 T^{\rm res}
 =T^{\rm phys,LC}+U^{\rm phys,spin}-T_b
   +\frac{\Lambda_{\rm ref}-\Lambda_{\rm int}}{\kappa_x}\,g.
 \label{x_reciprocidad:vii:eq:residual-composicion-cartan}
\end{equation}
En efecto, sustituir
\(\mathsf G+\Lambda_{\rm int}g
  =\kappa_x(T^{\rm phys,LC}+U^{\rm phys,spin})\)
en \eqref{x_reciprocidad:vii:eq:residual-definicion} da la igualdad término a término.
Cuando el transductor material publica precisamente
\(T_b=T^{\rm phys,LC}\) y \(\Lambda_{\rm ref}=\Lambda_{\rm int}\),
se obtiene \(T^{\rm res}=U^{\rm phys,spin}\).
La condición de conservación separada de \(T_b\) del
teorema~\ref{x_reciprocidad:vii:thm:residual-unicidad-balance} conserva su función
propia: \eqref{x_reciprocidad:vii:eq:balance-metrico-total} establece el balance de la
fuente total, y el reparto entre sus componentes depende de la acción
material utilizada.

\section{Sector de traza y aceleración de las cartas seleccionadas}
\label{x_reciprocidad:vii:subsec:residual-aceleracion-ds}

El lector TRIT del estado \(x\) publica \(s=\tau(x)\in\{+1,0,-1\}\),
y su reloj positivo publica
\begin{equation}
 H_*=\frac{2\pi}{108t_0(x)}>0.
 \label{x_reciprocidad:vii:eq:reloj-carta-ds}
\end{equation}
Sea \(h_s\) una métrica tridimensional de curvatura seccional constante
\(s\). Las tres cartas de la realización seleccionada se escriben
\begin{equation}
 g_s=-\mathrm dt^2+a_s(t)^2h_s,\qquad
 a_+(t)=H_*^{-1}\cosh(H_*t),\quad
 a_0(t)=\exp(H_*t),\quad
 a_-(t)=H_*^{-1}\sinh(H_*t).
 \label{x_reciprocidad:vii:eq:cartas-ds-residual}
\end{equation}
Las cartas cerrada y plana se consideran para \(t\in\mathbb R\);
la carta abierta, para \(t>0\). El signo \(s\) y el reloj \(H_*\)
mantienen sus antecedentes en el estado que selecciona la realización.

\begin{lemma}[Curvatura de las tres cartas del mismo reloj]
\label{x_reciprocidad:vii:lem:curvatura-cartas-ds}
Para las funciones de \eqref{x_reciprocidad:vii:eq:cartas-ds-residual},
\begin{equation}
 \frac{\ddot a_s}{a_s}=H_*^2,\qquad
 \left(\frac{\dot a_s}{a_s}\right)^2+\frac{s}{a_s^2}=H_*^2,
 \qquad
 \operatorname{Ric}[g_s]=3H_*^2g_s.
 \label{x_reciprocidad:vii:eq:curvatura-cartas-ds}
\end{equation}
En consecuencia, \(R[g_s]=12H_*^2\) y
\(\mathsf G[g_s]=-3H_*^2g_s\).
\end{lemma}
\begin{proof}
Las dos primeras igualdades se obtienen diferenciando las tres funciones:
en las cartas cerrada y abierta se utiliza
\(\cosh^2z-\sinh^2z=1\), y en la plana
\(\dot a_0=H_*a_0\). Para verificar la identificación tensorial,
los coeficientes de conexión que involucran la coordenada temporal son
\[
 \Gamma^0{}_{ij}=a_s\dot a_s(h_s)_{ij},\qquad
 \Gamma^i{}_{0j}=(\dot a_s/a_s)\delta^i_j,
\]
y los puramente espaciales coinciden con los de \(h_s\).
La contracción de la curvatura da
\[
 \operatorname{Ric}_{00}=-3\ddot a_s/a_s,\qquad
 \operatorname{Ric}_{0i}=0,\qquad
 \operatorname{Ric}_{ij}
  =(a_s\ddot a_s+2\dot a_s^2+2s)(h_s)_{ij}.
\]
Sustituir las dos primeras identidades de
\eqref{x_reciprocidad:vii:eq:curvatura-cartas-ds} produce
\(\operatorname{Ric}_{00}=-3H_*^2\) y
\(\operatorname{Ric}_{ij}=3H_*^2a_s^2(h_s)_{ij}\).
Tomar la traza y formar el tensor de Einstein prueba las otras
afirmaciones, con la convención de curvatura de
\eqref{x_reciprocidad:vii:eq:ecuacion-metrica-efectiva}.
\end{proof}

\begin{theorem}[Densidad, presión y signo de la aceleración seleccionada]
\label{x_reciprocidad:vii:thm:residual-ds-aceleracion}
En las cartas anteriores, para \(\Lambda_{\rm ref}=0\) y \(T_b=0\),
el tensor residual es
\begin{equation}
 T^{\rm res}=-\rho_*g_s,\qquad
 \rho_*:=\frac{3H_*^2}{8\pi G_{\rm int}}>0,
 \qquad T^\Lambda=T^{\rm res},\quad T^0=0.
 \label{x_reciprocidad:vii:eq:residual-ds-puro}
\end{equation}
Para un observador comóvil unitario \(u=\partial_t\), la densidad
y la presión isotrópica son \(\rho_*\) y \(p_*=-\rho_*\).
La ecuación de aceleración se evalúa como
\begin{equation}
 \frac{\ddot a_s}{a_s}
 =-\frac{4\pi G_{\rm int}}3(\rho_*+3p_*)=H_*^2>0.
 \label{x_reciprocidad:vii:eq:aceleracion-residual-ds}
\end{equation}
\end{theorem}
\begin{proof}
El lema anterior y \eqref{x_reciprocidad:vii:eq:residual-definicion} dan
\(T^{\rm res}=-(3H_*^2/\kappa_x)g_s\).
Como \(g_s(u,u)=-1\),
\(T^{\rm res}_{\mu\nu}u^\mu u^\nu=\rho_*\).
El proyector espacial \(q^{\mu\nu}=g_s^{\mu\nu}+u^\mu u^\nu\)
satisface \(q^{\mu\nu}(g_s)_{\mu\nu}=3\), y por ello
\[
 p_*:=\tfrac13q^{\mu\nu}T^{\rm res}_{\mu\nu}=-\rho_*.
\]
La traza vale \(-4\rho_*\), de modo que
\eqref{x_reciprocidad:vii:eq:descomposicion-traza-residual} produce los dos sectores
de \eqref{x_reciprocidad:vii:eq:residual-ds-puro}. Finalmente,
\(\rho_*+3p_*=-2\rho_*\), y la sustitución en el segundo miembro
de \eqref{x_reciprocidad:vii:eq:aceleracion-residual-ds} da
\(\kappa_x\rho_*/3=H_*^2\), igual al valor geométrico ya calculado.
La positividad utiliza exactamente \(G_{\rm int}>0\) y \(H_*>0\).
\end{proof}

En esta realización, el sector de traza publica una densidad constante
y una presión opuesta, y el reloj interno fija la aceleración positiva.
En una carta general, \(T^\Lambda\) describe la componente isotrópica
puntual, mientras \(T^0\) conserva el residual de traza nula y ambos
obedecen el intercambio de \eqref{x_reciprocidad:vii:eq:balances-sectores-residuales}.
La identificación de estos sectores con fuentes oscuras observacionales
requiere que una misma realización reproduzca conjuntamente expansión,
lentes, rotación, crecimiento de estructura y fondo cósmico. La
construcción presente fija el tensor, sus balances y la aceleración de
las cartas seleccionadas que intervienen en esa comparación.

% END OWNER


% BEGIN OWNER /Users/ruben/Documents/New project/output/EXTREMA_MEDIA_RAZON_RECIPROCIDAD_ES_EN_20260918/ARTICULO/ES/source/antecedentes_vii/70_horizontes_e_informacion.tex
% Recuperacion de 09c_entropia_sectorial.tex, 11_confluencia.tex y
% 85_agujero_negro_doble_hoja.tex; originales en fuentes_conservadas/
% horizontes_informacion. Procedencia: desarrollo/HORIZONTES_INFORMACION_FUENTES.md.
% Las referencias a antecedentes se acuerdan con el editor de VII.
% No incorpora una corriente material ni evalua la amplitud torsional.

\chapter{Estados bipartitos puros, entropía reducida y leyes de área del horizonte}
\label{x_reciprocidad:vii:sec:horizontes-informacion}

El estado enriquecido conserva las rutas y sus registros durante la
evolución; una lectura de subsistema determina qué parte de esa información
queda accesible. La construcción de esta sección utiliza los canales
angulares de la sección~\ref{x_reciprocidad:iv:angular:section}, la incidencia y el
operador de área de la sección~\ref{x_reciprocidad:v:ant:sec:area}, y las unidades
de acción, área y temperatura de las
secciones~\ref{x_reciprocidad:v:ant:sec:gravity}
y~\ref{x_reciprocidad:v:ant:sec:ii-kb-aut-desarrollo}. Son publicaciones anteriores de la
misma construcción APP--TRIT--TPK. Cada una transmite aquí su dominio,
orientación, normalización y condiciones de realización.

Se distinguen tres espacios: el registro binario de hoja, los factores
tríticos de la incidencia sectorial y la realización geométrica de los
exteriores de un horizonte. La purificación, el transporte de observables
y la conversión dimensional precisan las relaciones entre ellos. Esta
distinción conserva la estructura que acompaña a cada valor entrópico.

\section{Registro de hoja y purificación cuántica bipartita}
\label{x_reciprocidad:vii:subsec:hojas-purificacion}

Escribamos \(a=A_{\rm rad}\) y \(y=C^*_{\rm rad}\) para las
coordenadas angulares ya construidas. Los canales y sus probabilidades son
\begin{equation}
 \mathfrak q_\pm=e^{-a\mp y},\qquad
 p_\pm=\frac{\mathfrak q_\pm}{\mathfrak q_++\mathfrak q_-}
       =\frac{e^{\mp y}}{2\cosh y}.
 \label{x_reciprocidad:vii:eq:probabilidades-hoja}
\end{equation}
Los signos indexan estos dos canales. En
\(\mathcal H_{\rm h}=\operatorname{span}\{|+\rangle,|-\rangle\}\)
definimos
\[
 \rho_{\rm h}(y)=p_+|+\rangle\langle+|+p_-|-\rangle\langle-|.
\]
El espacio auxiliar
\(\mathcal H_{\rm ar}=\operatorname{span}
\{|\mathrm{adv}\rangle,|\mathrm{ret}\rangle\}\)
porta dos etiquetas ortonormales. Fijada una fase \(\chi_0\), el estado
conjunto es
\begin{equation}
 |\Psi_{\rm ar}\rangle
 =\sqrt{p_+}|+\rangle\otimes|\mathrm{adv}\rangle
  +e^{i\chi_0}\sqrt{p_-}|-\rangle\otimes|\mathrm{ret}\rangle.
 \label{x_reciprocidad:vii:eq:purificacion-hoja}
\end{equation}

\begin{proposition}[Información accesible en el registro de hoja]
\label{x_reciprocidad:vii:prop:purificacion-hoja}
El vector~\eqref{x_reciprocidad:vii:eq:purificacion-hoja} está normalizado y su traza
parcial sobre \(\mathcal H_{\rm ar}\) es \(\rho_{\rm h}(y)\).
Su entropía de entrelazamiento, expresada en nats, es
\begin{equation}
 s_{\rm h}(y)=\log(2\cosh y)-y\tanh y.
 \label{x_reciprocidad:vii:eq:entropia-hoja}
\end{equation}
El intercambio de las etiquetas \(+\) y \(-\) conjuga
\(\rho_{\rm h}(y)\) con \(\rho_{\rm h}(-y)\), y conserva
\(s_{\rm h}\).
\end{proposition}
\begin{proof}
La norma al cuadrado es \(p_++p_-=1\). La ortogonalidad de las
etiquetas auxiliares elimina los términos cruzados de la traza parcial.
Los coeficientes de Schmidt al cuadrado son precisamente \(p_\pm\).
Por~\eqref{x_reciprocidad:vii:eq:probabilidades-hoja},
\(p_+-p_-=-\tanh y\); sustituir
\(\log p_\pm=\mp y-\log(2\cosh y)\) en
\(-\sum_\pm p_\pm\log p_\pm\) demuestra la fórmula.
El intercambio de signos permuta los dos autovalores y conserva la entropía.
\end{proof}

El factor común \(e^{-a}\) se cancela al normalizar el registro binario.
La fase relativa \(\chi_0\) permanece en el estado conjunto aunque esta
reducción diagonal no la publique. La reversibilidad lógica se refiere
al transporte del registro completo, con sus etiquetas y correlaciones;
el retorno de fase conserva el avance de memoria descrito en la
sección~\ref{x_reciprocidad:vii:sec:retorno-memoria-gravedad}. Una implementación térmica
tiene además su propia ley de energía y disipación. En particular,
\(s_{\rm h}\) mide este subsistema binario y su bipartición concreta.

\section{Purificación cuántica de los sectores de incidencia}
\label{x_reciprocidad:vii:subsec:purificacion-sectorial}

La incidencia de borde construida anteriormente distingue dos conjuntos
de dieciocho enlaces, con etiquetas sectoriales \(m=90,120\). Cada enlace
porta \(\mathbb C^3\), con base \(|0\rangle,|1\rangle,|2\rangle\)
y operador de ocupación \(\operatorname{diag}(0,1,2)\). Sean
\(\mathsf N_m\) sus sumas sectoriales, y
\begin{equation}
 \mathsf N=\mathsf N_{90}+\mathsf N_{120},\qquad
 \mathsf D_\partial=90\mathsf N_{90}+120\mathsf N_{120}.
 \label{x_reciprocidad:vii:eq:ocupaciones-sectoriales}
\end{equation}
En la rama positiva, \(\gamma_{\rm BI}\) denota el parámetro de
Barbero--Immirzi ya obtenido, y \(a_0>0\) su factor areal angular.
Esta notación distingue el parámetro de la conexión y de sus holonomías.
El área se publica como
\begin{equation}
 \widehat{\mathcal A}_m=\ell_P^2\gamma_{\rm BI}a_0\mathsf N_m,
 \qquad
 \widehat{\mathcal A}=\widehat{\mathcal A}_{90}
                         +\widehat{\mathcal A}_{120}.
 \label{x_reciprocidad:vii:eq:area-sectores}
\end{equation}

Para el parámetro finito \(\Delta=\Delta_{\rm mod}\in\mathbb R\) del
estado reducido, ponemos
\begin{align}
 x_m&=e^{-m\Delta},& z_m&=1+x_m+x_m^2,&
 Z(\Delta)&=z_{90}^{18}z_{120}^{18},\\
 \rho_\partial(\Delta)&=Z(\Delta)^{-1}e^{-\Delta\mathsf D_\partial}.
 \label{x_reciprocidad:vii:eq:estado-sectorial}
\end{align}
La traza en esta sección es la traza matricial ordinaria y los estados
tienen traza uno. El parámetro \(\Delta\) conserva la selección y la
normalización del estado antecedente.

\begin{proposition}[Purificación tensorial de la frontera]
\label{x_reciprocidad:vii:prop:purificacion-sectorial}
Para un peso real finito \(w\), sea
\begin{equation}
 |\operatorname{TFD}(w)\rangle
 =\frac{|0\rangle|0\rangle+e^{-w/2}|1\rangle|1\rangle
                   +e^{-w}|2\rangle|2\rangle}
        {\sqrt{1+e^{-w}+e^{-2w}}}.
 \label{x_reciprocidad:vii:eq:tfd-local}
\end{equation}
El producto
\begin{equation}
 |\Psi_\Delta\rangle
 =|\operatorname{TFD}(90\Delta)\rangle^{\otimes18}
  \otimes|\operatorname{TFD}(120\Delta)\rangle^{\otimes18}
 \label{x_reciprocidad:vii:eq:tfd-sectorial}
\end{equation}
es una purificación normalizada de \(\rho_\partial(\Delta)\).
\end{proposition}
\begin{proof}
En cada enlace, la suma de los cuadrados de los tres coeficientes es uno.
La traza parcial elimina los términos con índices distintos y devuelve
\(\operatorname{diag}(1,e^{-w},e^{-2w})/(1+e^{-w}+e^{-2w})\).
La traza parcial conmuta con el producto tensorial. Multiplicar las
treinta y seis reducciones da el exponencial y la función de partición
de~\eqref{x_reciprocidad:vii:eq:estado-sectorial}.
\end{proof}

\section{Ecuación informacional y fluctuaciones areales}
\label{x_reciprocidad:vii:subsec:ecuacion-informacional}

Las medias y varianzas de ocupación por enlace son
\begin{equation}
 f_m=\frac{x_m+2x_m^2}{z_m},\qquad
 v_m=\frac{x_m+4x_m^2}{z_m}-f_m^2.
 \label{x_reciprocidad:vii:eq:momentos-sectoriales}
\end{equation}
Escribimos \(s(\rho)=-\operatorname{Tr}(\rho\log\rho)\), con
\(0\log0=0\). En la purificación anterior, \(s(\rho_\partial)\)
es su entropía de entrelazamiento.

\begin{proposition}[Área media, entropía y respuesta sectorial]
\label{x_reciprocidad:vii:prop:respuesta-sectorial}
Para \(\Delta\) real y finito se cumple
\begin{align}
 \langle\mathsf N\rangle_\Delta&=18(f_{90}+f_{120}),\\
 \langle\mathsf D_\partial\rangle_\Delta&=18(90f_{90}+120f_{120}),\\
 \langle\widehat{\mathcal A}\rangle_\Delta
 &=18\ell_P^2\gamma_{\rm BI}a_0(f_{90}+f_{120}),\\
 s(\Delta)&=\log Z(\Delta)+\Delta\langle\mathsf D_\partial\rangle_\Delta,
 \label{x_reciprocidad:vii:eq:entropia-sectorial}\\
 \frac{\mathrm d\langle\mathsf D_\partial\rangle_\Delta}{\mathrm d\Delta}
 &=-\operatorname{Var}_\Delta(\mathsf D_\partial)
  =-18(90^2v_{90}+120^2v_{120}),\\
 \frac{\mathrm ds}{\mathrm d\Delta}
 &=-\Delta\operatorname{Var}_\Delta(\mathsf D_\partial),\\
 \frac{\mathrm d\langle\widehat{\mathcal A}\rangle_\Delta}{\mathrm d\Delta}
 &=-18\ell_P^2\gamma_{\rm BI}a_0(90v_{90}+120v_{120}).
 \label{x_reciprocidad:vii:eq:respuestas-area-entropia}
\end{align}
\end{proposition}
\begin{proof}
Los pesos locales son \((1,x_m,x_m^2)/z_m\), cuyos dos primeros
momentos dan \(f_m\) y \(v_m\). La aditividad de los operadores
ocupación y la factorización del estado dan las tres primeras fórmulas.
Tomar la traza de
\(-\rho_\partial\log\rho_\partial
=\rho_\partial(\log Z+\Delta\mathsf D_\partial)\)
da la entropía. La diferenciación de las sumas finitas produce
\(\partial_\Delta\log Z=-\langle\mathsf D_\partial\rangle_\Delta\)
y \(\partial_\Delta^2\log Z
=\operatorname{Var}_\Delta(\mathsf D_\partial)\).
Los factores independientes suman sus varianzas. En la derivada de
\eqref{x_reciprocidad:vii:eq:entropia-sectorial} se cancelan los dos términos de
primer momento. Finalmente, \(\partial_\Delta f_m=-mv_m\)
da la respuesta del área.
\end{proof}

La pérdida de uniformidad se cuantifica en dimensión \(d_\partial=3^{36}\):
\begin{equation}
 I_\partial(\Delta)
 =36\log3-s(\Delta)
 =\mathcal D\bigl(\rho_\partial(\Delta)\Vert I/d_\partial\bigr)\ge0.
 \label{x_reciprocidad:vii:eq:deficit-informacional}
\end{equation}
Aquí \(\mathcal D\) denota la entropía relativa, definida y justificada
abajo. Sustituir \(\log(I/d_\partial)=-36\log3\,I\) prueba la
identidad. La convexidad de \(t\log t\) da la positividad; la igualdad
corresponde a la distribución uniforme, que en esta familia ocurre
exactamente en \(\Delta=0\). El déficit mide información adimensional;
la escala \(\gamma_{\rm BI}a_0\ell_P^2\) mide área por ocupación.

\begin{theorem}[Inversión de ocupaciones y área complementaria]
\label{x_reciprocidad:vii:thm:inversion-sectorial}
Sea \(\mathcal R_\partial\) la involución que intercambia
\(|0\rangle\) y \(|2\rangle\) y fija \(|1\rangle\) en cada
factor. Entonces
\begin{align}
 \mathcal R_\partial\mathsf N\mathcal R_\partial^*&=72I-\mathsf N,\\
 \mathcal R_\partial\mathsf D_\partial\mathcal R_\partial^*
 &=7560I-\mathsf D_\partial,\\
 \rho_\partial(-\Delta)&=\mathcal R_\partial\rho_\partial(\Delta)
                          \mathcal R_\partial^*,\\
 s(-\Delta)&=s(\Delta),\\
 \langle\widehat{\mathcal A}\rangle_{-\Delta}
 +\langle\widehat{\mathcal A}\rangle_\Delta
 &=72\gamma_{\rm BI}a_0\ell_P^2,\\
 \mathcal D\bigl(\rho_\partial(\Delta)\Vert\rho_\partial(-\Delta)\bigr)
 &=\Delta\bigl(7560-2\langle\mathsf D_\partial\rangle_\Delta\bigr).
 \label{x_reciprocidad:vii:eq:inversion-sectorial}
\end{align}
\end{theorem}
\begin{proof}
Cada ocupación se transforma por \(n\mapsto2-n\). En cada sector,
\(\mathsf N_m\mapsto36I-\mathsf N_m\), por lo que el término
constante del lector ponderado es \(36(90+120)=7560\).
Además, \(z_m(-\Delta)=e^{2m\Delta}z_m(\Delta)\) y
\(Z(-\Delta)=e^{7560\Delta}Z(\Delta)\). La sustitución en el
estado normalizado prueba la conjugación. Su espectro conserva la entropía;
la transformación de \(\mathsf N\) da el área complementaria. Restar
los logaritmos de los dos estados y tomar la esperanza en
\(\rho_\partial(\Delta)\) da la última igualdad.
\end{proof}

En \(\Delta=0\), \(s=36\log3\). Cuando \(\Delta\to+\infty\),
la ocupación nula concentra el estado y la entropía tiende a cero;
la inversión lleva ese límite a la ocupación máxima. Para \(\Delta>0\),
la entropía decrece estrictamente porque el lector ponderado tiene más
de un autovalor y el estado tiene rango completo. La configuración
completa, el área y la entropía conservan así funciones diferenciadas.

\section{Ley modular y variaciones finitas de área}
\label{x_reciprocidad:vii:subsec:ley-modular-finita}

Fijemos \(\rho_0=\rho_\partial(\Delta_0)\), con
\(\Delta_0\) real finito, y mantengamos
\(\Delta_0,\gamma_{\rm BI},a_0,\ell_P\) constantes. Las áreas
normalizadas \(\mathcal A_m^{\rm n}=\gamma_{\rm BI}a_0\mathsf N_m\)
expresan el operador modular como
\begin{equation}
 \mathsf K_0=-\log\rho_0
 =c_0I+\sum_{m=90,120}\beta_m\mathcal A_m^{\rm n},
 \quad \beta_m=\frac{m\Delta_0}{\gamma_{\rm BI}a_0},
 \quad c_0=\log Z(\Delta_0).
 \label{x_reciprocidad:vii:eq:operador-modular-areal}
\end{equation}
Se tiene \(\beta_{120}=\tfrac43\beta_{90}\), también cuando ambos
son cero. El cociente es \(4/3\) cuando \(\Delta_0\ne0\).
\(\mathsf K_0\) es el logaritmo de este estado de referencia; su tipo
lo distingue del transporte temporal y de la conexión gravitatoria.

\begin{theorem}[Primera ley modular y corrección finita]
\label{x_reciprocidad:vii:thm:ley-modular-finita}
Para cualquier estado normalizado \(\sigma\) del mismo espacio, sea
\(\Delta_\sigma\langle B\rangle
=\operatorname{Tr}[(\sigma-\rho_0)B]\). Entonces
\begin{equation}
 s(\sigma)-s(\rho_0)
 =\sum_{m=90,120}\beta_m\Delta_\sigma\langle\mathcal A_m^{\rm n}\rangle
  -\mathcal D(\sigma\Vert\rho_0),
 \label{x_reciprocidad:vii:eq:ley-modular-finita}
\end{equation}
donde
\(\mathcal D(\sigma\Vert\rho_0)
=\operatorname{Tr}[\sigma(\log\sigma-\log\rho_0)]\)
es finita y no negativa. Si \(\sigma(\varepsilon)\) es diferenciable,
normalizada y \(\sigma(0)=\rho_0\), se obtiene
\begin{equation}
 \left.\frac{\mathrm ds(\sigma(\varepsilon))}{\mathrm d\varepsilon}\right|_0
 =\sum_{m=90,120}\beta_m
   \left.\frac{\mathrm d\langle\mathcal A_m^{\rm n}\rangle_{
                 \sigma(\varepsilon)}}{\mathrm d\varepsilon}\right|_0.
 \label{x_reciprocidad:vii:eq:primera-ley-modular}
\end{equation}
\end{theorem}
\begin{proof}
Por definición,
\(\mathcal D(\sigma\Vert\rho_0)=-s(\sigma)
+\operatorname{Tr}(\sigma\mathsf K_0)\), mientras
\(s(\rho_0)=\operatorname{Tr}(\rho_0\mathsf K_0)\).
La resta y~\eqref{x_reciprocidad:vii:eq:operador-modular-areal} prueban la identidad
finita, pues la normalización cancela el término \(c_0I\).

Para la positividad, diagonalicemos \(\sigma\) con autovalores
\(p_i\) y vectores \(u_i\), y \(\rho_0\) con autovalores
\(q_j>0\) y vectores \(v_j\). Si
\(r_i=\langle u_i,\rho_0u_i\rangle>0\), la concavidad del
logaritmo da
\[
 \sum_j|\langle u_i,v_j\rangle|^2\log q_j\le\log r_i.
\]
En consecuencia,
\(\mathcal D(\sigma\Vert\rho_0)
\ge\sum_i p_i\log(p_i/r_i)\ge0\).
La última desigualdad es la convexidad de \(t\log t\) con pesos
\(r_i\), cuya suma es uno. La fidelidad de \(\rho_0\) mantiene
finitos sus logaritmos; los autovalores nulos de \(\sigma\) se
interpretan mediante \(0\log0=0\).

Para la fórmula infinitesimal, en el entorno del estado fiel \(\rho_0\),
la derivada de la traza de \(-\sigma\log\sigma\) es
\(-\operatorname{Tr}[\dot\sigma(\log\rho_0+I)]\).
Esta identidad se obtiene de la derivada espectral de la traza.
La condición \(\operatorname{Tr}\dot\sigma=0\) deja
\(\operatorname{Tr}(\dot\sigma\mathsf K_0)\), y sustituir
\eqref{x_reciprocidad:vii:eq:operador-modular-areal} concluye la prueba.
\end{proof}

Para las áreas físicas, los coeficientes son \(\beta_m/\ell_P^2\).
La aplicación del transductor común \(k_B\) da \(S=k_Bs\) y
conserva íntegro el término \(k_B\mathcal D\). La derivada en
\eqref{x_reciprocidad:vii:eq:respuestas-area-entropia} varía \(\Delta\); la ley
\eqref{x_reciprocidad:vii:eq:primera-ley-modular} mantiene fijo \(\Delta_0\) y
varía el estado alrededor de la referencia. Son operaciones distintas.

\section{Acción, energía de decisión y unidad areal}
\label{x_reciprocidad:vii:subsec:conversion-informacion-area}

Fijamos una misma sección orientada de acción y sus secciones gravitatoria
y térmica, en la rama positiva. Los antecedentes proporcionan
\begin{equation}
 \begin{aligned}
 \ell_P^2&=\frac{\hbar G}{c^3},& t_P&=\frac{\ell_P}{c},\\
 T_{108}&=2\pi t_P,& h&=2\pi\hbar,\\
 E_P&=\frac{\hbar}{t_P}=k_B\Theta_{\rm clk}\log3.&&
 \end{aligned}
 \label{x_reciprocidad:vii:eq:unidad-accion-informacion}
\end{equation}
La igualdad térmica conserva la pareja y la normalización previamente
construidas; al cambiar su base dimensional se transportan conjuntamente
\(k_B\) y \(\Theta_{\rm clk}\). El incremento elemental de área es
\(\delta\mathcal A=\gamma_{\rm BI}a_0\ell_P^2\).

\begin{proposition}[Coeficiente de conversión entre decisión y área]
\label{x_reciprocidad:vii:prop:tension-informacion-area}
El cociente de la energía de decisión y del incremento areal satisface
\begin{equation}
 \begin{aligned}
 \mathcal T_{I/A}:=\frac{E_P}{\delta\mathcal A}
 &=\frac{k_B\Theta_{\rm clk}\log3}{\gamma_{\rm BI}a_0\ell_P^2}\\
 &=\frac{c^3}{\gamma_{\rm BI}a_0G t_P}
  =\frac{c^4}{\gamma_{\rm BI}a_0G\ell_P}.
 \end{aligned}
 \label{x_reciprocidad:vii:eq:tension-informacion-area}
\end{equation}
Su dimensión es energía por área.
\end{proposition}
\begin{proof}
La sustitución de \(E_P=\hbar/t_P\) y
\(\ell_P^2=\hbar G/c^3\) cancela la sección de acción común.
La identidad \(\ell_P=ct_P\) da la última expresión. La positividad
de los factores asegura que todos los cocientes están definidos.
\end{proof}

La capacidad \(81\) del corte cúbico y la bipartición de los treinta y
seis enlaces sectoriales corresponden a dominios diferentes. En el
grafo cúbico \(\{0,\ldots,N-1\}^3\), las \(N^2\) rectas
paralelas entre dos caras opuestas son disjuntas por enlaces: todo corte
intercepta las \(N^2\), y una capa transversal alcanza esa cota.
Para \(N=9\), el estado producto uniforme de los trits de esa capa
tiene \(s_{\rm corte}=81\log3\). La asignación de una decisión del
mismo ciclo a cada enlace publica
\begin{equation}
 S_{\rm corte}^{\rm fis}=81k_B\log3,\qquad
 E_{\rm corte}^{(1)}=81E_P
                  =\Theta_{\rm clk}S_{\rm corte}^{\rm fis}.
 \label{x_reciprocidad:vii:eq:corte-energia-informacion}
\end{equation}
La prueba combina el corte mínimo, la aditividad de la entropía del
estado producto y~\eqref{x_reciprocidad:vii:eq:unidad-accion-informacion}. En el estado
sectorial, la ecuación finita pertinente sigue siendo
\eqref{x_reciprocidad:vii:eq:ley-modular-finita}, con su entropía relativa.

\section{Normalización cosmológica del horizonte}
\label{x_reciprocidad:vii:subsec:balance-horizonte}

Para un radio areal \(R>0\), la realización cosmológica considerada
utiliza la tríada
\begin{equation}
 E_\Lambda(R)=\frac{c^4R}{2G},\qquad
 S_{\rm H}(R)=\frac{k_B\pi R^2}{\ell_P^2},\qquad
 T_{\rm cos}(R)=\frac{\hbar c}{2\pi k_BR}.
 \label{x_reciprocidad:vii:eq:horizonte-cosmologico}
\end{equation}
El factor \(2\pi\) pertenece a esta normalización de temperatura.
El área \(\mathcal A_{\rm H}=4\pi R^2\) expresa la misma entropía
como \(S_{\rm H}/k_B=\mathcal A_{\rm H}/(4\ell_P^2)\).

\begin{theorem}[Celda gravitatoria de media acción]
\label{x_reciprocidad:vii:thm:celda-media-accion}
En~\eqref{x_reciprocidad:vii:eq:horizonte-cosmologico},
\begin{align}
 T_{\rm cos}(R)S_{\rm H}(R)&=E_\Lambda(R),\\
 T_{\rm cos}(\ell_P)S_{\rm H}(\ell_P)
 &=\frac{E_P}{2}=\frac{k_B\Theta_{\rm clk}\log3}{2},\\
 E_\Lambda(\ell_P)t_P&=\frac{\hbar}{2},&
 E_\Lambda(\ell_P)T_{108}&=\frac h2.
 \label{x_reciprocidad:vii:eq:media-accion-horizonte}
\end{align}
Manteniendo fijas las constantes de la carta y variando \(R\),
\begin{equation}
 T_{\rm cos}\,\mathrm dS_{\rm H}=2\,\mathrm dE_\Lambda,
 \qquad S_{\rm H}\,\mathrm dT_{\rm cos}=-\mathrm dE_\Lambda.
 \label{x_reciprocidad:vii:eq:diferencial-horizonte}
\end{equation}
\end{theorem}
\begin{proof}
Multiplicar temperatura y entropía da
\(\hbar cR/(2\ell_P^2)=c^4R/(2G)\).
En \(R=\ell_P\), el resultado es
\(\hbar/(2t_P)=E_P/2\). Las dos duraciones de
\eqref{x_reciprocidad:vii:eq:unidad-accion-informacion} dan las relaciones de acción.
Para los diferenciales,
\(\mathrm dE_\Lambda=c^4\mathrm dR/(2G)\),
\(\mathrm dS_{\rm H}=2\pi k_BR\mathrm dR/\ell_P^2\) y
\(\mathrm dT_{\rm cos}=-T_{\rm cos}\mathrm dR/R\).
La sustitución prueba ambas identidades.
\end{proof}

El producto \(E_\Lambda=T_{\rm cos}S_{\rm H}\) expresa una
identidad de estado. Sus variaciones radiales conservan los dos términos
de~\eqref{x_reciprocidad:vii:eq:diferencial-horizonte}; una interpretación termodinámica
completa especifica además el generador temporal, el trabajo y la
realización material correspondientes. En esta carta, acción y
gravitación fijan la unidad \(\ell_P^2\), mientras Boltzmann convierte
información en entropía dimensional. El transporte de secciones con
\(G\hbar\) fijo conserva esa unidad cuando la sección \(c\) se
mantiene fija, y conserva \(S_{\rm H}/k_B\) a radio fijo.

En la celda \(R=\ell_P\), \(S_{\rm H}/k_B=\pi\), de modo que
la multiplicidad efectiva es \(\Omega_{\rm eff}=e^\pi\).
Para la involución hiperbólica trítica \(H_{\rm trit}\), con espectro
\(\{1,-1\}\), la diagonalización da
\(\operatorname{spec}(e^{\pi H_{\rm trit}})=\{e^\pi,e^{-\pi}\}\).
Esta relación espectral especifica el valor expansivo; una identificación
microscópica de ambos espacios conserva adicionalmente su mapa de
realización y el estado.

\section{Carta de Kruskal e intercambio de los exteriores}
\label{x_reciprocidad:vii:subsec:kruskal-dos-hojas}

Fijemos un estado publicado \(x\), su radio \(R=R(x)>0\) y las
secciones positivas \(c,G\). En la carta que sigue, \(x,R,c,G\)
son parámetros fijos; \(t,r\) son las coordenadas espacio-temporales.
La realización exterior de las dos hojas se especifica por
\begin{equation}
 f(r)=1-\frac Rr,\qquad
 \mathrm ds^2=-f(r)c^2\mathrm dt^2+f(r)^{-1}\mathrm dr^2
                 +r^2\mathrm d\Omega_2^2,\qquad r>R.
 \label{x_reciprocidad:vii:eq:metrica-exterior}
\end{equation}
Para el lector de hoja \(\sigma\in\{+1,-1\}\), definimos
\begin{align}
 r_*&=r+R\log\left(\frac rR-1\right),\\
 \mathcal U&=-\sigma\exp\left(-\frac{ct-r_*}{2R}\right),&
 \mathcal V&=\sigma\exp\left(\frac{ct+r_*}{2R}\right).
 \label{x_reciprocidad:vii:eq:carta-kruskal}
\end{align}
Los símbolos \(\mathcal U,\mathcal V\) son coordenadas y se
distinguen de los operadores de transporte de rutas.

\begin{theorem}[Realización geométrica e involución de hoja]
\label{x_reciprocidad:vii:thm:kruskal-hojas}
La carta~\eqref{x_reciprocidad:vii:eq:carta-kruskal} realiza las dos hojas como los
dos exteriores \(\mathcal U\mathcal V<0\) de la extensión de
Kruskal. Se cumple
\begin{equation}
 -\mathcal U\mathcal V=\left(\frac rR-1\right)e^{r/R},\qquad
 \mathrm ds^2=-\frac{4R^3}{r}e^{-r/R}\,
                    \mathrm d\mathcal U\,\mathrm d\mathcal V
                    +r^2\mathrm d\Omega_2^2.
 \label{x_reciprocidad:vii:eq:metrica-kruskal}
\end{equation}
La involución
\begin{equation}
 \mathcal C_{\rm H}:(\sigma,\mathcal U,\mathcal V)
 \longmapsto(-\sigma,-\mathcal U,-\mathcal V)
 \label{x_reciprocidad:vii:eq:involucion-horizonte}
\end{equation}
es isométrica, fija el radio areal, intercambia los exteriores y preserva
el conjunto de horizontes \(\mathcal U\mathcal V=0\).
La métrica es regular en \(r=R\) y satisface las ecuaciones de vacío
en el dominio regular \(r>0\).
\end{theorem}
\begin{proof}
Multiplicar las coordenadas da la primera identidad. Como
\(\mathrm dr_*/\mathrm dr=f^{-1}\),
\[
 \mathrm d\mathcal U=\frac{\mathcal U}{2R}
            (-c\,\mathrm dt+f^{-1}\mathrm dr),\qquad
 \mathrm d\mathcal V=\frac{\mathcal V}{2R}
            ( c\,\mathrm dt+f^{-1}\mathrm dr).
\]
Su producto reproduce la parte radial y temporal de la métrica.
La función \(z\mapsto(z-1)e^z\) tiene derivada \(ze^z>0\)
para \(z>0\), lo que determina \(r\) a partir del producto.
En el exterior se recuperan
\(\sigma=\operatorname{sign}\mathcal V\) y
\(t=(R/c)\log(-\mathcal V/\mathcal U)\).
En \(r=R\), el coeficiente radial de la métrica extendida es
\(-4R^2/e\), finito y distinto de cero. La inversión conserva el
producto, el radio y \(\mathrm d\mathcal U\mathrm d\mathcal V\),
e intercambia las dos elecciones de signo.

Para comprobar el vacío en la carta exterior, las componentes mixtas
independientes del tensor de Einstein se reducen a
\[
 \mathsf G^t{}_t=\mathsf G^r{}_r
 =\frac{rf'+f-1}{r^2},\qquad
 \mathsf G^\theta{}_\theta=\mathsf G^\phi{}_\phi
 =\frac{f'}r+\frac{f''}2.
\]
Con \(f'=R/r^2\) y \(f''=-2R/r^3\), todas se anulan.
La expresión regular~\eqref{x_reciprocidad:vii:eq:metrica-kruskal} prolonga el
resultado a través de los horizontes en el dominio \(r>0\).
\end{proof}

Esta realización conserva
\(E_{\rm H}=c^4R/(2G)\) y
\(S_{\rm H}=k_B\pi R^2/\ell_P^2\), que dependen del radio areal.
La temperatura \(T_{\rm cos}\) de
\eqref{x_reciprocidad:vii:eq:horizonte-cosmologico} pertenece a la realización
cosmológica allí declarada. La carta exterior
\eqref{x_reciprocidad:vii:eq:metrica-exterior} conserva su propia normalización temporal.
El mapa geométrico establece el intercambio de exteriores; la comparación
con formación por colapso, espectros de radiación o modos cuasinormales
utiliza además sus condiciones físicas respectivas.

\section{Transporte de estados, observables y refinamientos}
\label{x_reciprocidad:vii:subsec:transporte-informacion-horizonte}

Denotemos por \(F_x\) la carta de los dos exteriores recién construida.
En cada carta angular regular, su jacobiano temporal y radial es
\begin{equation}
 \left|\frac{\partial(\mathcal U,\mathcal V)}{\partial(t,r)}\right|
 =\frac{cr}{2R^3}e^{r/R}>0.
 \label{x_reciprocidad:vii:eq:jacobiano-kruskal}
\end{equation}
La identidad resulta de los dos diferenciales de la prueba anterior.
Los restantes registros del estado se transportan como coordenadas de
fibra; la carta conserva su identidad y su memoria.

\begin{proposition}[Equivalencia unitaria de las lecturas de hoja y exterior]
\label{x_reciprocidad:vii:prop:transporte-unitario-horizonte}
Sea \(\mu\) una medida de hoja y \(\nu=(F_x)_*\mu\) su medida
transportada. El cambio de variables define una unitaria
\begin{equation}
 W_{\rm H}:L^2(\mu)\longrightarrow L^2(\nu),\qquad
 (W_{\rm H}\psi)(z)=\psi(F_x^{-1}z).
 \label{x_reciprocidad:vii:eq:unitaria-horizonte}
\end{equation}
Con densidades respecto de medidas coordenadas fijadas, la misma
aplicación incorpora el factor de jacobiano correspondiente.
La conjugación
\begin{equation}
 \mathfrak I_{\rm H}(B)=W_{\rm H}BW_{\rm H}^*
 \label{x_reciprocidad:vii:eq:mapa-observables-horizonte}
\end{equation}
es un \(*\)-isomorfismo unital completamente positivo. Transporta el
estado, sus esperanzas y su cálculo funcional modular. Para la medida
simétrica de las dos hojas, entrelaza también sus involuciones.
\end{proposition}
\begin{proof}
La definición de medida imagen da
\(\int|\psi\circ F_x^{-1}|^2\,\mathrm d\nu
=\int|\psi|^2\,\mathrm d\mu\).
La carta es invertible en los exteriores, de modo que la isometría es
sobreyectiva. La conjugación unitaria preserva producto, adjunto e
identidad; en cada ampliación matricial conjuga por
\(I_n\otimes W_{\rm H}\), lo que prueba la positividad completa.
Para un estado \(\rho\), la ciclicidad de la traza da
\(\operatorname{Tr}(W_{\rm H}\rho W_{\rm H}^*
\,\mathfrak I_{\rm H}(B))=\operatorname{Tr}(\rho B)\).
El cálculo funcional implica
\(g(W_{\rm H}\rho W_{\rm H}^*)=W_{\rm H}g(\rho)W_{\rm H}^*\)
en su dominio, en particular para el logaritmo y el flujo modular de
un estado fiel en su soporte. La igualdad
\(F_x\circ\mathcal C_{\rm hoja}=\mathcal C_{\rm H}\circ F_x\)
prueba el entrelazamiento de involuciones; la medida simétrica las realiza
unitariamente.
\end{proof}

% REMATE_LOCAL_X_HORIZONTE
\Needspace{25\baselineskip}
Si los refinamientos conservan \((\sigma,t,r,R)\) y utilizan la misma
isometría sobre los registros de fibra, las inclusiones
\(j_m^{\rm hoja}\) y \(j_m^{\rm ext}\) satisfacen
\begin{equation}
 W_{{\rm H},m+1}j_m^{\rm hoja}
 =j_m^{\rm ext}W_{{\rm H},m}.
 \label{x_reciprocidad:vii:eq:naturalidad-horizonte}
\end{equation}
En efecto, el cambio de variables actúa sobre las coordenadas geométricas
conservadas y la inclusión sobre las mismas coordenadas de fibra; evaluar
ambas composiciones da la misma función y la misma densidad transportada.
La igualdad extiende el mapa al sistema refinado sin reiniciar los registros.

Para una bipartición \(\mathcal H_A\otimes\mathcal H_B\), un
transporte factorizado \(U_A\otimes U_B\) cumple
\[
 \rho'_A=U_A\rho_AU_A^*,\qquad s(\rho'_A)=s(\rho_A).
\]
La primera identidad sigue de la traza parcial y la segunda del espectro.
Para una unitaria general, la formulación equivalente transporta también
el álgebra del subsistema:
\(\mathcal B(\mathcal H_A)\otimes I\mapsto
U(\mathcal B(\mathcal H_A)\otimes I)U^*\).
Así se conserva la misma pregunta operacional acerca de la información
accesible. Las igualdades de entropía se aplican a estados con las entropías
finitas pertinentes.

La relación entre información y gravitación queda expresada mediante
operaciones distintas y composables: el estado y la bipartición fijan la
información accesible; la incidencia y Barbero--Immirzi fijan la resolución
areal; acción y gravitación construyen la unidad \(\ell_P^2\);
Boltzmann realiza dimensionalmente la entropía; la carta de horizonte
transporta la geometría y sus observables. Las leyes finitas conservan la
entropía relativa, y las realizaciones geométricas conservan sus medidas,
su normalización temporal y sus dominios.

% END OWNER


% BEGIN OWNER /Users/ruben/Documents/New project/output/EXTREMA_MEDIA_RAZON_RECIPROCIDAD_ES_EN_20260918/ARTICULO/ES/source/antecedentes_vii/80_observador_y_respuesta_relacional.tex
% Recuperación y articulación de los propietarios documentados en
% PROCEDENCIA_OBSERVADOR_RESPUESTA.md. Originales conservados íntegros.
% Antecedentes de integración: núcleo APP--TRIT--TPK, realización del
% transporte de borde, lectores areal/volumétrico y reloj energético.
% Las condiciones de representación se enuncian antes de cada uso.

\chapter{Observador, memoria y respuesta relacional}
\label{x_reciprocidad:vii:sec:observador-respuesta}

La operación de lectura actúa sobre el transporte que publica el estado
enriquecido APP--TRIT--TPK. Su dominio conserva ruta, hoja, orientación,
acarreo y memoria; el resultado accesible al observador es una aplicación
de ese dominio. Esta distinción permite describir conjuntamente tres
operaciones: la elección de un representante orientado del transporte,
la inscripción de su registro en un aparato y la respuesta local inducida
por la frontera global. Cada operación transmite una información distinta.
La construcción siguiente especifica sus mapas y las condiciones que
permiten componerlos.

\section{Transporte de borde y lectura intrínseca}
\label{x_reciprocidad:vii:sec:observador-ciclo}

La publicación finita del transporte proporciona un ciclo orientado
\(C_n\), con \(n\geq1\), vértices \(0,\ldots,n-1\) y aristas
\(e_j:j\longrightarrow j+1\pmod n\). Fijamos la realización reticular
\(L\) del registro, su espacio vectorial
\(W=L\otimes_{\mathbb Z}\mathbb R\), un producto interno positivo y
el representante primitivo no nulo \(v\in L\) de la clase residual
distinguida, en la cámara declarada. Las hojas suma y producto de APP
aportan las polarizaciones del transporte; la orientación del TRIT y la
actualización del TPK preceden a esta publicación lineal. El ciclo,
la realización reticular y la clase residual son los datos que esta
sección recibe del transporte construido anteriormente.

Una \(0\)-cocadena es una función sobre los vértices con valores en
\(W\); una \(1\)-cocadena asigna un vector a cada arista orientada.
El coborde actúa por
\[
 (d\Phi)(e_j)=\Phi(j+1)-\Phi(j).
\]
Denotamos por \(\mathbf1_E\) la cocadena escalar constante en las
aristas. Para un transporte \(\omega\in C^1(C_n;W)\), su publicación
acumulada \(\Gamma_\omega\) se obtiene por sumas parciales e
interpolación afín; su defecto de cierre es
\begin{equation}
 B(\omega)=\Gamma_\omega(1)-\Gamma_\omega(0)
 =\sum_{j=0}^{n-1}\omega(e_j).
 \label{x_reciprocidad:vii:eq:observador-defecto}
\end{equation}

\begin{theorem}[Descomposición exacta y modo armónico del ciclo]
\label{x_reciprocidad:vii:thm:observador-descomposicion}
Todo transporte tiene una única descomposición
\[
 \omega=d\Phi+u\mathbf1_E,\qquad \Phi(0)=0,
 \qquad u=\frac1n\sum_j\omega(e_j).
\]
En particular, \(B(d\Phi)=0\) y \(B(\omega)=nu\).
\end{theorem}

\begin{proof}
Sea \(y_j=\omega(e_j)-u\). La suma de los \(y_j\) es cero.
Definimos \(\Phi(k)=\sum_{j=0}^{k-1}y_j\); el valor para \(k=n\)
coincide con \(\Phi(0)=0\), de modo que la función está bien definida
en el ciclo, incluida su arista final. Se cumple \(d\Phi(e_j)=y_j\).
Si \(d\Phi+u\mathbf1_E=d\Psi+w\mathbf1_E\), la suma de aristas da
\(nu=nw\). Entonces \(d(\Phi-\Psi)=0\); la conectividad del ciclo
y la normalización en el vértice inicial implican \(\Phi=\Psi\).
La identidad \(\sum_j(\Phi(j+1)-\Phi(j))=0\) prueba directamente
la cancelación telescópica.
\end{proof}

\begin{definition}[Transporte admisible y lectura normalizada]
\label{x_reciprocidad:vii:def:observador-admisible}
El transporte es admisible cuando su modo armónico pertenece a
\(\mathbb Rv\). Sobre el espacio \(\mathcal T_{\rm adm}\) de esos
transportes definimos la lectura global y su densidad por fase:
\begin{equation}
 \mathfrak a_v(\omega)
 =\frac{\langle B(\omega),v\rangle}{\langle v,v\rangle},
 \qquad
 \mu_v(\omega)=\frac{\mathfrak a_v(\omega)}n.
 \label{x_reciprocidad:vii:eq:observador-lectura}
\end{equation}
Se tiene \(B(\omega)=\mathfrak a_v(\omega)v\).
\end{definition}

\begin{theorem}[Determinación de la lectura intrínseca]
\label{x_reciprocidad:vii:thm:observador-unicidad}
El cociente
\(\mathcal T_{\rm adm}/dC^0(C_n;W)\) es unidimensional, generado
por la clase de \(v\mathbf1_E\). La aplicación \(\mu_v\) es el
único funcional lineal sobre \(\mathcal T_{\rm adm}\) que satisface
\[
 F(\omega+d\Phi)=F(\omega),\qquad F(v\mathbf1_E)=1.
\]
\end{theorem}

\begin{proof}
La descomposición anterior escribe cada transporte admisible como
\(d\Phi+\lambda v\mathbf1_E\). Su clase módulo cocadenas exactas
queda determinada por \(\lambda\). La clase de \(v\mathbf1_E\)
es no nula porque su suma es \(nv\neq0\). Por linealidad e
invariancia, todo funcional de las características indicadas vale
\(F(\omega)=\lambda F(v\mathbf1_E)=\lambda\). La fórmula
\eqref{x_reciprocidad:vii:eq:observador-lectura} da precisamente \(\mu_v=\lambda\).
La lectura global es \(\mathfrak a_v=n\mu_v\).
\end{proof}

\begin{definition}[Observador de borde]
\label{x_reciprocidad:vii:def:observador-borde}
Fijada la polarización del transporte publicado, un observador de borde
es un par \(\mathcal O=(\varepsilon,\Psi)\), donde
\(\varepsilon\in\{+1,-1\}\) y \(\Psi\in C^0(C_n;W)\).
Su acción de lectura es
\begin{equation}
 \omega^{\mathcal O}=\varepsilon\omega+d\Psi.
 \label{x_reciprocidad:vii:eq:observador-accion}
\end{equation}
El signo fija la orientación y la cocadena fija la trivialización.
Dos elecciones de polarización se comparan mediante su transporte
publicado; la relación \eqref{x_reciprocidad:vii:eq:observador-accion} describe
exactamente los representantes que difieren por orientación y coborde.
\end{definition}

\begin{proposition}[Covarianza de la lectura]
\label{x_reciprocidad:vii:prop:observador-covarianza}
La acción conserva la admisibilidad y verifica
\[
 B(\omega^{\mathcal O})=\varepsilon B(\omega),\quad
 \mathfrak a_v(\omega^{\mathcal O})=\varepsilon\mathfrak a_v(\omega),
 \quad \mu_v(\omega^{\mathcal O})=\varepsilon\mu_v(\omega).
\]
Las trivializaciones con igual orientación publican el mismo cierre.
\end{proposition}

\begin{proof}
La suma del coborde \(d\Psi\) es cero. Por tanto, el modo armónico
se transforma en \(\varepsilon u\), que sigue perteneciendo a
\(\mathbb Rv\), y el defecto se transforma en \(\varepsilon B\).
Las dos igualdades escalares se obtienen por la proyección normalizada.
\end{proof}

La unidad de esta lectura es el representante residual \(v\), con su
normalización interna. La libertad de trivialización actúa sobre el
representante de la cocadena; la información de cierre reside en su
clase. La descripción permite comparar observadores manteniendo
separados el transporte genealógico, su publicación lineal y el escalar
que evalúa la clase residual.

\section{Identidad de borde e interior celular}
\label{x_reciprocidad:vii:sec:observador-stokes}

Sea \(K\) un complejo celular finito orientado de dimensión dos,
provisto de una \(2\)-cadena fundamental \(c_K\) tal que
\(\partial c_K=[C_n]\). Esta condición expresa que las incidencias
interiores se cancelan con sus multiplicidades y orientaciones.
Sea \(\Omega\in C^1(K;W)\) una extensión de \(\omega\) y
\(F_\Omega=d\Omega\) su curvatura discreta aditiva.

\begin{theorem}[Identidad holográfica celular de la lectura]
\label{x_reciprocidad:vii:thm:observador-stokes}
Con los datos anteriores,
\begin{equation}
 B(\omega)=\langle\Omega,\partial c_K\rangle
          =\langle d\Omega,c_K\rangle,
 \qquad
 \mathfrak a_v(\omega)
 =\frac{\langle\langle d\Omega,c_K\rangle,v\rangle}
        {\langle v,v\rangle}.
 \label{x_reciprocidad:vii:eq:observador-stokes}
\end{equation}
En particular, \(F_\Omega=0\) implica \(\mathfrak a_v(\omega)=0\).
Si el cierre es no nulo, toda extensión del transporte en ese complejo
tiene curvatura discreta no nula.
\end{theorem}

\begin{proof}
El emparejamiento entre cadenas y cocadenas satisface
\(\langle d\Omega,c_K\rangle=\langle\Omega,\partial c_K\rangle\).
Al expandirlo, cada arista interior aparece con la suma de los números
de incidencia de las celdas contiguas; la condición sobre
\(\partial c_K\) deja exactamente el ciclo orientado de borde.
La restricción de \(\Omega\) a ese ciclo es \(\omega\), cuya suma
es \(B(\omega)\). La proyección sobre \(v\) da la segunda fórmula
y la implicación de planitud.
\end{proof}

La igualdad compara dos evaluaciones de la misma cocadena extendida:
el defecto acumulado del borde y el coborde integrado sobre el interior.
Su curvatura toma valores en \(W\) y usa el coborde celular aditivo.
La realización por una conexión geométrica conserva, adicionalmente,
su propio espacio de coeficientes y su ley de transporte.

\section{Inscripción isométrica del registro en el aparato}
\label{x_reciprocidad:vii:sec:observador-aparato}

Sea \(\mathfrak M_n\) un conjunto finito no vacío de registros
publicados a nivel \(n\), con resultado, ruta, acarreo, hoja y memoria.
Su espacio de registros es
\(\mathcal H_n=\ell^2(\mathfrak M_n)\), con base \(|m\rangle\).
La actualización visible \(F_n\) tiene un levantamiento de registro
\[
 \widehat F_n(m)=(F_n(m),m).
\]
Tomamos como registros del nivel siguiente un conjunto que contiene esas
parejas. La segunda coordenada conserva el antecedente y hace inyectivo
el levantamiento. Esta construcción especifica la memoria retenida por
la inscripción, también cuando la actualización visible reúne valores.

En un espacio finito de aparato \(\mathcal K_{A,n}\), fijamos una
familia ortonormal \((|a_m\rangle)_{m\in\mathfrak M_n}\) y definimos
\[
 \iota_n:\mathcal H_n\longrightarrow\mathcal K_{A,n},
 \quad \iota_n|m\rangle=|a_m\rangle,
 \qquad
 V_n:\mathcal H_n\longrightarrow\mathcal H_{n+1},
 \quad V_n|m\rangle=|\widehat F_n(m)\rangle.
\]
Sea \(\mathcal P_n=\iota_n\mathcal H_n\) el subespacio de puntero.

\begin{theorem}[Compatibilidad entre actualización e inscripción]
\label{x_reciprocidad:vii:thm:observador-entrelazador}
Los mapas \(\iota_n\) y \(V_n\) son isometrías. Existe una única
isometría \(W_n^0:\mathcal P_n\longrightarrow\mathcal K_{A,n+1}\)
que verifica
\begin{equation}
 W_n^0\iota_n=\iota_{n+1}V_n.
 \label{x_reciprocidad:vii:eq:observador-cuadrado}
\end{equation}
Si \(\dim\mathcal K_{A,n+1}\geq\dim\mathcal K_{A,n}\),
se prolonga a una isometría de todo \(\mathcal K_{A,n}\).
Mediante complementos auxiliares que igualen las dimensiones, se obtiene
una prolongación unitaria entre los espacios ampliados.
\end{theorem}

\begin{proof}
Las imágenes de las bases por \(\iota_n\) son ortonormales; las
imágenes por \(V_n\) también lo son por la inyectividad de
\(\widehat F_n\). Para \(z=\iota_n\psi\) definimos
\(W_n^0z=\iota_{n+1}V_n\psi\). La representación de \(z\) es única,
y los tres mapas conservan su norma. La fórmula determina el mapa en
una base de \(\mathcal P_n\) y prueba el cuadrado.
Si la dimensión de llegada es suficiente, el complemento ortogonal de
\(W_n^0\mathcal P_n\) contiene una familia ortonormal del tamaño de
\(\mathcal P_n^\perp\); enviando una base de este último a aquella
se construye la prolongación. Finalmente se amplían ambos espacios a
una misma dimensión y se completan las dos familias prescritas a bases
ortonormales. La aplicación entre esas bases es unitaria.
\end{proof}

Cuando el registro incluye rótulos de un instrumento finito, esta misma
inscripción realiza su ambiente: si los operadores
\(K_{y,a}:\mathcal H_S\to\mathcal H'_S\) satisfacen
\(\sum_{y,a}K_{y,a}^*K_{y,a}=I\), la aplicación
\[
 V_{\rm inst}\psi=\sum_{y,a}K_{y,a}\psi\otimes|y,a\rangle
\]
es isométrica, pues su norma al cuadrado es
\(\sum_{y,a}\langle\psi,K_{y,a}^*K_{y,a}\psi\rangle=\|\psi\|^2\).
La inscripción \(\iota\) transporta cada rótulo al puntero que le
corresponde. El modelo fija así el registro de la premedición; un
dispositivo material requiere además su interacción y su calibración.


% BEGIN OWNER /Users/ruben/Documents/New project/output/EXTREMA_MEDIA_RAZON_RECIPROCIDAD_ES_EN_20260918/ARTICULO/ES/source/antecedentes_vii/80b_salida_y_condicionamiento.tex
\section{Aplicación de salida y condicionamiento del registro}
\label{x_reciprocidad:vii:sec:salida-condicionamiento}

La inscripción anterior se compone con el lector del aparato sobre el
estado genealógico completo. El acoplamiento conserva la operación que
produce la correlación; la aplicación de salida publica el resultado; el
condicionamiento restringe las prolongaciones compatibles con ese registro.

Sean \((\Omega_S,\Sigma_S)\) y \((\Omega_M,\Sigma_M)\) los espacios
medibles de estados genealógicos del sistema y del aparato, obtenidos de
sus cilindros y registros compatibles. Un contexto de medición se escribe
\(\mathfrak M_a=(U_a,q_a,\mathcal C_a,m_0)\): \(m_0\) es la
preparación del aparato, \(\mathcal C_a\) fija las compatibilidades y el
horizonte de prolongación, y
\[
 U_a:\Omega_S\times\Omega_M\longrightarrow\Omega_S\times\Omega_M,
 \qquad q_a:\Omega_M\longrightarrow Y_a
\]
son, respectivamente, el acoplamiento medible y el lector medible hacia
el espacio de resultados \((Y_a,\Sigma_{Y_a})\). En un régimen reversible,
\(U_a\) es biyectivo y bimedible; en la realización de Hilbert se exige
asimismo la unitariedad correspondiente. La salida individual es la composición
\begin{equation}
 \operatorname{Out}_a(x)
 =q_a\bigl(\operatorname{pr}_M U_a(x,m_0)\bigr).
 \label{x_reciprocidad:vii:eq:out-genealogico}
\end{equation}
Queda determinada por el estado completo, la preparación del aparato y
el contexto. La descripción estadística de una preparación constituye
otra lectura del conjunto de estados compatibles.

\begin{proposition}[Medida de resultados y restricción a la fibra]
\label{x_reciprocidad:vii:prop:medida-salida}
Sea \(\Omega_a\in\Sigma_S\) el conjunto de estados compatibles con la
preparación y el contexto, provisto de una probabilidad \(\mu_a\) sobre
\((\Omega_a,\Sigma_S|_{\Omega_a})\).
La aplicación \(\operatorname{Out}_a|_{\Omega_a}\) es medible y produce
la probabilidad de resultados
\begin{equation}
 \nu_a(B)=\mu_a\bigl(\operatorname{Out}_a^{-1}(B)\cap\Omega_a\bigr),
 \qquad B\in\Sigma_{Y_a}.
 \label{x_reciprocidad:vii:eq:medida-imagen-salida}
\end{equation}
Para un resultado \(y\) con singleton medible, la fibra
\(\Omega_{a,y}=\{x\in\Omega_a:\operatorname{Out}_a(x)=y\}\)
es medible. Cuando \(\mu_a(\Omega_{a,y})>0\),
\begin{equation}
 \mu_{a,y}(E)=
 \frac{\mu_a(E\cap\Omega_{a,y})}{\mu_a(\Omega_{a,y})}
 \label{x_reciprocidad:vii:eq:condicionamiento-salida}
\end{equation}
para \(E\in\Sigma_S|_{\Omega_a}\) define una probabilidad concentrada
en esa fibra.
\end{proposition}
\begin{proof}
La inclusión \(x\mapsto(x,m_0)\), el acoplamiento, la proyección sobre
el aparato y su lector son medibles; su composición también lo es.
Las preimágenes conservan uniones disjuntas y el conjunto total, por lo
que \eqref{x_reciprocidad:vii:eq:medida-imagen-salida} es numerablemente aditiva y tiene
masa uno. La fibra es la preimagen del singleton \(\{y\}\) dentro de
\(\Omega_a\). La restricción de \(\mu_a\) a esa fibra es una medida
positiva y numerablemente aditiva; dividir por su masa positiva la normaliza.
\end{proof}

Para un prefijo \(s\) de profundidad \(n\), las prolongaciones posteriores
al registro forman
\[
 \operatorname{Fut}_{a,y}(s)
 =\{x\in\Omega_{a,y}:x|_n=s\}.
\]
La restricción conserva el prefijo ya recorrido y el registro del
acoplamiento; modifica la familia de prolongaciones admitidas. La salida
\eqref{x_reciprocidad:vii:eq:out-genealogico} actúa sobre estados completos; la ley de
resultados \eqref{x_reciprocidad:vii:eq:medida-imagen-salida} se obtiene transportando
la medida de preparación por esa aplicación. La representación por una
matriz de densidad expresa el nivel
estadístico de la realización cuántica y conserva sus condiciones de
correspondencia con las fibras genealógicas.

% END OWNER


\section{Álgebra accesible y complemento del puntero}
\label{x_reciprocidad:vii:sec:observador-expectativa}

Definimos \(P_m=|a_m\rangle\langle a_m|\),
\(P=\sum_mP_m\) y \(P_\perp=I-P\). El álgebra del subespacio de
puntero es \(P\mathcal B(\mathcal K_{A,n})P\), cuya unidad es \(P\).
En ella, el lector diagonal es
\[
 \mathcal D_P(A)=\sum_mP_mAP_m.
\]

\begin{proposition}[Expectativa del registro y extensión unital]
\label{x_reciprocidad:vii:prop:observador-pinzamiento}
La aplicación \(\mathcal D_P\) es una expectativa condicional
completamente positiva, preservadora de traza e idempotente sobre el
álgebra del puntero, con \(\mathcal D_P(P)=P\). Su imagen es
\(\bigoplus_m\mathbb CP_m\).
Sobre el álgebra del aparato completo, la aplicación
\begin{equation}
 \mathcal E_P(A)=\sum_mP_mAP_m+P_\perp AP_\perp
 \label{x_reciprocidad:vii:eq:observador-complemento}
\end{equation}
es una expectativa condicional unital, completamente positiva y
preservadora de traza. Su imagen es
\[
 \left(\bigoplus_m\mathbb CP_m\right)
 \oplus\mathcal B(\mathcal P_n^\perp).
\]
Ambas aplicaciones son proyecciones ortogonales para el producto de
Hilbert--Schmidt. Sus espectros están contenidos en \(\{0,1\}\).
Cuando existen componentes eliminadas, la separación entre el
subespacio fijo y el núcleo es uno.
\end{proposition}

\begin{proof}
Los proyectores \(P_m\) son mutuamente ortogonales. En el aparato
completo, la familia que añade \(P_\perp\) suma \(I\). Cada sumando
\(A\mapsto P_mAP_m\) es completamente positivo: tras cualquier
ampliación matricial actúa por congruencia con \(I\otimes P_m\).
Lo mismo ocurre para \(P_\perp\). La suma conserva esa propiedad.
La ciclicidad de la traza y la suma de los proyectores prueban la
preservación de traza en cada dominio. La ortogonalidad anula los
términos cruzados al iterar el mapa, por lo que es idempotente.
Sus elementos fijos son precisamente los bloques anunciados y la
aplicación es bimodular respecto de esa subálgebra.
Además,
\(\operatorname{Tr}(A^*P_mBP_m)
 =\operatorname{Tr}((P_mAP_m)^*B)\),
de donde resulta la autoadjunción para Hilbert--Schmidt. Una proyección
autoadjunta tiene únicamente los valores espectrales cero y uno cuando
ambos subespacios son no nulos. Si sólo queda el subespacio fijo,
la aplicación es la identidad.
\end{proof}

La suma \(\sum_mP_mAP_m\), aplicada sin el complemento a todo el
aparato, envía \(I\) a \(P\). Su unidad natural pertenece al álgebra
del puntero. La fórmula \eqref{x_reciprocidad:vii:eq:observador-complemento} conserva
también el bloque auxiliar; una resolución ortonormal de ese bloque
proporciona, cuando se desea, la diagonalización del aparato entero.

Si \(U\) es una transformación unitaria del aparato y
\(P'_m=UP_mU^*\), entonces
\[
 \mathcal E_{P'}(UAU^*)=U\mathcal E_P(A)U^*.
\]
La igualdad se obtiene sustituyendo cada proyector en la suma.
Asimismo, para \(A\in\mathcal B(\mathcal P_n)\), el cuadrado
\eqref{x_reciprocidad:vii:eq:observador-cuadrado} implica
\[
 \mathcal D_{P_{n+1}}(W_n^0A(W_n^0)^*)
 =W_n^0\mathcal D_{P_n}(A)(W_n^0)^*.
\]
En efecto, cada unidad matricial \(|a_m\rangle\langle a_{m'}|\)
se transporta a la unidad correspondiente a los registros distintos
\(\widehat F_n(m),\widehat F_n(m')\); el lector diagonal conserva
exactamente los términos con \(m=m'\).

\section{Factorización de la respuesta por la frontera efectiva}
\label{x_reciprocidad:vii:sec:observador-factorizacion}

Sea \(\mathcal X_\infty^{\rm enr}\) el dominio de estados compatibles
enriquecidos y sea
\(b:\mathcal X_\infty^{\rm enr}\longrightarrow B_\infty\) su traza
global de frontera. Denotamos por \(B_{\rm ef}=\operatorname{im}b\)
la frontera efectivamente publicada. Para un espacio de respuestas
\(\mathcal V_v\), consideramos
\(I_v:\mathcal X_\infty^{\rm enr}\longrightarrow\mathcal V_v\).

\begin{theorem}[Criterio de respuesta relacional]
\label{x_reciprocidad:vii:thm:observador-factorizacion}
Existe una única aplicación \(\overline I_v:B_{\rm ef}\to\mathcal V_v\)
tal que \(I_v=\overline I_v\circ b\) si y sólo si
\begin{equation}
 b(x)=b(x')\quad\Longrightarrow\quad I_v(x)=I_v(x').
 \label{x_reciprocidad:vii:eq:observador-fibra}
\end{equation}
La unicidad se extiende a todo \(B_\infty\) cuando \(b\) es
sobreyectiva.
\end{theorem}

\begin{proof}
Una composición por \(b\) es constante en cada fibra. Recíprocamente,
para \(\beta\in B_{\rm ef}\), se define
\(\overline I_v(\beta)=I_v(x)\) con \(b(x)=\beta\).
La condición \eqref{x_reciprocidad:vii:eq:observador-fibra} garantiza independencia
del representante; todo valor de \(\overline I_v\) queda determinado
por algún estado de esa fibra. Si \(b\) es sobreyectiva,
\(B_{\rm ef}=B_\infty\). En otro caso, la ecuación de composición
determina únicamente los valores sobre \(B_{\rm ef}\).
\end{proof}

Este criterio es conjuntista. Para una factorización continua, se
conservan además las topologías y la condición de aplicación cociente
de \(b:\mathcal X_\infty^{\rm enr}\to B_{\rm ef}\); con ella, la
continuidad de \(I_v\) equivale a la de su factor. Esta precisión
separa la determinación de una lectura de sus propiedades de regularidad.

Sea \(\ell_v:\mathcal X_\infty^{\rm enr}\to\mathcal L_v\) la
restricción accesible a una vecindad local. El carácter global efectivo
se verifica mediante dos estados que satisfagan
\begin{equation}
 \ell_v(x)=\ell_v(x'),\qquad I_v(x)\neq I_v(x').
 \label{x_reciprocidad:vii:eq:observador-testigo-global}
\end{equation}
Si \(I_v\) factoriza por \(b\), estos dos estados tienen fronteras
distintas. Además, la respuesta no puede escribirse como función de
\(\ell_v\), pues esa función asignaría el mismo valor a ambos estados.
La diferencia de fronteras adquiere así un efecto operatorio comprobable.

\section{Pesos de ambivalencia y respuesta autoinercial}
\label{x_reciprocidad:vii:sec:observador-ambivalencia}

Para la fórmula general consideramos lectores positivos
\(A_\partial,V_\partial\) y un lector energético \(Q\)
definidos sobre el dominio efectivo de frontera declarado.
A continuación se explicita una realización sobre
\(B_{\rm av}\subseteq B_{\rm ef}\). La memoria permanece como coordenada
\(\mathsf M_\partial(\beta)\) de la frontera. La bisagra de
ambivalencia proporciona las dos orientaciones
\[
 r_{\rm av}=\frac{\pi_{\rm HMT}}{\varphi_{\rm HMT}^2},\qquad
 r_{\rm av}^{J}=\frac{\varphi_{\rm HMT}}{\pi_{\rm HMT}^2}.
\]
El superíndice \(J\) expresa intercambio del par generador.
\subsubsection{Realización normalizada de los lectores de frontera.}
La incidencia modal APP--TRIT--TPK, desarrollada en la construcción
regional, proporciona el envolvente de caminos de matriz
\[
 F_{\rm av}=\begin{pmatrix}0&1\\1&1\end{pmatrix}.
\]
El vector positivo \(v=(1,\varphi_{\rm HMT})^{\mathsf T}\)
satisface \(F_{\rm av}v=\varphi_{\rm HMT}v\). Su transición
normalizada y sus pesos estacionarios son
\begin{equation}
 P_{ij}=\frac{(F_{\rm av})_{ij}v_j}{\varphi_{\rm HMT}v_i},
 \qquad
 P=\begin{pmatrix}0&1\\
 \varphi_{\rm HMT}^{-2}&\varphi_{\rm HMT}^{-1}\end{pmatrix},
 \qquad
 (\mu_{\rm ar},\mu_{\rm vol})
 =\frac{(1,\varphi_{\rm HMT}^{2})}{1+\varphi_{\rm HMT}^{2}}.
 \label{x_reciprocidad:vii:eq:lectores-parry}
\end{equation}
Las filas de \(P\) suman uno por
\(\varphi^{-2}+\varphi^{-1}=1\); la ecuación estacionaria
\(\mu_{\rm ar}=\mu_{\rm vol}\varphi^{-2}\), junto con la
normalización, determina los pesos escritos. Esta medida de Parry
pondera los caminos del envolvente; la frecuencia cronológica
\((1/3,2/3)\) cuenta los instantes del calendario modal y conserva
su función distinta.

La realización de amplitud común de memoria se define sobre el sector
\(\mathcal X_{\rm av}\subset\mathcal X_\infty^{\rm enr}\) en el que
los lectores de las dos hojas tienen la factorización
\begin{equation}
 \mathcal L_{\rm vol}(x)=\mu_{\rm vol}\mathcal M(x),\qquad
 \mathcal L_{\rm ar}(x)=\pi_{\rm HMT}\mu_{\rm ar}\mathcal M(x),
 \qquad \mathcal M(x)>0.
 \label{x_reciprocidad:vii:eq:lectores-amplitud-comun}
\end{equation}
Ambos valores pertenecen a la misma recta positiva de amplitud de
memoria. La publicación areal lleva la marca regional
\(\pi_{\rm HMT}\) del retorno. La lectura adimensional se obtiene
dividiendo ambas secciones por \(\mathcal L_{\rm vol}(x)\).
De este modo las sumas de los pesos se efectúan entre magnitudes del
mismo tipo; una eventual realización como áreas y volúmenes físicos
se compone después con sus respectivos transductores dimensionales.

\begin{proposition}[Descenso de los lectores normalizados y del reloj energético]
\label{x_reciprocidad:vii:prop:lectores-frontera-explicitos}
Sobre \(B_{\rm av}=b(\mathcal X_{\rm av})\), las reglas
\begin{equation}
 A_\partial(b(x))
 =\frac{\mathcal L_{\rm ar}(x)}{\mathcal L_{\rm vol}(x)}
 =r_{\rm av},\qquad V_\partial(b(x))=1
 \label{x_reciprocidad:vii:eq:lectores-normalizados}
\end{equation}
definen dos lectores positivos independientes del representante.
Fijada una orientación \(a\) de la sección de acción, la energía de
decisión de \eqref{x_reciprocidad:v:ant:eq:ii-kb-aut-energia} determina sobre estados
\begin{equation}
 q_a(x)=\frac{h_a(x)}{108t_0(x)}
       =\frac{\hbar_a(x)}{t_P(x)}\in\mathcal L_E^+,
 \qquad t_P=\frac{54}{\pi_{\rm HMT}}t_0.
 \label{x_reciprocidad:vii:eq:lector-q-accion}
\end{equation}
Escribimos
\[
 b_{\rm av}:=b|_{\mathcal X_{\rm av}}:
 \mathcal X_{\rm av}\longrightarrow B_{\rm av}.
\]
Existe un único lector \(Q:B_{\rm av}\to\mathcal L_E^+\) tal que
\(q_a|_{\mathcal X_{\rm av}}=Q\circ b_{\rm av}\)
si y sólo si, para cualesquiera \(x,x'\in\mathcal X_{\rm av}\),
\begin{equation}
 b(x)=b(x')\ \Longrightarrow\
 \frac{h_a(x)}{t_0(x)}=\frac{h_a(x')}{t_0(x')}.
 \label{x_reciprocidad:vii:eq:descenso-reloj}
\end{equation}
En particular, una carta de acción y duración fijas satisface esta
condición y publica \(Q=h_a/(108t_0)\).
\end{proposition}
\begin{proof}
El cociente de \eqref{x_reciprocidad:vii:eq:lectores-amplitud-comun} cancela la
amplitud común y vale
\(\pi_{\rm HMT}\mu_{\rm ar}/\mu_{\rm vol}
=\pi_{\rm HMT}/\varphi_{\rm HMT}^2\).
El resultado es el mismo para todos los representantes de una fibra
de \(b_{\rm av}\) dentro de \(\mathcal X_{\rm av}\) y es positivo.
Dividir ambos lectores por el mismo factor positivo
conserva los cocientes ponderados en los que intervienen.
Para el reloj, el criterio de factorización del
teorema~\ref{x_reciprocidad:vii:thm:observador-factorizacion}, aplicado a \(q_a\),
equivale exactamente a \eqref{x_reciprocidad:vii:eq:descenso-reloj}. Las secciones
fijas lo satisfacen; también lo satisfacen secciones variables cuyo
cociente descienda por \(b_{\rm av}\). La unicidad resulta de la definición
de \(B_{\rm av}\) como imagen efectiva.
\end{proof}

El valor del lector \(Q\) pertenece a la recta de energía
\(\mathcal L_E=\mathcal L_{\rm act}\otimes\mathcal L_T^*\).
Con \(E_1,E_2\) en esa misma recta,
\(E_1E_2/Q^2\) es adimensional e independiente de la base energética
positiva. El símbolo \(Q\) de esta sección designa exclusivamente
esa energía de reloj; la cantidad de calor de la realización de
Landauer pertenece al balance térmico correspondiente.
La factorización de amplitud común especifica una realización
concreta de los lectores generales. Cada otro dominio conserva sus
aplicaciones de frontera y su propia comprobación del descenso.
La memoria completa sigue en la historia enriquecida y en su
publicación de frontera, aunque se cancele en el cociente normalizado.


Definimos
\[
 w_A(\beta)=
 \frac{A_\partial(\beta)r_{\rm av}}
 {A_\partial(\beta)r_{\rm av}+V_\partial(\beta)r_{\rm av}^{J}},
 \qquad w_V(\beta)=1-w_A(\beta).
\]
En la realización \eqref{x_reciprocidad:vii:eq:lectores-normalizados}, los pesos
admiten la evaluación exacta
\begin{equation}
 w_A=\frac{r_{\rm av}^{\,2}}{r_{\rm av}^{\,2}+r_{\rm av}^{J}}
     =\frac{\pi_{\rm HMT}^{4}}
            {\pi_{\rm HMT}^{4}+\varphi_{\rm HMT}^{5}},\qquad
 w_V=\frac{\varphi_{\rm HMT}^{5}}
            {\pi_{\rm HMT}^{4}+\varphi_{\rm HMT}^{5}}.
 \label{x_reciprocidad:vii:eq:pesos-amplitud-comun}
\end{equation}
La sustitución conserva los dos factores de la definición: el lector
areal normalizado y su ponderación posterior por \(r_{\rm av}\).
La generación y el marcado regional del retorno se realizan una vez.
Multiplicar numerador
y denominador por \(\varphi_{\rm HMT}^4\pi_{\rm HMT}^2\)
da la última expresión. Estos pesos permanecen constantes en el
sector de amplitud común. Con proyectores y energías fijos, una
variación de la respuesta puede proceder del lector \(Q\);
el testigo de dependencia global requiere además los estados con
igual restricción local y distinta respuesta de
\eqref{x_reciprocidad:vii:eq:observador-testigo-global}.


Sean \(P_A,P_V\) proyecciones ortogonales complementarias de un
espacio de respuesta \(\mathcal H_v\), y sean \(E_1,E_2>0\)
dos energías internas fijadas para la comparación. La respuesta es
\begin{equation}
 \overline I_v(\beta)=\frac{E_1E_2}{Q(\beta)^2}
       \bigl(w_A(\beta)P_A+w_V(\beta)P_V\bigr),
 \qquad I_v=\overline I_v\circ b.
 \label{x_reciprocidad:vii:eq:observador-respuesta-ponderada}
\end{equation}
También se admiten energías que sean lectores explícitos de
\(\beta\); en ese caso se conservan como tales en la fórmula.

\begin{proposition}[Positividad y covarianza de la respuesta]
\label{x_reciprocidad:vii:prop:observador-respuesta-positiva}
La respuesta \eqref{x_reciprocidad:vii:eq:observador-respuesta-ponderada} es
autoadjunta, positiva definida y constante en las fibras de \(b\).
Sus valores propios en los sectores no nulos son
\[
 \lambda_A=E_1E_2Q^{-2}w_A,\qquad
 \lambda_V=E_1E_2Q^{-2}w_V.
\]
El intercambio completo
\((A_\partial,r_{\rm av},P_A)
 \leftrightarrow(V_\partial,r_{\rm av}^{J},P_V)\)
permuta los dos sumandos y conserva el operador. Un transporte
unitario del espacio de respuesta conjuga ese operador y conserva
su espectro con multiplicidades.
\end{proposition}

\begin{proof}
La positividad de los lectores da \(0<w_A,w_V<1\). Para
\(z=z_A+z_V\), con \(z_A=P_Az\), \(z_V=P_Vz\),
\[
 \langle z,\overline I_vz\rangle
 =\frac{E_1E_2}{Q^2}
   (w_A\|z_A\|^2+w_V\|z_V\|^2)>0
\]
si \(z\neq0\). La descomposición ortogonal prueba las expresiones
espectrales. La dependencia exclusiva de \(\beta\) da constancia
en las fibras y la factorización. El intercambio de las ternas
intercambia numeradores, pesos y proyectores a la vez, por lo que
conserva la suma. Para un transporte unitario \(U\), la sustitución
\(P_A\mapsto UP_AU^*\), \(P_V\mapsto UP_VU^*\) produce
\(U\overline I_vU^*\), con el mismo espectro.
\end{proof}

Un intercambio implementado dentro de un mismo espacio por una unitaria
que envíe \(P_A\) a \(P_V\) requiere que ambos sectores tengan
igual dimensión. El intercambio de las dos ternas como rótulos de la
fórmula es válido con las dimensiones propias de cada sector.

La condición \eqref{x_reciprocidad:vii:eq:observador-testigo-global} exige comparar
las respuestas. Por ejemplo, multiplicar simultáneamente
\(A_\partial,V_\partial\) por un mismo factor positivo conserva
ambos pesos; si también se conservan \(E_1,E_2,Q\), la respuesta
permanece idéntica. En cambio, con iguales pesos y energías, sustituir
\(Q\) por \(qQ\), \(q>0\), multiplica la respuesta por \(q^{-2}\).
Si dos estados con igual restricción local realizan ese cambio con
\(q\neq1\), proporcionan el testigo global. Estas son condiciones
operatorias de separación; la existencia de los estados se comprueba
en el dominio de frontera que se esté realizando.

\section{Conexión del estado conjunto y lectura de aceleración}
\label{x_reciprocidad:vii:sec:observador-autoinercia}

En una realización diferenciable del estado conjunto
\(\Gamma=(x,\mathcal O,m)\), sean \(N\) su espacio de realización,
\(\nabla^{\rm tot}\) su conexión y \(\pi_S:N\to N_S\) una
proyección diferenciable hacia el subsistema, dotado de conexión
\(\nabla^S\). La conexión y la proyección conservan la procedencia
de los transportes y lectores que realizan. El marco observado se
obtiene mediante
\[
 \mathfrak F_{\mathcal O}(\Gamma)
 =\operatorname{Res}_{\mathcal O}(\Gamma,\nabla^{\rm tot}).
\]
Para una trayectoria dos veces diferenciable \(\Gamma(t)\),
escribimos \(\gamma_S=\pi_S\circ\Gamma\).

\begin{definition}[Autoinercia y defecto de proyección]
\label{x_reciprocidad:vii:def:observador-autoinercia}
Una trayectoria es autoinercial respecto de la conexión declarada
cuando \(\nabla^{\rm tot}_{\dot\Gamma}\dot\Gamma=0\).
El defecto de su proyección es
\begin{equation}
 \mathcal K_{\pi_S}(\dot\Gamma,\dot\Gamma)
 =\nabla^S_{\dot\gamma_S}\dot\gamma_S
  -d\pi_S\bigl(\nabla^{\rm tot}_{\dot\Gamma}\dot\Gamma\bigr).
 \label{x_reciprocidad:vii:eq:observador-defecto-proyeccion}
\end{equation}
\end{definition}

\begin{proposition}[Aceleración del subsistema]
\label{x_reciprocidad:vii:prop:observador-aceleracion}
Para una trayectoria autoinercial,
\[
 a_S=\nabla^S_{\dot\gamma_S}\dot\gamma_S
     =\mathcal K_{\pi_S}(\dot\Gamma,\dot\Gamma).
\]
En coordenadas \(x^B\) de \(N\) y \(y^a\) de \(N_S\), el
defecto es cuadrático en la velocidad, con componentes
\begin{equation}
 (\mathcal K_{\pi_S})^a{}_{BC}
 =\partial_B\partial_C\pi_S^a
  +(\Gamma^S)^a{}_{bc}\,
        \partial_B\pi_S^b\partial_C\pi_S^c
  -\partial_D\pi_S^a(\Gamma^{\rm tot})^D{}_{BC}.
 \label{x_reciprocidad:vii:eq:observador-defecto-coordenadas}
\end{equation}
\end{proposition}

\begin{proof}
La regla de la cadena da
\(\ddot\gamma_S^a=\partial_B\pi_S^a\ddot\Gamma^B
 +\partial_B\partial_C\pi_S^a\dot\Gamma^B\dot\Gamma^C\).
Sumando la contribución de \(\nabla^S\) y restando la imagen de
la aceleración total, los términos que contienen \(\ddot\Gamma\)
se cancelan. Los restantes son
\eqref{x_reciprocidad:vii:eq:observador-defecto-coordenadas} contraídos con
\(\dot\Gamma^B\dot\Gamma^C\). Al anularse la aceleración total
resulta la igualdad anunciada.
\end{proof}

La respuesta geométrica factoriza por \(b\) cuando la conexión,
la proyección y el dato cinemático utilizado en la evaluación se
determinan por \(\beta=b(x)\), o cuando esos datos cinemáticos se
mantienen fijados en la comparación. Con esas condiciones,
\eqref{x_reciprocidad:vii:eq:observador-defecto-coordenadas} da el mismo resultado
para cualesquiera representantes de una fibra; el
teorema~\ref{x_reciprocidad:vii:thm:observador-factorizacion} construye su factor
único sobre \(B_{\rm ef}\). Si varía el dato cinemático, éste
forma parte del dominio de la respuesta que se compara.

La formulación autoinercial describe una trayectoria geodésica del
estado conjunto y una aceleración efectiva de su subsistema debida
al defecto de proyección. La realización gravitatoria identifica
después esa conexión, el marco y la respuesta con sus observables
físicos. La covarianza del observador, la inscripción de memoria y
la factorización por la frontera conservan así funciones complementarias:
la primera determina la clase de lectura, la segunda retiene la
genealogía y la tercera especifica la dependencia relacional de la
respuesta.

% END OWNER


% BEGIN OWNER /Users/ruben/Documents/New project/output/EXTREMA_MEDIA_RAZON_RECIPROCIDAD_ES_EN_20260918/ARTICULO/ES/source/antecedentes_vii/81_propagacion_bidireccional.tex
% Recuperación de 13_accion_interaccion_absorbedor y83_operador_absorbedor.
% Desarrollo del canal cronológico desde el transporte de rutas.
% Procedencia y condiciones en PROCEDENCIA_PROPAGACION_BIDIRECCIONAL.md.
% Los originales y la sección80 permanecen íntegros. Sin contenido de radión.

\chapter{Propagación bidireccional y condición de absorción en la frontera}
\label{x_reciprocidad:vii:sec:propagacion-bidireccional}

El transporte APP--TRIT--TPK conserva la orientación de la ruta y la memoria
que distingue su ida de su retorno. Sobre ese antecedente se construyen tres
operaciones sucesivas: la acción de una historia enriquecida, la propagación
de una fuente a través de los transportes realizados y la imposición de una
condición común en la frontera. La inscripción del registro en un espacio de
aparato, establecida en la sección~\ref{x_reciprocidad:vii:sec:observador-aparato}, permite
expresar esas operaciones mediante aplicaciones lineales sin confundir la
publicación observable con el estado que conserva sus antecedentes.

La dirección temporal, la conjugación de hojas y la reflexión espacial
tienen dominios distintos. Sus relaciones se fijan primero sobre el registro;
después se especifican los entrelazamientos que utiliza una realización
dinámica. Este orden determina qué información permanece en los propagadores
y qué condición introduce la elección de una frontera absorbente.

\section{Acción de rutas e involuciones del registro}
\label{x_reciprocidad:vii:subsec:accion-bidireccional}

Una ruta finita contiene sus posiciones, direcciones cardinales
\((d_1,\ldots,d_N)\), hojas \((\mu_1,\ldots,\mu_N)\), con
\(\mu_t\in\{\Sigma,\Pi\}\), y los incrementos de su registro
enriquecido. Para \(\lambda\geq0\), definimos
\begin{equation}
 \mathcal A_\lambda(\gamma)
 =\sum_{t=2}^N
 \bigl(1-\delta_{d_t,d_{t-1}}
       +\lambda(1-\delta_{\mu_t,\mu_{t-1}})\bigr).
 \label{x_reciprocidad:vii:eq:accion-ruta-bidireccional}
\end{equation}
El primer sumando cuenta cambios de dirección; el segundo cuenta cambios
de hoja con el peso declarado. Para \(N\leq1\), la suma es cero.
La acción general del sector conserva además sus términos de frontera y
memoria:
\begin{equation}
 \mathcal S[\gamma]
 =\sum_{a\subset\gamma}\mathcal L_a
  +\mathcal B_{\rm in}(\partial_-\gamma)
  +\mathcal B_{\rm out}(\partial_+\gamma)
  +\mathcal M(\gamma).
 \label{x_reciprocidad:vii:eq:accion-enriquecida-frontera}
\end{equation}
Los términos de esta expresión especifican una clase variacional con
extremos tipados y memoria conservada. La dinámica del sector determina
\(\mathcal L_a\), las condiciones admisibles y el término \(\mathcal M\).

Sobre los registros definimos \(\mathcal C_{\rm r}\) como el intercambio
global de hojas; \(\mathcal P_{\rm r}\) como una reflexión de la carta
espacial y de sus direcciones; y \(\mathcal T_{\rm r}\) como la inversión
del orden temporal, la sustitución de cada dirección por su opuesta y el
intercambio de extremos. La reflexión cardinal es una permutación involutiva.
La inversión del registro firmado es
\(C_M(a_1,\ldots,a_N)=(-a_N,\ldots,-a_1)\).
Las hojas, los cocientes y la memoria se transportan con la ruta completa;
la lectura escalar de una palabra constituye una operación posterior.

\begin{proposition}[Invariancia de la acción local de ruta]
\label{x_reciprocidad:vii:prop:invariancia-ruta-bidireccional}
En el dominio de registros anterior,
\[
 \mathcal A_\lambda(\gamma)
 =\mathcal A_\lambda(\mathcal C_{\rm r}\gamma)
 =\mathcal A_\lambda(\mathcal P_{\rm r}\gamma)
 =\mathcal A_\lambda(\mathcal T_{\rm r}\gamma)
 =\mathcal A_\lambda(\mathcal C_{\rm r}\mathcal P_{\rm r}
                              \mathcal T_{\rm r}\gamma).
\]
\end{proposition}
\begin{proof}
Una permutación global del alfabeto de direcciones conserva la igualdad
entre direcciones consecutivas. El intercambio de las dos hojas conserva
la igualdad entre hojas consecutivas. La inversión del orden establece una
biyección entre pares adyacentes, y la sustitución por direcciones opuestas
también conserva sus igualdades. Cada operación conserva los dos recuentos
de \eqref{x_reciprocidad:vii:eq:accion-ruta-bidireccional}; su composición los conserva.
\end{proof}

La proposición determina la simetría del funcional local. Para que una
transformación actúe además sobre un problema dinámico con fronteras fijas,
debe transportar su clase de prolongaciones admisibles y los otros tres
términos de \eqref{x_reciprocidad:vii:eq:accion-enriquecida-frontera}. En particular, dos
rutas con igual publicación final pueden conservar distintas clases
variacionales por su memoria.

Una interacción entre dos secciones realizadas en una frontera \(B\)
utiliza un mapa de acoplamiento
\[
 \mathcal I_B:E_B^{(1)}\otimes E_B^{(2)}\longrightarrow E_B^{(12)}
\]
y un término \(\mathcal S_{\rm int}\) de la misma acción. Su covariancia
se expresa por
\[
 \rho_{12}(g)\mathcal I_B
 =\mathcal I_B\bigl(\rho_1(g)\otimes\rho_2(g)\bigr).
\]
Al componer las historias, el registro de la frontera compartida se cuenta
una vez; las memorias de ambos lados se transportan a la sección compuesta.
El acoplamiento y esta condición de covariancia son datos de la interacción
concreta, distintos de las involuciones del registro.

\section{Reversión de la evolución realizada}
\label{x_reciprocidad:vii:subsec:reversion-evolucion}

La realización del registro proporciona un espacio de Hilbert real
\(\mathcal H_{\mathbb R}\), una estructura compleja ortogonal
\(J_E\), con \(J_E^2=-I\), y un Hamiltoniano autoadjunto \(H\).
Su dominio \(\mathcal D(H)\) es invariante por \(J_E\); se supone
\(HJ_E=J_EH\) en ese dominio. Sea \(\mathcal C\) una involución
ortogonal que conserva \(\mathcal D(H)\) y satisface
\begin{equation}
 \mathcal C J_E\mathcal C^{-1}=-J_E,\qquad
 \mathcal C H\mathcal C^{-1}=H.
 \label{x_reciprocidad:vii:eq:dominio-reversion-evolucion}
\end{equation}
La constante interna de acción \(\hbar_{\rm int}>0\) conserva su
generación anterior; aquí normaliza el grupo de evolución.

\begin{theorem}[Conjugación temporal del grupo]
\label{x_reciprocidad:vii:thm:reversion-evolucion}
El generador \(A=-J_EH/\hbar_{\rm int}\), con dominio \(\mathcal D(H)\),
es antiadjunto. Su grupo ortogonal, unitario respecto de \(J_E\), cumple
\begin{equation}
 U(t)=\exp(-J_EHt/\hbar_{\rm int}),\qquad
 \mathcal C U(t)\mathcal C^{-1}=U(-t).
 \label{x_reciprocidad:vii:eq:reversion-grupo-bidireccional}
\end{equation}
\end{theorem}
\begin{proof}
La ortogonalidad y \(J_E^2=-I\) implican \(J_E^*=-J_E\).
La invariancia del dominio y la conmutación con \(H\) permiten tratar
\(H\) como operador autoadjunto complejo en el espacio definido por
\(J_E\). Por cálculo funcional, \(-J_EH/\hbar_{\rm int}\) es
antiadjunto y genera el grupo anunciado. En dimensión finita, las mismas
afirmaciones siguen de la serie exponencial y de \(A^*=-A\).
Las hipótesis sobre \(\mathcal C\) dan
\(\mathcal C A\mathcal C^{-1}=-A\). Conjugar el grupo, o su cálculo
funcional, produce \(\exp(-tA)=U(-t)\).
\end{proof}

El mapa que realiza la involución firmada \(C_M\) como \(\mathcal C\)
conserva la orientación temporal y el dominio anterior. La conjugación de
carga de una representación de campos conserva, además, sus propios
caracteres y operadores. La igualdad
\eqref{x_reciprocidad:vii:eq:reversion-grupo-bidireccional} fija la reversión temporal
del grupo bajo las condiciones expresadas.

\section{Propagadores del canal cronológico de una ruta}
\label{x_reciprocidad:vii:subsec:green-cronologico}

La inscripción isométrica del registro admite las prolongaciones unitarias
del teorema~\ref{x_reciprocidad:vii:thm:observador-entrelazador}. Fijemos una realización
bilateral de un canal cronológico con fibras \(\mathcal H_n\),
\(n\in\mathbb Z\), y transportes unitarios
\(U_n:\mathcal H_n\to\mathcal H_{n+1}\).
La elección incluye los espacios auxiliares cuando se requieren; conserva
la ruta y sus registros anteriores. Denotamos por \(U(n,k)\) el transporte
de \(k\) a \(n\), con
\[
 U(k,k)=I,\qquad U(n,k)U(k,l)=U(n,l),\qquad
 U(n,k)^*=U(k,n).
\]
Para \(n>k\), es el producto \(U_{n-1}\cdots U_k\); para \(n<k\),
es el inverso del transporte de \(n\) a \(k\).

El espacio de fuentes \(\mathcal J_c\) contiene las secciones de soporte
finito. El espacio de respuestas \(\mathcal F\) contiene todas las
secciones de esas fibras. La acción cuadrática de la ruta, en unidades
internas de esta realización, utiliza
\[
 (D\psi)_n=\psi_{n+1}-U_n\psi_n,\qquad
 E[\psi]=\sum_n\|(D\psi)_n\|^2
\]
para secciones de soporte finito. El operador cinético con la convención
de signo \(L=-D^*D\) es
\begin{equation}
 (L\psi)_n=U_n^*\psi_{n+1}-2\psi_n+U_{n-1}\psi_{n-1}.
 \label{x_reciprocidad:vii:eq:diferencia-covariante-bidireccional}
\end{equation}
En efecto, \((D^*z)_n=z_{n-1}-U_n^*z_n\), y la sustitución produce
esa fórmula. La variación de \(E\) es
\(-2\operatorname{Re}\langle\delta\psi,L\psi\rangle\).

\begin{theorem}[Green exactos de la diferencia covariante]
\label{x_reciprocidad:vii:thm:green-cronologico}
Escribiendo \(x_+=\max(x,0)\), las aplicaciones
\begin{align}
 (G_{\rm ret}j)_n&=\sum_k(n-k)_+U(n,k)j_k,\label{x_reciprocidad:vii:eq:green-ret-cronologico}\\
 (G_{\rm adv}j)_n&=\sum_k(k-n)_+U(n,k)j_k\label{x_reciprocidad:vii:eq:green-adv-cronologico}
\end{align}
de \(\mathcal J_c\) en \(\mathcal F\) satisfacen
\[
 LG_{\rm ret}j=LG_{\rm adv}j=j.
\]
La respuesta retardada se anula antes y en la primera posición de la
fuente; la avanzada se anula después y en la última. Son las únicas
soluciones que se anulan suficientemente lejos en la dirección respectiva.
Para \(\psi\) de soporte finito también se cumple
\(G_{\rm ret}L\psi=G_{\rm adv}L\psi=\psi\).
Respecto del emparejamiento con fuentes de soporte finito,
\(G_{\rm adv}=G_{\rm ret}^{\dagger}\).
\end{theorem}
\begin{proof}
Cada suma tiene tantos términos como posiciones de la fuente.
La composición de los transportes reduce la aplicación de \(L\) a
la identidad escalar
\[
 (n+1-k)_+-2(n-k)_++(n-1-k)_+=\delta_{nk}.
\]
La función \((k-n)_+\) satisface la misma identidad. Esto prueba
las dos ecuaciones y los soportes. Para la unicidad, una solución
homogénea que se anula en dos posiciones consecutivas queda determinada
como cero en la siguiente por \eqref{x_reciprocidad:vii:eq:diferencia-covariante-bidireccional};
la recurrencia, usando transportes invertibles, se prolonga en ambas
direcciones. La diferencia de dos soluciones con el mismo soporte causal
se anula suficientemente lejos y, por tanto, es cero.
Aplicar este argumento a \(\psi\) y a \(G_{\rm ret}L\psi\), o al
caso avanzado, da las identidades por la derecha sobre soporte finito.
Finalmente, para dos fuentes compactas, la suma doble del emparejamiento
es finita y
\[
 \bigl((k-n)_+U(k,n)\bigr)^*=(k-n)_+U(n,k).
\]
Intercambiar los índices demuestra la adjunción indicada.
\end{proof}

Las respuestas pertenecen a \(\mathcal F\); en general crecen
linealmente fuera de la fuente y no pertenecen a \(\ell^2\).
El dominio de soporte finito y su emparejamiento especifican la adjunción
de Green utilizada aquí. El resultado construye la propagación del canal
cronológico desde su acción; una realización de campos espaciales conserva
además su operador, geometría y dominio de propagación propios.

Sea \(C_0\) una involución isométrica, lineal o antilineal, de
\(\mathcal H_0\). El transporte desde la fibra de referencia construye
\[
 C_n=U(-n,0)C_0U(0,n):\mathcal H_n\longrightarrow\mathcal H_{-n},
 \qquad (\mathcal C\psi)_n=C_{-n}\psi_{-n}.
\]
Se tiene \(C_{-n}C_n=I\) y
\(C_{n+1}U_n=U_{-n-1}^*C_n\), por composición.
En particular, \(\mathcal C^2=I\), y el cambio \((n,k)\mapsto(-n,-k)\)
en las sumas anteriores prueba
\begin{equation}
 \mathcal C L=L\mathcal C,\qquad
 \mathcal C G_{\rm ret}\mathcal C^{-1}=G_{\rm adv}.
 \label{x_reciprocidad:vii:eq:intercambio-green-cronologico}
\end{equation}
Cuando \(C_0\) es la realización declarada de la involución del registro,
estas fórmulas transportan esa involución a todo el canal, sin identificar
el retorno de fase con un reinicio de memoria.

\section{Intercambio causal y descomposición de la respuesta}
\label{x_reciprocidad:vii:subsec:green-realizacion-campos}

En una realización de campos, sea
\(L:\mathcal D(L)\subset\mathcal H_\Phi\to\mathcal H_J\)
un operador con Green retardado y avanzado únicos en los dominios de
fuentes admisibles. Denotamos por \(J^+(S)\) y \(J^-(S)\) los
dominios causales futuro y pasado de un soporte \(S\). Las condiciones son
\[
 \begin{gathered}
 LG_{\rm ret}=LG_{\rm adv}=I,\\
 \operatorname{supp}(G_{\rm ret}j)\subset J^+(\operatorname{supp}j),\\
 \operatorname{supp}(G_{\rm adv}j)\subset J^-(\operatorname{supp}j).
 \end{gathered}
\]
Sean \(C_\Phi,C_J\) involuciones isométricas, ambas lineales o ambas
antilineales, que transportan esos dominios,
invierten sus orientaciones causales y satisfacen
\(LC_\Phi=C_JL\). Se conserva un emparejamiento de Green para el cual
\(G_{\rm adv}=G_{\rm ret}^{\dagger}\).

\begin{proposition}[Conjugación de los Green y paridades]
\label{x_reciprocidad:vii:prop:green-paridades}
Con las condiciones anteriores,
\[
 C_\Phi G_{\rm ret}C_J^{-1}=G_{\rm adv}.
\]
Sobre el espacio real de operadores de respuesta, definamos
\[
 \Theta(K)=C_\Phi K C_J^{-1},\qquad
 \mathsf P_\pm(K)=\tfrac12(K\pm\Theta(K)).
\]
Son proyecciones complementarias para la involución \(\Theta\), y
\begin{equation}
 G_{\rm sym}=\tfrac12(G_{\rm ret}+G_{\rm adv}),\qquad
 G_{\rm rad}=\tfrac12(G_{\rm ret}-G_{\rm adv})
 \label{x_reciprocidad:vii:eq:descomposicion-green}
\end{equation}
satisfacen \(\Theta(G_{\rm sym})=G_{\rm sym}\),
\(\Theta(G_{\rm rad})=-G_{\rm rad}\),
\(G_{\rm sym}^{\dagger}=G_{\rm sym}\) y
\(LG_{\rm rad}=0\). El soporte de la respuesta simétrica está contenido
en la unión de los dos dominios causales.
\end{proposition}
\begin{proof}
El operador conjugado resuelve la ecuación de fuente, pues
\[
 LC_\Phi G_{\rm ret}C_J^{-1}=C_JLG_{\rm ret}C_J^{-1}=I.
\]
La inversión de la orientación lleva su soporte al dominio avanzado;
su unicidad lo identifica con \(G_{\rm adv}\). Las involuciones dan
\(\Theta^2=I\); expandir \((I\pm\Theta)^2\) y
\((I+\Theta)(I-\Theta)\) demuestra las afirmaciones sobre proyecciones.
La adjunción de Green demuestra la simetría y la diferencia de las dos
ecuaciones de fuente da \(LG_{\rm rad}=0\). La suma de dos respuestas
tiene soporte contenido en la unión de sus soportes.
\end{proof}

Si las involuciones son antilineales, \(\Theta\) es lineal sobre
escalares reales; ése es el espacio operatorio usado en la proposición.
La conjugación del grupo temporal, la condición de soporte y la
adjunción de Green son tres propiedades distintas, transmitidas por los
mapas de realización correspondientes.

La reducción causal de la actualización bilateral, desarrollada en la
sección sobre conservación y recuperación, identifica otro origen preciso
del retardo: la eliminación de una componente de memoria que sigue
participando en la evolución conjunta. Su núcleo discreto conserva el
orden de las interacciones y la contribución de la memoria inicial.
Los operadores de Green aquí construidos resuelven, en cambio, una
ecuación de propagación con condiciones de soporte. Ambas construcciones
pueden componerse cuando se fija el mapa de realización que identifica
sus dominios y operadores; sus respectivas condiciones de causalidad
permanecen explícitas.

\section{Proyección de corrientes absorbidas y variación de la acción}
\label{x_reciprocidad:vii:subsec:proyeccion-absorcion}

Sea \(\gamma_\partial:\mathcal H_\Phi\to\mathcal H_\partial\)
la lectura de frontera, distinta de la ruta \(\gamma\). Definimos
\[
 \mathcal A_\partial=\gamma_\partial G_{\rm rad}.
\]
La ventana finita fija el espacio de fuentes y la lectura de frontera;
los Green conservan su dominio de respuestas. En un nivel finito con
productos internos fijados, su inversa de
Moore--Penrose determina
\begin{equation}
 P_{\rm abs}=I-\mathcal A_\partial^+\mathcal A_\partial.
 \label{x_reciprocidad:vii:eq:proyector-absorcion}
\end{equation}

\begin{theorem}[Condición de absorción y respuesta de borde]
\label{x_reciprocidad:vii:thm:proyeccion-absorcion}
\(P_{\rm abs}\) es la proyección ortogonal sobre
\(\ker\mathcal A_\partial\). Para cualquier fuente \(j\), la fuente
\(j_{\rm abs}=P_{\rm abs}j\) cumple
\begin{equation}
 \gamma_\partial G_{\rm ret}j_{\rm abs}
 =\gamma_\partial G_{\rm adv}j_{\rm abs}.
 \label{x_reciprocidad:vii:eq:igualdad-borde-absorbido}
\end{equation}
Es la fuente absorbida más próxima a \(j\) para el producto interno
declarado. Existen fuentes absorbidas no nulas exactamente cuando
\(\operatorname{rank}\mathcal A_\partial<\dim\mathcal H_J\).
\end{theorem}
\begin{proof}
La descomposición ortogonal
\(\mathcal H_J=\ker\mathcal A_\partial\oplus
 \operatorname{im}\mathcal A_\partial^*\) muestra que
\(\mathcal A_\partial^+\mathcal A_\partial\) es la identidad
en el segundo sumando y cero en el primero. Su complemento es el
proyector anunciado. Por definición,
\(2\mathcal A_\partial j_{\rm abs}
 =\gamma_\partial(G_{\rm ret}-G_{\rm adv})j_{\rm abs}=0\).
Para \(k\in\ker\mathcal A_\partial\), la descomposición ortogonal da
\[
 \|j-k\|^2=\|(I-P_{\rm abs})j\|^2+\|P_{\rm abs}j-k\|^2.
\]
El mínimo único corresponde a \(k=P_{\rm abs}j\).
La equivalencia final es el teorema de rango y nulidad.
\end{proof}

En el emparejamiento de fuente y campo, la acción simétrica es
\(\mathscr S[j]=\tfrac12\langle j,G_{\rm sym}j\rangle\).
La adjunción declarada hace real esa cantidad. Para variaciones reales,
\[
 \mathrm d\mathscr S[j](h)
 =\operatorname{Re}\langle h,G_{\rm sym}j\rangle.
\]
La acción restringida a fuentes absorbidas se escribe sobre el espacio
ambiente como
\begin{equation}
 \mathscr S_{\rm abs}[j]
 =\tfrac12\langle P_{\rm abs}j,G_{\rm sym}P_{\rm abs}j\rangle,
 \quad
 \mathrm d\mathscr S_{\rm abs}[j](h)
 =\operatorname{Re}\langle P_{\rm abs}h,G_{\rm sym}P_{\rm abs}j\rangle.
 \label{x_reciprocidad:vii:eq:variacion-restringida-absorcion}
\end{equation}
Cuando el emparejamiento identifica campo y fuente por Riesz, el gradiente
ambiente es \(P_{\rm abs}G_{\rm sym}P_{\rm abs}j\).
La respuesta de campo sigue siendo \(G_{\rm sym}j_{\rm abs}\); el
proyector del gradiente expresa la restricción de las variaciones.
Las fórmulas se demuestran expandiendo el término cuadrático de
\(j+th\) y utilizando la autoadjunción. Un término material
\(\mathscr S_{\rm matter}\) añade su propio diferencial.

\section{Condición absorbente explícita en una ventana cronológica}
\label{x_reciprocidad:vii:subsec:absorcion-momentos}

El canal de la subsección~\ref{x_reciprocidad:vii:subsec:green-cronologico} permite
evaluar esa condición sin introducir una respuesta objetivo. Sea una
fuente apoyada en \(0,\ldots,N\), con \(N\geq1\). Transportemos
sus componentes a la fibra de referencia:
\[
 \widehat j_k=U(0,k)j_k,\qquad
 M_0(j)=\sum_{k=0}^N\widehat j_k,\qquad
 M_1(j)=\sum_{k=0}^Nk\widehat j_k.
\]
La identidad \((n-k)_+-(k-n)_+=n-k\) da
\begin{equation}
 (G_{\rm ret}-G_{\rm adv})j\big|_n
 =U(n,0)\bigl(nM_0(j)-M_1(j)\bigr).
 \label{x_reciprocidad:vii:eq:diferencia-green-momentos}
\end{equation}

\begin{proposition}[Absorción y dos momentos transportados]
\label{x_reciprocidad:vii:prop:absorcion-momentos}
Si la traza evalúa el campo en dos posiciones distintas \(a,b\),
la igualdad de trazas avanzada y retardada equivale a
\(M_0(j)=M_1(j)=0\). En una fibra de dimensión \(d\), el espacio
de tales fuentes tiene dimensión \((N-1)d\). Para \(N\geq2\),
contiene las fuentes no nulas con componentes transportadas
\((z,-2z,z,0,\ldots,0)\), \(z\neq0\).
\end{proposition}
\begin{proof}
La invertibilidad de \(U(a,0)\) y \(U(b,0)\) reduce la condición
de borde a \(aM_0-M_1=0\) y \(bM_0-M_1=0\).
Restar las ecuaciones da \((a-b)M_0=0\), y después \(M_1=0\).
El mapa de los dos momentos es sobreyectivo: para valores prescritos
\((z_0,z_1)\), basta tomar
\(\widehat j_0=z_0-z_1\), \(\widehat j_1=z_1\) y las demás
componentes cero. Su rango es \(2d\); la nulidad es
\((N+1)d-2d\). La fuente indicada tiene ambos momentos nulos.
\end{proof}

En el marco transportado, escribamos
\[
 B_N=\begin{pmatrix}1&1&\cdots&1\\0&1&\cdots&N\end{pmatrix}.
\]
El proyector absorbente queda evaluado por
\begin{equation}
 P_N=\bigl[I-B_N^{\mathsf T}(B_NB_N^{\mathsf T})^{-1}B_N\bigr]
       \otimes I_{\mathcal H_0}.
 \label{x_reciprocidad:vii:eq:proyector-momentos-cronologicos}
\end{equation}
Las dos filas de \(B_N\) son independientes para \(N\geq1\), por lo
que la inversa existe. El transporte unitario entre marcos lleva este
proyector a las fibras originales. Para \(N=2\),
\[
 P_2=\frac16
 \begin{pmatrix}1&-2&1\\-2&4&-2\\1&-2&1\end{pmatrix}
 \otimes I_{\mathcal H_0}.
\]
Esta construcción produce una clase absorbida no nula y su operación de
proyección en el canal cronológico. Su amplitud y su identificación con una
corriente material conservan las reglas del sector físico correspondiente.

\section{Refinamiento y continuidad de la condición de borde}
\label{x_reciprocidad:vii:subsec:refinamiento-absorcion}

Sean \(R_m^\Phi,R_m^J,R_m^\partial\) los mapas de refinamiento de
campos, fuentes y fronteras. Supongamos que entrelazan los dos Green y la
traza. Por composición se obtiene
\begin{equation}
 R_m^\Phi G_{{\rm sym},m}=G_{{\rm sym},m+1}R_m^J,\qquad
 R_m^\partial\mathcal A_{\partial,m}
 =\mathcal A_{\partial,m+1}R_m^J.
 \label{x_reciprocidad:vii:eq:refinamiento-lectura-absorcion}
\end{equation}
La segunda igualdad lleva una fuente absorbida a otra fuente absorbida.
La naturalidad de los proyectores conserva adicionalmente la estructura
ortogonal de las realizaciones.

\begin{proposition}[Naturalidad de la proyección absorbente]
\label{x_reciprocidad:vii:prop:naturalidad-proyeccion-absorcion}
Para una isometría de fuentes \(R=R_m^J\), se tiene
\(P_{{\rm abs},m+1}R=RP_{{\rm abs},m}\) exactamente cuando
\[
 R(\ker\mathcal A_{\partial,m})
       \subset\ker\mathcal A_{\partial,m+1},\qquad
 R((\ker\mathcal A_{\partial,m})^\perp)
       \subset(\ker\mathcal A_{\partial,m+1})^\perp.
\]
Con esas condiciones en todos los niveles, los proyectores definen una
proyección ortogonal en la completación del límite inductivo isométrico
de los espacios de fuentes.
\end{proposition}
\begin{proof}
Se descompone cada fuente en sus componentes en el núcleo y en su
complemento ortogonal. Las dos inclusiones implican que aplicar la
proyección antes o después del refinamiento conserva respectivamente la
primera componente y anula la segunda. Recíprocamente, la igualdad de
operadores aplicada a cada componente da las dos inclusiones.
En el límite inductivo, la igualdad permite definir un operador sobre la
unión de los niveles sin depender del representante. Todos los proyectores
tienen norma a lo sumo uno, de modo que el operador se extiende a la
completación. Sus identidades de autoadjunción e idempotencia, válidas en
la unión densa, pasan por continuidad a la completación.
\end{proof}

Para prolongar también una respuesta de Green a una completación se
conservan sus dominios y las estimaciones de continuidad correspondientes.
La propagación cronológica de soporte finito y la factorización de la
respuesta por la imagen de borde, demostrada en
\ref{x_reciprocidad:vii:thm:observador-factorizacion}, mantienen así sus dominios
respectivos durante el refinamiento.

La realización física de emisión y absorción compone estos operadores con
las corrientes materiales, la acción y las condiciones causales del sector.
El entrelazamiento debe conservar la involución de ruta, el emparejamiento
de fuente y campo, la acción transportada y las condiciones de frontera a
cada nivel. La construcción precedente aporta la reversión del grupo, los
Green del canal cronológico, la descomposición temporal y una proyección
absorbente explícita; las realizaciones espaciales emplean sus mapas y
condiciones de causalidad propios.

% END OWNER

\chapter{Realización de la acción y comparación metrológica}

% BEGIN OWNER /Users/ruben/Documents/New project/output/EXTREMA_MEDIA_RAZON_RECIPROCIDAD_ES_EN_20260918/ARTICULO/ES/source/antecedentes_i/sections/revision_planck_contraste.tex
\section{Contraste de la constante de Planck reducida}
\label{x_reciprocidad:sec:revision-planck-contraste}

El Sistema Internacional vigente desde el 20 de mayo de 2019 fija
exactamente \(h=6.62607015\times10^{-34}\,\mathrm{J\,s}\)
\cite{x_reciprocidad:bipm2018}. Por tanto, su constante reducida es
\[
 \hbar_{\mathrm{SI}}=\frac{6.62607015}{2\pi}\,10^{-34}\,
 \mathrm{J\,s}
 =1.0545718176461563912\ldots\times10^{-34}\,\mathrm{J\,s}.
\]
Los puntos suspensivos corresponden a la expansión de una expresión
exacta, no a una incertidumbre experimental de \(h\) en el SI actual.
Este valor de referencia se utiliza aquí después de la evaluación de
los lectores de HMT; su función es comparativa.

Las dos coordenadas del retorno del artículo son
\begin{align*}
 H_5&=1.0545718176461544696\ldots,\\
 \eta_{\mathrm{ret}}&=1.0545718169150390611\ldots.
\end{align*}
En la identificación de \(\mathcal U_S\) con \(\mathrm{J\,s}\), la
diferencia relativa de la sección anterior respecto de
\(\hbar_{\mathrm{SI}}\) es aproximadamente \(-1.82\times10^{-15}\).
La de la sección posterior es aproximadamente \(-6.93\times10^{-10}\).
La segunda diferencia contiene el efecto del retorno que utiliza el
corrector electrónico; no equivale al error de evaluación de la
primera sección.

La precisión de estas comparaciones permite reconocer la proximidad
cuantitativa de la construcción de \(H_5\) y, simultáneamente, distinguir
proximidad de identidad exacta. En el presente manuscrito las expresiones
exhibidas no proporcionan una igualdad exacta con \(h/(2\pi)\).
El resultado demostrado sobre \(\Sigma_H\) es su década \(-34\);
el resultado analítico sobre \(H_5\) es la evaluación del carácter
declarado. La identificación física se examina en esta carta y con las
diferencias explícitas que acabamos de dar. El valor SI no interviene
en la definición de la norma determinantal ni en la recurrencia regional.

% END OWNER

